% !TeX program = xelatex
% The article class is initialized at 12pt; the body size is set to exactly
% 13pt by the explicit \normalsize definition below.
\documentclass[b5paper,12pt,twoside]{article}

\usepackage{fix-cm}
\usepackage[T1]{fontenc}
\usepackage{newpxtext}
\usepackage[
  b5paper,
  includeheadfoot,
  left=0.7cm,
  right=0.7cm,
  top=1.0cm,
  bottom=1.0cm,
  headheight=16pt,
  headsep=0.35cm,
  footskip=0.65cm
]{geometry}
\usepackage{microtype}
\usepackage{amsmath,amssymb,amsthm,mathtools}
\usepackage{newpxmath}
\usepackage{array}
\usepackage{enumitem}
\usepackage[numbers,sort&compress]{natbib}
\usepackage{fancyhdr}
\usepackage{xcolor}
\definecolor{LinkBlue}{HTML}{1F5A8A}
\definecolor{CiteGreen}{HTML}{2D6A4F}
\definecolor{UrlPurple}{HTML}{6F3C8C}
\usepackage[
  unicode,
  linktoc=all,
  colorlinks=true,
  linkcolor=LinkBlue,
  citecolor=CiteGreen,
  urlcolor=UrlPurple,
  anchorcolor=LinkBlue
]{hyperref}

\hypersetup{
  pdftitle={On Rings of Operators. Reduction Theory},
  pdfauthor={John von Neumann},
  pdfsubject={Faithful LaTeX transcription of the 1938 manuscript published in 1949},
  pdfstartview=FitH,
  bookmarksnumbered=true
}
\setcitestyle{numbers,open={(},close={)},comma}
\makeatletter
\renewcommand{\@biblabel}[1]{#1.}
\makeatother
\linespread{1.1}
\makeatletter
\renewcommand{\normalsize}{%
  \@setfontsize\normalsize{13pt}{15.6pt}%
  \abovedisplayskip 12pt plus 3pt minus 7pt%
  \abovedisplayshortskip 0pt plus 3pt%
  \belowdisplayshortskip 7pt plus 3pt minus 4pt%
  \belowdisplayskip \abovedisplayskip%
  \let\@listi\@listI
}
\renewcommand\section{\@startsection{section}{1}{0pt}%
  {0.9\baselineskip plus 0.25\baselineskip minus 0.15\baselineskip}%
  {0.35\baselineskip}%
  {\normalfont\normalsize\bfseries}}
\makeatother
\AtBeginDocument{\normalsize}

\setlength{\parindent}{1.5em}
\setlength{\parskip}{0pt}
\setlength{\emergencystretch}{4em}
\setlist{nosep,leftmargin=2.2em}
\allowdisplaybreaks[1]
\numberwithin{equation}{section}

\pagestyle{fancy}
\fancyhf{}
\fancyhead[CE,CO]{\small\textsc{On Rings of Operators. Reduction Theory}}
\fancyhead[LE,RO]{\small\thepage}
\renewcommand{\headrulewidth}{0pt}
\renewcommand{\footrulewidth}{0pt}
\fancypagestyle{plain}{%
  \fancyhf{}
  \fancyfoot[C]{\raisebox{1.2mm}[0pt][0pt]{\thepage}}
  \renewcommand{\headrulewidth}{0pt}%
}

\newtheoremstyle{vnplain}
  {0.65\baselineskip}{0.45\baselineskip}
  {\itshape}{}% body font, indent
  {\normalfont\bfseries}{.}{0.5em}{}
\newtheoremstyle{vndefinition}
  {0.65\baselineskip}{0.45\baselineskip}
  {\itshape}{}% body font, indent
  {\normalfont\bfseries}{.}{0.5em}{}
\theoremstyle{vnplain}
\newtheorem{theorem}{Theorem}
\renewcommand{\thetheorem}{\Roman{theorem}}
\newtheorem{lemma}{Lemma}
\theoremstyle{vndefinition}
\newtheorem{definition}{Definition}
\newtheorem{remark}{Remark}
\renewcommand{\proofname}{\textnormal{\textit{Proof}}}

\newcommand{\secref}[1]{\hyperref[#1]{\S\ref*{#1}}}
\newcommand{\tocsec}[1]{\hyperref[sec:#1]{\S#1}}
\newcommand{\vnpart}[2]{%
  \par\addvspace{1.2\baselineskip}%
  \phantomsection\label{part:#1}%
  \addcontentsline{toc}{part}{Part #1: #2}%
  \begin{center}\large\bfseries\textsc{Part #1: #2}\end{center}%
  \nobreak\vspace{0.35\baselineskip}%
}
\newcommand{\partref}[1]{\hyperref[part:#1]{Part~#1}}
\newcommand{\defref}[1]{\hyperref[#1]{Definition~\ref*{#1}}}
\newcommand{\lemref}[1]{\hyperref[#1]{Lemma~\ref*{#1}}}
\newcommand{\thmref}[1]{\hyperref[#1]{Theorem~\ref*{#1}}}
\newcommand{\origpage}[1]{\phantomsection\label{origpage:#1}}
\newcommand{\esssup}{\operatorname*{ess\,sup}}
\newcommand{\essinf}{\operatorname*{ess\,inf}}
\newcommand{\vnmod}[1]{\mathrel{\cdots}\operatorname{mod}\,#1}
\DeclareMathOperator{\RePart}{Re}
\DeclareMathOperator{\ImPart}{Im}

\renewcommand{\refname}{Bibliography}
\begin{document}
\setcounter{page}{1}

% PDF pages 1--28 (original journal pages 400--427).

\thispagestyle{plain}
% With includeheadfoot the running header is part of the 1 cm top margin.
% The first page has no running header, so reclaim that reserved space here.
\vspace*{-0.67cm}
\noindent
{\fontsize{9pt}{10.5pt}\selectfont
  \textsc{\bfseries Annals of Mathematics}\par
  \noindent\textbf{Vol. 50, No. 2, April, 1949}\par
}
\vspace{0.15cm}

\begin{center}
  {\normalsize\bfseries ON RINGS OF OPERATORS.\enspace REDUCTION THEORY\par}
  \vspace{0.10\baselineskip}
  {\fontsize{12pt}{14.4pt}\selectfont By \textsc{John von Neumann}\par}
  \vspace{0.10\baselineskip}
  {\fontsize{11pt}{13.2pt}\selectfont (Received June 7, 1948)\par}
\end{center}

\vspace{0.25\baselineskip}
\noindent\textit{Note.}\quad
This paper was written in 1937--38. Various other commitments prevented the
author from effecting some changes, which he had intended to carry out before
publishing the paper. This delayed the publication until the present time.

The paper is now published in the 1938 form, with only minor modifications:
\begin{enumerate}[label=\arabic*)]
\item \hyperref[fn:17]{Footnote 17} following
\hyperref[def:9]{Definition 9} and \hyperref[fn:18]{footnote 18} following
\hyperref[thm:IX]{Theorem IX},
both in \secref{sec:22}, are new.
\item The second reference 12 is new.
\item \hyperref[lem:22]{Lemma 22} and the derivation of the properties of $E\mathbin{\#}F$ based on it differs
from the 1938 version. This Lemma was, however, used by the author in his
Princeton Lectures on Operator Theory in 1935.
\item The second part of ($\delta$) in \secref{sec:24} differs from the 1938
version.
\end{enumerate}

\begin{center}
  {\large\scshape Analytic Table of Contents}
\end{center}

\noindent\textsc{Part I: Generalized direct sums.}
\begin{description}
\item[\tocsec{1}] Description of unitary spaces.
\item[\tocsec{2}] $N$-functions $\rho(\lambda)$. Stieltjes-measure
for $\sigma(\lambda)$. $\sigma(\lambda)$-total continuity of $\rho(\lambda)$,
the integral representation. The same for
$\rho_f(\lambda)=\lVert E(\lambda)f\rVert^2$;
$\sigma(\lambda)$-total-continuity of a resolution of unity $E(\lambda)$.
\item[\tocsec{3}] Definition and discussion of the notion of a
generalized direct sum.
\begin{center}\hyperref[def:1]{\textbf{Definition 1.}}\quad
\hyperref[rem:3.1]{\textbf{Remark 1.}}\quad
\hyperref[rem:3.2]{\textbf{Remark 2.}}\quad
\hyperref[rem:3.3]{\textbf{Remark 3.}}\end{center}
\item[\tocsec{4}] Discussion of the special case of a totally
discontinuous $\sigma(\lambda)$, that is of a discrete (common) direct sum.
\item[\tocsec{5}] Elementary properties of the
$\displaystyle\int_{-} f(\lambda)\sqrt{d\sigma(\lambda)}$.
(Addition, convergence, certain measurability properties.)
\begin{center}\hyperref[lem:1]{\textbf{Lemma 1.}}\quad
\hyperref[lem:2]{\textbf{Lemma 2.}}\quad
\hyperref[lem:3]{\textbf{Lemma 3.}}\end{center}
\item[\tocsec{6}] Construction of the complete normalised orthogonal
systems $\psi_n(\lambda)$ in $\mathfrak H^{(\lambda)}$. Definition of
measurable families. Existence and (``up to measurable unitary
transformations'') uniqueness of measurable families. Definition of
generalized direct sum, with the help of measurable families.
\begin{center}\hyperref[def:2]{\textbf{Definition 2.}}\quad
\hyperref[rem:6.1]{\textbf{Remark.}}\quad
\hyperref[thm:I]{\textbf{Theorem I.}}\quad
\hyperref[thm:II]{\textbf{Theorem II.}}\end{center}
\item[\tocsec{7}] Construction of the operators $a(\Lambda)$, in
particular $1_K(\Lambda)$, and more particularly
$E(\mu)=1_K(\Lambda)$, where $\lambda\leq\mu\Longleftrightarrow\lambda\in K$.
The $E(\mu)$ form the resolution of unity which belongs to the decomposition
$\mathfrak H^{(\lambda)}$ of $\mathfrak H$. Characterization of the
decomposition $\mathfrak H^{(\lambda)}$ of $\mathfrak H$ with the $E(\lambda)$
and $\sigma(\lambda)$, and conversely, existence and uniqueness theorems.
\begin{center}\hyperref[rem:7.1]{\textbf{Remark.}}\quad
\hyperref[thm:III]{\textbf{Theorem III.}}\end{center}
\end{description}

\noindent\textsc{Part II: Abelian rings and the generalized direct sums.}
\begin{description}
\item[\tocsec{8}] Effect of a change of $\sigma(\lambda)$, if the
$E(\lambda)$ remain the same.
\item[\tocsec{9}] Definition of the equivalence. Elementary
properties.
\begin{center}\hyperref[def:3]{\textbf{Definition 3.}}\quad
\hyperref[rem:9.1]{\textbf{Remark 1.}}\quad
\hyperref[rem:9.2]{\textbf{Remark 2.}}\end{center}
\item[\tocsec{10}] Discussion of the discrete special case
(cf. \secref{sec:4}).
\item[\tocsec{11}] $E(\lambda)$ constant outside of
$-2<\lambda<-1$ can always be secured by equivalence. $P=$ set of all
$a(\Lambda)$, and is an equivalence-invariant.
\begin{center}\hyperref[lem:4]{\textbf{Lemma 4.}}\end{center}
\item[\tocsec{12}] $P$ can be any Abelian ring containing $1$, but
no other ring; $P$ determines $\mathfrak H^{(\lambda)}$, $\sigma(\lambda)$
precisely up to equivalence.
\begin{center}\hyperref[thm:IV]{\textbf{Theorem IV.}}\end{center}
\item[\tocsec{13}] Discussion of the decomposition of an operator
$A\in P'$: $A\sim\sum A(\lambda)$. Necessary and sufficient condition for $A$
and for the $A(\lambda)$, uniqueness of $A$ and of the $A(\lambda)$.
\begin{center}\hyperref[def:4]{\textbf{Definition 4.}}\quad
\hyperref[thm:V]{\textbf{Theorem V.}}\end{center}
\item[\tocsec{14}] Computations rules for
$A\sim\sum A(\lambda)$. Convergence criteria. Operator functions. Behavior
under equivalences.
\end{description}

\noindent\textsc{Part III: Decomposition of rings.}
\begin{description}
\item[\tocsec{15}] $\sigma(\lambda)$-measurable dependence of
$\lambda$: for $f(\lambda)$, $A(\lambda)$, $M(\lambda)$. Discussion of
$f(\lambda)$, $A(\lambda)$.
\begin{center}\hyperref[def:5]{\textbf{Definition 5.}}\quad
\hyperref[rem:15.1]{\textbf{Remark 1.}}\quad
\hyperref[rem:15.2]{\textbf{Remark 2.}}\end{center}
\item[\tocsec{16}] The operator $\lVert\!\lvert A\rvert\!\rVert\leq1$ in the weak
topology is a complete and separable metric space. The possibility of choice in
an analytic set $A$ in a complete and separable space $S$: If $F(a)$ is a
continuous and (real) numerical function in $A$, then a
$\sigma(\lambda)$-measurable $f(\lambda)$ ($\lambda$ a real number,
$\lambda\in K$, $K=$ set of all $F(a)$, $f(\lambda)\in S$) with
$F(f(\lambda))=\lambda$ exists.
\begin{center}\hyperref[lem:5]{\textbf{Lemma 5.}}\end{center}
\item[\tocsec{17}] $\sigma(\lambda)$-measurable dependence and the
operations $M(\lambda)'$, $M_1(\lambda)\cdot M_2(\lambda)\cdots$,
$R(M_1(\lambda),M_2(\lambda),\ldots)$, the rings $0(\lambda)$, $1(\lambda)$,
and the relations $=$, $\ne$, $\subset$, $\supset$, $\subsetneqq$,
$\supsetneqq$. If $M(\lambda)\subset N(\lambda)$ existence of a projection
$E=\sum E(\lambda)$, such that $M(\lambda)\ne N(\lambda)$ implies
$E(\lambda)\notin M(\lambda)$, $E(\lambda)\in N(\lambda)$.
\begin{center}\hyperref[lem:6]{\textbf{Lemma 6.}}\quad
\hyperref[lem:7]{\textbf{Lemma 7.}}\quad
\hyperref[lem:8]{\textbf{Lemma 8.}}\quad
\hyperref[lem:9]{\textbf{Lemma 9.}}\quad
\hyperref[lem:10]{\textbf{Lemma 10.}}\quad
\hyperref[rem:17]{\textbf{Remark.}}\end{center}
\item[\tocsec{18}] Definition and discussion of the relation
$M\approx\sum M(\lambda)$. It is an (essentially) one-to-one decomposition resp.
composition for all rings $M$ with $P\subset M\subset P'$ (and only these
rings).
\begin{center}\hyperref[def:6]{\textbf{Definition 6.}}\quad
\hyperref[lem:11]{\textbf{Lemma 11.}}\quad
\hyperref[lem:12]{\textbf{Lemma 12.}}\quad
\hyperref[thm:VI]{\textbf{Theorem VI.}}\end{center}
\item[\tocsec{19}] Behavior of the decomposition
$M\approx\sum M(\lambda)$ with respect to the operations $M'$, $M_1\cdot
M_2\cdots$, $R(M_1,M_2,\ldots)$, the relation $M\subset N$, and under
equivalences.
\begin{center}\hyperref[lem:13]{\textbf{Lemma 13.}}\quad
\hyperref[lem:14]{\textbf{Lemma 14.}}\quad
\hyperref[lem:15]{\textbf{Lemma 15.}}\end{center}
\end{description}

\noindent\textsc{Part IV: Decomposition of rings with respect to their center.}
\begin{description}
\item[\tocsec{20}] Possible decompositions
$M\approx\sum M(\lambda)$ of a given $M$: The possible $P$'s are all
$P\subset M\cdot M'$. $P=M\cdot M'$ is characterized by the $M(\lambda)$'s
being factors: $M(\lambda)\cdot M(\lambda)'=0(\lambda)$.
\begin{center}\hyperref[thm:VII]{\textbf{Theorem VII.}}\end{center}
\item[\tocsec{21}] Relative dimensionalities
$d_{M(\lambda)}(E)$, weight functions $\Delta_{M(\lambda)}(E)$. The
$\sigma(\lambda)$-measurable $d_{M(\lambda)}(E)$,
$\Delta_{M(\lambda)}(E)$. The set $\widetilde N$.
\begin{center}\hyperref[def:7]{\textbf{Definition 7.}}\quad
\hyperref[def:8]{\textbf{Definition 8.}}\quad
\hyperref[lem:16]{\textbf{Lemma 16.}}\quad
\hyperref[lem:17]{\textbf{Lemma 17.}}\quad
\hyperref[lem:18]{\textbf{Lemma 18.}}\quad
\hyperref[lem:19]{\textbf{Lemma 19.}}\quad
\hyperref[thm:VIII]{\textbf{Theorem VIII.}}\end{center}
\item[\tocsec{22}] The classes $C_c$ where
$c=(I_m,I_n)$, $m,n=1,2,\ldots,\infty$; $(II_\alpha,II_\beta)$,
$\alpha,\beta=1,\infty$; $(III,III)$. Their $\sigma(\lambda)$-measurabilities:
All $C_c$ for $c\ne(II_\infty,II_\infty),(III,III)$, also
$C_{(II_\infty,II_\infty)}+C_{(III,III)}$, while
$C_{(II_\infty,II_\infty)}$, $C_{(III,III)}$ remain undecided. The role of
normal forms.
\begin{center}\hyperref[def:9]{\textbf{Definition 9.}}\quad
\hyperref[lem:20]{\textbf{Lemma 20.}}\quad
\hyperref[thm:IX]{\textbf{Theorem IX.}}\end{center}
\item[\tocsec{23}] Choice theorems for various partially isometric
mappings in the $M(\lambda)$ ($\sigma(\lambda)$-measurability of their
dependence on $\lambda$). Construction of all weight functions $\Delta_M(E)$
in $M$ by those $\Delta_{M(\lambda)}(E)$ in $M(\lambda)$.
\begin{center}\hyperref[lem:21]{\textbf{Lemma 21.}}\quad
\hyperref[lem:22]{\textbf{Lemma 22.}}\quad
\hyperref[thm:X]{\textbf{Theorem X.}}\end{center}
\item[\tocsec{24}] Determination of $a(\lambda)$ by
$\Delta_M(E)=\int a(\lambda)d_{M(\lambda)}(E(\lambda))\,d\sigma(\lambda)$.
$r-\widetilde N$ and $C_{(III,III)}$. Role of condition (iii) in
\hyperref[def:7]{Definition 7}.
Unitary $U$'s and partially isometric $V$'s. Role of condition (ii) eod.:
Countable and finite additivity. The Abelian $M$'s and Lebesgue measure.
Criteria for $E\lesssim F\vnmod{M}$ and $E\sim F\vnmod{M}$ for a general $M$.
\end{description}

\renewcommand{\refname}{Bibliography}
\begin{thebibliography}{99}
\bibitem{ref1} J. v. Neumann: \textit{Allgemeine Eigenwerttheorie
Hermitescher Funktionaloperatoren}. Math. Ann., Vol. 102, pp. 49--131 (1929).
\bibitem{ref2} J. v. Neumann: \textit{Zur Algebra der
Funktionaloperatoren und Theorie der normalen Operatoren}. Math. Ann., Vol. 102,
pp. 370--427 (1929).
\bibitem{ref3} J. v. Neumann: \textit{Zur Theorie der unbeschr\"ankten
Matrizen}. J. Reine Angew. Math., Vol. 161, pp. 208--236 (1929).
\bibitem{ref4} J. v. Neumann: \textit{\"Uber Funktionen von
Funktionaloperatoren}. Ann. of Math., Vol. 32, pp. 191--226 (1931).
\bibitem{ref5} J. v. Neumann: \textit{\"Uber adjungierte
Funktionaloperatoren}. Ann. of Math., Vol. 33, pp. 294--310 (1932).
\bibitem{ref6} J. v. Neumann: \textit{Einige S\"atze \"uber messbare
Abbildungen}. Ann. of Math., Vol. 33, pp. 574--586 (1932).
\bibitem{ref7} J. v. Neumann: \textit{Zur Operatormethode in der
klassichen Mechanik}. Ann. of Math., Vol. 33, pp. 587--642 (1932).
\bibitem{ref8} J. v. Neumann: \textit{Die Einf\"uhrung analytischer
Parameter in topologischen Gruppen}. Ann. of Math., Vol. 34, pp. 170--190
(1933).
\bibitem{ref9} J. v. Neumann: \textit{Almost periodic functions
in a group}. Trans. Amer. Math. Soc., Vol. 36, pp. 445--492 (1934).
\bibitem{ref10} F. J. Murray and J. v. Neumann: \textit{On Rings of
Operators}. Ann. of Math., Vol. 37, pp. 116--229 (1935).
\bibitem{ref11} M. H. Stone: \textit{Linear Transformations in Hilbert
Spaces}. New York, Amer. Math. Soc. Colloquium Publications Vol. XV, 1932.
\bibitem{ref12} F. Maeda: \textit{Theory of Vector Valued Set
Functions}. J. Sci. Hiroshima Univ., A4, pp. 57--91 (1934) and A4,
pp. 141--169 (1939).
\bibitem{ref13} C. Carath\'eodory: \textit{Reelle Funktionen}.
Leipzig and Berlin, Teubner, 1918.
\bibitem{ref14} N. Lusin: \textit{Le\c{c}ons sur les ensembles analytiques et
leurs applications}. Paris, Gauthier-Villars, 1930.
\bibitem{ref15} F. Hausdorff: \textit{Mengenlehre}. Berlin, Walter de
Gruyter and Co., 1927.
\bibitem{ref16} N. Lusin: \textit{Sur les propri\'et\'es des fonctions
mesurables}. C. R. Acad. Sci. Paris, Vol. 154, pp. 1688--1690 (1912).
\bibitem{ref17} W. Sierpinski: \textit{D\'emonstration de quelques
th\'eor\`emes fondamentaux sur les fonctions mesurables}. Fund. Math., Vol. 3,
pp. 314--321 (1922).
\end{thebibliography}

\vnpart{I}{Generalized direct sums}

\section{Description of unitary spaces}\label{sec:1}

We will consider spaces which fulfill the postulates $A,B,C,E$ of
\cite[pp. 64--66]{ref1}. They were discussed in
\cite[Chapter I]{ref10}. Such a space is either a finite-dimensional
Euclidean space, or a Hilbert space. Loc. cit. \cite{ref10} they were
called ``linear spaces with an inner product'', but we now prefer to follow the
terminology of \cite[pp. 16--18]{ref11}, and call them ``unitary spaces''.
(It would perhaps be more consequent, to require $A,B$ only for ``unitary''
spaces, and to add the further specifications ``separable, complete'' for $C,E$.
In the present context, however, we need not pay much attention to these
nuances. Omission of $C$, separability, would leave most of our results
unaffected, but we propose to discuss this problem, as well as various other ones
connected with the generalization of the results of \cite{ref10} to the
non-separable case, at another occasion.)

Unitary spaces will always be denoted, as in \cite{ref10}, by the letter
$\mathfrak H$. If several unitary spaces occur simultaneously, we will
distinguish them by indices.

The 0-dimensional space (with the sole element 0) will again be explicitly
excluded from the considerations of this paper, as it was the case in
\cite{ref10} too. Thus the only possible dimensions for unitary spaces
are: $1,2,\ldots$ (Euclidean spaces), and $+\infty$ (Hilbert space).

The fundamental operations of a unitary space $\mathfrak H$ are the ``linear
operations'' of $\alpha f$, $f+g$ ($\alpha$ a complex number,
$f,g\in\mathfrak H$, and $\alpha f,f+g\in\mathfrak H$) and the ``inner
product'' $(f,g)$ ($f,g\in\mathfrak H$ and $(f,g)$ a complex number). From
these the ``absolute value'' $\lVert f\rVert=\sqrt{(f,f)}$ and the ``distance''
$\lVert f-g\rVert$ are derived (cf. \cite[pp. 64--65]{ref1}). Even when
we consider several unitary spaces $\mathfrak H$ simultaneously, we will use
these symbols $\alpha f$, $f+g$, $(f,g)$, $\lVert f\rVert$ irrespectively of the
$\mathfrak H$ in which the $f,g$ are. This will not lead to misunderstandings,
because it will never happen that an $f$ belongs to several $\mathfrak H$'s with
different definitions of these notions.

\section{\texorpdfstring{$N$-functions, Stieltjes measure, and total
continuity}{N-functions, Stieltjes measure, and total continuity}}\label{sec:2}

Another essential element of our considerations are the functions
$\rho(\lambda)$ which satisfy the following conditions:
\begin{enumerate}[label=(\roman*)]
\item $\rho(\lambda)$ is real valued, defined for all
$\lambda>-\infty$, $<+\infty$;
\item $\rho(\lambda)$ is monotonous non-decreasing.
($\lambda_1\leq\lambda_2$ implies
$\rho(\lambda_1)\leq\rho(\lambda_2)$);
\item $\rho(\lambda)$ is right semi-continuous.
($\lambda\geq\lambda_0$, $\lambda\to\lambda_0$ imply
$\rho(\lambda)\to\rho(\lambda_0)$, or:
$\rho(\lambda_0+0)=\rho(\lambda_0)$);
\item $\rho(\lambda)$ is bounded. (A constant $C$ with
$|\rho(\lambda)|\leq C$ for all $\lambda>-\infty$, $<+\infty$ exists.)
\end{enumerate}
Such a $\rho(\lambda)$ will be called an $N$-function.

This definition differs from the one given in \cite[p. 193]{ref4}
for $M$-functions only by the domain of $\rho(\lambda)$: it was then
$0\leq\lambda\leq a$ ($<+\infty$), now it is
$-\infty<\lambda<+\infty$. However these functions too were introduced in
\cite[p. 199]{ref4}. The restriction of boundedness, (iv), could be
omitted without essentially affecting the use which we wish to make of the
$\rho(\lambda)$, and either procedure has its specific advantages. Our choice,
requiring boundedness, seems however to be technically preferable.

We will need the notion of Lebesgue-Stieltjes integration with respect to
$N$-functions $\rho(\lambda)$. We will use this theory in the form in which it
appears in \cite[pp. 198--221]{ref11}. For the convenience of the reader we
list below the essential notions, definitions and facts which we wish to use, as
well as some minor deviations in terminology.

\noindent\textit{($\alpha$)} Some subsets $K$ of
$-\infty<\lambda<+\infty$ are called $\rho(\lambda)$-\emph{measurable} or
\emph{measurable with respect to $\rho(\lambda)$}, cf.
\cite[pp. 203--205]{ref11}. For them a $\rho(\lambda)$-measure is defined,
which we will denote by $\mu_{\rho(\lambda)}(K)$. The situation which we
discuss is somewhat simpler than the one loc. cit., because our
$\rho(\lambda)$ are real and monotonous, while the $\rho(\lambda)$ loc. cit.
are complex and of bounded variation. In the notation loc. cit. we have
$\rho(\lambda)=\rho_1(\lambda)$, while
$\rho_2(\lambda)=\rho_3(\lambda)=\rho_4(\lambda)=0$. Consequently we have
$\mu_{\rho(\lambda)}(K)=M(K)=M_1(K)$. (We write $K$ for the set $E$ loc. cit.,
as we wish to use the letter $E$ for projections only.) As discussed loc. cit.,
the notion of $\rho(\lambda)$-measurability and the
$\rho(\lambda)$-measure $\mu_{\rho(\lambda)}(K)$ share all essential properties
with Lebesgue-measurability and measure, except its invariance under
translations. They are somewhat more convenient insofar as the
$\rho(\lambda)$-measure of the entire interval $-\infty<\lambda<+\infty$ is
finite: $\rho(+\infty)-\rho(-\infty)$.

\noindent\textit{($\beta$)} Based on the notions ($\alpha$) the
$\rho(\lambda)$-integral or \emph{Lebesgue-Stieltjes integral with respect to
$\rho(\lambda)$} can be defined, cf. \cite[pp. 205--213]{ref11}; an
alternative procedure is described in \cite[pp. 198--202]{ref4}.
A complex-valued function $f(\lambda)$ is $\rho(\lambda)$-measurable, if the set
\[
K_f\binom{a_2,b_2}{a_1,b_1}:\quad
a_1<\Re f(\lambda)\leq a_2,\qquad b_1<\Im f(\lambda)\leq b_2
\]
is $\rho(\lambda)$-measurable for every choice of $a_1,a_2,b_1,b_2$. If
$f(\lambda)$ is $\rho(\lambda)$-measurable and $f(\lambda)\geq0$, then the
integral $\int f(\lambda)\,d\rho(\lambda)$ is defined, its value being
$\geq0$, $\leq+\infty$ (it may be $=+\infty$!); if $f(\lambda)$ is
$\rho(\lambda)$-measurable and complex-valued, then
$\int f(\lambda)\,d\rho(\lambda)$ is only defined if
$\int|f(\lambda)|\,d\rho(\lambda)$ is finite (this applies to real-valued
$f(\lambda)$'s too, if they are not $\geq0$!). In this case $f(\lambda)$ is
$\rho(\lambda)$-summable, and $\int f(\lambda)\,d\rho(\lambda)$ is complex,
but necessarily finite.

If $K$ is a $\rho(\lambda)$-measurable set, then
\[
1_K(\lambda)=
\begin{cases}
1,&\lambda\in K,\\
0,&\lambda\notin K
\end{cases}
\]
is a $\rho(\lambda)$-measurable function. The notions of
$\rho(\lambda)$-measurability and $\rho(\lambda)$-summability of $f(\lambda)$
over $K$ and of $\int_K f(\lambda)\,d\rho(\lambda)$ arise from those of
$\rho(\lambda)$-measurability and $\rho(\lambda)$-summability of $f(\lambda)$
and of $\int f(\lambda)\,d\rho(\lambda)$ by replacing $f(\lambda)$ by
$1_K(\lambda)f(\lambda)$.

\noindent\textit{($\gamma$)} Let $\rho(\lambda),\sigma(\lambda)$ be two
$N$-functions. $\rho(\lambda)$ is $\sigma(\lambda)$-\emph{totally continuous}
if $\mu_{\sigma(\lambda)}(K)=0$ implies
$\mu_{\rho(\lambda)}(K)=0$ for Borel-sets $K$. (We formulate this for Borel
sets, because they are always $\sigma(\lambda)$- and
$\rho(\lambda)$-measurable.) By \cite[p. 204, criterium (5)]{ref11}, this
implies, that every $\sigma(\lambda)$-measurable $K$ with
$\mu_{\sigma(\lambda)}(K)=0$ is $\rho(\lambda)$-measurable with
$\mu_{\rho(\lambda)}(K)=0$; and now the same criterium gives, that
$\sigma(\lambda)$-measurability implies $\rho(\lambda)$-measurability in
general.

\emph{The necessary and sufficient condition for a $\rho(\lambda)$ to be
$\sigma(\lambda)$-totally-continuous is, that it be possible to write it in the
form}
\begin{equation}\label{eq:2.1}
\rho(\lambda)=a+\int_{\mu\leq\lambda}f(\mu)\,d\sigma(\mu),
\end{equation}
\emph{where $a$ is a real constant, and $f(\mu)$ a
$\sigma(\lambda)$-summable function, $f(\mu)\geq0$.
$\rho(\lambda),\sigma(\lambda)$ determine $a$ uniquely and $f(\lambda)$
uniquely up to an arbitrary set of $\sigma(\lambda)$-measure 0. $a$ is equal
to $\lim_{\lambda\to-\infty}\rho(\lambda)=\rho(-\infty)$,
$f(\lambda)$ will be denoted by $d\rho(\lambda)/d\sigma(\lambda)$.}
(In \cite{ref11} the $\sigma(\lambda)$-total continuity of
$\rho(\lambda)$ is denoted by $\rho<\sigma$, and the roles of the definition
and the necessary and sufficient criterium are interchanged. The decisive fact
is stated on p. 215, Lemma 6.5, (1).)

It is worth while to state the following consequence of \eqref{eq:2.1}, which
holds for all $\sigma(\lambda)$-measurable sets $K$:
\begin{align*}
\mu_{\rho(\lambda)}(K)
 &=\int_K1\,d\rho(\lambda)=\int1_K(\lambda)\,d\rho(\lambda)\\
 &=\int1_K(\lambda)\,d\left[\int_{\mu\leq\lambda}
        f(\mu)\,d\sigma(\mu)\right]
  =\int1_K(\lambda)f(\lambda)\,d\sigma(\lambda)\\
 &=\int_Kf(\lambda)\,d\sigma(\lambda)
  =\int_K\frac{d\rho(\lambda)}{d\sigma(\lambda)}\,d\sigma(\lambda).
\end{align*}

\noindent\textit{($\delta$)} Let $E(\lambda)$,
$-\infty<\lambda<+\infty$, be a resolution of unity in the unitary space
$\mathfrak H$, cf. \cite[p. 91]{ref1}, or
\cite[p. 174]{ref11}. One verifies immediately, that
\[
\rho_f(\lambda)=(E(\lambda)f,f)=\lVert E(\lambda)f\rVert^2
\]
is an $N$-function for every $f\in\mathfrak H$. $E(\lambda)$ is
$\sigma(\lambda)$-\emph{totally-continuous} if every $\rho_f(\lambda)$
($f\in\mathfrak H$) is $\sigma(\lambda)$-totally-continuous.

If $\varphi_1,\varphi_2,\ldots$ is a complete, normalized, orthogonal set in
$\mathfrak H$, and $K$ a Borel set with
$\mu_{\rho_{\varphi_n}(\lambda)}(K)=0$ for $n=1,2,\ldots$, then
$\mu_{\rho_f(\lambda)}(K)=0$ for every $f$, by
\cite[p. 227, Theorem 6.3]{ref11}. Therefore all $\rho_f(\lambda)$ are
$\sigma(\lambda)$-totally-continuous, if the
$\rho_{\varphi_n}(\lambda)$, $n=1,2,\ldots$, are such. In other words:

\emph{$E(\lambda)$ is $\sigma(\lambda)$-totally-continuous if and only if all
$\rho_{\varphi_n}(\lambda)$, $n=1,2,\ldots$, are
$\sigma(\lambda)$-totally-continuous.}

Choose a sequence $a_1,a_2,\ldots$ of numbers $>0$, with a finite
$\sum_n a_n$. Considering
$\rho_{\varphi_n}(\lambda)=\lVert E(\lambda)\varphi_n\rVert^2\geq0$ and
$\leq\lVert\varphi_n\rVert^2=1$, we see, that
$\sum_n a_n\rho_{\varphi_n}(\lambda)$ is (absolutely) majorised by
$\sum_n a_n$. Thus
\begin{equation}\label{eq:2.2}
\sigma(\lambda)=\sum_n a_n\rho_{\varphi_n}(\lambda)
\end{equation}
is an $N$-function. By \cite[p. 215, end of Lemma 6.5]{ref11}, all
$\rho_{\varphi_n}(\lambda)$ are $\sigma(\lambda)$-totally continuous.
Therefore $E(\lambda)$ is $\sigma(\lambda)$-totally continuous.

Conversely: If $E(\lambda)$ is $\tau(\lambda)$-totally-continuous for some
$N$-function $\tau(\lambda)$, then all $\rho_{\varphi_n}(\lambda)$ are
$\tau(\lambda)$-totally continuous, and with them
$\sigma(\lambda)=\sum_n a_n\rho_{\varphi_n}(\lambda)$ is it too. (Verify this
by using the normal form \eqref{eq:2.1}!) So we see:

\emph{$E(\lambda)$ is $\tau(\lambda)$-totally continuous if and only if
$\sigma(\lambda)$ is $\tau(\lambda)$-totally continuous. Every $N$-function
$\sigma(\lambda)$ with this property is equivalent to $E(\lambda)$.}

Before we close this section we wish to emphasize, that it is important for our
further purposes, to have defined the $\rho(\lambda)$-measures in such a way,
that the possibility of a one-point-sets having a $\rho(\lambda)$-measure
$\ne0$ is not excluded. The $\rho(\lambda)$-measure of the set with the unique
element $\lambda_0$ is obviously
\[
\lim_{\lambda>\lambda_0,\,\lambda\to\lambda_0}\rho(\lambda)
-\lim_{\lambda<\lambda_0,\,\lambda\to\lambda_0}\rho(\lambda)
=\rho(\lambda_0+0)-\rho(\lambda_0-0)
=\rho(\lambda_0)-\rho(\lambda_0-0),
\]
so it is $\ne0$ if and only if $\rho(\lambda)$ is discontinuous at
$\lambda=\lambda_0$. This is the reason, why we required only
right-semi-continuity, and not continuity, of our $N$-functions.

\section{Definition and discussion of the generalized direct sum}\label{sec:3}

After these preliminaries we are able to define the notion of the
\emph{generalized direct sum}, which is the basis of all our subsequent
discussions. It is related to the common notion of direct sum (of linear spaces)
in a similar way, as the general resolutions of unity (which may include
continuous spectra) to those with pure point spectra. Therefore it is the
appropriate tool for an ``additive decomposition'' of Hilbert space.

In this connection we will have to introduce postulatorially a certain Stieltjes
integral in Hilbert space, which is similar to some integrals introduced by F.
Maeda, \cite{ref12}.

We will first define the generalized notion of direct sum, and discuss only
afterwards, how the common notion of direct sum arises as a special case from it.

\begin{definition}\label{def:1}
Let $\mathfrak H$ be a unitary space, $\mathfrak H^{(\lambda)}$ a unitary space
for each $\lambda>-\infty$, $<+\infty$, and $\sigma(\lambda)$ an $N$-function.
$\mathfrak H$ is the \emph{generalized direct sum of the
$\mathfrak H^{(\lambda)}$ with the weight-function $\sigma(\lambda)$}, if the
following is the case:

Consider all functions $f(\lambda)$ which are defined for all
$\lambda>-\infty$, $<+\infty$, and the values of which are subject to this
restriction: $f(\lambda)\in\mathfrak H^{(\lambda)}$.

For these the notion of $\sigma(\lambda)$-summability is defined in some
way\footnote{This notion will not be defined explicitly at present, all we
require is, that it must be defined in some way that satisfies the postulates
which follow.}, and whenever $f(\lambda)$ is $\sigma(\lambda)$-summable, its
$\sigma(\lambda)$-integral
$\displaystyle\int_{-}f(\lambda)\sqrt{d\sigma(\lambda)}$ is
defined, being necessarily $\in\mathfrak H$. The following postulates are
satisfied:
\begin{enumerate}[label=(\roman*)]
\item In order that a given $f(\lambda)$ be $\sigma(\lambda)$-summable, it is
necessary and sufficient, that it fulfill these two conditions:
\begin{enumerate}[label=(i)$_{\arabic*}$]
\item For every $\sigma(\lambda)$-summable $g(\lambda)$ the (complex-valued,
numerical) function of $\lambda$, $(f(\lambda),g(\lambda))$ is
$\sigma(\lambda)$-measurable.
\item If the (non-negative, numerical) function of $\lambda$
$\lVert f(\lambda)\rVert^2$ is $\sigma(\lambda)$-measurable at
all\footnote{We will see in \hyperref[rem:3.1]{Remark 1} in this section, that it is
$\sigma(\lambda)$-measurable, provided that (i)$_1$ holds.}, then it is
$\sigma(\lambda)$-summable too. In other words: Then
$\int\lVert f(\lambda)\rVert^2\,d\sigma(\lambda)$ is finite.
\end{enumerate}
\item If $f(\lambda),g(\lambda)$ are both $\sigma(\lambda)$-summable, and if
the (complex-valued, numerical) function of $\lambda$
$(f(\lambda),g(\lambda))$ is summable at all\footnote{We will see in
\hyperref[rem:3.1]{Remark 1}
in this section, that it is $\sigma(\lambda)$-summable, as a consequence of
(i). (More precisely: Of the necessary character of the conditions
(i)$_1$, (i)$_2$.)}, then we have for
\[
f=\int_{-}f(\lambda)\sqrt{d\sigma(\lambda)},\qquad
g=\int_{-}g(\lambda)\sqrt{d\sigma(\lambda)}
\]
(observe $f,g\in\mathfrak H$):
\[
(f,g)=\int(f(\lambda),g(\lambda))\,d\sigma(\lambda).
\]
\item Every $f\in\mathfrak H$\footnote{It would be sufficient to require,
that the set $\mathfrak S$ of all $f\in\mathfrak H$ which can be put in this
form determines the linear, closed set $\mathfrak H$. In other words: that the
(finite) linear aggregates formed from elements of $\mathfrak S$ be everywhere
dense. This is so, because we will prove in \hyperref[lem:1]{Lemmas 1} and
\hyperref[lem:2]{2} in \secref{sec:5},
that $\mathfrak S$ is a linear, closed set, irrespectively of whether we require
(iii) or not.} can be put in the form
\[
f=\int_{-}f(\lambda)\sqrt{d\sigma(\lambda)}
\]
with some $\sigma(\lambda)$-summable $f(\lambda)$.
\end{enumerate}
\end{definition}

We will now prove the statements made in footnotes 2 and 3, and in
\secref{sec:5} the statement of footnote 4. For this reason we will avoid the use
of (iii) in the rest of this section as well as in \secref{sec:5}.

\medskip\noindent\phantomsection\label{rem:3.1}\textbf{Remark 1.}
Assume that $f(\lambda)$ satisfies (i)$_1$. Then it is either
$\sigma(\lambda)$-summable, or it is not. If it is, then we can apply (i)$_1$ to
$g(\lambda)=f(\lambda)$ obtaining the $\sigma(\lambda)$-measurability of
$(f(\lambda),f(\lambda))=\lVert f(\lambda)\rVert^2$. If it is not, then
(i)$_2$ must be violated, which means that $\lVert f(\lambda)\rVert^2$ is
$\sigma(\lambda)$-measurable (but not $\sigma(\lambda)$-summable). Thus
$\lVert f(\lambda)\rVert^2$ is $\sigma(\lambda)$-measurable in any event,
proving the statement of footnote 2. Now (i)$_1$, (i)$_2$ imply, that
$\lVert f(\lambda)\rVert^2$ is $\sigma(\lambda)$-summable, whenever
$f(\lambda)$ is $\sigma(\lambda)$-summable.

If $f(\lambda),g(\lambda)$ are both $\sigma(\lambda)$-summable, then
$(f(\lambda),g(\lambda))$ is $\sigma(\lambda)$-measurable by (i)$_1$, and it
is even $\sigma(\lambda)$-summable by what we proved above, considering
\[
|(f(\lambda),g(\lambda))|
 \leq\tfrac12\lVert f(\lambda)\rVert^2
      +\tfrac12\lVert g(\lambda)\rVert^2.
\]
This proves the statement of footnote 3, and the unrestricted validity of the
formula in (ii).

\medskip\noindent\phantomsection\label{rem:3.2}\textbf{Remark 2.}
If $f(\lambda)$ fulfills (i)$_1$, (i)$_2$, $\alpha f(\lambda)$ does too. If
$f(\lambda),g(\lambda)$ fulfill (i)$_1$, $f(\lambda)+g(\lambda)$ does it too,
and owing to
\[
\lVert f(\lambda)+g(\lambda)\rVert^2
 \leq2\lVert f(\lambda)\rVert^2+2\lVert g(\lambda)\rVert^2
\]
the same is the case for (i)$_2$ (remembering
\hyperref[rem:3.1]{Remark 1}). Thus (i) implies:
If $f(\lambda),g(\lambda)$ are $\sigma(\lambda)$-summable,
$\alpha f(\lambda)$, $f(\lambda)+g(\lambda)$ are
$\sigma(\lambda)$-summable too.

Consider now
\[
\int_{-}\alpha f(\lambda)\sqrt{d\sigma(\lambda)}
 -\alpha\int_{-}f(\lambda)\sqrt{d\sigma(\lambda)}
\]
and
\[
\int_{-}(f(\lambda)+g(\lambda))\sqrt{d\sigma(\lambda)}
 -\int_{-}f(\lambda)\sqrt{d\sigma(\lambda)}
 -\int_{-}g(\lambda)\sqrt{d\sigma(\lambda)}.
\]
By (ii), they are orthogonal (in $\mathfrak H$) to every
$\int_{-}h(\lambda)\sqrt{d\sigma(\lambda)}$. (iii), or even its weakened form
from footnote 4 would give at once, that they vanish, but we can prove this
without using (iii) at all: Both expressions are linear aggregates from
$\mathfrak S$ (cf. footnote 4), and orthogonal to all elements of
$\mathfrak S$, thus they are orthogonal to all linear aggregates from
$\mathfrak S$ too---therefore to themselves, and so they must both vanish. So
we have proved:
\[
\int_{-}\alpha f(\lambda)\sqrt{d\sigma(\lambda)}
 =\alpha\int_{-}f(\lambda)\sqrt{d\sigma(\lambda)},
\]
\[
\int_{-}(f(\lambda)+g(\lambda))\sqrt{d\sigma(\lambda)}
 =\int_{-}f(\lambda)\sqrt{d\sigma(\lambda)}
  +\int_{-}g(\lambda)\sqrt{d\sigma(\lambda)}.
\]

\medskip\noindent\phantomsection\label{rem:3.3}\textbf{Remark 3.}
Put in the formula of (ii) $f(\lambda)=g(\lambda)$, and then replace
$f(\lambda)$ by $f(\lambda)-g(\lambda)$, using the formula of
\hyperref[rem:3.2]{Remark 2}. Then
\[
\left\lVert
 \int_{-}f(\lambda)\sqrt{d\sigma(\lambda)}
 -\int_{-}g(\lambda)\sqrt{d\sigma(\lambda)}
\right\rVert^2
=\int\lVert f(\lambda)-g(\lambda)\rVert^2\,d\sigma(\lambda)
\]
results. This proves: We have
\[
\int_{-}f(\lambda)\sqrt{d\sigma(\lambda)}
=\int_{-}g(\lambda)\sqrt{d\sigma(\lambda)}
\]
if and only if $f(\lambda)=g(\lambda)$ everywhere, except perhaps for a
$\lambda$-set of $\sigma(\lambda)$-measure 0. (The $\lambda$-set where
$f(\lambda)\ne g(\lambda)$ is $\sigma(\lambda)$-measurable in any event, as it
is characterized by $\lVert f(\lambda)-g(\lambda)\rVert^2\ne0$, and the left
side is a $\sigma(\lambda)$-measurable function of $\lambda$.)

If we change a $\sigma(\lambda)$-summable $f(\lambda)$ on a $\lambda$-set of
$\sigma(\lambda)$-measure 0, then the validity of (i)$_1$, (i)$_2$, and thus
the $\sigma(\lambda)$-summability of $f(\lambda)$ is unimpaired. So we have
proved:

If $f(\lambda)$ is $\sigma(\lambda)$-summable, then all other
$\sigma(\lambda)$-summable $h(\lambda)$'s with
\[
\int_{-}f(\lambda)\sqrt{d\sigma(\lambda)}
=\int_{-}h(\lambda)\sqrt{d\sigma(\lambda)}
\]
are obtained, by changing $f(\lambda)$ arbitrarily on any $\lambda$-set of
$\sigma(\lambda)$-measure 0.

This proves, that from the point of view $\sigma(\lambda)$-summability and of
the value of $\int_{-}f(\lambda)\sqrt{d\sigma(\lambda)}$ (in $\mathfrak H$)
$f(\lambda)$ may even be undefined on a $\lambda$-set of
$\sigma(\lambda)$-measure 0. Therefore even the
$\mathfrak H^{(\lambda)}$'s themselves may be undefined for a fixed
$\lambda$-set of $\sigma(\lambda)$-measure 0.

\section{Totally discontinuous weight-functions and the discrete direct
sum}\label{sec:4}

It is instructive to discuss the behavior of our
\hyperref[def:1]{Definition 1} of generalized
direct sums for \emph{totally discontinuous} $\sigma(\lambda)$'s, before we
continue the general discussion. These $\sigma(\lambda)$'s arise in the
following way:

As the $\sigma(\lambda)$-measure of $-\infty<\lambda<+\infty$ is finite, the
number of points $\lambda$ with a $\sigma(\lambda)$-measure $\geq1/i$,
$i=1,2,\ldots$, is finite, and therefore the number of points $\lambda$ with a
$\sigma(\lambda)$-measure $\ne0$ is finite or countably infinite. Therefore we
may write them as a (finite or infinite) sequence
$\lambda_1,\lambda_2,\ldots$. Let their respective
$\sigma(\lambda)$-measures be $a_1,a_2,\ldots$ (by assumption all $a_n>0$).
As we saw at the end of \secref{sec:2},
\[
a_n=\sigma(\lambda_n+0)-\sigma(\lambda_n-0)
   =\sigma(\lambda_n)-\sigma(\lambda_n-0).
\]

We say that $\sigma(\lambda)$ is totally discontinuous, if the complementary
set to $(\lambda_1,\lambda_2,\ldots)$ has the $\sigma(\lambda)$-measure 0. This
means, that $-\infty<\lambda<+\infty$ has the same
$\sigma(\lambda)$-measure as $(\lambda_1,\lambda_2,\ldots)$, that is
\begin{equation}\label{eq:4.1}
\sigma(+\infty)-\sigma(-\infty)=a_1+a_2+\cdots.
\end{equation}
Total discontinuity implies, that the $\mu$-set $\mu\leq\lambda$ (for any
given $\lambda$) has the same $\sigma(\lambda)$-measure as the set of all
$\lambda_n$ with $\lambda_n\leq\mu$. That is:
$\sigma(\lambda)-\sigma(-\infty)=\sum_{\lambda_n\leq\lambda}a_n$, and
therefore
\begin{equation}\label{eq:4.2}
\sigma(\lambda)=a+\sum_{\lambda_n\leq\lambda}a_n.
\end{equation}
Conversely: If $a,a_1,a_2,\ldots$ are real constants, $a_n>0$ and
$\sum_n a_n$ is finite, then \eqref{eq:4.2} defines an $N$-function, the
discontinuities of which are at $\lambda_1,\lambda_2,\ldots$ with the respective
``jumps'' $a_1,a_2,\ldots$. (It suffices to verify
\[
\sigma(\lambda+0)=\sigma(\lambda)
 =a+\sum_{\lambda_n\leq\lambda}a_n,\qquad
\sigma(\lambda-0)=a+\sum_{\lambda_n<\lambda}a_n.)
\]
One sees, that \eqref{eq:4.1} holds. So we see:

\emph{All totally discontinuous $N$-functions $\sigma(\lambda)$ can be
obtained by \eqref{eq:4.2}, if $a,a_1,a_2,\ldots$ are any real constants, with
$a_n>0$ and $\sum_n a_n$ finite.}

If the number of $\lambda_n$'s is finite, then $\sigma(\lambda)$ is everywhere
interval-wise constant, except at the points $\lambda_n$; the same is the case
for an infinite but isolated set of $\lambda_n$'s. If the $\lambda_n$'s have
condensation points, and in particular if they are everywhere dense, the
outward appearance of $\sigma(\lambda)$ is very different, but all this is
unimportant from the point of view of total discontinuity.

Assume now that $\sigma(\lambda)$ is totally discontinuous. Every subset of
$(\lambda_1,\lambda_2,\ldots)$ is finite or countably infinite, therefore a
Borel-set, therefore $\sigma(\lambda)$-measurable. Every subset of the
complement of $(\lambda_1,\lambda_2,\ldots)$ is a subset of a set of
$\sigma(\lambda)$-measure 0, and therefore $\sigma(\lambda)$-measurable too.
Thus every set $K$ is $\sigma(\lambda)$-measurable, and therefore every
(numerical) function $f(\lambda)$ is such too. (This, by the way, is
characteristic for total discontinuity, but we need not convince ourselves of
it now.)

Now we verify easily:
\[
\mu_{\sigma(\lambda)}(K)=\sum_{\lambda_n\in K}a_n;
\]
if $f(\lambda)\geq0$ then
\[
\int f(\lambda)\,d\sigma(\lambda)=\sum_n f(\lambda_n)a_n;
\]
if $f(\lambda)$ is complex-valued, then it is $\sigma(\lambda)$-summable, if
$\int|f(\lambda)|\,d\sigma(\lambda)=\sum_n|f(\lambda_n)|a_n$ is finite, and
in this case
$\int f(\lambda)\,d\sigma(\lambda)=\sum_n f(\lambda_n)a_n$.

Thus all statements of \hyperref[def:1]{Definition 1} which refer to
$\sigma(\lambda)$-measurability, become void; in particular (i)$_1$ drops out
completely. Owing to \hyperref[rem:3.3]{Remark 3}, we need only to consider the
$\mathfrak H^{(\lambda_n)}$ and the $f(\lambda_n)$, $n=1,2,\ldots$. So
\hyperref[def:1]{Definition 1} becomes this:

$\int_{-}f(\lambda)\sqrt{d\sigma(\lambda)}$ exists (that is: $f(\lambda)$ is
$\sigma(\lambda)$-summable) if and only if
$\sum_n\lVert f(\lambda_n)\rVert^2a_n$ is finite (by (i)$_2$). These
$\int_{-}f(\lambda)\sqrt{d\sigma(\lambda)}$ exhaust $\mathfrak H$ by (iii),
and we have for them
\[
\left(\int_{-}f(\lambda)\sqrt{d\sigma(\lambda)},
      \int_{-}g(\lambda)\sqrt{d\sigma(\lambda)}\right)
=\sum_n(f(\lambda_n),g(\lambda_n))a_n.
\]

These statements can be transformed and simplified further. Choose an
$f^{(m)}\in\mathfrak H^{(\lambda_m)}$, and define:
\begin{equation}\label{eq:4.3}
f(\lambda_n)=
\begin{cases}
\dfrac{1}{\sqrt{a_m}}f^{(m)},&n=m,\\[4pt]
0,&n\ne m.
\end{cases}
\end{equation}
Then $\sum_n\lVert f(\lambda_n)\rVert^2a_n$ is finite, and thus
$\int_{-}f(\lambda)\sqrt{d\sigma(\lambda)}$ exists, denote it by
$[f^{(m)}]_m$ ($\in\mathfrak H$). Now the inner-product-formula becomes
\[
([f^{(m)}]_m,[g^{(l)}]_l)=
\begin{cases}
(f^{(m)},g^{(m)}),&m=l,\\
0,&m\ne l.
\end{cases}
\]
Denote the set of all $[f^{(m)}]_m$, $f^{(m)}\in\mathfrak H^{(\lambda_m)}$,
by $\mathfrak H_m^*$. Then the above formula states for $m=l$:
$f^{(m)}\longleftrightarrow[f^{(m)}]_m$ is a one-to-one isomorphic mapping of
$\mathfrak H^{(\lambda_m)}$ on the subset $\mathfrak H_m^*$ of $\mathfrak H$.
As $\mathfrak H^{(\lambda_m)}$ is a linear, complete space,
$\mathfrak H_m^*$ is too; that is, it is a linear closed set in $\mathfrak H$.
For $m\ne l$ the formula gives: $\mathfrak H_m^*$,
$\mathfrak H_l^*$ are orthogonal. Thus the
$\mathfrak H_1^*,\mathfrak H_2^*,\ldots$ are mutually orthogonal linear,
closed sets in $\mathfrak H$.

The existence of $\int_{-}f(\lambda)\sqrt{d\sigma(\lambda)}$, that is the
finiteness of
$\sum_n\lVert f(\lambda_n)\rVert^2a_n
=\sum_n\lVert[f(\lambda_n)]_n\rVert^2a_n$ is obviously equivalent to the
convergence of the (orthogonal) series
$\sum_n[f(\lambda_n)]_n\sqrt{a_n}$ in $\mathfrak H$.
(Cf. \cite[p. 9]{ref11}, or \cite[p. 67]{ref1}.) Therefore it is natural to
consider the correspondence
\[
\int_{-}f(\lambda)\sqrt{d\sigma(\lambda)}
\longleftrightarrow\sum_n[f(\lambda_n)]_n\sqrt{a_n}.
\]
The inner-product-formula shows, that this leaves inner products invariant.
Thus it is a one-to-one isomorphic mapping of a subset $\mathfrak H_1$ of
$\mathfrak H$ on another one $\mathfrak H_2$. By (iii)
$\mathfrak H_1=\mathfrak H$; and as
$[f(\lambda_n)]_n\in\mathfrak H_n^*$,
$\mathfrak H_2\subset[\mathfrak H_1^*,\mathfrak H_2^*,\ldots]=$ the linear,
closed set determined by $\mathfrak H_1^*,\mathfrak H_2^*,\ldots$. Now the two
sides of this correspondence are clearly equal, if $f(\lambda)$ is the
expression from \eqref{eq:4.3}. Thus the correspondence is identical in each
$\mathfrak H_m^*$, therefore in
$[\mathfrak H_1^*,\mathfrak H_2^*,\ldots]$, and a fortiori in
$\mathfrak H_2$. It maps $\mathfrak H_2$ on itself: Therefore
$\mathfrak H_1=\mathfrak H_2$, thus
$[\mathfrak H_1^*,\mathfrak H_2^*,\ldots]=\mathfrak H$ and
\begin{equation}\label{eq:4.4}
\int_{-}f(\lambda)\sqrt{d\sigma(\lambda)}
=\sum_n[f(\lambda_n)]_n\sqrt{a_n}.
\end{equation}

Conversely: Let $\mathfrak H$ be decomposed into a (finite or infinite)
sequence of mutually orthogonal linear, closed sets
$\mathfrak H_1^*,\mathfrak H_2^*,\ldots$, which determine the linear, closed
set $\mathfrak H$. Let $\mathfrak H_n^*$ be isomorphic to some unitary space
$\mathfrak H^{(\lambda_n)}$ by an isomorphism
$f^{(n)}\longleftrightarrow[f^{(n)}]_n$
($f^{(n)}\in\mathfrak H^{(\lambda_n)}$,
$[f^{(n)}]_n\in\mathfrak H_n^*\subset\mathfrak H$). Define
$\int_{-}f(\lambda)\sqrt{d\sigma(\lambda)}$ by equation \eqref{eq:4.4} (so
that it is defined if and only if the right side of \eqref{eq:4.4} converges,
that is $\sum_n\lVert f(\lambda_n)\rVert^2a_n
=\sum_n\lVert[f(\lambda_n)]_n\rVert^2a_n$ is finite). Then one verifies
immediately that all parts of \hyperref[def:1]{Definition 1} hold, in the
modified form which we
gave them for the totally discontinuous $\sigma(\lambda)$.

Now $\mathfrak H$ is what one usually calls the direct sum of
$\mathfrak H_1^*,\mathfrak H_2^*,\ldots$, and these spaces are isomorphic to
$\mathfrak H^{(\lambda_1)},\mathfrak H^{(\lambda_2)},\ldots$ respectively.
Thus the totally discontinuous case permits us to replace the unitary spaces
$\mathfrak H^{(\lambda)}$ (of course only the essential ones in the sense of
\hyperref[rem:3.3]{Remark 3}: $\lambda=\lambda_1,\lambda_2,\ldots$) by
isomorphic linear, closed
subsets of $\mathfrak H$, and leads back to the common notion of the direct sum,
as we stated it at the beginning of \secref{sec:3}. (If the number of
$\lambda_n$'s is finite, we can even get rid of the $a_n$, by putting
$a_1=a_2=\cdots=1$. If it is infinite, this cannot be done, as
$\sum_n a_n$ must be finite, in order to safeguard the boundedness of
$\sigma(\lambda)$. Cf. \secref{sec:2}, condition (iv).)

If $\sigma(\lambda)$ is not totally discontinuous, we may expect new
phenomena. Then our use of unitary spaces $\mathfrak H^{(\lambda)}$ outside of
$\mathfrak H$ becomes essential, because there is no natural way, as there was
one in the totally discontinuous case, to replace them by isomorphic images
inside of $\mathfrak H$. We will now discuss the properties of this general case
more thoroughly.

\section{Elementary properties of generalized integrals}\label{sec:5}

From now on $\sigma(\lambda)$ is again an arbitrary $N$-function. As it was
pointed out in \secref{sec:3}, we will not use (iii) for a while.

Let, as in footnote 4, $\mathfrak S$ be the set of all $f\in\mathfrak H$ of the
form $f=\int_{-}f(\lambda)\sqrt{d\sigma(\lambda)}$ for some
$\sigma(\lambda)$-summable $f(\lambda)$. We will deduce some properties of
$\mathfrak S$ and of this representation. The latter ones are those, which are
important for our purposes, we state the former ones only, because they follow
from those easily. The analysis of $\mathfrak S$ would be altogether
unnecessary, if we wanted to use (iii) (this prescribes
$\mathfrak S=\mathfrak H$, and thus determines $\mathfrak S$ completely). But
we wish to establish the statement of footnote 4, according to which
$\mathfrak S$ is linear, closed in consequence of (i), (ii) alone, thus making
a seemingly weaker but really equivalent formulation of (iii) possible (cf.
there). This will be secured by the statements of
\hyperref[lem:1]{Lemmas 1} and \hyperref[lem:2]{2} on
$\mathfrak S$.

We will prove three Lemmas. They are as follows:

\begin{lemma}\label{lem:1}
$\mathfrak S$ is a linear set. If $f,g\in\mathfrak S$,
\[
f=\int_{-}f(\lambda)\sqrt{d\sigma(\lambda)},\qquad
g=\int_{-}g(\lambda)\sqrt{d\sigma(\lambda)},
\]
then
\[
\alpha f=\int_{-}\alpha f(\lambda)\sqrt{d\sigma(\lambda)},\qquad
f+g=\int_{-}(f(\lambda)+g(\lambda))\sqrt{d\sigma(\lambda)}.
\]
\end{lemma}
\begin{proof}
Immediately by \hyperref[rem:3.2]{Remark 3 in \S2}.
\end{proof}

\begin{lemma}\label{lem:2}
$\mathfrak S$ is a closed set. If $g_n\in\mathfrak S$, $n=1,2,\ldots$,
\[
g_n=\int_{-}g_n(\lambda)\sqrt{d\sigma(\lambda)},
\]
and if $\lim_{n\to\infty}g_n$ exists (in the metric, ``strong'', topology of
$\mathfrak H$), then a function $g(\lambda)$ (defined for all
$\lambda>-\infty$, $<+\infty$, $g(\lambda)\in\mathfrak H^{(\lambda)}$),
exists, which has this property:

\noindent\textup{(I)} A subsequence $p_1,p_2,\ldots$ exists, such that
$\lim_{r\to\infty}g_{p_r}(\lambda)=g(\lambda)$ (in the metric, ``strong'',
topology of $\mathfrak H^{(\lambda)}$) for all $\lambda>-\infty$,
$<+\infty$, except perhaps for a $\lambda$-set of $\sigma(\lambda)$-measure 0.

\textup{(I)} determines $g(\lambda)$ uniquely, except for an arbitrary
$\lambda$-set of $\sigma(\lambda)$-measure 0, irrespective of what the sequence
$p_1,p_2,\ldots$ happened to be. This $g(\lambda)$ is
$\sigma(\lambda)$-summable and
\[
\lim_{n\to\infty}g_n
=\int_{-}g(\lambda)\sqrt{d\sigma(\lambda)}.
\]
\end{lemma}
\begin{proof}
The statement on $\mathfrak S$ follows from the rest of the Lemma, because the
final equation shows, that if $\lim_{n\to\infty}g_n$ exists at all (all
$g_n\in\mathfrak S$!), then it belongs to $\mathfrak S$. Thus $\mathfrak S$ is
closed. We need therefore to consider only the other parts of the Lemma.

Assume now $g_n\in\mathfrak S$,
$g_n=\int_{-}g_n(\lambda)\sqrt{d\sigma(\lambda)}$,
$\lim_{n\to\infty}g_n$ exists. Then we have
$\lim_{m,n\to\infty}\lVert g_m-g_n\rVert=0$, or by
\hyperref[rem:3.3]{Remark 3 in \S2}
\[
\lim_{m,n\to\infty}\int
 \lVert g_m(\lambda)-g_n(\lambda)\rVert^2\,d\sigma(\lambda)=0.
\]
Thus we have a ``Convergence in the mean'', and we can proceed in the customary
way. The above equation means, that there exists for every $\epsilon>0$ an
$s=s(\epsilon)$, such that $m,n\geq s$ imply
\[
\int\lVert g_m(\lambda)-g_n(\lambda)\rVert^2\,d\sigma(\lambda)<\epsilon.
\]
Denote the set of all $\lambda$ with
$\lVert g_m(\lambda)-g_n(\lambda)\rVert\geq\delta>0$ by
$K_1(m,n,\delta)$. Then we see:
$\mu_{\sigma(\lambda)}(K_1(m,n,\delta))<\epsilon/\delta^2$ if
$m,n\geq s(\epsilon)$. Thus
$\mu_{\sigma(\lambda)}(K_1(m,n,1/2^r))<1/2^r$ if
$m,n\geq s(1/8^r)$. Choose now a sequence $p_1<p_2<\cdots$ with
$p_r\geq s(1/8^r)$; then we have
$\mu_{\sigma(\lambda)}(K_1(p_r,p_{r+1},1/2^r))<1/2^r$. Put
\[
K_2(r)=\sum_{t=r}^{\infty}K_1(p_t,p_{t+1},1/2^t),\qquad
K_3=\prod_{r=1}^{\infty}K_2(r),
\]
then
\[
\mu_{\sigma(\lambda)}(K_2(r))
 <\frac1{2^r}+\frac1{2^{r+1}}+\cdots=\frac1{2^{r-1}},\qquad
\mu_{\sigma(\lambda)}(K_3)
 \leq\mu_{\sigma(\lambda)}(K_2(r))<\frac1{2^{r-1}}
\]
for every $r=1,2,\ldots$ and thus $\mu_{\sigma(\lambda)}(K_3)=0$.

If $\lambda\notin K_3$, then there exists an $r$ with
$\lambda\notin K_2(r)$, and thus
$\lambda\notin K_1(p_t,p_{t+1},1/2^t)$ for all $t\geq r$; this means
$\lVert g_{p_t}(\lambda)-g_{p_{t+1}}(\lambda)\rVert<1/2^t$ for all
$t\geq r$. Then $\lim_{t\to\infty}g_{p_t}(\lambda)$ must exist (in the
metric, ``strong'', topology of $\mathfrak H^{(\lambda)}$). So we see:
$\lim_{r\to\infty}g_{p_r}(\lambda)$ exists for all $\lambda>-\infty$,
$<+\infty$, except perhaps for a $\lambda$-set of
$\sigma(\lambda)$-measure 0 (the above $K_3$). Define
$g(\lambda)=\lim_{r\to\infty}g_{p_r}(\lambda)$ for all $\lambda$'s for which
this limit exists, for all other $\lambda$ $g(\lambda)$ may be anything (their
set has the $\sigma(\lambda)$-measure 0).

So we found a $g(\lambda)$ which fulfills the condition (I).

Let us now consider any such $g(\lambda)$, and see whether the final equation of
our Lemma holds for it.

In what follows, we will always neglect $\lambda$-sets of
$\sigma(\lambda)$-measure 0. We have always
$\lim_{r\to\infty}g_{p_r}(\lambda)=g(\lambda)$, therefore
\[
\lim_{r\to\infty}(g_{p_r}(\lambda)-g_n(\lambda))
=g(\lambda)-g_n(\lambda)
\]
and thus
\[
\lim_{r\to\infty}\lVert g_{p_r}(\lambda)-g_n(\lambda)\rVert^2
=\lVert g(\lambda)-g_n(\lambda)\rVert^2.
\]
As
$\lVert g_{p_r}(\lambda)-g_n(\lambda)\rVert^2\geq0$,
$\lVert g(\lambda)-g_n(\lambda)\rVert^2\geq0$, this implies by a well-known
property of Lebesgue integrals\footnote{Cf. for instance \cite[p. 443,
Theorem 15]{ref13}. Using those notations, put $S(P)=0$, then this obtains: If
all $f_k(P)\leq0$, then
\[
\limsup_{k\to\infty}\left[\int f_k(P)\,dw\right]
 \leq\int\left[\limsup_{k\to\infty}f_k(P)\right]dw.
\]
This is even so if the case
$\limsup_{k\to\infty}[\int f_k(P)\,dw]=-\infty$ occurs, which was excluded
loc. cit., in fact then it is trivial. Replacing $f_k(P)$ by $-f_k(P)$ and
assuming the existence of $\lim_{k\to\infty}f_k(P)$, this becomes: If
$f_k(P)\geq0$, then
\[
\liminf_{k\to\infty}\left[\int f_k(P)\,dw\right]
 \geq\int\left[\lim_{k\to\infty}f_k(P)\right]dw.
\]
This is, disregarding the notation, the desired inequality. The entire argument
holds unchanged for Lebesgue-Stieltjes-integrals, or their representation by
Lebesgue integrals, cf. \cite[p. 198]{ref4}, may be used to carry it over.},
that
\[
\liminf_{r\to\infty}\int
 \lVert g_{p_r}(\lambda)-g_n(\lambda)\rVert^2\,d\sigma(\lambda)
\geq\int\lVert g(\lambda)-g_n(\lambda)\rVert^2\,d\sigma(\lambda).
\]

The left side is equal to
$\liminf_{r\to\infty}\lVert g_{p_r}-g_n\rVert^2$ by
\hyperref[rem:3.3]{Remark 3 in \S2}. If $n\geq s(\epsilon)$ then as for
$r\to\infty$ finally
$p_r\geq s(\epsilon)$ too, this expression is $\leq\epsilon$. So we see: For
$n\geq s(\epsilon)$ we have
\[
\int\lVert g(\lambda)-g_n(\lambda)\rVert^2\,d\sigma(\lambda)
 \leq\epsilon.
\]

If $f(\lambda)$ is $\sigma(\lambda)$-summable, then
$(g_m(\lambda),f(\lambda))$ is $\sigma(\lambda)$-measurable for all
$m=1,2,\ldots$ by (i)$_1$, thus
$(g_{p_r}(\lambda)-g_n(\lambda),f(\lambda))$ is
$\sigma(\lambda)$-measurable too, and so is its limit for $r\to\infty$,
$(g(\lambda)-g_n(\lambda),f(\lambda))$. So
$g(\lambda)-g_n(\lambda)$ fulfills (i)$_1$. As we saw it fulfills (i)$_2$ too
(if $n$ is great enough, for instance if $n\geq s(1)$). Thus (i) applies:
$g(\lambda)-g_n(\lambda)$ is $\sigma(\lambda)$-summable, and with it
$g(\lambda)=[g(\lambda)-g_n(\lambda)]+g_n(\lambda)$ too.

Put $g=\int_{-}g(\lambda)\sqrt{d\sigma(\lambda)}$. The inequality we had before
becomes now: For $n\geq s(\epsilon)$ we have
$\lVert g-g_n\rVert^2\leq\epsilon$. Thus
\[
\lim_{n\to\infty}g_n=g
=\int_{-}g(\lambda)\sqrt{d\sigma(\lambda)},
\]
establishing the final equation of our Lemma.

This equation, together with \hyperref[rem:3.3]{Remark 3 in \S2}, proves, that
$g(\lambda)$ is uniquely determined, except for an arbitrary $\lambda$-set of
$\sigma(\lambda)$-measure 0.
\end{proof}

\begin{lemma}\label{lem:3}
Let $f(\lambda)$ be an arbitrary function, defined for all
$\lambda>-\infty$, $<+\infty$, $f(\lambda)\in\mathfrak H^{(\lambda)}$. Let
$\mathfrak S'$ be the set of all those $g\in\mathfrak S$,
$g=\int_{-}g(\lambda)\sqrt{d\sigma(\lambda)}$, for which the (complex-valued,
numerical) function of $\lambda$ $(f(\lambda),g(\lambda))$ is
$\sigma(\lambda)$-measurable. (By \hyperref[rem:3.3]{Remark 3 in \S2} it does not
matter, which $g(\lambda)$ of a given $g$ we choose.) Then $\mathfrak S'$ is a
closed linear set.
\end{lemma}
\begin{proof}
Assume $g,h\in\mathfrak S'$. Then
\[
g=\int_{-}g(\lambda)\sqrt{d\sigma(\lambda)},\qquad
h=\int_{-}h(\lambda)\sqrt{d\sigma(\lambda)},
\]
and $(f(\lambda),g(\lambda))$, $(f(\lambda),h(\lambda))$ are
$\sigma(\lambda)$-measurable. From this
$\alpha g=\int_{-}\alpha g(\lambda)\sqrt{d\sigma(\lambda)}$,
$g+h=\int_{-}(g(\lambda)+h(\lambda))\sqrt{d\sigma(\lambda)}$ follow, and
\[
(f(\lambda),\alpha g(\lambda))
=\overline\alpha(f(\lambda),g(\lambda)),\qquad
(f(\lambda),g(\lambda)+h(\lambda))
=(f(\lambda),g(\lambda))+(f(\lambda),h(\lambda))
\]
are $\sigma(\lambda)$-measurable. Thus $\alpha g,g+h\in\mathfrak S'$, and so
$\mathfrak S'$ is a linear set.

Assume now that $g$ is a condensation point of $\mathfrak S'$. Choose
$g_1,g_2,\ldots\in\mathfrak S'$ with $\lim_{n\to\infty}g_n=g$. Put
\[
g_n=\int_{-}g_n(\lambda)\sqrt{d\sigma(\lambda)},\qquad
g=\int_{-}g(\lambda)\sqrt{d\sigma(\lambda)}.
\]
Then we have by \hyperref[lem:2]{Lemma 2} a subsequence
$p_1,p_2,\ldots$ of $1,2,\ldots$, so
that $\lim_{r\to\infty}g_{p_r}(\lambda)=g(\lambda)$, except perhaps for a
$\lambda$-set of $\sigma(\lambda)$-measure 0. Similarly
$\lim_{r\to\infty}(f(\lambda),g_{p_r}(\lambda))
=(f(\lambda),g(\lambda))$. As all functions of $\lambda$,
$(f(\lambda),g_n(\lambda))$ are $\sigma(\lambda)$-measurable, we have this now
for $(f(\lambda),g(\lambda))$ too. Therefore $g\in\mathfrak S'$. Thus
$\mathfrak S'$ is closed too.
\end{proof}

\section{Measurable families and existence and uniqueness of generalized
direct sums}\label{sec:6}

As the linear, closed set character of $\mathfrak S$ is established, we can be
satisfied concerning the observations of footnote 4, and we will use from now
on (iii): That is $\mathfrak S=\mathfrak H$.

Let $\varphi_1,\varphi_2,\ldots$ be a complete normalized orthogonal set in
$\mathfrak H$ (an everywhere dense sequence would do just as well). Put
\[
\varphi_n=\int_{-}\varphi_n(\lambda)\sqrt{d\sigma(\lambda)}.
\]
Consider now the sequence $\varphi_1(\lambda),\varphi_2(\lambda),\ldots$ in the
unitary space $\mathfrak H^{(\lambda)}$. Apply to it the ``orthogonalization
process'' of E. Schmidt, cf. \cite[p. 13]{ref11}, or \cite[p. 69]{ref1}. Then
a normalized orthogonal set
$\psi_1(\lambda),\psi_2(\lambda),\ldots$ (in
$\mathfrak H^{(\lambda)}$) arises. It has $k=k(\lambda)=0,1,2,\ldots,\infty$
elements, where $k$ is the ``rank'' of
$\varphi_1(\lambda),\varphi_2(\lambda),\ldots$ (the number of linearly
independent elements in that sequence). If $k\ne\infty$, we define
\[
\psi_{k+1}(\lambda)=\psi_{k+2}(\lambda)=\cdots=0,
\]
thus $\psi_1(\lambda),\psi_2(\lambda),\ldots$ is now an infinite sequence in
any event---it is an orthogonal set, $\psi_n(\lambda)$ being normalized if
$n\leq k$, and vanishing if $n>k$.

If one visualizes nature of the ``orthogonalization process'', one sees at
once, that the $(\varphi_m(\lambda),\psi_n(\lambda))$ and the
$\lVert\psi_n(\lambda)\rVert^2$ with $m,n\leq l$ are Baire-functions of the
$(\varphi_p(\lambda),\varphi_q(\lambda))$ with $p,q\leq l$. These are
$\sigma(\lambda)$-measurable, therefore the
$(\varphi_m(\lambda),\psi_n(\lambda))$, and the
$\lVert\psi_n(\lambda)\rVert^2$ are too, and so is
\[
k(\lambda)=\sum_{1}^{\infty}\lVert\psi_n(\lambda)\rVert^2.
\]

We now propose to prove:

\emph{If $f(\lambda)$ is a function, defined for all $\lambda>-\infty$,
$<+\infty$, $f(\lambda)\in\mathfrak H^{(\lambda)}$, and if all
(complex-valued-numerical) functions of $\lambda$
$(f(\lambda),\psi_n(\lambda))$, $n=1,2,\ldots$, are
$\sigma(\lambda)$-measurable, then $f(\lambda)$ fulfills (i)$_1$.}

As we know, $\varphi_p(\lambda)$ is a linear aggregate of those
$\psi_1(\lambda),\ldots,\psi_k(\lambda)$ ($k=k(\lambda)$) which have indices
$\leq p$, thus of $\psi_1(\lambda),\ldots,\psi_l(\lambda)$
($l=\min(k(\lambda),p)$). These form a normalized, orthogonal set, so
\[
\varphi_p(\lambda)=\sum_{r=1}^{l}
 (\varphi_p(\lambda),\psi_r(\lambda))\psi_r(\lambda).
\]
Replacing $\sum_{1}^{l}$ by $\sum_{1}^{p}$ adds only vanishing terms, so we
may do it. Now we have
\[
(f(\lambda),\varphi_p(\lambda))
=\sum_{r=1}^{p}(f(\lambda),\psi_r(\lambda))
 \overline{(\varphi_p(\lambda),\psi_r(\lambda))}
\]
and so it is $\sigma(\lambda)$-measurable too. Form the $\mathfrak S'$ (in the
sense of \hyperref[lem:3]{Lemma 3}) for $f(\lambda)$, then we have proved
$\varphi_p\in\mathfrak S'$, $p=1,2,\ldots$. By
\hyperref[lem:3]{Lemma 3} this implies
$\mathfrak S'=\mathfrak H$. Thus $(f(\lambda),g(\lambda))$ is
$\sigma(\lambda)$-measurable for every $\sigma(\lambda)$-summable
$g(\lambda)$, in other words: $f(\lambda)$ fulfills (i)$_1$. This completes
the proof.

As
\[
(\psi_m(\lambda),\psi_n(\lambda))=
\begin{cases}
\lVert\psi_m(\lambda)\rVert^2,&m=n,\\
0,&m\ne n,
\end{cases}
\]
it is $\sigma(\lambda)$-measurable in any event. Thus $\psi_m(\lambda)$
fulfills (i)$_1$. As $\lVert\psi_m(\lambda)\rVert^2\leq1$,
\[
\int\lVert\psi_m(\lambda)\rVert^2\,d\sigma(\lambda)
 \leq\int1\,d\sigma(\lambda)=\sigma(+\infty)-\sigma(-\infty)
\]
is finite, and so (i)$_2$ holds too. Therefore $\psi_m(\lambda)$ is
$\sigma(\lambda)$-summable by (i). Put
\[
\psi_m=\int_{-}\psi_m(\lambda)\sqrt{d\sigma(\lambda)},
\qquad m=1,2,\ldots.
\]

Let $K_4$ be the set of those $\lambda>-\infty$, $<+\infty$, for which the
non-vanishing $\psi_n(\lambda)$'s (that is: the $k(\lambda)$ first ones) do not
form a complete normalized orthogonal set in $\mathfrak H^{(\lambda)}$. For
every $\lambda\in K_4$ choose an $f(\lambda)$ which is orthogonal to all
$\psi_n(\lambda)$, and $\lVert f(\lambda)\rVert^2=1$. For all
$\lambda\notin K_4$ put $f(\lambda)=0$. Then we have always
$(f(\lambda),\psi_n(\lambda))=0$, so this is
$\sigma(\lambda)$-measurable. Thus $f(\lambda)$ fulfills (i)$_1$, and by
\hyperref[rem:3.1]{Remark 1 in \S2}, $\lVert f(\lambda)\rVert^2$ must be
$\sigma(\lambda)$-measurable. So $K_4$ is a $\sigma(\lambda)$-measurable set.
Now
\[
\int\lVert f(\lambda)\rVert^2\,d\sigma(\lambda)
 \leq\int1\,d\sigma(\lambda)=\sigma(+\infty)-\sigma(-\infty)
\]
is finite, so $f(\lambda)$ fulfills (i)$_2$ too. Therefore $f(\lambda)$ is
$\sigma(\lambda)$-summable by (i), put
$f=\int_{-}f(\lambda)\sqrt{d\sigma(\lambda)}$. The expression we used above
for $(f(\lambda),\varphi_p(\lambda))$ shows, that in the present case this
vanishes, because all $(f(\lambda),\psi_n(\lambda))$ vanish. Now apply (ii):
\[
(f,\varphi_p)=\int(f(\lambda),\varphi_p(\lambda))\,d\sigma(\lambda)=0,
\]
\[
\lVert f\rVert^2
=\int\lVert f(\lambda)\rVert^2\,d\sigma(\lambda)
=\int1_{K_4}(\lambda)\,d\sigma(\lambda)
=\mu_{\sigma(\lambda)}(K_4).
\]
$f$ is orthogonal to all $\varphi_1,\varphi_2,\ldots$, therefore $f=0$,
$\lVert f\rVert^2=0$, and $\mu_{\sigma(\lambda)}(K_4)=0$. So $K_4$ has the
$\sigma(\lambda)$-measure 0.

Thus the non-vanishing $\psi_n(\lambda)$'s (that is: the $k(\lambda)$ first
ones) form a complete normalized orthogonal set in
$\mathfrak H^{(\lambda)}$ for every $\lambda>-\infty$, $<+\infty$, except
perhaps for a $\lambda$-set of $\sigma(\lambda)$-measure 0. For these
$\lambda$'s however, we can redefine the $\psi_n(\lambda)$ and $k(\lambda)$ so
that this should be true for them too, while the $\psi_n(\lambda)$ with
$n>k(\lambda)$ again vanish. This will not effect the above derived criterium
for (i)$_1$ in terms of the $\psi_n(\lambda)$'s, nor the
$\sigma(\lambda)$-summability of $\psi_n(\lambda)$, nor the value of $\psi_n$.
And now $k(\lambda)$ is the dimensionality of the unitary space
$\mathfrak H^{(\lambda)}$.

The results obtained so far can be best formulated with the help of this
definition:

\begin{definition}\label{def:2}
Let $k(\lambda)$ be the dimensionality of $\mathfrak H^{(\lambda)}$. Denote
the set of all $\lambda>-\infty$, $<+\infty$ with $k(\lambda)=n$ by
$\Delta_n$, and the set of those with $k(\lambda)\geq n$ by $H_n$.
($n=1,2,\ldots,\infty$,
$H_n=\Delta_n+\Delta_{n+1}+\cdots+\Delta_\infty$.)

Assume that for every $\lambda>-\infty$, $<+\infty$ a complete normalized
orthogonal set $\psi_1(\lambda),\psi_2(\lambda),\ldots$ in
$\mathfrak H^{(\lambda)}$ is given (the length of this sequence is of course
$k(\lambda)$, that is $n$ for $\lambda\in\Delta_n$). Then every function
$f(\lambda)$ ($-\infty<\lambda<+\infty$,
$f(\lambda)\in\mathfrak H^{(\lambda)}$) can be written in the form
\[
f(\lambda)=\sum_m x_m(\lambda)\psi_m(\lambda),\qquad
x_m(\lambda)=(f(\lambda),\psi_m(\lambda));
\]
and the complex numbers $x_m(\lambda)$ are only subject to these two
restrictions:
\begin{enumerate}
\item[($\alpha$)] $x_m(\lambda)$'s domain is characterized by $m\leq k(\lambda)$, that is
by $\lambda\in H_m$,
\item[($\beta$)] $\sum_m|x_m(\lambda)|^2$ is finite for every $\lambda>-\infty$,
$<+\infty$.
\end{enumerate}

If the $\sigma(\lambda)$-summability of $f(\lambda)$ is equivalent to the two
further restrictions
\begin{enumerate}
\item[($\gamma$)] $x_m(\lambda)$ is a measurable function of $\lambda$ for every
$m=1,2,\ldots$,
\item[($\delta$)] $\displaystyle\sum_{m=1}^{\infty}\int_{H_m}
 |x_m(\lambda)|^2\,d\sigma(\lambda)$ is finite,
\end{enumerate}
then the systems $\psi_m(\lambda)$ form a \emph{measurable family}.
\end{definition}

\medskip\noindent\phantomsection\label{rem:6.1}\textbf{Remark.}
Considering the non-negativity of the integrands, we have
\begin{align*}
\sum_{m=1}^{\infty}\int_{H_m}|x_m(\lambda)|^2\,d\sigma(\lambda)
 &=\sum_{m=1}^{\infty}\sum_{n=m}^{\infty}
   \int_{\Delta_n}|x_m(\lambda)|^2\,d\sigma(\lambda)\\
 &=\sum_{n=1}^{\infty}\sum_{m=1}^{n}
   \int_{\Delta_n}|x_m(\lambda)|^2\,d\sigma(\lambda)\\
 &=\sum_{n=1}^{\infty}\int_{\Delta_n}
   \sum_{m=1}^{n}|x_m(\lambda)|^2\,d\sigma(\lambda)\\
 &=\int\sum_m|x_m(\lambda)|^2\,d\sigma(\lambda).
\end{align*}
Thus the finiteness of the expression at the extreme left implies the finiteness
of $\sum_m|x_m(\lambda)|^2$ for all $\lambda>-\infty$, $<+\infty$, except
perhaps for a set of $\sigma(\lambda)$-measure 0. But by
\hyperref[rem:3.3]{Remark 3} in
\secref{sec:3}, such a set does not matter. So we see, that ($\delta$) implies
($\beta$), and that the $\sigma(\lambda)$-summable $f(\lambda)$'s are
characterized by ($\alpha$), ($\gamma$), ($\delta$) alone.

We now prove:

\begin{theorem}\phantomsection\label{thm:I}
\begin{enumerate}[label=(\roman*)$^*$]
\item $k(\lambda)$ is a measurable function, all $\Delta_n,H_n$ are measurable
sets.
\item Measurable families $\psi_m(\lambda)$ exist.
\item If $\psi_m^0(\lambda)$ is a measurable family, then all others
$\psi_m(\lambda)$ are obtained by choosing for each $\lambda>-\infty$,
$<+\infty$ a unitary matrix $\{u_{mn}(\lambda)\}$ of degree $k(\lambda)$,
such that $u_{mn}(\lambda)$ is a measurable function of $\lambda$ for every
pair $m,n=1,2,\ldots$, and putting
\[
\psi_m(\lambda)=\sum_n u_{mn}(\lambda)\psi_n^0(\lambda).
\]
Then of course
\[
u_{mn}(\lambda)=(\psi_m(\lambda),\psi_n^0(\lambda)).
\]
\end{enumerate}
\end{theorem}
\begin{proof}
\textit{Ad (i)$^*$.} We defined at the beginning of this section a function
$k(\lambda)$, which was then shown to be measurable. We saw later, that it
coincides with the dimensionality of $\mathfrak H^{(\lambda)}$. $\Delta_n$ and
$H_n$ are the sets $k(\lambda)=n$ and $\geq n$, respectively.

\textit{Ad (ii)$^*$.} Consider the $\psi_n(\lambda)$ defined before
\hyperref[def:2]{Definition 2}. We proved that
\hyperref[def:1]{Definition 1}, (i)$_1$ is equivalent to
\hyperref[def:2]{Definition 2}, ($\gamma$), and
\hyperref[def:1]{Definition 1}, (i)$_2$ is clearly identical with
\hyperref[def:2]{Definition 2}, ($\delta$). Thus
\hyperref[def:1]{Definition 1}, (i) proves, that these $\psi_n(\lambda)$ form a
measurable family.

\textit{Ad (iii)$^*$. Sufficiency:} If the $\psi_n(\lambda)$ have this form,
then the two representations
\[
f(\lambda)=\sum_m x_m^0(\lambda)\psi_m^0(\lambda),\qquad
f(\lambda)=\sum_m x_m(\lambda)\psi_m(\lambda)
\]
are connected by the formulae
\[
x_m(\lambda)=\sum_n\overline{u_{mn}(\lambda)}x_n^0(\lambda),\qquad
x_n^0(\lambda)=\sum_m u_{mn}(\lambda)x_m(\lambda).
\]
As $u_{mn}(\lambda)$ is measurable, ($\gamma$) for $\psi_m(\lambda)$ is
equivalent to ($\gamma$) for $\psi_m^0(\lambda)$. As
\[
\lVert f(\lambda)\rVert^2
=\sum_m|x_m^0(\lambda)|^2=\sum_m|x_m(\lambda)|^2,
\]
the same is the case for ($\delta$). Thus $\psi_m(\lambda)$ is a measurable
family along with $\psi_m^0(\lambda)$.

\textit{Necessity:} Let $\psi_m(\lambda),\psi_m^0(\lambda)$ be two measurable
families. As $\psi_n(\lambda)$ ($n$ fixed, put $\psi_n(\lambda)=0$ for
$\lambda\notin H_n$) fulfills ($\gamma$), ($\delta$) for the family
$\psi_m(\lambda)$, it must also do this for the family
$\psi_m^0(\lambda)$. Thus
$u_{mn}(\lambda)=(\psi_m(\lambda),\psi_n^0(\lambda))$ is a measurable function
of $\lambda$ by ($\gamma$). The other statements follow from the fact, that
$\psi_m(\lambda)$ and $\psi_m^0(\lambda)$ are both complete normalized
orthogonal sets in $\mathfrak H^{(\lambda)}$.
\end{proof}

The theorem which follows may be considered as a converse of
\hyperref[thm:I]{Theorem I}:

\begin{theorem}\label{thm:II}
Let for every $\lambda>-\infty$, $<+\infty$ a unitary space
$\mathfrak H^{(\lambda)}$ be given, and a complete normalized orthogonal set
$\psi_1(\lambda),\psi_2(\lambda),\ldots$ in
$\mathfrak H^{(\lambda)}$. Let $k(\lambda)$ be the dimensionality of
$\mathfrak H^{(\lambda)}$, $\Delta_n$ and $H_n$ the $\lambda$-sets with
$k(\lambda)=n$ and $\geq n$, respectively ($n=1,2,\ldots,\infty$,
$H_n=\Delta_n+\Delta_{n+1}+\cdots+\Delta_\infty$).

Let $\sigma(\lambda)$ be an $N$-function, $k(\lambda)$ being
$\sigma(\lambda)$-measurable. Therefore the sets $\Delta_n,H_n$ are such, too.

Consider all complex-valued functions $x_m(\lambda)$ of the two variables
$m=1,2,\ldots$, $\lambda>-\infty$, $<+\infty$, which fulfill the conditions
($\alpha$), ($\gamma$), ($\delta$) of \hyperref[def:2]{Definition 2}. Defining for
$x\equiv x_m(\lambda)$, $y\equiv y_m(\lambda)$
\[
\alpha x\equiv\alpha x_m(\lambda),\qquad
x+y\equiv x_m(\lambda)+y_m(\lambda),
\]
\[
(x,y)=\sum_{m=1}^{\infty}\int_{H_m}
 x_m(\lambda)\overline{y_m(\lambda)}\,d\sigma(\lambda),
\]
we obtain a unitary space $\mathfrak H$\footnote{Cf.
\cite[pp. 23--32]{ref11} or \cite[pp. 108--111]{ref1}. The most convenient
way to look at this $\mathfrak H$ is, to consider the $x_m(\lambda)$ as
functions in the $m,\lambda$-space defined by $m\leq k(\lambda)$, that is by
$\lambda\in H_m$, in which
$\sum_{m=1}^{\infty}\int_{H_m}\cdots\,d\sigma(\lambda)$ plays the role of the
Lebesgue-(Stieltjes) integral.}.

If a function $f(\lambda)$ of $\lambda>-\infty$, $<+\infty$,
$f(\lambda)\in\mathfrak H^{(\lambda)}$ is given, define $x_m(\lambda)$ by
\[
f(\lambda)=\sum_m x_m(\lambda)\psi_m(\lambda),\qquad
x_m(\lambda)=(f(\lambda),\psi_m(\lambda)).
\]
Define $f(\lambda)$ to be $\sigma(\lambda)$-summable by
$x\equiv x_m(\lambda)\in\mathfrak H$, and then define
\[
\int_{-}f(\lambda)\sqrt{d\sigma(\lambda)}=x\equiv x_m(\lambda).
\]

With these definitions $\mathfrak H$ is the generalized direct sum of the
$\mathfrak H^{(\lambda)}$ with the weight-function $\sigma(\lambda)$, and the
systems $\psi_m(\lambda)$ form a measurable family.
\end{theorem}
\begin{proof}
We must first verify the conditions (i), (ii), (iii) of
\hyperref[def:1]{Definition 1}.

\textit{Ad (i):} Put
$f(\lambda)=\sum_n x_n(\lambda)\psi_n(\lambda)$ and similarly
$g(\lambda)=\sum_n y_n(\lambda)\psi_n(\lambda)$, the
$\sigma(\lambda)$-summability of $g(\lambda)$ means
$y\equiv y_m(\lambda)\in\mathfrak H$. (i)$_1$ requires the measurability of
\[
(f(\lambda),g(\lambda))
=\sum_n x_n(\lambda)\overline{y_n(\lambda)}
\]
for all $y\equiv y_m(\lambda)\in\mathfrak H$. As we can put
\[
y_n(\lambda)=
\begin{cases}
1,&n=m,\\
0,&n\ne m,
\end{cases}
\]
the measurability of each $x_m(\lambda)$ is necessary; as all $y_m(\lambda)$
are measurable, this is sufficient. (i)$_2$ states, that
\[
\int\lVert f(\lambda)\rVert^2\,d\sigma(\lambda)
=\int\sum_m|x_m(\lambda)|^2\,d\sigma(\lambda)
=\sum_m\int_{H_m}|x_m(\lambda)|^2\,d\sigma(\lambda)
\]
(cf. the \hyperref[rem:6.1]{Remark} after
\hyperref[def:2]{Definition 2}) is finite. So (i)$_1$, (i)$_2$ are
equivalent to \hyperref[def:2]{Definition 2}, ($\gamma$), ($\delta$), and as
\hyperref[def:2]{Definition 2},
($\alpha$) is automatically fulfilled, so $x\equiv x_m(\lambda)\in\mathfrak H$.
Thus they are equivalent to the $\sigma(\lambda)$-summability of $f(\lambda)$.

\textit{Ad (ii):} This follows from the computation
\begin{align*}
(x,y)
 &=\sum_{m=1}^{\infty}\int_{H_m}
 x_m(\lambda)\overline{y_m(\lambda)}\,d\sigma(\lambda)\\
 &=\int\sum_m x_m(\lambda)\overline{y_m(\lambda)}\,d\sigma(\lambda)\\
 &=\int(f(\lambda),g(\lambda))\,d\sigma(\lambda).
\end{align*}
The replacement of $\sum_m\int_{H_m}$ by $\int\sum_m$ is legitimate for
$x=y$, because everything is non-negative (cf. above); replacing $x$ by
$\frac12(x+y)$, $\frac12(x-y)$, and subtracting gives its real part in the case
of arbitrary $x,y$; replacing $x,y$ by $ix,y$ gives its imaginary part.

\textit{Ad (iii):} Obvious.

Now the application of \hyperref[def:2]{Definition 2} shows immediately that the
$\psi_m(\lambda)$ form a measurable family.
\end{proof}

\section{Operators, resolutions of unity, and characterization of the
decomposition}\label{sec:7}

We continue the analysis of an $\mathfrak H$, which is the generalized direct
sum of a family of $\mathfrak H^{(\lambda)}$'s, with the weight function
$\sigma(\lambda)$.

Let $\alpha(\lambda)$ be a numerical (complex-valued),
$\sigma(\lambda)$-measurable function of $\lambda$, which is bounded, if a
suitable set of $\sigma(\lambda)$-measure 0 is disregarded. Application of the
criteria of \hyperref[def:1]{Definition 1}, (i) (or just as well of
\hyperref[def:2]{Definition 2}, ($\alpha$),
($\gamma$), ($\delta$)), shows, that $\alpha(\lambda)f(\lambda)$ is
$\sigma(\lambda)$-summable if $f(\lambda)$ is such. Remembering
\hyperref[rem:3.3]{Remark 3} in
\secref{sec:3}, we see, that
$f=\int_{-}f(\lambda)\sqrt{d\sigma(\lambda)}$ determines
$f^*=\int_{-}\alpha(\lambda)f(\lambda)\sqrt{d\sigma(\lambda)}$ and thus an
operator $A$ is defined by $Af=f^*$. We denote this $A$ by $\alpha(\Lambda)$:
\[
\alpha(\Lambda)\int_{-}f(\lambda)\sqrt{d\sigma(\lambda)}
=\int_{-}\alpha(\lambda)f(\lambda)\sqrt{d\sigma(\lambda)}.
\]

\emph{This $\alpha(\Lambda)$ has the following properties:}
\begin{enumerate}[label=(\alph*)]
\item $\alpha(\Lambda)$ is a bounded operator.
\item $\alpha(\Lambda)$ is unaffected if $\alpha(\lambda)$ is changed on a set
of $\sigma(\lambda)$-measure 0.
\item $\alpha(\lambda)\equiv a=\text{Constant}$, or
$\equiv\overline{\beta(\lambda)}$, or
$\equiv\beta(\lambda)+\gamma(\lambda)$ or
$\equiv\beta(\lambda)\gamma(\lambda)$ implies
$\alpha(\Lambda)=a1$, or $=\beta(\Lambda)^*$ (the adjoint!), or
$=\beta(\Lambda)+\gamma(\Lambda)$ or
$=\beta(\Lambda)\gamma(\Lambda)$, respectively.
\item If $\lim_{n\to\infty}\alpha_n(\lambda)=\alpha(\lambda)$ for every
$\lambda$, and if the $\alpha_n(\lambda)$ are uniformly bounded for
$n=1,2,\ldots$, then
$\lim_{n\to\infty}\alpha_n(\Lambda)=\alpha(\Lambda)$. (In the strong sense,
that is $\lim_{n\to\infty}\lVert\alpha_n(\Lambda)f-
\alpha(\Lambda)f\rVert=0$ for every $f\in\mathfrak H$, cf.
\cite[pp. 381--382]{ref2}.)
\end{enumerate}

(a) results from \hyperref[def:1]{Definition 1}, (ii): If
$|\alpha(\lambda)|\leq C$ disregarding
a set of $\sigma(\lambda)$-measure 0, then
\[
\lVert\alpha(\Lambda)f\rVert^2
=\int|\alpha(\lambda)|^2\lVert f(\lambda)\rVert^2\,d\sigma(\lambda)
\leq C^2\int\lVert f(\lambda)\rVert^2\,d\sigma(\lambda)
=C^2\lVert f\rVert^2.
\]
(b) is obvious by \hyperref[rem:3.3]{Remark 3} in \secref{sec:3}. The second
statement of (c) again follows from \hyperref[def:1]{Definition 1}, (ii):
\begin{align*}
(\beta(\Lambda)^*f,g)
 &=(f,\beta(\Lambda)g)
  =\int\overline{\beta(\lambda)}(f(\lambda),g(\lambda))\,d\sigma(\lambda)\\
 &=\int\alpha(\lambda)(f(\lambda),g(\lambda))\,d\sigma(\lambda)
  =(\alpha(\Lambda)f,g);
\end{align*}
the other statements of (c) are obvious. (d) again follows from
\hyperref[def:1]{Definition 1},
(ii), as
\[
\lVert\alpha_n(\Lambda)f-\alpha(\Lambda)f\rVert^2
=\int|\alpha_n(\lambda)-\alpha(\lambda)|^2
 \lVert f(\lambda)\rVert^2\,d\sigma(\lambda).
\]

\noindent\textit{(e)} Let $\alpha(\lambda),\gamma(\lambda)$ be
$\sigma(\lambda)$-measurable, $\gamma(\lambda)$ real, let $\beta(x)$ have the
measurability-properties described loc. cit. below, and let it be bounded in
every finite $x$-interval. If $A$ is a bounded Hermitean operator,
$\beta(A)$ denotes the operator-function in the sense of \cite[p. 241]{ref11}
or \cite[pp. 202--203]{ref4}. Then
$\alpha(\lambda)\equiv\beta(\gamma(\lambda))$ implies
$\alpha(\Lambda)=\beta(\gamma(\Lambda))$.

This is obvious for $\beta(x)=x$, and so by (c) for every polynomial
$\beta(x)$. Now (d) extends it to every continuous- and then to every
Baire-function $\beta(x)$, and then (b) to every $\beta(x)$.
(Cf. \cite[p. 234]{ref11} or \cite[p. 196]{ref4}.)

\medskip\noindent\phantomsection\label{rem:7.1}\textbf{Remark.}
If we could put $\gamma(\lambda)=\lambda$, (e) would show, that
$\alpha(\Lambda)$ is the operator-function of $\Lambda$ in the sense loc. cit.
But $\gamma(\lambda)=\lambda$ will in general violate our
boundedness-requirements, and thereby exclude $\gamma(\lambda)=\lambda$. It
would be easy, however, to extend our definitions to this case, and then the
theory of functions of unbounded operators (cf. loc. cit. \cite{ref11}) would
be particularly useful. If $\sigma(\lambda)$ is constant outside of a finite
$\lambda$-interval, then the complement of this interval has the
$\sigma(\lambda)$-measure 0. Thus $\gamma(\lambda)=\lambda$ can be used, and
we need even not change our present set-up if we desire to interpret
$\alpha(\Lambda)$ as an operator-function. Of this we will make use in the
proof of \hyperref[thm:IV]{Theorem IV} (cf. the part of the proof dealing with the ``Relationship
to equivalence'').

If $K$ is a $\sigma(\lambda)$-measurable set, define
\[
1_K(\lambda)=
\begin{cases}
1,&\lambda\in K,\\
0,&\lambda\notin K.
\end{cases}
\]
Then $1_K(\Lambda)$ is a projection by (c); and if $K'$ is a subset of $K''$,
then we have $1_{K'}(\Lambda)1_{K''}(\Lambda)=1_{K'}(\Lambda)$, that is
$1_{K'}(\Lambda)\leq1_{K''}(\Lambda)$ by (c). If $K',K'',\ldots$ is a
sequence of sets with $K'\supset K''\supset\cdots$ and if
$K'\cdot K''\cdots=K$, then clearly
$\lim_{n\to\infty}1_{K^{(n)}}(\lambda)=1_K(\lambda)$, and
$|1_{K^{(n)}}(\lambda)|\leq1$ so by (d)
$\lim_{n\to\infty}1_{K^{(n)}}(\Lambda)=1_K(\Lambda)$.

If $K$ is the set of all $\lambda>-\infty$, $\leq\mu$, then denote
$1_K(\Lambda)$ by $E(\mu)$. The above statements prove that $E(\mu)$,
$-\infty<\mu<+\infty$ is a resolution of unity. We have:
\[
E(\mu)\int_{-}f(\lambda)\sqrt{d\sigma(\lambda)}
=\int_{-}f_\mu(\lambda)\sqrt{d\sigma(\lambda)},
\]
where
\[
f_\mu(\lambda)=
\begin{cases}
f(\lambda),&\lambda\leq\mu,\\
0,&\lambda>\mu.
\end{cases}
\]
We call these $E(\mu)$ the \emph{resolution of unity which belongs to our
generalized direct sum}.

We prove next:

\begin{theorem}\label{thm:III}
A given unitary space $\mathfrak H$ can be represented as a generalized direct
sum to which a given weight function $\sigma(\lambda)$ and a given resolution
of unity $E(\lambda)$ belong, if and only if $\sigma(\lambda)$ is equivalent to
$E(\lambda)$ (cf. \secref{sec:2}). In this case this representation is uniquely
determined (up to replacements of the $\mathfrak H^{(\lambda)}$ by isomorphic
images, and up to an arbitrary $\lambda$-set of $\sigma(\lambda)$-measure 0),
by $\sigma(\lambda)$ and $E(\lambda)$.
\end{theorem}
\begin{proof}
\textit{Necessity of the condition:} Assume that $\mathfrak H$ is the
generalized direct sum of the $\mathfrak H^{(\lambda)}$ with the weight function
$\sigma(\lambda)$, and let $E(\lambda)$ be the resolution of unity which
belongs to it. Form the $\rho_f(\mu)$ of \secref{sec:2}. Put
$f=\int_{-}f(\lambda)\sqrt{d\sigma(\lambda)}$, then
\[
\rho_f(\mu)=\lVert E(\mu)f\rVert^2
=\int\lVert f_\mu(\lambda)\rVert^2\,d\sigma(\lambda)
=\int_{\lambda\leq\mu}\lVert f(\lambda)\rVert^2\,d\sigma(\lambda).
\]
Thus every $\rho_f(\lambda)$ is $\sigma(\lambda)$-totally continuous by the
formula \eqref{eq:2.1} in \secref{sec:2}. Therefore if $\sigma(\lambda)$ is
$\tau(\lambda)$-totally-continuous, then the same is true for every
$\rho_f(\lambda)$, and so for $E(\lambda)$ too. Conversely: Choose
$f(\lambda)$ with $\lVert f(\lambda)\rVert=1$ (for instance
$f(\lambda)=\psi_1(\lambda)$, that is
$x_m(\lambda)=1$ for $m=1$ and $0$ for $m\ne1$); then our above formula gives
$\rho_f(\mu)=\sigma(\mu)$. So if $E(\lambda)$ is
$\tau(\lambda)$-totally-continuous, then the same is true for
$\sigma(\lambda)$. Thus we have proved that $\sigma(\lambda)$ is equivalent to
$E(\lambda)$.

\textit{Sufficiency of the condition:} Assume that an $N$-function
$\sigma(\lambda)$, and an equivalent resolution of unity $E(\lambda)$ are
given. Under these assumptions a (finite or infinite) sequence of mutually
orthogonal linear closed sets $\mathfrak M_1,\mathfrak M_2,\ldots$, which
together determine the linear closed set $\mathfrak H$, can be found, so that
the following requirement is fulfilled: Each $\mathfrak M_m$ reduces all
$E(\lambda)$, and in each $\mathfrak M_m$ the resolution of unity $E(\lambda)$
is simple. Simplicity is defined in \cite[p. 275, Definition 7.2]{ref11}, for
self-adjoint transformations, and we will call a resolution of unity simple,
if it belongs to a simple self-adjoint transformation. A definition of
simplicity in terms of $E(\lambda)$ themselves results by combining Theorems
7.2 and 7.9 (\cite[pp. 243, 275]{ref11}, respectively). The above statement is
contained in Theorem 7.4 (\cite[p. 247]{ref11}), if we denote the
$\mathfrak M(\varphi_1),\mathfrak M(\varphi_2),\ldots,
\mathfrak M(g_1),\mathfrak M(g_2),\ldots$ of that theorem (in some order) by
$\mathfrak M_1,\mathfrak M_2,\ldots$, and remember, that every
$\mathfrak M(f)$ reduces the $E(\lambda)$ by a remark in
\cite[p. 244]{ref11}, and that in every $\mathfrak M(f)$ the $E(\lambda)$ are
simple by Theorem 7.9, quoted above. In each $\mathfrak M_m$ we apply Theorem
6.2 (\cite[p. 226]{ref11}): We have
$\mathfrak M_m=\mathfrak M(f_m)$ for a suitable $f_m\in\mathfrak M_m$, and
$\mathfrak M_m$ is isomorphic to the unitary space $\mathfrak H_m^*$ which is
defined as follows: $\rho_m(\lambda)=\rho_{f_m}(\lambda)
=\lVert E(\lambda)f_m\rVert^2$ is an $N$-function, consider all
$\rho_m(\lambda)$-measurable (complex-valued) functions $\xi(\lambda)$, for
which $\int|\xi(\lambda)|^2\,d\rho_m(\lambda)$ is finite, and define for
$\xi\equiv\xi(\lambda)$, $\eta\equiv\eta(\lambda)$
\[
\alpha\xi\equiv\alpha\xi(\lambda),\qquad
\xi+\eta\equiv\xi(\lambda)+\eta(\lambda),
\]
\[
(\xi,\eta)=\int\xi(\lambda)\overline{\eta(\lambda)}\,d\rho_m(\lambda).
\]
The set $\mathfrak H_m^*$ of all these $\xi\equiv\xi(\lambda)$ is then a
unitary space. (Cf. loc. cit. footnote 5.)

As $\rho_m(\lambda)$ is $\sigma(\lambda)$-totally continuous, we can form
$d\rho_m(\lambda)/d\sigma(\lambda)$ from \secref{sec:2}, ($\gamma$). Let
$N_m$ be the set of all $\lambda$ with
$d\rho_m(\lambda)/d\sigma(\lambda)=0$, owing to the final formula loc. cit.
$\mu_{\rho_m(\lambda)}(N_m)=0$. Thus we can restrict the $\xi(\lambda)$ of
$\mathfrak H_m^*$ to the $\lambda\notin N_m$. Consider now all
$\sigma(\lambda)$-measurable (complex-valued) functions $x(\lambda)$, defined
for the $\lambda\notin N_m$, for which
$\int_{\lambda\notin N_m}|x(\lambda)|^2\,d\sigma(\lambda)$ is finite, and
define for $x\equiv x(\lambda)$, $y\equiv y(\lambda)$
\[
\alpha x\equiv\alpha x(\lambda),\qquad
x+y\equiv x(\lambda)+y(\lambda),
\]
\[
(x,y)=\int_{\lambda\notin N_m}x(\lambda)\overline{y(\lambda)}
 \,d\sigma(\lambda).
\]
The set $\mathfrak H_m^{**}$ of all these $x\equiv x(\lambda)$ is again a
unitary space. (Cf. as above.) But now the requirement $\lambda\notin N_m$
may be essential, because $\mu_{\sigma(\lambda)}(N_m)$ may be $>0$. The
correspondence
\[
\xi(\lambda)\longleftrightarrow x(\lambda)
=\sqrt{\frac{d\rho_m(\lambda)}{d\sigma(\lambda)}}\,\xi(\lambda)
\]
is clearly an isomorphism of $\mathfrak H_m^*$ and $\mathfrak H_m^{**}$
(cf. the final formula of \secref{sec:2}, ($\gamma$)).

In the $\mathfrak M_m\longleftrightarrow\mathfrak H_m^*$ isomorphism this is
true (cf. Theorem 6.2 loc. cit. above, and a remark in
\cite[p. 243]{ref11}): If $f\in\mathfrak M_m$ corresponds to
$\xi(\lambda)\in\mathfrak H_m^*$ then $E(\mu)f$ corresponds to
$1_\mu(\lambda)\xi(\lambda)$, where
\[
1_\mu(\lambda)=
\begin{cases}
1,&\lambda\leq\mu,\\
0,&\lambda>\mu.
\end{cases}
\]
Obviously the same holds in $\mathfrak H_m^{**}$: $x(\lambda)$ becomes
$1_\mu(\lambda)x(\lambda)$. So we have in the
$\mathfrak M_m\longleftrightarrow\mathfrak H_m^{**}$-isomorphism:
\[
E(\mu)x(\lambda)=
\begin{cases}
x(\lambda),&\lambda\leq\mu,\\
0,&\lambda>\mu.
\end{cases}
\]

Let us now return to $\mathfrak H$ itself. Every $f\in\mathfrak H$ can be
written uniquely as a sum $f=f_1+f_2+\cdots$, $f_m\in\mathfrak M_m$ (clearly
$f_m=P_{\mathfrak M_m}f$). Conversely such a sum converges if and only if
$\sum_m\lVert f_m\rVert^2$ is finite (cf. \cite[p. 67]{ref1}, put in Theorem
5, $\varphi_m=(1/\lVert f_m\rVert)f_m$ for every $f_m\ne0$). At the same time
$f=f_1+f_2+\cdots$, $g=g_1+g_2+\cdots$, $f_m,g_m\in\mathfrak M_m$, imply
\[
\alpha f=\alpha f_1+\alpha f_2+\cdots,\qquad
f+g=(f_1+g_1)+(f_2+g_2)+\cdots,
\]
\[
(f,g)=\sum_m(f_m,g_m).
\]

Now use for each $f_m$ the
$\mathfrak H_m^*\longleftrightarrow\mathfrak H_m^{**}$-isomorphism, replacing
it by $x_m(\lambda)$. $x_m(\lambda)$ can be looked at as a complex-valued
function of the two variables $m=1,2,\ldots$, $\lambda>-\infty$,
$<+\infty$, subject to these two restrictions: (a) The inclusive domain
$1,2,\ldots$ of $m$ may have to be replaced by a finite set $1,\ldots,m_0$.
(b) $\lambda\notin N_m$. We can omit (a), by putting $N_m$ equal to the entire
interval $\lambda>-\infty$, $<+\infty$ for the forbidden $m$'s (i.e. for
$m=m_0+1,m_0+2,\ldots$). Denoting the complement of $N_m$ by $H_m$, we
obtain the condition $\lambda\in H_m$. Thus $\mathfrak H$ is isomorphic to the
set of all $x_m(\lambda)$ satisfying the conditions ($\alpha$), ($\gamma$),
($\delta$), in \hyperref[def:2]{Definition 2} and in the sense of
\hyperref[thm:II]{Theorem II} (in ($\alpha$) of
course $\lambda\in H_m$ alone must be used), with the definitions of
$\alpha x,x+y,(x,y)$, given in \hyperref[thm:II]{Theorem II}. At the same time
\[
E(\mu)x_m(\lambda)=
\begin{cases}
x_m(\lambda),&\lambda\leq\mu,\\
0,&\lambda>\mu.
\end{cases}
\]

Assume now, that a $\sigma(\lambda)$-measurable function
$k(\lambda)=1,2,\ldots,\infty$ exists, so that $k(\lambda)\geq m$ is
equivalent to $\lambda\in H_m$. Form then for each $\lambda>-\infty$,
$<+\infty$ a $k(\lambda)$-dimensional unitary space
$\mathfrak H^{(\lambda)}$, and choose in it a complete normalized orthogonal
set $\psi_1(\lambda),\psi_2(\lambda),\ldots$. Then our above remarks
establish, with the help of \hyperref[thm:II]{Theorem II}, that $\mathfrak H$
is the generalized
direct sum of the $\mathfrak H^{(\lambda)}$ with the weight-function
$\sigma(\lambda)$, and with the help of the remark which precedes the present
theorem, that the resolution of unity which belongs to this direct sum consists
of our $E(\mu)$. Thus our only remaining task is to construct this $k(\lambda)$.
But this requires certain changes in the $H_m$.

For each $\lambda>-\infty$, $<+\infty$ denote the number of $m$'s with
$\lambda\in H_m$ by $k_0(\lambda)=0,1,2,\ldots,\infty$, let them be
$m_1(\lambda)<m_2(\lambda)<\cdots$. Denote the domain of our
$x_m(\lambda)$'s by $\Delta$. It consists of all pairs $(m,\lambda)$ with
$\lambda\in H_m$, that is $m=m_n(\lambda)$,
$n=1,2,\ldots,k_0(\lambda)$. Replace each $(m,\lambda)$ by $(n,\lambda)$,
where $m=m_n(\lambda)$. Thus the set $\Delta_0$ of all $(n,\lambda)$ with
$n=1,2,\ldots,k_0(\lambda)$ results, that is with $\lambda\in H_{0n}$, where
$H_{0n}$ is defined by $k_0(\lambda)\geq n$. We claim: This transformation
carries every $\sigma(\lambda)$-measurable $x_m(\lambda)$ into a
$\sigma(\lambda)$-measurable $x_{0n}(\lambda)$, and the
$\sum_{m=1}^{\infty}\int_{H_m}\cdots\,d\sigma(\lambda)$ go over into
$\sum_{n=1}^{\infty}\int_{H_{0n}}\cdots\,d\sigma(\lambda)$. This becomes
obvious if we observe, that our transformation does not affect $\lambda$ at
all, and that the set of those $\lambda$'s, for which a given $m$ goes over
into a given $n$ is $\sigma(\lambda)$-measurable. (It is the set
\[
H_m\cdot\sum_{m_1<\cdots<m_{n-1}<m}
H_{m_1}\cdots H_{m_{n-1}}.)
\]
which is $\sigma(\lambda)$-measurable because all the $H_l$ are such.) Thus
all our previous statements remain true, if we replace our $H_m$ by these
$H_{0n}$, and then $k_0(\lambda)$ has all properties we required for
$k(\lambda)$, except for the possibility of $k_0(\lambda)=0$.

$k_0(\lambda)=0$ means that $\lambda\in H_m$ holds for no
$m=1,2,\ldots$, that is
$\lambda\in\operatorname{Complement}(H_1+H_2+\cdots)=N$. Choose any
$f\in\mathfrak H$, it corresponds to an $x_m(\lambda)$, which may be defined
to be 0 for $\lambda\notin H_m$, so that we can replace the integrals
$\int_{\lambda\in H_m}\cdots\,d\sigma(\lambda)$ by
$\int\cdots\,d\sigma(\lambda)$. In particular $x_m(\lambda)=0$ for
$\lambda\in N$. Now we have for
$\rho_f(\lambda)=\lVert E(\lambda)f\rVert^2$:
\begin{align*}
\mu_{\rho_f(\lambda)}(N)
 &=\int_Nd(\lVert E(\mu)f\rVert^2)
  =\int1_N(\mu)\,d(\lVert E(\mu)f\rVert^2)\\
 &=\int1_N(\mu)\,d\left(
   \sum_{m=1}^{\infty}\int_{\lambda\leq\mu}
   |x_m(\lambda)|^2\,d\sigma(\lambda)\right)\\
 &=\sum_{m=1}^{\infty}\int1_N(\mu)\,d\left(
   \int_{\lambda\leq\mu}|x_m(\lambda)|^2\,d\sigma(\lambda)\right)\\
 &=\sum_{m=1}^{\infty}\int1_N(\lambda)|x_m(\lambda)|^2
   \,d\sigma(\lambda)=0,
\end{align*}
as the integrand $1_N(\lambda)|x_m(\lambda)|^2$ is everywhere 0.

Now form the $N$-function
$\tau(\lambda)=\sum_n a_n\rho_{\varphi_n}(\lambda)$ from
\secref{sec:2}, ($\delta$) ($\varphi_1,\varphi_2,\ldots$ a complete,
normalized, orthogonal set, $a_1,a_2,\ldots>0$, $\sum_n a_n$ finite). We have
\[
\mu_{\tau(\lambda)}(N)
=\sum_n a_n\mu_{\rho_{\varphi_n}(\lambda)}(N)=0,
\]
but $E(\lambda)$ is $\tau(\lambda)$-totally-continuous and equivalent to
$\sigma(\lambda)$ (cf. loc. cit.), so $\sigma(\lambda)$ is
$\tau(\lambda)$-totally-continuous, and therefore
$\mu_{\sigma(\lambda)}(N)=0$. Thus $k_0(\lambda)=0$ occurs only on a set of
$\sigma(\lambda)$-measure 0, and therefore it can be neglected. (We can add
$N$ to $H_1$ without changing anything essential, and nevertheless excluding
$k_0(\lambda)=0$.)

\textit{Uniqueness of the representation:} Assume that $\mathfrak H$ is at
the same time generalized direct sum of the $\mathfrak H'^{(\lambda)}$ and of
the $\mathfrak H''^{(\lambda)}$, in both cases with the weight-function
$\sigma(\lambda)$, and that the same resolution of unity $E(\lambda)$ belongs
to both of them. We will distinguish the notions relative to
$\mathfrak H'^{(\lambda)}$ and to $\mathfrak H''^{(\lambda)}$, by affixing one
and two primes: $'$, $''$, respectively.

The $f\in\mathfrak H$ are then in a one-to-one correspondence with the
$'\sigma(\lambda)$-summable functions $f'(\lambda)$,
\[
f=\int_{-}'f'(\lambda)\sqrt{d\sigma(\lambda)},\qquad
f'(\lambda)\in\mathfrak H'^{(\lambda)},
\]
and similarly with the $''\sigma(\lambda)$-summable functions $f''(\lambda)$,
\[
f=\int_{-}''f''(\lambda)\sqrt{d\sigma(\lambda)},\qquad
f''(\lambda)\in\mathfrak H''^{(\lambda)}.
\]
Thus a one-to-one correspondence between the $'\sigma(\lambda)$-summable
functions $f'(\lambda)$ and the $''\sigma(\lambda)$-summable functions
$f''(\lambda)$ results:
\[
\int_{-}'f'(\lambda)\sqrt{d\sigma(\lambda)}
=\int_{-}''f''(\lambda)\sqrt{d\sigma(\lambda)},\qquad
f'(\lambda)\in\mathfrak H'^{(\lambda)},\quad
f''(\lambda)\in\mathfrak H''^{(\lambda)}.
\]
Of course in all these correspondences sets of $\sigma(\lambda)$-measure 0 are
disregarded.

If $f'(\lambda)$ and $g'(\lambda)$ correspond in this way to $f''(\lambda)$
and $g''(\lambda)$ respectively, put
\[
f=\int_{-}'f'(\lambda)\sqrt{d\sigma(\lambda)}
 =\int_{-}''f''(\lambda)\sqrt{d\sigma(\lambda)},
\]
\[
g=\int_{-}'g'(\lambda)\sqrt{d\sigma(\lambda)}
 =\int_{-}''g''(\lambda)\sqrt{d\sigma(\lambda)}.
\]
Considering that the resolution of unity $E(\lambda)$ belongs to both
representations of $\mathfrak H$ as generalized direct sum, this implies
\[
(E(\mu)f,g)=\int_{\lambda\leq\mu}(f'(\lambda),g'(\lambda))
 \,d\sigma(\lambda)
=\int_{\lambda\leq\mu}(f''(\lambda),g''(\lambda))
 \,d\sigma(\lambda)
\]
and so
\[
\int_{\lambda\leq\mu}
[(f'(\lambda),g'(\lambda))-(f''(\lambda),g''(\lambda))]
\,d\sigma(\lambda)=0
\]
for every $\mu>-\infty$, $<+\infty$. Thus
$(f'(\lambda),g'(\lambda))=(f''(\lambda),g''(\lambda))$ except perhaps for a
$\lambda$-set of $\sigma(\lambda)$-measure 0\footnote{Put
\[
\nu(\mu)=\int_{\lambda\leq\mu}
[(f'(\lambda),g'(\lambda))-(f''(\lambda),g''(\lambda))]
\,d\sigma(\lambda)
\]
 and form for any $\sigma(\lambda)$-measurable function $\epsilon(\mu)$ the
$\int\epsilon(\mu)\,d\nu(\mu)$. This is 0 and at the same time
\[
\int\epsilon(\lambda)
[(f'(\lambda),g'(\lambda))-(f''(\lambda),g''(\lambda))]
\,d\sigma(\lambda),
\]
thus this expression must vanish. Choose
\[
\epsilon(\lambda)=
\begin{cases}
|Z|/Z,&Z\ne0,\\
0,&Z=0,
\end{cases}
\]
where $Z$ is the $[\cdots]$ expression, then
$\int|(f'(\lambda),g'(\lambda))-(f''(\lambda),g''(\lambda))|
\,d\sigma(\lambda)=0$ results, and thus the statement in the text.\label{fn:7}}, which
may depend on the $f'(\lambda),g'(\lambda),f''(\lambda),g''(\lambda)$.

Let now $\varphi_m'(\lambda)$ be a measurable family in the
$\mathfrak H'^{(\lambda)}$, defined for $m\leq k'(\lambda)$; and
$\psi_n''(\lambda)$ a measurable family in the
$\mathfrak H''^{(\lambda)}$, defined for $n\leq k''(\lambda)$. Define
\[
f_m'(\lambda)=
\begin{cases}
\varphi_m'(\lambda),&m\leq k'(\lambda),\\
0,&m>k'(\lambda),
\end{cases}
\]
this function $f_m'(\lambda)$ is clearly $'\sigma(\lambda)$-summable ($m$
fixed!), let the corresponding $''$-function be $f_m''(\lambda)$. Define
similarly
\[
g_n''(\lambda)=
\begin{cases}
\psi_n''(\lambda),&n\leq k''(\lambda),\\
0,&n>k''(\lambda),
\end{cases}
\]
this function $g_n''(\lambda)$ is clearly $''\sigma(\lambda)$-summable ($n$
fixed!), let the corresponding $'$-function be $g_n'(\lambda)$. Then we have:
\begin{align*}
\tag{1}\label{eq:7.1}
(f_m''(\lambda),f_p''(\lambda))
 &=(f_m'(\lambda),f_p'(\lambda))
 =\begin{cases}
1,&m=p\leq k'(\lambda),\\
0,&\text{otherwise};
\end{cases}\\
\tag{2}\label{eq:7.2}
(g_n'(\lambda),g_q'(\lambda))
 &=(g_n''(\lambda),g_q''(\lambda))
 =\begin{cases}
 1,&n=q\leq k''(\lambda),\\
0,&\text{otherwise};
\end{cases}\\
\tag{3}\label{eq:7.3}
(f_m''(\lambda),g_n''(\lambda))
 &=(f_m'(\lambda),g_n'(\lambda))=u_{mn}(\lambda),
\end{align*}
except when $\lambda\in N$, where $N$ is the sum of the exceptional sets of
all these equations. Thus $N$ is the sum of a countably infinite number of
sets, each of $\sigma(\lambda)$-measure 0, therefore
\[
\mu_{\sigma(\lambda)}(N)=0.
\]

Equation \eqref{eq:7.1} proves that the $f_m''(\lambda)$, $m\leq k'(\lambda)$, form a
normalized orthogonal set in $\mathfrak H''^{(\lambda)}$, while the
$g_n''(\lambda)$, $n\leq k''(\lambda)$, form a complete normalized orthogonal
set in $\mathfrak H''^{(\lambda)}$. Therefore
$u_{mn}(\lambda)=(f_m''(\lambda),g_n''(\lambda))$ in \eqref{eq:7.3} leads to
\[
\tag{4}\label{eq:7.4}
\sum_n u_{mn}(\lambda)\overline{u_{pn}(\lambda)}
=\begin{cases}
1,&m=p\leq k'(\lambda),\\
0,&\text{otherwise}.
\end{cases}
\]
Similarly equation \eqref{eq:7.2} proves that the $g_n'(\lambda)$,
$n\leq k''(\lambda)$, form a normalized orthogonal set in
$\mathfrak H'^{(\lambda)}$, while the $f_m'(\lambda)$,
$m\leq k'(\lambda)$, form a complete normalized orthogonal set in
$\mathfrak H'^{(\lambda)}$. Therefore
$u_{mn}(\lambda)=(f_m'(\lambda),g_n'(\lambda))$ in \eqref{eq:7.3} leads to
\[
\tag{5}\label{eq:7.5}
\sum_m u_{mn}(\lambda)\overline{u_{mq}(\lambda)}
=\begin{cases}
1,&n=p\leq k''(\lambda),\\
0,&\text{otherwise}.
\end{cases}
\]
\eqref{eq:7.4}, \eqref{eq:7.5} state together, that $k'(\lambda)=k''(\lambda)$ (both are
$=\sum_{m,n}|u_{mn}(\lambda)|^2$) and that the matrix
$\{u_{mn}(\lambda)\}_{m,n=1,2,\ldots}$
($m,n\leq k'(\lambda)=k''(\lambda)$) is unitary. This holds for all
$\lambda\notin N$. As $N$ has the $\sigma(\lambda)$-measure 0, we can change
the $\mathfrak H'^{(\lambda)}$, $\lambda\in N$, arbitrarily. We make use of
 this liberty by replacing these $\mathfrak H'^{(\lambda)}$ by the corresponding
 $\mathfrak H''^{(\lambda)}$, so that $k'(\lambda)=k''(\lambda)$, and choose
 \[
 f_m'(\lambda)=f_m''(\lambda)=g_m'(\lambda)=g_m''(\lambda)=
 \begin{cases}
 \psi_m''(\lambda),&m\leq k'(\lambda),\\
 0,&m>k'(\lambda),
 \end{cases}
 \]
 so that
\[
u_{mn}(\lambda)=
\begin{cases}
1,&m=n,\\
0,&m\ne n.
\end{cases}
\]
(All this for $\lambda\in N$ only!) Thus we have
$k'(\lambda)=k''(\lambda)$ and $\{u_{mn}(\lambda)\}$ unitary without
exceptions.

Owing to its definition in \eqref{eq:7.3} $u_{mn}(\lambda)$ is a
$\sigma(\lambda)$-measurable function. Besides \eqref{eq:7.3} gives
\[
f_m''(\lambda)=\sum_n u_{mn}(\lambda)g_n''(\lambda),\qquad
m,n\leq k'(\lambda).
\]
As the $g_n''(\lambda)=\psi_n''(\lambda)$, $n\leq k'(\lambda)$, form a
measurable family in the $''$-sense, we can apply
\hyperref[thm:I]{Theorem I}, (iii)$^*$, and
thus see, that the $f_m''(\lambda)$, $m\leq k'(\lambda)$, too form a measurable
family in the $''$-sense. On the other hand the $f_m'(\lambda)$,
$m\leq k'(\lambda)$, form a measurable family in the $'$-sense.

$\sum_m x_mf_m'(\lambda)$ is the general element of
$\mathfrak H'^{(\lambda)}$, $\sum_m x_mf_m''(\lambda)$ is the general element
of $\mathfrak H''^{(\lambda)}$ ($\sum_m|x_m|^2$ finite), if we let them
correspond to each other, an isomorphism
$\mathfrak H'^{(\lambda)}\longleftrightarrow
\mathfrak H''^{(\lambda)}$ results.

Assume now
\[
\int_{-}'h'(\lambda)\sqrt{d\sigma(\lambda)}
=\int_{-}''h''(\lambda)\sqrt{d\sigma(\lambda)}
\]
and put
\[
h'(\lambda)=\sum_m x_m(\lambda)f_m'(\lambda),\qquad
h''(\lambda)=\sum_m y_m(\lambda)f_m''(\lambda).
\]
Then we have, as we proved at the beginning of this discussion, always with
the possible exception of a set of $\sigma(\lambda)$-measure 0:
$(h'(\lambda),f_m'(\lambda))=(h''(\lambda),f_m''(\lambda))$,
$x_m(\lambda)=y_m(\lambda)$, and so
$h'(\lambda)\longleftrightarrow h''(\lambda)$ in the above isomorphism
$\mathfrak H'^{(\lambda)}\longleftrightarrow
\mathfrak H''^{(\lambda)}$. Thus the isomorphisms
$\mathfrak H'^{(\lambda)}\longleftrightarrow
\mathfrak H''^{(\lambda)}$ carry all $'$-notions into the corresponding
$''$-notions, and conversely.
\end{proof}

\vnpart{II}{Abelian rings and the generalized direct sums}

\section{Effect of a change of the weight-function with the resolution of
unity fixed}\label{sec:8}

We have seen that a representation of a unitary space $\mathfrak H$ as the
generalized direct sum of the unitary spaces $\mathfrak H^{(\lambda)}$ is
essentially (that is, up to replacements of the $\mathfrak H^{(\lambda)}$ by
isomorphic images and up to an arbitrary $\lambda$-set of
$\sigma(\lambda)$-measure 0) characterized by its weight-function
$\sigma(\lambda)$ and its resolution of unity $E(\lambda)$. The only
restriction we had to impose on these was that $\sigma(\lambda)$ must be
equivalent to $E(\lambda)$ (\hyperref[thm:III]{Theorem III}). Our next task is
to show that even
less than this (that is, than $\sigma(\lambda)$ and $E(\lambda)$) suffices to
determine the $\mathfrak H^{(\lambda)}$ up to unessential transformations.

We first get rid of $\sigma(\lambda)$. Let $E(\lambda)$ be a resolution of
unity, and $\sigma_1(\lambda),\sigma_2(\lambda)$ two $N$-functions, both
equivalent to $E(\lambda)$. Then $\sigma_2(\lambda)$ is
$\sigma_1(\lambda)$-totally-continuous and conversely, therefore we can form
$d\sigma_2(\lambda)/d\sigma_1(\lambda)$ and
$d\sigma_1(\lambda)/d\sigma_2(\lambda)$ by \secref{sec:2}, ($\gamma$). Now
we have by \cite{ref11},
\[
\frac{d\sigma_2(\lambda)}{d\sigma_1(\lambda)}
\frac{d\sigma_1(\lambda)}{d\sigma_2(\lambda)}=1,
\]
except perhaps for a set of $\sigma(\lambda)$-measure 0. But on this set we
may redefine both differential quotients, giving them the value 1. Then we have
the displayed identity everywhere, and thus both factors are $\ne0$
everywhere.

Let now $\mathfrak H^{(\lambda)}$ be a representation of $\mathfrak H$ as a
generalized direct sum with the weight-function $\sigma_1(\lambda)$ and the
resolution of unity $E(\lambda)$. We speak correspondingly of
$\sigma_1(\lambda)$-summability and of
$f=\int_{-}f(\lambda)\sqrt{d\sigma_1(\lambda)}$. We now define: A function
$f(\lambda)\in\mathfrak H^{(\lambda)}$ is $\sigma_2(\lambda)$-summable if the
function
\[
\sqrt{\frac{d\sigma_2(\lambda)}{d\sigma_1(\lambda)}}\,f(\lambda)
\quad(\in\mathfrak H^{(\lambda)})
\]
is $\sigma_1(\lambda)$-summable, and then (again by definition!)
\[
\int_{-}f(\lambda)\sqrt{d\sigma_2(\lambda)}
=\int_{-}\sqrt{\frac{d\sigma_2(\lambda)}{d\sigma_1(\lambda)}}
 f(\lambda)\sqrt{d\sigma_1(\lambda)}.
\]
One verifies easily that $\mathfrak H$ is the generalized direct sum of the
$\mathfrak H^{(\lambda)}$ with the weight-function $\sigma_2(\lambda)$ using
these definitions, because it is the generalized direct

% pp. 428--455 of the original journal pagination.
% This portion begins in the final paragraph of Section 8.

% p. 428
sum of the $\mathfrak H^{(\lambda)}$ with the weight-function $\sigma _1(\lambda)$
using the original definitions.  The essential fact to be used, in order to
establish this, is
\begin{align*}
 &\int\left(
   \sqrt{\frac{d\sigma _2(\lambda)}{d\sigma _1(\lambda)}}f(\lambda),
   \sqrt{\frac{d\sigma _2(\lambda)}{d\sigma _1(\lambda)}}g(\lambda)
 \right)d\sigma _1(\lambda)\\
 &\qquad=\int (f(\lambda),g(\lambda))
       \frac{d\sigma _2(\lambda)}{d\sigma _1(\lambda)}d\sigma _1(\lambda)\\
 &\qquad=\int (f(\lambda),g(\lambda))d\sigma _2(\lambda).
\end{align*}
Denoting the resolution of unity which belongs to the
$\sigma _2(\lambda)$-decomposition of $\mathfrak H$ by $E'(\lambda)$ we see that
we have again
\[
 E'(\mu)\int_{-}f(\lambda)\sqrt{d\sigma _2(\lambda)}
   =\int_{-}f_\mu(\lambda)\sqrt{d\sigma _2(\lambda)},
\]
where
\[
 f_\mu(\lambda)=
 \begin{cases}
  f(\lambda),&\lambda\leq\mu,\\
  0,&\lambda>\mu,
 \end{cases}
\]
because the factor
$\sqrt{d\sigma _2(\lambda)/d\sigma _1(\lambda)}$ does not affect the
corresponding rule which holds for the $\sigma _1(\lambda)$-decomposition.
Thus $E'(\mu)=E(\mu)$.  In other words: With the new definitions $\mathfrak H$
is the generalized direct sum of the $\mathfrak H^{(\lambda)}$ with the
weight-function $\sigma _2(\lambda)$ and the resolution of unity $E(\lambda)$.

Remembering Theorem~\ref{thm:III}, we can therefore state:

\emph{Let $E(\lambda)$ be a resolution of unity, and $\sigma _1(\lambda)$,
$\sigma _2(\lambda)$ two $N$-functions equivalent to $E(\lambda)$.  Let
$\mathfrak H$ be the generalized direct sum of the
$\mathfrak H_1^{(\lambda)}$ resp. of the $\mathfrak H_2^{(\lambda)}$ with the
weight-function $\sigma _1(\lambda)$ resp. $\sigma _2(\lambda)$, and in both
cases with the resolution of unity $E(\lambda)$.  Then these two
representations become identical if we replace every
$f(\lambda)\in\mathfrak H_1^{(\lambda)}$ by
$\sqrt{d\sigma _2(\lambda)/d\sigma _1(\lambda)}f(\lambda)$ and apply to this a
suitable isomorphism
$\mathfrak H_1^{(\lambda)}\rightleftarrows\mathfrak H_2^{(\lambda)}$
(depending on $\lambda$ only), and disregard a (fixed) $\lambda$-set of
$\sigma(\lambda)$-measure $0$.}

This result justifies our symbolic notation
$\int_{-}f(\lambda)\sqrt{d\sigma(\lambda)}$: If we do not express the
isomorphism
$\mathfrak H_1^{(\lambda)}\rightleftarrows\mathfrak H_2^{(\lambda)}$ in our
formulae (that is, if we use the same symbol for the corresponding elements of
$\mathfrak H_1^{(\lambda)}$ and $\mathfrak H_2^{(\lambda)}$) it implies the
relation
\[
 \int_{-}\sqrt{\frac{d\sigma _2(\lambda)}{d\sigma _1(\lambda)}}
 f(\lambda)\sqrt{d\sigma _1(\lambda)}
 =\int_{-}f(\lambda)\sqrt{d\sigma _2(\lambda)}.
\]

The above result shows that if we have once given the resolution of unity
$E(\lambda)$ then it is not very important which particular $\sigma(\lambda)$
(equivalent to $E(\lambda)$!) we choose: This choice does not affect the
$\mathfrak H^{(\lambda)}$'s essentially, and within each
$\mathfrak H^{(\lambda)}$ it has only the effect of multiplying everything with
a constant factor
$\sqrt{d\sigma _2(\lambda)/d\sigma _1(\lambda)}$ if we pass from
$\sigma(\lambda)=\sigma _1(\lambda)$ to
$\sigma(\lambda)=\sigma _2(\lambda)$.

\section{Definition of the equivalence. Elementary properties}
\label{sec:9}

We next define a notion of equivalence for two generalized direct
sum-representations of a unitary space $\mathfrak H$.

% p. 429
\begin{definition}
\label{def:3}
Assume that $\mathfrak H$ is at the same time generalized direct sum of the
$\mathfrak H_1^{(\lambda)}$ with the weight-function $\sigma _1(\lambda)$ and of
the $\mathfrak H_2^{(\lambda)}$ with the weight-function $\sigma _2(\lambda)$.
We call these two representations \emph{equivalent} if this is the case:

There exists a (real-valued) function $m(\lambda)$, a positive function
$\gamma(\lambda)$, a set $N_1$ with the $\sigma _1(\lambda)$-measure $0$, and a
set $N_2$ with the $\sigma _2(\lambda)$-measure $0$, which have the following
properties:
\begin{enumerate}
\renewcommand{\labelenumi}{(\arabic{enumi})}
\item $\lambda\rightleftarrows m(\lambda)$ is a one-to-one mapping of the
  complement of $N_1$ on the complement of $N_2$.
\item If $\lambda\notin N_1$ then there exist an isomorphic mapping of
  $\mathfrak H_1^{(\lambda)}$ on $\mathfrak H_2^{(m(\lambda))}$.  We denote this
  isomorphism
  $\mathfrak H_1^{(\lambda)}\rightleftarrows
   \mathfrak H_2^{(m(\lambda))}$ by $g^{(\lambda)}$.
\item $f_2(\lambda)$ ($\in\mathfrak H_2^{(\lambda)}$) is
  $\sigma _2(\lambda)$-summable if and only if
  \[
    f_1(\lambda)=\gamma(\lambda)(g^{(\lambda)})^{-1}f_2(m(\lambda))
  \]
  is $\sigma _1(\lambda)$-summable, and then
  \[
   \int_{-}f_1(\lambda)\sqrt{d\sigma _1(\lambda)}
    =\int_{-}f_2(\lambda)\sqrt{d\sigma _2(\lambda)}.
  \]
\end{enumerate}
We call $\lambda\rightleftarrows m(\lambda)$ the \emph{mapping} and
$\gamma(\lambda)$ the \emph{factor} of this equivalence.
\end{definition}

\medskip\noindent\textbf{Remark 1.}\phantomsection\label{rem:9.1}
Let $K_1$ be an arbitrary set, and $K_2$ its
$\lambda\rightleftarrows m(\lambda)$ image.  Assume that $K_2$ is
$\sigma _2(\lambda)$-measurable.  Choose a measurable family
$\varphi _{2m}(\lambda)$ in $\mathfrak H_2^{(\lambda)}$, and define
\[
 f_2(\lambda)=
 \begin{cases}
  \varphi _{21}(\lambda),&\lambda\in K_2,\\
  0,&\lambda\notin K_2.
 \end{cases}
\]
This $f_2(\lambda)$ is clearly $\sigma _2(\lambda)$-summable; therefore
$f_1(\lambda)=\gamma(\lambda)(g^{(\lambda)})^{-1}f_2(m(\lambda))$ is
$\sigma _1(\lambda)$-summable.  Thus the function
\[
 \lVert f_1(\lambda)\rVert
  =\gamma(\lambda)\lVert f_2(m(\lambda))\rVert
  =\begin{cases}
    \gamma(\lambda),&\lambda\in K_1,\\
    0,&\lambda\notin K_1
   \end{cases}
\]
(we now disregard sets of $\sigma _1(\lambda)$-measure $0$) is
$\sigma _1(\lambda)$-measurable, and so is the set of all $\lambda$ with
$\lVert f_1(\lambda)\rVert\ne0$: $K_1$.  Conversely: If $K_1$ is
$\sigma _1(\lambda)$-measurable, then $K_2$ is
$\sigma _2(\lambda)$-measurable, because the inverse $m^{-1}(\lambda)$ of
$m(\lambda)$ (define $m^{-1}(\lambda)$ as the inverse of $m(\lambda)$ in the
complement of $N_2$, and, as $N_2$ has the $\sigma _2(\lambda)$-measure $0$,
arbitrarily in $N_2$) has the same properties as $m(\lambda)$, if we interchange
$\mathfrak H_1^{(\lambda)},\sigma _1(\lambda)$ and
$\mathfrak H_2^{(\lambda)},\sigma _2(\lambda)$.

So we see: $K_1$ is $\sigma _1(\lambda)$-measurable if and only if $K_2$ is
$\sigma _2(\lambda)$-measurable.

With the above notations we have
\[
 \left\lVert\int_{-}f_1(\lambda)\sqrt{d\sigma _1(\lambda)}\right\rVert^2
  =\int\lVert f_1(\lambda)\rVert^2d\sigma _1(\lambda)
  =\int_{K_1}\gamma(\lambda)d\sigma _1(\lambda),
\]
\[
 \left\lVert\int_{-}f_2(\lambda)\sqrt{d\sigma _2(\lambda)}\right\rVert^2
  =\int\lVert f_2(\lambda)\rVert^2d\sigma _2(\lambda)
  =\int_{K_2}d\sigma _2(\lambda)
  =\mu_{\sigma _2(\lambda)}(K_2),
\]
and as the left sides are equal,
\begin{equation}
 \mu_{\sigma _2(\lambda)}(K_2)
  =\int_{K_1}\gamma(\lambda)d\sigma _1(\lambda).
 \label{eq:9.1}
\end{equation}

% p. 430
So $\mu_{\sigma _1(\lambda)}(K_1)=0$ implies
$\mu_{\sigma _2(\lambda)}(K_2)=0$, and of course conversely.  Furthermore one
easily proves from \eqref{eq:9.1} that
\begin{equation}
 \int\omega(\lambda)d\sigma _2(\lambda)
  =\int\omega(m(\lambda))\gamma(\lambda)d\sigma _1(\lambda)
 \label{eq:9.2}
\end{equation}
for every $\sigma _2(\lambda)$-summable numerical function $\omega(\lambda)$.
(If $\omega(\lambda)\geq0$, then apply the definition
\[
 \int\omega(\lambda)d\sigma(\lambda)
  =\int_0^{+\infty}x\,d[\mu_{\sigma(\lambda)}
      (\text{set of all $\lambda$ with $\omega(\lambda)\leq x$})]
\]
of the Lebesgue--Stieltjes integral, and then extend to complex-valued
$\omega(\lambda)$'s.)

\medskip\noindent\textbf{Remark 2.}\phantomsection\label{rem:9.2}
That our notion of equivalence is reflexive and symmetric is obvious.  The
results of \hyperref[rem:9.1]{Remark 1} concerning the images of sets of measure $0$
prove that it is transitive too.  Thus our notion of equivalence possesses the
three essential properties of an equivalence.  Therefore all possible
representations of $\mathfrak H$ as a generalized direct sum (of any
$\mathfrak H^{(\lambda)}$'s with any weight-function $\sigma(\lambda)$) can be
classified in the customary way: Form for each representation the set of all
those which are equivalent to it, and call this set an \emph{equivalence-class}.
Then any two equivalence-classes are either identical or they have no common
element, and so every representation belongs to one and only one
equivalence-class.

The result of \secref{sec:8} can now be formulated as follows:

\emph{All representations of $\mathfrak H$ as a generalized direct sum to which
the same resolution of unity belongs, are equivalent.  These equivalences have
all the identical mapping $\lambda\rightleftarrows\lambda$.  The factor is
$\gamma(\lambda)=\sqrt{d\sigma _2(\lambda)/d\sigma _1(\lambda)}$ if the
weight-function $\sigma _1(\lambda)$ is changed into $\sigma _2(\lambda)$.}

\section{Discussion of the discrete special case
(cf. \texorpdfstring{\hyperref[sec:4]{\S4}}{Section 4})}
\label{sec:10}

We repeat the procedure of \secref{sec:4} after Definition~\ref{def:1}: Before
we continue the systematic investigation we first discuss the meaning of the
new notion (now it is the equivalence, introduced by Definition~\ref{def:3}) for
totally discontinuous $\sigma(\lambda)$'s, that is for the case in which our
generalized direct sum becomes a common direct sum (cf. \secref{sec:4}).

Consider two equivalent representations of $\mathfrak H$ as a direct sum:
$\mathfrak H_1^{(\lambda)},\sigma _1(\lambda)$ and
$\mathfrak H_2^{(\lambda)},\sigma _2(\lambda)$.  Assume that
$\sigma _1(\lambda)$ is totally discontinuous.  Let, as in \secref{sec:4},
$\lambda _1,\lambda _2,\ldots$ be the (finite or infinite) sequence of all
points $\lambda$ with a $\sigma _1(\lambda)$-measure $\ne0$.  Apply now
Definition~\ref{def:3}.  As $\mu_{\sigma _1(\lambda)}(N_1)=0$, no $\lambda_n$
can belong to $N_1$.  Thus the $\lambda _1,\lambda _2,\ldots$ have distinct
$\lambda\rightleftarrows m(\lambda)$-images
$\mu _1,\mu _2,\ldots$, $\mu_n=m(\lambda_n)$.  By \hyperref[rem:9.1]{Remark 1} in
\secref{sec:9} each $\mu_n$ has a $\sigma _2(\lambda)$-measure $\ne0$.  As the
complement of $(\lambda _1,\lambda _2,\ldots)$ has the
$\sigma _1(\lambda)$-measure $0$, the same \hyperref[rem:9.1]{Remark} yields that its
$\lambda\rightleftarrows m(\lambda)$-image has the
$\sigma _2(\lambda)$-measure $0$.  Adding to it $N_2$, we see: The complement
of $(\mu _1,\mu _2,\ldots)$ has the $\sigma _2(\lambda)$-measure $0$.  Thus
$\sigma _2(\lambda)$, too, is totally discontinuous.

Forming the
$a_n=\mu_{\sigma _1(\lambda)}(\lambda_n)>0$ and the
$b_n=\mu_{\sigma _2(\lambda)}(\mu_n)>0$, as in \secref{sec:4}, the

% p. 431
\hyperref[rem:9.2]{Remark 2} of \secref{sec:9} gives for the sets
$K_1=(\lambda_n)$, $K_2=(\mu_n)$ that
$a_n=(\gamma(\lambda_n))^2b_n$.  That is:
\[
 \gamma(\lambda_n)=\sqrt{\frac{a_n}{b_n}}.
\]

\emph{Summing up, we can state that an equivalence-class contains no
representations with a totally discontinuous $\sigma(\lambda)$, or only such
ones.  In the latter case the common direct sums which arise (in the sense of
\secref{sec:4}) differ only by permutations of their $\lambda_n$'s and the
numerical values of their $a_n$'s.  The subspaces
$\mathfrak H_1^*,\mathfrak H_2^*,\ldots$ of $\mathfrak H$
(cf. \secref{sec:4}) are invariant.}

Thus for the common direct sums our notion of equivalence states precisely the
invariance of the essential features.

\section{\texorpdfstring{$E(\lambda)$}{E(lambda)} constant outside of
\texorpdfstring{$-2<\lambda<-1$}{-2<lambda<-1} can always be secured by
equivalence. \texorpdfstring{$P$}{P} is the set of all
\texorpdfstring{$\alpha(\Lambda)$}{alpha(Lambda)}, and is an equivalence-invariant}
\label{sec:11}

We now resume the general discussion.

Let $\mathfrak H$ be the direct sum of the $\mathfrak H^{(\lambda)}$ with the
weight-function $\sigma(\lambda)$, and let the resolution of unity $E(\lambda)$
belong to this representation.  Form the ring of bounded operators which the
$E(\lambda)$, $-\infty<\lambda<+\infty$, generate: $P$,
cf. \cite[p.~388]{ref2}.  We call $P$ \emph{the ring which belongs} to the above
representation of $\mathfrak H$.

As $E(\lambda)\in P$ and $E(\lambda)$ converges (strongly) to $1$ for
$\lambda\to+\infty$, we have $1\in P$.  As the $E(\lambda)$ are all Hermitean
operators and as they all commute, the set of all $E(\lambda)$,
$-\infty<\lambda<+\infty$, is Abelian, and therefore the ring $P$ is Abelian
too, cf. \cite[p.~389]{ref2}.  So we have:

\emph{The ring $P$ which belongs to a generalized direct sum is always Abelian
and contains $1$.}

Consider now the operators $\alpha(\Lambda)$ defined in \secref{sec:7}.  We
ask: For which functions $\alpha(\lambda)$ is $\alpha(\Lambda)\in P$?  For
\[
 \alpha(\lambda)=e_\mu(\lambda)=
 \begin{cases}
  1,&\lambda\leq\mu,\\
  0,&\lambda>\mu,
 \end{cases}
\]
we have $\alpha(\Lambda)=E(\mu)\in P$, thus
$\alpha(\lambda)=e_\mu(\lambda)$ is such.  Using (c) in \secref{sec:7} we see
that
\[
 \alpha(\lambda)=e_\mu(\lambda)-e_\nu(\lambda)=
 \begin{cases}
  1,&\nu<\lambda\leq\mu,\\
  0,&\text{otherwise}
 \end{cases}
 \qquad(\nu<\mu),
\]
too, is such, and every finite linear aggregate of functions
$\alpha(\lambda)=e_\mu(\lambda)-e_\nu(\lambda)$ and
$\alpha(\lambda)=1$ similarly,---that is every finite-valued intervalwise
constant function.  Now by (d) in \secref{sec:7} every continuous function
$\alpha(\lambda)$ is such, as it is a (uniform) limit of the above functions.
Renewed application of (d) in \secref{sec:7} extends this to all
Baire-functions $\alpha(\lambda)$ and (b) in \secref{sec:7} to all
$\sigma(\lambda)$-measurable $\alpha(\lambda)$'s.  So we see: The ring $P$
contains all $\alpha(\Lambda)$'s.  Conversely: If a ring contains all
$\alpha(\Lambda)$'s, then it contains all $e_\mu(\Lambda)=E(\mu)$ and therefore
$P$ must be a subset of it.  Thus $P$ is the ring generated by all
$\alpha(\Lambda)$'s.

Consider now two equivalent representations of $\mathfrak H$:
$\mathfrak H_1^{(\lambda)},\sigma _1(\lambda)$ and
$\mathfrak H_2^{(\lambda)},\sigma _2(\lambda)$.

% p. 432
Denote their rings by $P_1$ and $P_2$ respectively and their
$\alpha(\Lambda)$'s by $\alpha(\Lambda_1)$ and $\beta(\Lambda_2)$ respectively.
If $\beta(\lambda)$ is $\sigma _2(\lambda)$-measurable, then we know from
\hyperref[rem:9.1]{Remark 1}, after Definition~\ref{def:3}, that
$\alpha(\lambda)=\beta(m(\lambda))$ is $\sigma _1(\lambda)$-measurable, and the
definitions of \secref{sec:7} make it evident that
$\alpha(\Lambda_1)=\beta(\Lambda_2)$.  Thus the set of all
$\alpha(\Lambda_1)$'s is a subset of the set of all $\beta(\Lambda_2)$'s.
Interchanging $\mathfrak H_1^{(\lambda)},\sigma _1(\lambda)$ and
$\mathfrak H_2^{(\lambda)},\sigma _2(\lambda)$, we see that the converse too is
true, thus these two sets are identical.  Therefore they generate the same
ring, that is: $P_1=P_2$.

In other words:

\emph{The set of all $\alpha(\Lambda)$'s is an equivalence-invariant, and so is
the ring $P$ which it generates.}

\medskip
Let $m(\lambda)$ be a function which maps the interval
$-2<\lambda<-1$ in a one-to-one a bicontinuous way on the interval
$-\infty<\lambda<+\infty$ (For instance:
\[
 m(\lambda)=\frac{2\lambda+3}{1-|2\lambda+3|}.)
\]
If the $E(\lambda),\sigma(\lambda)$ are given, define
\[
 E'(\lambda)=
 \begin{cases}
  0,&\lambda\leq-2,\\
  E(m(\lambda)),&-2<\lambda<-1,\\
  1,&\lambda\geq-1,
 \end{cases}
\qquad
 \sigma'(\lambda)=
 \begin{cases}
  \sigma(-\infty),&\lambda\leq-2,\\
  \sigma(m(\lambda)),&-2<\lambda<-1,\\
  \sigma(+\infty),&\lambda\geq-1.
 \end{cases}
\]
Then $E'(\lambda)$ is a resolution of unity, $\sigma'(\lambda)$ is an
$N$-function, and one sees that $\sigma'(\lambda)$ is equivalent to
$E'(\lambda)$ because $\sigma(\lambda)$ is equivalent to $E(\lambda)$.  The set
of all $\lambda$ with $\lambda\leq-2$ or $\lambda\geq-1$ has the
$\sigma'(\lambda)$-measure $0$.  Thus $\lambda\rightleftarrows m(\lambda)$ is a
one-to-one mapping of the complement of a set of $\sigma'(\lambda)$-measure
$0$ on the complement of a set of $\sigma(\lambda)$-measure $0$ (of
$-2<\lambda<-1$ on $-\infty<\lambda<+\infty$) which maps
$\sigma'(\lambda)$-measurable sets on $\sigma(\lambda)$-measurable sets and
conversely, transforming the $\sigma'(\lambda)$ measure into the
$\sigma(\lambda)$ measure.

Define next $\mathfrak H'^{(\lambda)}=\mathfrak H^{(m(\lambda))}$ (for
$-2<\lambda<-1$, the $\lambda\leq-2$ or $\geq-1$ do not matter, as their set
has the $\sigma'(\lambda)$-measure $0$).  Define $\sigma'(\lambda)$-summability
for a function $f'(\lambda)$ ($\in\mathfrak H'^{(\lambda)}$) by the
$\sigma(\lambda)$-summability of the $f(\lambda)$
($\in\mathfrak H^{(\lambda)}$) for which
\[
 f(m(\lambda))=f'(\lambda)
 \quad(\in\mathfrak H'^{(\lambda)}=\mathfrak H^{(m(\lambda))}),
\]
and then
\[
 \int_{-}f'(\lambda)\sqrt{d\sigma'(\lambda)}
  =\int_{-}f(\lambda)\sqrt{d\sigma(\lambda)}.
\]

Clearly $\mathfrak H$ is the generalized direct sum of the
$\mathfrak H'^{(\lambda)}$ with the weight-function

% p. 433
$\sigma'(\lambda)$ because it is so for $\mathfrak H^{(\lambda)}$ and
$\sigma(\lambda)$.  And these two representations are equivalent if we use the
mapping $\lambda\rightleftarrows m(\lambda)$ and the factor $1$.

The resolution of unity belonging to
$\mathfrak H'^{(\lambda)},\sigma'(\lambda)$ is $E'(\lambda)$.

Summing up we can state:

\emph{Every equivalence-class of representations of $\mathfrak H$ as a
generalized direct sum contains at least one, for the resolution of unity of
which $E(\lambda)=0$, $\sigma(\lambda)=\sigma(-\infty)$ for
$\lambda\leq-2$ and $E(\lambda)=1$, $\sigma(\lambda)=\sigma(+\infty)$ for
$\lambda\geq-1$.}

We prove next:

\begin{lemma}
\label{lem:4}
$P$ is the set of all $\alpha(\Lambda)$ where $\alpha(\lambda)$ is any
$\sigma(\lambda)$-measurable function which is bounded if a suitable set of
$\sigma(\lambda)$-measure $0$ is disregarded.  One may restrict
$\alpha(\lambda)$ to bounded Baire-functions, without affecting the result.
\end{lemma}

\begin{proof}
Let us first consider the first part of the Lemma.  We have already seen that
$P$ contains the set of all $\alpha(\Lambda)$, thus it suffices to prove that it
is contained in it.  As both sets are equivalence-invariant, therefore the
above result implies that we need to consider only the case where
$E(\lambda)=0$, $\sigma(\lambda)=\sigma(-\infty)$ for $\lambda\leq-2$ and
$E(\lambda)=1$, $\sigma(\lambda)=\sigma(+\infty)$ for $\lambda\geq-1$.

In this case the \hyperref[rem:7.1]{Remark} in \secref{sec:7} applies: The operator $\Lambda$
exists and is bounded, and $\alpha(\Lambda)$ is the operator-function of
$\Lambda$ in the usual sense.  Thus by the theorem of \cite[p.~213]{ref4}, the
ring generated by $\Lambda$ is contained in the set of all
$\alpha(\Lambda)$.

$E(\lambda)$ is clearly the resolution of unity of $\Lambda$, thus $\Lambda$
generates the same ring as the $E(\lambda)$ with $\lambda<0$ and the
$1-E(\lambda)$ with $\lambda\geq0$, by \cite[pp.~389--390]{ref2}.  As
$E(\lambda)=1$ for all $\lambda\geq-1$, we may as well speak of this ring
generated by all $E(\lambda)$, $-\infty<\lambda<+\infty$, that is $P$.  Thus
$P$ is contained in the set of all $\alpha(\Lambda)$.  This completes the proof
of the first half.  The second half follows from (b) in \secref{sec:7} together
with \cite[p.~213]{ref4}.
\end{proof}

\section{\texorpdfstring{$P$}{P} can be any Abelian ring containing 1, but no
other ring; \texorpdfstring{$P$}{P} determines the generalized direct sum
precisely up to equivalence}
\label{sec:12}

The foregoing results put us in the position to prove this theorem:

\begin{theorem}
\label{thm:IV}
A given unitary space $\mathfrak H$ can be represented as a generalized direct
sum to which a given ring $P$ belongs, if and only if $P$ is Abelian and
contains $1$.  In this case these representations form precisely one
equivalence-class.
\end{theorem}

\begin{proof}
\emph{Necessity of the condition:} We proved at the beginning of
\secref{sec:11} that every ring $P$ which belongs to a generalized direct sum
is Abelian and contains $1$.

\emph{Sufficiency of the condition:} Let $P$ be an Abelian ring containing
$1$.  By \cite[pp.~401, 403]{ref2}, a resolution of unity $E(\lambda)$ exists,
so that $E(\lambda)$ with $\lambda<0$ and the $1-E(\lambda)$ with
$\lambda\geq0$ generate $P$.  As $1\in P$, we can add the $1$ to these.  If
$1$ is included, we can replace the $1-E(\lambda)$ by the $E(\lambda)$ (for
$\lambda\geq0$), and then, as $\lim_{\lambda\to+\infty}E(\lambda)=1$, we may
omit $1$.  Thus $P$ is generated by all $E(\lambda)$,
$-\infty<\lambda<+\infty$.

Now choose an $N$-function $\sigma(\lambda)$ equivalent to $E(\lambda)$, and
apply Theorem~\ref{thm:III} to $E(\lambda),\sigma(\lambda)$.  Then a
representation of $\mathfrak H$ as a generalized direct sum results, to which
the resolution of unity $E(\lambda)$, and thus the ring $P$ belongs.

\emph{Relationship to equivalence:} Let $P$ be an Abelian ring containing
$1$, and $\mathfrak H_1^{(\lambda)}$ a representation of $\mathfrak H$ as a
generalized direct sum with the weight function $\sigma _1(\lambda)$ and the
ring $P$.  Let $\mathfrak H_2^{(\lambda)}$ be another representation of
$\mathfrak H$, with the weight-

% p. 434
function $\sigma _2(\lambda)$.  If these two representations are equivalent,
they must have, as we proved in \secref{sec:11}, the same rings; therefore
$\mathfrak H_2^{(\lambda)},\sigma _2(\lambda)$ has the ring $P$.  If we can
prove conversely, that whenever
$\mathfrak H_2^{(\lambda)},\sigma _2(\lambda)$ has the ring $P$ then
$\mathfrak H_1^{(\lambda)},\sigma _1(\lambda)$ and
$\mathfrak H_2^{(\lambda)},\sigma _2(\lambda)$ are equivalent, then we see:
The $\mathfrak H_2^{(\lambda)},\sigma _2(\lambda)$ with the ring $P$ form
precisely the equivalence-class of
$\mathfrak H_1^{(\lambda)},\sigma _1(\lambda)$.

Assume therefore that $\mathfrak H_1^{(\lambda)},\sigma _1(\lambda)$ has the
resolution of unity $E_1(\lambda)$, that
$\mathfrak H_2^{(\lambda)},\sigma _2(\lambda)$ has the resolution of unity
$E_2(\lambda)$, and that both have the ring $P$.  We want to prove their
equivalence.  Owing to the results of \secref{sec:11} we can restrict ourselves
to the case where $E_1(\lambda)$ and $E_2(\lambda)$ are $0$ for $\lambda<-2$
and $1$ for $\lambda\geq-1$.  And remembering those at the end of
\secref{sec:9} we can replace $\sigma _1(\lambda)$ and $\sigma _2(\lambda)$ by
any two $N$-functions $\tau _1(\lambda)$ and $\tau _2(\lambda)$ equivalent to
$E_1(\lambda)$ and $E_2(\lambda)$ respectively.  Let
$\varphi _1,\varphi _2,\ldots$ be a complete normalized orthogonal set in
$\mathfrak H$, and $a_1,a_2,\ldots$ a sequence of numbers $>0$ with a finite
$\sum_n a_n$, then we can put by ($\delta$) in \secref{sec:2}
\[
 \tau _1(\lambda)=\sum_n a_n\lVert E_1(\lambda)\varphi_n\rVert^2,
 \qquad
 \tau _2(\lambda)=\sum_n a_n\lVert E_2(\lambda)\varphi_n\rVert^2.
\]
$\tau _1(\lambda),\tau _2(\lambda)$ are constant for $\lambda<-2$ and for
$\lambda>-1$.

We form now the operators $\alpha(\Lambda)$ for both resolutions of unity
$E_1(\lambda)$ and $E_2(\lambda)$, denoting them by $\alpha(\Lambda_1)$ and
$\beta(\Lambda_2)$ respectively.  As $\tau _1(\lambda),\tau _2(\lambda)$ are
constant outside of the finite interval $-2\leq\lambda\leq-1$, the \hyperref[rem:7.1]{Remark} in
\secref{sec:7} applies: The operators $\Lambda_1,\Lambda_2$ exist and are
bounded, and $\alpha(\Lambda_1),\beta(\Lambda_2)$ are their respective
operator-functions in the usual sense.  Now apply Lemma~\ref{lem:4}:
$\alpha(\lambda)=\lambda$ is bounded if we disregard the set $\lambda<-2$ or
$\lambda>-1$, which has $\tau _1(\lambda)$- and $\tau _2(\lambda)$-measure
$0$, so $\Lambda_1,\Lambda_2\in P$.  Therefore
$\Lambda_1=m(\Lambda_2)$, where $m(\lambda)$ is a bounded Baire-function, and
similarly $\Lambda_2=n(\Lambda_1)$, where $n(\lambda)$ is a bounded
Baire-function.

So $\Lambda_1=n(m(\Lambda_1))$, therefore the set $N'_1$ of all $\lambda$ with
$n(m(\lambda))\ne\lambda$ has the $\tau _1(\lambda)$-measure $0$ (this is best
seen by applying $\Lambda_1=n(m(\Lambda_1))$ to a
$\int_{-}f(\lambda)\sqrt{d\tau _1(\lambda)}$ for which
$f(\lambda)\ne0$ for all $\lambda$, for instance
$f(\lambda)=\varphi _1(\lambda)$, where $\varphi_m(\lambda)$ is a measurable
family) and similarly the set $N'_2$ of all $\lambda$ with
$m(n(\lambda))\ne\lambda$ has the $\tau _2(\lambda)$-measure $0$.

If $n(m(\lambda))=\lambda$, then $m(n(m(\lambda)))=m(\lambda)$, that is:
$\lambda\notin N'_1$ implies $m(\lambda)\notin N'_2$, thus
$\lambda\rightleftarrows m(\lambda)$ maps the complement of $N'_1$ on a subset
of the complement of $N'_2$, and has there the inverse
$\lambda\rightleftarrows n(\lambda)$.  The same is true if we interchange
$N'_1,m(\lambda)$ with $N'_2,n(\lambda)$.  Therefore
$\lambda\rightleftarrows m(\lambda)$ is a one-to-one mapping of the complement
of $N'_1$ on the complement of $N'_2$ and has there the inverse
$\lambda\rightleftarrows n(\lambda)$.

Consider a non-negative, bounded Baire-function $\gamma(\lambda)$.  We have
\[
 \int\gamma(\lambda)d\tau _1(\lambda)
 =\sum_n a_n\int\gamma(\lambda)
       d\bigl(\lVert E_1(\lambda)\varphi_n\rVert^2\bigr)
 =\sum_n a_n(\gamma(\Lambda_1)\varphi_n,\varphi_n)
\]
\[
 \left(\text{observe that }
 \int\gamma(\lambda)d\bigl(\lVert E_1(\lambda)\varphi_n\rVert^2\bigr)
 =\int\gamma(\lambda)d\left[
     \int_{\mu\leq\lambda}\lVert\varphi_n(\mu)\rVert^2d\tau _1(\mu)
   \right]
 \right.
\]
\[
 \left.
 =\int\gamma(\lambda)\lVert\varphi_n(\lambda)\rVert^2d\tau _1(\lambda)
 =\int(\gamma(\lambda)\varphi_n(\lambda),\varphi_n(\lambda))d\tau _1(\lambda)
 =(\gamma(\Lambda_1)\varphi_n,\varphi_n)\right)
\]

% p. 435
and similarly
\[
 \int\gamma(\lambda)d\tau _2(\lambda)
 =\sum_n a_n\int\gamma(\lambda)
       d\bigl(\lVert E_2(\lambda)\varphi_n\rVert^2\bigr)
 =\sum_n a_n(\gamma(\Lambda_2)\varphi_n,\varphi_n).
\]
As $\gamma(\Lambda_1)=\gamma(n(\Lambda_2))$, these formulae imply
\[
 \int\gamma(\lambda)d\tau _1(\lambda)
  =\int\gamma(n(\lambda))d\tau _2(\lambda).
\]

Now let $K_1$ be a Borel-set,
\[
 \gamma(\lambda)=1_{K_1}(\lambda)=
 \begin{cases}
  1,&\lambda\in K_1,\\
  0,&\lambda\notin K_1.
 \end{cases}
\]
Then
\[
 \gamma(n(\lambda))=1_{K'_2}(\lambda)=
 \begin{cases}
  1,&\lambda\in K'_2,\\
  0,&\lambda\notin K'_2,
 \end{cases}
\]
where $K'_2$ is the inverse-image of $K_1$ under the mapping
$\lambda\rightleftarrows n(\lambda)$.
\footnote{Being an inverse-image of a Borel-set with respect to a
Baire-function, $K'_2$ is necessarily a Borel-set too.  Observe that this is
not true in general for images!\label{fn:8}}
So we have proved
$\mu_{\tau _1(\lambda)}(K_1)=\mu_{\tau _2(\lambda)}(K'_2)$.  Assume now
$K_1\subset\operatorname{Complement}N'_1$.  Then
$K_2=K'_2\cdot\operatorname{Complement}N'_2$ is a Borel set, because
$K'_2,N'_2$ are such sets.  $K_2$ is the
$\lambda\rightleftarrows m(\lambda)$-image of $K_1$ by the properties of
$m(\lambda),n(\lambda)$ derived above.  As
$\mu_{\tau _2(\lambda)}(N'_2)=0$, we have
\begin{equation}
 \mu_{\tau _1(\lambda)}(K_1)=\mu_{\tau _2(\lambda)}(K_2).
 \label{eq:12.1}
\end{equation}

If $K_1$ is an arbitrary Borel-set, then its
$\lambda\rightleftarrows m(\lambda)$-image is the sum of those of
$K_1\cdot N'_1$ and of $K_1\cdot\operatorname{Complement}N'_1$.  The former is
a subset of $N'_2$ and thus has the $\tau _2(\lambda)$ measure $0$.  The latter
is the image of a Borel-subset of $\operatorname{Complement}N'_1$, thus a Borel
set, with the $\tau _2(\lambda)$-measure
$\mu_{\tau _1(\lambda)}(K_1\cdot\operatorname{Complement}N'_1)
=\mu_{\tau _1(\lambda)}(K_1)$.  Thus the total image $K_2$ of $K_1$ is
$\tau _2(\lambda)$-measurable, and \eqref{eq:12.1} holds again.  If $K_1$ is any
set with $\mu_{\tau _1(\lambda)}(K_1)=0$, then it is contained in a Borel-set
$K_1^0$ with $\mu_{\tau _1(\lambda)}(K_1^0)=0$.  If their
$\lambda\rightleftarrows m(\lambda)$-images are $K_2$ and $K_2^0$ respectively,
then $K_2$ is a subset of $K_2^0$, and by \eqref{eq:12.1}
$\mu_{\tau _2(\lambda)}(K_2^0)=0$, therefore
$\mu_{\tau _2(\lambda)}(K_2)=0$.  Thus \eqref{eq:12.1} holds for sets $K_1$ of
$\tau _1(\lambda)$-measure $0$, too.  Combining this with the result for
Borel-sets, we see: If $K_1$ is $\tau _1(\lambda)$-measurable, then its
$\lambda\rightleftarrows m(\lambda)$-image $K_2$ is
$\tau _2(\lambda)$-measurable, and \eqref{eq:12.1} is true.

Consideration of the mapping $\lambda\rightleftarrows n(\lambda)$ proves that
the converse, too, holds: $K_1$ is $\tau _1(\lambda)$-measurable if and only if
$K_2$ is $\tau _2(\lambda)$-measurable, and in this case \eqref{eq:12.1} holds.
The same is true if we make the replacement
\[
 (\tau _1(\lambda),\tau _2(\lambda),m(\lambda))
 \quad\text{by}\quad
 (\tau _2(\lambda),\tau _1(\lambda),n(\lambda)).
\]

Let now $\varphi_m(\lambda)$ be a measurable family in the
$E_1(\lambda),\tau _1(\lambda)$-representation of $\mathfrak H$.
$\varphi_m(\lambda)$ is defined for $m\leq k_1(\lambda)$
($k_1(\lambda)$ is the dimensionality of $\mathfrak H_1^{(\lambda)}$), define
for $m>k_1(\lambda)$ $\varphi_m(\lambda)=0$ (as in \secref{sec:6}).
$\varphi_m(\lambda)$ is clearly $\tau _1(\lambda)$-summable, define
\[
 f_m=\int_{-}\varphi_m(\lambda)\sqrt{d\tau _1(\lambda)},
\]
and then in the $E_2(\lambda),\tau _2(\lambda)$-representation of
$\mathfrak H$
\[
 f_m=\int_{-}\psi_m(\lambda)\sqrt{d\tau _2(\lambda)},
\]

% p. 436
We have for any bounded Baire-function $\alpha(\lambda)$
\[
 \begin{aligned}
 (\alpha(\Lambda _1)f_m,f_n)
 &=\left(\int_{-}\alpha(\lambda)\varphi_m(\lambda)
          \sqrt{d\tau _1(\lambda)},
          \int_{-}\varphi_n(\lambda)\sqrt{d\tau _1(\lambda)}\right)\\
 &=\int(\alpha(\lambda)\varphi_m(\lambda),\varphi_n(\lambda))
       d\tau _1(\lambda)\\
 &=\int\alpha(\lambda)(\varphi_m(\lambda),\varphi_n(\lambda))
       d\tau _1(\lambda)\\
 &=\begin{cases}
    0,&m\ne n,\\
    \displaystyle\int\alpha(\lambda)d\tau _1(\lambda),&m=n,
   \end{cases}
 \end{aligned}
\]
and at the same time
\[
 \begin{aligned}
 (\alpha(\Lambda _2)f_m,f_n)
 &=\left(\int_{-}\alpha(\lambda)\psi_m(\lambda)
          \sqrt{d\tau _2(\lambda)},
          \int_{-}\psi_n(\lambda)\sqrt{d\tau _2(\lambda)}\right)\\
 &=\int(\alpha(\lambda)\psi_m(\lambda),\psi_n(\lambda))
       d\tau _2(\lambda)\\
 &=\int\alpha(\lambda)(\psi_m(\lambda),\psi_n(\lambda))
       d\tau _2(\lambda).
 \end{aligned}
\]
Considering $\alpha(\Lambda _2)=\alpha(m(\Lambda _1))$ this means:
\[
 \begin{aligned}
 &\int\alpha(\lambda)(\psi_m(\lambda),\psi_n(\lambda))d\tau _2(\lambda)\\[-1mm]
 &\quad=\begin{cases}
   0,&m\ne n,\\[2mm]
   \displaystyle\begin{aligned}
   &\int_{\{k_1(\lambda)\geq m\}}\alpha(m(\lambda))d\tau _1(\lambda)\\
   &\quad=\int_{\{k_1(n(\lambda))\geq m\}}\alpha(\lambda)d\tau _2(\lambda)
   \end{aligned},
   &m=n.
  \end{cases}
 \end{aligned}
\]
Thus
\[
 \begin{aligned}
 &\int\alpha(\lambda)\bigl((\psi_m(\lambda),\psi_n(\lambda))
             -e_{mn}(\lambda)\bigr)d\tau _2(\lambda)=0,\\
 &e_{mn}(\lambda)=
  \begin{cases}
   1,&m=n\leq k_1(n(\lambda)),\\
   0,&\text{otherwise},
  \end{cases}
 \end{aligned}
\]
for every bounded Baire-function $\alpha(\lambda)$, and therefore for every
bounded $\tau _2(\lambda)$-measurable $\alpha(\lambda)$.  This implies
$(\psi_m(\lambda),\psi_n(\lambda))-e_{mn}(\lambda)=0$, except for a set
$N_2^{\prime\prime\prime(m,n)}$ of $\tau _2(\lambda)$-measure $0$.
\footnote{Choose $\alpha_1(\lambda)=\epsilon(Z)$ for the $\epsilon(Z)$ of
footnote~\ref{fn:7} and with $Z$ equal to the
$\{\cdots\}$-expression, then
$\int|\,(\psi_m(\lambda),\psi_n(\lambda))-e_{mn}(\lambda)\,|
d\tau _2(\lambda)=0$ results, and thus the statement in the text.\label{fn:9}}
The sum of all $N_2^{\prime\prime\prime(m,n)}$, say $N'''_2$, is a fixed set of
$\tau _2(\lambda)$-measure $0$, and with the same property.  So for
$\lambda\notin N'''_2$ the $\psi_m(\lambda)$,
$m<k_1(n(\lambda))$ form a normalized orthogonal set in
$\mathfrak H_2^{(\lambda)}$, therefore the dimensionality $k_2(\lambda)$ of
$\mathfrak H_2^{(\lambda)}$ is then $\geq k_1(n(\lambda))$.

Similarly there exists a set $N'''_1$ of $\tau _1(\lambda)$-measure $0$, so
that $\lambda\notin N'''_1$ implies
$k_1(\lambda)\geq k_2(m(\lambda))$.  Let now $N_1$ be the sum of
$N'_1$, $N'''_1$, and the $\lambda\rightleftarrows n(\lambda)$-image of
$N'''_2$, and $N_2$ the sum of $N'_2,N'''_2$, and the
$\lambda\rightleftarrows m(\lambda)$-image of $N'''_1$.  Clearly $N_1,N_2$ are
respective sets of $\tau _1(\lambda),\tau _2(\lambda)$-measure $0$, and
$\lambda\rightleftarrows m(\lambda)$ is a one-to-one mapping of the complement
of $N_1$ on the complement of $N_2$, its inverse

% p. 437
being there $\lambda\rightleftarrows n(\lambda)$.  Furthermore we have for
$\lambda\notin N_1$, $k_1(\lambda)\geq k_2(m(\lambda))$ and for
$\lambda\notin N_2$, $k_2(\lambda)\geq k_1(n(\lambda))$.
$\lambda\notin N_1$ implies $m(\lambda)\notin N_2$,
$k_2(m(\lambda))\geq k_1(n(m(\lambda)))=k_1(\lambda)$ too, so
$k_1(\lambda)=k_2(m(\lambda))$.

Choose a measurable family $\varphi_{1m}(\lambda)$ for the
$E_1(\lambda),\tau _1(\lambda)$-representation,\linebreak and one
$\varphi_{2m}(\lambda)$ for the $E_2(\lambda),\tau _2(\lambda)$-representation.
Define the isomorphism $g^{(\lambda)}$ of $\mathfrak H_1^{(\lambda)}$ on
$\mathfrak H_2^{(m(\lambda))}$, for $\lambda\notin N_1$, by
\[
 g^{(\lambda)}\left(\sum_m x_m\varphi_{1m}(\lambda)\right)
   =\sum_m x_m\varphi_{2m}(m(\lambda))
 \qquad(\text{remember }k_1(\lambda)=k_2(m(\lambda))!).
\]
Introduce now the functions $x_{1m}(\lambda)$ and $x_{2m}(\lambda)$ in the
sense of Definition~\ref{def:2} for the respective
$E_1(\lambda),\tau _1(\lambda)$- and
$E_2(\lambda),\tau _2(\lambda)$-representation.  Then the connection between
$f_1(\lambda)$ and $f_2(\lambda)$ as described in (3) in
Definition~\ref{def:3}, and putting $\gamma(\lambda)=1$, is equivalent to
$x_{1m}(\lambda)=x_{2m}(m(\lambda))$ (disregarding a $\lambda$-set of
$\tau _1(\lambda)$-measure $0$).  Now the equivalence of the
$\tau _1(\lambda)$-summability-condition of ($\gamma$), ($\delta$) in
Definition~\ref{def:2}, in terms of the $x_{1m}(\lambda)$ and of the
$\tau _2(\lambda)$-summability condition in the same sense, in terms of the
$x_{2m}(\lambda)$, can be established at once: It follows from the fact that if
$K_1$ is the $\lambda\rightleftarrows m(\lambda)$-inverse-image of $K_2$ then
the $\tau _1(\lambda)$-measurability of $K_1$ and the
$\tau _2(\lambda)$-measurability of $K_2$ are equivalent, and
$\mu_{\tau _1(\lambda)}(K_1)=\mu_{\tau _2(\lambda)}(K_2)$.

Summing up: Our $N_1,N_2,m(\lambda)$ and $\gamma(\lambda)=1$ fulfill the
conditions (1)--(3) in Definition~\ref{def:3}.  Therefore we have established
the equivalence of the $E_1(\lambda),\tau _1(\lambda)$- and of the
$E_2(\lambda),\tau _2(\lambda)$-representations of $\mathfrak H$, and thus
completed our proof.
\end{proof}

\medskip
By the \hyperref[thm:IV]{preceding theorem} we have found what is usually termed a ``complete
set of invariants'' for the representations of a given unitary space as a
generalized direct sum, with respect to the notion of equivalence, as defined
in Definition~\ref{def:3} and \hyperref[rem:9.2]{Remark 2} in \secref{sec:9}.  It
consists of the one invariant $P$, the ring which belongs to the representation,
and the range of $P$ is precisely the set of all Abelian rings containing $1$.

\section{Discussion of the decomposition of an operator
\texorpdfstring{$A\in P'$}{A in P'}:
\texorpdfstring{$A\sim\sum A(\lambda)$}{A as a sum of A(lambda)}. Necessary and
sufficient conditions and uniqueness}
\label{sec:13}

Let $\mathfrak H$ be the generalized direct sum of the
$\mathfrak H^{(\lambda)}$ with the weight-function $\sigma(\lambda)$, the
resolution of unity $E(\lambda)$, and the ring $P$.  Form the ring $P'$ of all
those bounded linear operators in $\mathfrak H$ which commute with every
element of $P$, cf. \cite[p.~388]{ref2}.

Consider an operator $A\in P'$, put
\[
 \mathop{\mathrm{l.u.b.}}_{f\ne0}\frac{\lVert Af\rVert}{\lVert f\rVert}
  =\lvert\!\lvert\!\lvert A\rvert\!\rvert\!\rvert=C.
\]
Let $\varphi_m(\lambda)$, $m\leq k(\lambda)$, be a measurable family, and
define again $\varphi_m(\lambda)=0$ for $m>k(\lambda)$.
$\varphi_m(\lambda)$ is $\sigma(\lambda)$-summable, put
$\omega_m=\int_{-}\varphi_m(\lambda)\sqrt{d\sigma(\lambda)}$.  Form now the
$f_m(\lambda)$ with
\[
 A^*\omega_m=\int_{-}f_m(\lambda)\sqrt{d\sigma(\lambda)}
\]
($A^*$ is the adjoint of $A$, cf. \cite[p.~112]{ref1} and
\cite[p.~295]{ref5}), and define
$a_{mn}(\lambda)=(f_m(\lambda),\varphi_n(\lambda))$, so that
$f_m(\lambda)=\sum_n a_{mn}(\lambda)\varphi_n(\lambda)$.

% p. 438
Now choose an arbitrary $f\in\mathfrak H$, and put
\[
 f=\int_{-}\bar f(\lambda)\sqrt{d\sigma(\lambda)},
 \qquad
 Af=\int_{-}\bar{\bar f}(\lambda)\sqrt{d\sigma(\lambda)},
\]
then
\[
 (f,A^*\omega_m)=(Af,\omega_m),
\]
and so
\[
 \int\left(\bar f(\lambda),
       \sum_n\overline{a_{mn}(\lambda)}\varphi_n(\lambda)\right)
       d\sigma(\lambda)
  =\int(\bar{\bar f}(\lambda),\varphi_m(\lambda))d\sigma(\lambda).
\]
Replace $f$ by
\[
 E(\mu)f=\int_{-}\bar f_\mu(\lambda)\sqrt{d\sigma(\lambda)},
 \qquad
 \bar f_\mu(\lambda)=
 \begin{cases}
  \bar f(\lambda),&\lambda\leq\mu,\\
  0,&\lambda>\mu.
 \end{cases}
\]
As $E(\mu)\in P$, $A\in P'$, so that $E(\mu)$ and $A$ commute, this has the
effect of replacing $Af$ by
\[
 AE(\mu)f=E(\mu)Af
 =\int_{-}\bar{\bar f}_\mu(\lambda)\sqrt{d\sigma(\lambda)},
 \qquad
 \bar{\bar f}_\mu(\lambda)=
 \begin{cases}
  \bar{\bar f}(\lambda),&\lambda\leq\mu,\\
  0,&\lambda>\mu.
 \end{cases}
\]
Thus the above formula becomes
\[
 \int_{\lambda\leq\mu}\left(\bar f(\lambda),
       \sum_n\overline{a_{mn}(\lambda)}\varphi_n(\lambda)\right)
       d\sigma(\lambda)
  =\int_{\lambda\leq\mu}(\bar{\bar f}(\lambda),\varphi_m(\lambda))
       d\sigma(\lambda),
\]
\[
 \int_{\lambda\leq\mu}\left[
  \sum_n a_{mn}(\lambda)(\bar f(\lambda),\varphi_n(\lambda))
  -(\bar{\bar f}(\lambda),\varphi_m(\lambda))\right]d\sigma(\lambda)=0.
\]
Thus
\[
 \sum_n a_{mn}(\lambda)(\bar f(\lambda),\varphi_n(\lambda))
   =(\bar{\bar f}(\lambda),\varphi_m(\lambda)),
\]
except perhaps for a set $\widetilde N^{(m)}$ of $\sigma(\lambda)$-measure $0$
(cf. the argument of footnote~\ref{fn:7}).  The sum of all $\widetilde N^{(m)}$, say
$\widetilde N$, is a fixed set of $\sigma(\lambda)$-measure $0$ (depending of course
on the $\bar f(\lambda)$ and $\bar{\bar f}(\lambda)$) and with the same
property.  If we put
\[
 \bar f(\lambda)=\sum_m x_m(\lambda)\varphi_m(\lambda),
 \qquad
 \bar{\bar f}(\lambda)=\sum_m y_m(\lambda)\varphi_m(\lambda),
\]
then we can formulate the above result as follows:
\begin{equation}
 y_m(\lambda)=\sum_n a_{mn}(\lambda)x_n(\lambda),
 \label{eq:13.1}
\end{equation}
except for a set of $\sigma(\lambda)$-measure $0$.  (The series on the right
side converges absolutely in any event, because
\[
 \sum_n|a_{mn}(\lambda)|^2=\lVert f_m(\lambda)\rVert^2,
 \qquad
 \sum_n|x_n(\lambda)|^2=\lVert\bar f(\lambda)\rVert^2
\]
are finite.)  Observe that while the $a_{mn}(\lambda)$ are determined numbers
(depending only on $A$), the exceptional set in \eqref{eq:13.1} depends on the
choice of the $\bar f(\lambda)$ and $\bar{\bar f}(\lambda)$ too.  If, however,
the $\bar f(\lambda)$ (that is, the $x_n(\lambda)$) are given, then
\eqref{eq:13.1} can be used to define the $\bar{\bar f}(\lambda)$ (that is the
$y_n(\lambda)$), and then no exceptions occur at all.

Consider now the function
\[
 \mathop{\mathrm{l.u.b.}}_{\substack{\sum_m|x_m|^2<+\infty\\
                                     \sum_m|x_m|^2>0}}
 \frac{\sum_m\left|\sum_n a_{mn}(\lambda)x_n\right|^2}
      {\sum_m|x_m|^2}=C(\lambda).
\]

% p. 439
The $\sum_n$ in the numerator are absolutely convergent (because
$\sum_n|a_{mn}(\lambda)|^2$, $\sum_n|x_n|^2$ are finite), the $\sum_m$ in the
numerator has non-negative terms and thus is either absolutely convergent or
properly divergent to $+\infty$, the $\sum_m$ in the denominator is absolutely
convergent and $>0$.  Thus $C(\lambda)$ is $\geq0$, $\leq+\infty$.

Let $S$ be the set of all $\lambda$ with $C(\lambda)>C^2$, and for a given
sequence $x_1,x_2,\ldots$ with $\sum_m|x_m|^2=1$ let
$S_p(x_1,x_2,\ldots)$ be the set of all $\lambda$ with $k(\lambda)=p$ and
\[
 \sum_m\left|\sum_n a_{mn}(\lambda)x_n\right|^2>C^2.
\]
Form all sequences $x_1,x_2,\ldots$ of length of $p$ with
$\sum_m|x_m|^2=1$, in which only a finite number of $x_n\ne0$ occur, and in
which the ratios
$\Re x_1:\Im x_1:\Re x_2:\Im x_2:\cdots$ are all rational.  They form a
countable set, and we can therefore enumerate them:
$x_1^{(\nu)},x_2^{(\nu)},\ldots$ for $\nu=1,2,\ldots$.  Now clearly
\[
 \begin{aligned}
 &\mathop{\mathrm{l.u.b.}}_{\substack{\sum_m|x_m|^2<+\infty\\
                                     \sum_m|x_m|^2>0}}
 \frac{\sum_m\left|\sum_n a_{mn}(\lambda)x_n\right|^2}
      {\sum_m|x_m|^2}\\
 &\quad=
 \mathop{\mathrm{l.u.b.}}_{\sum_m|x_m|^2=1}
 \sum_m\left|\sum_n a_{mn}(\lambda)x_n\right|^2\\
 &\quad=
 \mathop{\mathrm{l.u.b.}}_{\nu=1,2,\ldots}
 \sum_m\left|\sum_n a_{mn}(\lambda)x_n^{(\nu)}\right|^2.
 \end{aligned}
\]
Therefore
\[
 S=\sum_{p,\nu=1,2,\ldots}S_p(x_1^{(\nu)},x_2^{(\nu)},\ldots).
\]
We will prove below
$\mu_{\sigma(\lambda)}(S_p(x_1,x_2,\ldots))=0$ in general, then this relation
gives $\mu_{\sigma(\lambda)}(S)=0$.

Assume in fact that a sequence $x_1,x_2,\ldots$ with
$\sum_m|x_m|^2=1$ is given.  The set $S_p(x_1,x_2,\ldots)$ is
$\sigma(\lambda)$-measurable (because the functions $k(\lambda)$,
$a_{mn}(\lambda)$ are such), and we can therefore define
\[
 \bar f(\lambda)=
 \begin{cases}
  \sum_m x_m\varphi_m(\lambda),&\lambda\in S_p(x_1,x_2,\ldots),\\
  0,&\text{otherwise}.
 \end{cases}
\]
We have $\lVert\bar f(\lambda)\rVert^2=\sum_m|x_m|^2=1$ resp. $0$.  This and
the $\sigma(\lambda)$-measurability of $S_p(x_1,x_2,\ldots)$, give, by
($\gamma$), ($\delta$) in Definition~\ref{def:2}, at once that
$\bar f(\lambda)$ is $\sigma(\lambda)$-summable.  Now put
$f=\int_{-}\bar f(\lambda)\sqrt{d\sigma(\lambda)}$.  Thus
$\lVert f\rVert^2=\int\lVert\bar f(\lambda)\rVert^2d\sigma(\lambda)
=\mu_{\sigma(\lambda)}(S_p(x_1,x_2,\ldots))$.  For
$Af=\int_{-}\bar{\bar f}(\lambda)\sqrt{d\sigma(\lambda)}$ we have
\[
 \bar{\bar f}(\lambda)=
 \begin{cases}
  \sum_m y_m(\lambda)\varphi_m(\lambda),\quad
  y_m(\lambda)=\sum_n a_{mn}(\lambda)x_n,
    &\lambda\in S_p(x_1,x_2,\ldots),\\
  0,&\text{otherwise}.
 \end{cases}
\]
Thus $\lVert\bar{\bar f}(\lambda)\rVert^2
=\sum_m|y_m(\lambda)|^2
=\sum_m|\sum_n a_{mn}(\lambda)x_n|^2>C^2$ resp. $=0$, and so if
$\mu_{\sigma(\lambda)}(S_p(x_1,x_2,\ldots))>0$, then
$\lVert Af\rVert^2=\int\lVert\bar{\bar f}(\lambda)\rVert^2d\sigma(\lambda)
>C^2\mu_{\sigma(\lambda)}(S_p(x_1,x_2,\ldots))$.  This gives
$\lVert Af\rVert^2>C^2\lVert f\rVert^2$, therefore $f\ne0$ and
$\lVert Af\rVert/\lVert f\rVert>C$.  Considering the definition of $C$, this
is impossible.  Thus

% p. 440
there must be $\mu_{\sigma(\lambda)}(S_p(x_1,x_2,\ldots))=0$, and this proves,
as we saw above, $\mu_{\sigma(\lambda)}(S)=0$.

We have proved $C(\lambda)\leq C^2$, except on a set of
$\sigma(\lambda)$-measure $0$.  As such a set does not matter, we can replace
in it all $a_{mn}(\lambda)$'s by $0$, thus securing $C(\lambda)\leq C^2$
without exception.  In other words:
\[
 \frac{\sum_m\left|\sum_n a_{mn}(\lambda)x_n\right|^2}
      {\sum_m|x_m|^2}\leq C^2
 \quad\text{if}\quad
 \sum_m|x_m|^2
 \begin{cases}
  <+\infty,\\
  >0.
 \end{cases}
\]
This enables us to define an operation $A(\lambda)$ in
$\mathfrak H^{(\lambda)}$ by
\[
 f(\lambda)=\sum_m x_m\varphi_m(\lambda),
 \qquad
 A(\lambda)f(\lambda)=\sum_m y_m\varphi_m(\lambda),
 \qquad
 y_m=\sum_n a_{mn}(\lambda)x_n.
\]
Our inequality proves first of all that the series for $A(\lambda)f(\lambda)$
converges, so that $A(\lambda)$ is an everywhere defined operator, obviously
linear.  Second, it proves
$\lVert A(\lambda)f(\lambda)\rVert/\lVert f(\lambda)\rVert\leq C$, thus
establishing the boundedness of $A(\lambda)$, and even
\[
 \mathop{\mathrm{l.u.b.}}_{f(\lambda)\ne0}
 \frac{\lVert A(\lambda)f(\lambda)\rVert}{\lVert f(\lambda)\rVert}
 =\lvert\!\lvert\!\lvert A(\lambda)\rvert\!\rvert\!\rvert\leq C
 \qquad(\text{cf. \cite[p.~385]{ref2}}).
\]
Our previous results concerning $f$ and $Af$ can now be formulated as follows:
\[
 \text{If }f=\int_{-}f(\lambda)\sqrt{d\sigma(\lambda)},
 \quad\text{then}\quad
 Af=\int_{-}A(\lambda)f(\lambda)\sqrt{d\sigma(\lambda)}.
\]

\medskip
We now define:

\begin{definition}
\label{def:4}
Let $\mathfrak H$ be the generalized direct sum of the
$\mathfrak H^{(\lambda)}$, with the weight function $\sigma(\lambda)$, the
resolution of unity $E(\lambda)$, and the ring $P$.  Let $A$ be a linear,
bounded operator in $\mathfrak H$, and $A(\lambda)$ one in
$\mathfrak H^{(\lambda)}$, for each $\lambda>-\infty$, $<+\infty$.  We say
that $A$ is the \emph{generalized direct sum} of the $A(\lambda)$, in symbols
$A\sim\sum A(\lambda)$, if
\[
 A\int_{-}f(\lambda)\sqrt{d\sigma(\lambda)}
  =\int_{-}A(\lambda)f(\lambda)\sqrt{d\sigma(\lambda)}
\]
holds for all $\sigma(\lambda)$-summable $f(\lambda)$'s.
\end{definition}

The results obtained in the first part of this section put us in the position
to prove:

\begin{theorem}
\label{thm:V}
If $A$ is given, then the necessary and sufficient condition for the existence
of a system of $A(\lambda)$'s with $A\sim\sum A(\lambda)$ is $A\in P'$.  Then
these $A(\lambda)$ are uniquely determined (up to an arbitrary $\lambda$-set
of $\sigma(\lambda)$-measure $0$).

If the $A(\lambda)$ are given, then the necessary and sufficient conditions for
the existence of a $A$ with $A\sim\sum A(\lambda)$ are these:
\begin{enumerate}
\renewcommand{\labelenumi}{(\roman{enumi})}
\item Choose a measurable family $\varphi_m(\lambda)$.
Each (numerical) function
\[
 (A(\lambda)\varphi_m(\lambda),\varphi_n(\lambda))
\]
is $\sigma(\lambda)$-measurable.
\item The numbers
$\lvert\!\lvert\!\lvert A(\lambda)\rvert\!\rvert\!\rvert$ are bounded, if a
suitable $\lambda$-set of $\sigma(\lambda)$-measure $0$ is disregarded.
\end{enumerate}
Then this $A$ is uniquely determined.
\end{theorem}

% p. 441
\begin{proof}
\emph{$A$ is given: Necessity of the condition:}

The formula
\[
 A\int_{-}f(\lambda)\sqrt{d\sigma(\lambda)}
 =\int_{-}A(\lambda)f(\lambda)\sqrt{d\sigma(\lambda)}
\]
implies
\[
 \begin{aligned}
 A\alpha(\Lambda)\int_{-}f(\lambda)\sqrt{d\sigma(\lambda)}
 &=A\int_{-}\alpha(\lambda)f(\lambda)\sqrt{d\sigma(\lambda)}\\
 &=\int_{-}A(\lambda)\alpha(\lambda)f(\lambda)
       \sqrt{d\sigma(\lambda)}\\
 &=\int_{-}\alpha(\lambda)A(\lambda)f(\lambda)
       \sqrt{d\sigma(\lambda)}\\
 &=\alpha(\Lambda)\int_{-}A(\lambda)f(\lambda)
       \sqrt{d\sigma(\lambda)}\\
 &=\alpha(\Lambda)A\int_{-}f(\lambda)\sqrt{d\sigma(\lambda)},
 \end{aligned}
\]
that is $A\alpha(\Lambda)f=\alpha(\Lambda)Af$.  Thus $A$ commutes with every
$\alpha(\Lambda)$, that is with every element of $P$.  Therefore $A\in P'$.

\emph{Sufficiency of the condition:} This was proved in the first part of this
section.

\emph{Uniqueness:} Assume $A\sim\sum A(\lambda)$ and
$A\sim\sum B(\lambda)$.  Then we have for every $\sigma(\lambda)$-summable
$f(\lambda)$
\[
 A\int_{-}f(\lambda)\sqrt{d\sigma(\lambda)}
 =\int_{-}A(\lambda)f(\lambda)\sqrt{d\sigma(\lambda)}
 =\int_{-}B(\lambda)f(\lambda)\sqrt{d\sigma(\lambda)},
\]
and so $A(\lambda)f(\lambda)=B(\lambda)f(\lambda)$ except for a $\lambda$-set
of $\sigma(\lambda)$-measure $0$ (depending on the $f(\lambda)$).  Let
$\varphi_m(\lambda)$ be a measurable family, extend again the definition of
these $\varphi_m(\lambda)$, $m\leq k(\lambda)$:
$\varphi_m(\lambda)=0$ for $m>k(\lambda)$.  $\varphi_m(\lambda)$ is summable,
therefore we see: $A(\lambda)\varphi_m(\lambda)=B(\lambda)\varphi_m(\lambda)$
except for a set $N_m^*$ of $\sigma(\lambda)$-measure $0$.  So the sum of all
$N_m^*$, $N^{**}$, is a fixed set of $\sigma(\lambda)$-measure $0$, and for
$\lambda\notin N^{**}$ we have
$A(\lambda)\varphi_m(\lambda)=B(\lambda)\varphi_m(\lambda)$ for all
$m=1,2,\ldots$.  But the $\varphi_m(\lambda)$, $m\leq k(\lambda)$ form already
a complete normalized orthogonal set in $\mathfrak H^{(\lambda)}$, and so we
have even $A(\lambda)=B(\lambda)$ for $\lambda\notin N^{**}$.

\emph{The $A(\lambda)$ are given: Necessity of the condition:}

Assume $A\sim\sum A(\lambda)$, and form the $\varphi_m(\lambda)$ as above,
define $\omega_m=\int_{-}\varphi_m(\lambda)\sqrt{d\sigma(\lambda)}$.  Then as
$A\omega_m=\int_{-}A(\lambda)\varphi_m(\lambda)\sqrt{d\sigma(\lambda)}$, both
$A(\lambda)\varphi_m(\lambda)$ and $\varphi_n(\lambda)$ are
$\sigma(\lambda)$-summable, and so
$(A(\lambda)\varphi_m(\lambda),\varphi_n(\lambda))$ is
$\sigma(\lambda)$-measurable, fulfilling condition (i).

Constructing the $A(\lambda)$ of $A$ with the method of the first part of this
section, we obtain
$\lvert\!\lvert\!\lvert A(\lambda)\rvert\!\rvert\!\rvert
\leq\lvert\!\lvert\!\lvert A\rvert\!\rvert\!\rvert$ everywhere.  Therefore
we have in general
$\lvert\!\lvert\!\lvert A(\lambda)\rvert\!\rvert\!\rvert
\leq\lvert\!\lvert\!\lvert A\rvert\!\rvert\!\rvert$, except for a set of
$\sigma(\lambda)$-measure $0$.

\emph{Sufficiency of the condition:} All functions
$(A(\lambda)\varphi_m(\lambda),\varphi_n(\lambda))$ are
$\sigma(\lambda)$-measurable, and
$\lvert\!\lvert\!\lvert A(\lambda)\rvert\!\rvert\!\rvert\leq C$, except for a
set of $\sigma(\lambda)$-measure $0$.  Let $f(\lambda)$ be
$\sigma(\lambda)$-summable, then all
$(f(\lambda),\varphi_m(\lambda))$ are $\sigma(\lambda)$-measurable, and so
\[
 \begin{aligned}
 (A(\lambda)f(\lambda),\varphi_n(\lambda))
 &=(f(\lambda),A^*(\lambda)\varphi_n(\lambda))\\
 &=\sum_m(f(\lambda),\varphi_m(\lambda))
          (\varphi_m(\lambda),A^*(\lambda)\varphi_n(\lambda))\\
 &=\sum_m(f(\lambda),\varphi_m(\lambda))
          (A(\lambda)\varphi_m(\lambda),\varphi_n(\lambda))
 \end{aligned}
\]
is $\sigma(\lambda)$-measurable too.  Furthermore
$\int\lVert f(\lambda)\rVert^2d\sigma(\lambda)$ is finite, and so
\[
 \int\lVert A(\lambda)f(\lambda)\rVert^2d\sigma(\lambda)
 \leq\int C^2\lVert f(\lambda)\rVert^2d\sigma(\lambda)
 =C^2\int\lVert f(\lambda)\rVert^2d\sigma(\lambda)
\]
is finite too.

% p. 442
Thus $A(\lambda)f(\lambda)$ is $\sigma(\lambda)$-summable, and we can define
an operator $A$ by
\[
 A\int_{-}f(\lambda)\sqrt{d\sigma(\lambda)}
 =\int_{-}A(\lambda)f(\lambda)\sqrt{d\sigma(\lambda)}.
\]
This $A$ is clearly linear, and as
\[
 \begin{aligned}
 \left\lVert A\int_{-}f(\lambda)\sqrt{d\sigma(\lambda)}\right\rVert^2
 &=\left\lVert\int_{-}A(\lambda)f(\lambda)
       \sqrt{d\sigma(\lambda)}\right\rVert^2\\
 &=\int\lVert A(\lambda)f(\lambda)\rVert^2d\sigma(\lambda)\\
 &\leq C^2\int\lVert f(\lambda)\rVert^2d\sigma(\lambda)\\
 &=C^2\left\lVert\int_{-}f(\lambda)\sqrt{d\sigma(\lambda)}\right\rVert^2,
 \end{aligned}
\]
that is $\lVert Af\rVert^2\leq C^2\lVert f\rVert^2$,
$\lVert Af\rVert\leq C\lVert f\rVert$, therefore $A$ is bounded too.
$A\sim\sum A(\lambda)$ is obvious from $A$'s definition.

\emph{Uniqueness:} Obvious, as the $A(\lambda)$ determine each
$A\int_{-}f(\lambda)\sqrt{d\sigma(\lambda)}$, that is each $Af$.
\end{proof}

\section{Computation rules for \texorpdfstring{$A\sim\sum A(\lambda)$}{operator
decompositions}. Convergence criteria. Operator functions. Behavior under
equivalences}
\label{sec:14}

We now derive some properties of the decomposition $A\sim\sum A(\lambda)$,
which are immediate consequences of the existence-and-uniqueness-criteria of
the \hyperref[thm:V]{preceding theorem}.

In the fourth and fifth part of the proof of \hyperref[thm:V]{this theorem} we saw, that
$\lvert\!\lvert\!\lvert A(\lambda)\rvert\!\rvert\!\rvert
\leq\lvert\!\lvert\!\lvert A\rvert\!\rvert\!\rvert$, except for a set of
$\sigma(\lambda)$-measure $0$, and that if
$\lvert\!\lvert\!\lvert A(\lambda)\rvert\!\rvert\!\rvert\leq C$, except for a
set of $\sigma(\lambda)$-measure $0$, then
$\lvert\!\lvert\!\lvert A\rvert\!\rvert\!\rvert\leq C$.  This means:

\emph{$\lvert\!\lvert\!\lvert A\rvert\!\rvert\!\rvert$ is the smallest number
$C$ for which
$\lvert\!\lvert\!\lvert A(\lambda)\rvert\!\rvert\!\rvert\leq C$ holds,
disregarding a suitable set of $\sigma(\lambda)$-measure $0$.  In other words:
$\lvert\!\lvert\!\lvert A\rvert\!\rvert\!\rvert$ is the ``least upper bound up
to a set of $\sigma(\lambda)$-measure $0$'' of the
$\lvert\!\lvert\!\lvert A(\lambda)\rvert\!\rvert\!\rvert$.}

If $A\sim\sum\alpha(\lambda)1$ ($1$ is the ``unit operator'': $1f=f$, this time in
$\mathfrak H^{(\lambda)}$), then we see:
$(\alpha(\lambda)\varphi_1(\lambda),\varphi_1(\lambda))=\alpha(\lambda)$ is
$\sigma(\lambda)$-measurable,
$\lvert\!\lvert\!\lvert\alpha(\lambda)1\rvert\!\rvert\!\rvert
=|\alpha(\lambda)|$ is bounded up to a set of $\sigma(\lambda)$-measure $0$, so
$\alpha(\Lambda)$ can be formed.  Comparison of the respective definitions
shows, that $A=\alpha(\Lambda)$.  Conversely $A=\alpha(\Lambda)$ obviously
implies $A\sim\sum\alpha(\lambda)1$.  As $P$ is the set of all
$\alpha(\Lambda)$, we have:

\emph{$A\in P$ means (observe that we consider only the $A\in P'$, and that
$P\subset P'$, because $P$ is Abelian), that
$A\sim\sum\alpha(\lambda)1$ for some (complex-valued) function
$\alpha(\lambda)$.}

Finally we prove:

\emph{If $A\sim\sum A(\lambda)$, $B\sim\sum B(\lambda)$, then
$\alpha A\sim\sum\alpha A(\lambda)$,
$A^*\sim\sum A^*(\lambda)$,
$AB\sim\sum A(\lambda)B(\lambda)$.  Any one of the following properties is
possessed by $A$ if and only if it is possessed by all $A(\lambda)$, except for
a set of $\sigma(\lambda)$-measure $0$: To be a projection, to be definite, to
 be Hermitean, to be normal, to be unitary, to be partially unitary.
(Cf. \cite[pp.~74, 74, 70, 115, 71]{ref1}, and \cite[p.~141]{ref10}.)}

The statements concerning $\alpha A,A+B,AB$ are obvious.  The one concerning $A^*$
(the adjoint) is proved as follows: As
\[
 (A(\lambda)^*\varphi_m(\lambda),\varphi_n(\lambda))
 =\overline{(A(\lambda)\varphi_n(\lambda),\varphi_m(\lambda))},
 \qquad
 \lvert\!\lvert\!\lvert A(\lambda)^*\rvert\!\rvert\!\rvert
 =\lvert\!\lvert\!\lvert A(\lambda)\rvert\!\rvert\!\rvert,
\]

% p. 443
there exists an $A_1\sim\sum A(\lambda)^*$ owing to the existence of
$A\sim\sum A(\lambda)$.  Now one verifies immediately $A_1=A^*$.  All other
statements follow by applying these results to the equations which characterize
the respective notions.  These are respectively
$A^*=A$, $A=A^2$, an $X$ with $A=X^*X$ exists,
\footnote{$A=X^*X$ is clearly definite: $(Af,f)=(X^*Xf,f)=(Xf,Xf)\geq0$.
If $A$ is definite we can form the operator function $X=\sqrt A$, then
$X=X^*$, $A=X^2$, so $A=X^*X$.\label{fn:10}}
$A^*=A$, $A^*A=AA^*$, $A^*A=AA^*=1$, $AA^*A=A$.

The convergence-properties of the decomposition $A\sim\sum A(\lambda)$ are
contained in this criterium:

\emph{If $A\sim\sum A(\lambda)$,
$A_n\sim\sum A_n(\lambda)$, then $A$ is the strong limit of
$A_1,A_2,\ldots$ (cf. \cite[pp.~381--382]{ref2}) if and only if the two
following conditions are fulfilled:
\begin{enumerate}
\item[(I)] The
$\lvert\!\lvert\!\lvert A_n(\lambda)\rvert\!\rvert\!\rvert$ are uniformly
bounded, disregarding a suitable $\lambda$-set of $\sigma(\lambda)$-measure
$0$.
\item[(II)] Every subsequence $n_1,n_2,\ldots$ of $1,2,\ldots$ has a
subsequence $m_1,m_2,\ldots$, such that $A(\lambda)$ is the strong limit of
$A_{m_1}(\lambda),A_{m_2}(\lambda),\ldots$, disregarding a suitable
$\lambda$-set of $\sigma(\lambda)$-measure $0$ (that depends on
$n_1,n_2,\ldots$).
\end{enumerate}}

If $A$ is the strong limit of $A_1,A_2,\ldots$, then the
$\lvert\!\lvert\!\lvert A_n\rvert\!\rvert\!\rvert$ are uniformly bounded
\cite[p.~382]{ref2},
$\lvert\!\lvert\!\lvert A_n\rvert\!\rvert\!\rvert\leq C$.  Therefore
$\lvert\!\lvert\!\lvert A_n(\lambda)\rvert\!\rvert\!\rvert\leq C$, except for
a set of $\sigma(\lambda)$-measure $0$, and so (I) is necessary.  Thus we have
to prove, that (II) is characteristic for the strong convergence of
$A_1,A_2,\ldots$ to $A$ if
$\lvert\!\lvert\!\lvert A_n\rvert\!\rvert\!\rvert\leq C$,
$\lvert\!\lvert\!\lvert A_n(\lambda)\rvert\!\rvert\!\rvert\leq C$
(exceptions as above) are assumed.  In this situation strong convergence of
$A_n$ to $A$ is equivalent to strong limit $A_nf=Af$ and strong convergence of
$A_n(\lambda)$ to $A(\lambda)$ is equivalent to strong limit
$A_n(\lambda)f(\lambda)=A(\lambda)f(\lambda)$, for all elements $f$ or
$f(\lambda)$ of a complete normalized orthogonal system in $\mathfrak H$ or
$\mathfrak H^{(\lambda)}$ \cite[pp.~385--386]{ref2}.  Let the
$\varphi_m(\lambda)$ be a measurable family, and put
$\varphi_m=\int_{-}\varphi_m(\lambda)\sqrt{d\sigma(\lambda)}$; let $\psi_m$ be
a complete normalized orthogonal set in $\mathfrak H$, and put
$\psi_m=\int_{-}\psi_m(\lambda)\sqrt{d\sigma(\lambda)}$.  Let the
$f_p=\int_{-}f_p(\lambda)\sqrt{d\sigma(\lambda)}$ be a sequence which contains
these $\varphi_m$ and $\psi_m$.  Then we can replace strong limit $A_n=A$ by
strong limit $A_nf_p=Af_p$ for all $p=1,2,\ldots$; and strong limit
$A_n(\lambda)=A(\lambda)$ by strong limit
$A_n(\lambda)f_p(\lambda)=A(\lambda)f_p(\lambda)$ for all $p=1,2,\ldots$.
So we can make this substitution both in the statement of our theorem and in
(II).  (But we continue assuming
$\lvert\!\lvert\!\lvert A_n\rvert\!\rvert\!\rvert\leq C$,
$\lvert\!\lvert\!\lvert A_n(\lambda)\rvert\!\rvert\!\rvert\leq C$.)  Observe,
that the sequence $m_1,m_2,\ldots$ (and the set of exceptions) in (II) may even
depend on $f_p$: we can always get rid of this dependence, and obtain a new
sequence $m_1,m_2,\ldots$ which works for all $f_p$, $p=1,2,\ldots$,
simultaneously, by applying the diagonal process.  Now we will prove the
equivalence in question for each $f_p$ separately, in other words: If
$\lvert\!\lvert\!\lvert A_n\rvert\!\rvert\!\rvert\leq C$,
$\lvert\!\lvert\!\lvert A_n(\lambda)\rvert\!\rvert\!\rvert\leq C$, then $Af$
is the strong limit of $A_1f,A_2f,\ldots$, if and only if every subsequence
$n_1,n_2,\ldots$ of $1,2,\ldots$ has a subsequence $m_1,m_2,\ldots$, such that
$A(\lambda)f(\lambda)$ is the strong limit of
$A_{m_1}(\lambda)f(\lambda),A_{m_2}(\lambda)f(\lambda),\ldots$
($f=\int_{-}f(\lambda)\sqrt{d\sigma(\lambda)}$), disregarding a suitable
$\lambda$-set of $\sigma(\lambda)$-measure $0$ (depending on
$n_1,n_2,\ldots$ and $f$).

% p. 444
Observe that
\[
 Af=\int_{-}A(\lambda)f(\lambda)\sqrt{d\sigma(\lambda)},
 \qquad
 A_nf=\int_{-}A_n(\lambda)f(\lambda)\sqrt{d\sigma(\lambda)}.
\]
Application of Lemma~\ref{lem:2} in \secref{sec:5}, to the
$A_{m_1}f,A_{m_2}f,\ldots$ and $Af$ (for $g_1,g_2,\ldots$ and
$\lim_{n\to\infty}g_n$) shows, that our condition is necessary.  In order to
prove its sufficiency, assume that it is fulfilled.  We have then
\[
 \lim_{r\to\infty}
 \lVert A_{m_r}(\lambda)f(\lambda)-A(\lambda)f(\lambda)\rVert^2=0,
\]
except for a $\lambda$-set of $\sigma(\lambda)$-measure $0$, and
\[
 \lVert A_{m_r}(\lambda)f(\lambda)-A(\lambda)f(\lambda)\rVert^2
 \leq(2C\lVert f(\lambda)\rVert)^2,
\]
with the same exception.  As
\[
 \int(2C\lVert f(\lambda)\rVert)^2d\sigma(\lambda)
 =4C^2\int\lVert f(\lambda)\rVert^2d\sigma(\lambda)
 =4C^2\lVert f\rVert^2
\]
is finite, this implies
\[
 \lim_{r\to\infty}\int
 \lVert A_{m_r}(\lambda)f(\lambda)-A(\lambda)f(\lambda)\rVert^2
 d\sigma(\lambda)=0,
 \qquad
 \lim_{r\to\infty}\lVert A_{m_r}f-Af\rVert^2=0.
\]
Now if $Af$ were not the strong limit of $A_1f,A_2f,\ldots$, there would exist
an $\epsilon>0$, such that $\lVert A_nf-Af\rVert^2\geq\epsilon$ occurs for
infinitely many $n$'s.  Then a subsequence $n_1,n_2,\ldots$ of
$1,2,\ldots$ with $\lVert A_{n_r}f-Af\rVert^2\geq\epsilon$ would exist, and so
there could not be
$\lim_{r\to\infty}\lVert A_{m_r}f-Af\rVert^2=0$ for any subsequence
$m_1,m_2,\ldots$ of $n_1,n_2,\ldots$.  Thus the proof is complete.

\emph{Replacing (II) by the stronger condition}

\emph{(II$'$) $A(\lambda)$ is the strong limit of
$A_1(\lambda),A_2(\lambda),\ldots$, disregarding a suitable $\lambda$-set of
$\sigma(\lambda)$-measure $0$,}

\emph{makes our conditions sufficient (but not necessary).}

We are now in a position to prove:

\emph{If $A\sim\sum A(\lambda)$ and $\alpha(A)$ is an operator-function of $A$
(in the sense of \cite[p.~204]{ref4}, or \cite[p.~241]{ref11}), then
$\alpha(A)\sim\sum\alpha(A(\lambda))$.}

This is obvious for $\alpha(x)=1$ and $\alpha(x)=x$.  Our rules of computation
concerning $A,A+B,AB$ extend it to all polynomials $\alpha(x)$; the above
convergence criterium (preferably in its sufficient form (I), (II$'$)) to all
continuous functions, and then to all Baire-functions $\alpha(x)$.  For a
general $\alpha(x)$ apply \cite[p.~231]{ref11}: Two Baire-function
$\alpha_1(x),\alpha_2(x)$ can be found, for which
$\alpha_1(x)\leq\alpha(x)\leq\alpha_2(x)$, and which differ from each other
(and thus from $\alpha(x)$) only on a set, which, using the terminology of
\cite[p.~227]{ref11}, is a ``null set with respect to $A$''.  Thus
$\alpha_1(A)=\alpha_2(A)=\alpha(A)$.  Besides
$\alpha_1(A)\sim\sum\alpha_1(A(\lambda))$,
$\alpha_2(A)\sim\sum\alpha_2(A(\lambda))$, therefore both
$\alpha(A)\sim\sum\alpha_1(A(\lambda))$ and
$\alpha(A)\sim\sum\alpha_2(A(\lambda))$.  Thus
$\alpha_1(A(\lambda))=\alpha_2(A(\lambda))$, except for a set of
$\sigma(\lambda)$-measure $0$.  For these $\lambda$'s therefore
$\alpha_1(x),\alpha_2(x)$ differ from each other only on a ``null set with
respect to $A(\lambda)$'' (cf. above).  As
$\alpha_1(x)\leq\alpha(x)\leq\alpha_2(x)$, this means that
$\alpha(A(\lambda))$ too can be formed, and that
$\alpha_1(A(\lambda))=\alpha(A(\lambda))$.  Thus
$\alpha(A)\sim\sum\alpha(A(\lambda))$, completing the proof.

\medskip
The behavior of the decomposition $A\sim\sum A(\lambda)$ under equivalence is
this:

\emph{If $A\sim\sum A(\lambda)$, and the representation of $\mathfrak H$ as a
generalized direct sum is replaced
% p. 445
by an equivalent one, with the mapping
$\lambda\rightleftarrows m(\lambda)$, then the decomposition of $A$ in this new
representation will be $A\sim\sum B(\lambda)$, where $B(\lambda)$ is defined
as follows: $B(m(\lambda))$ is the operator which corresponds to $A(\lambda)$
under the isomorphism $g^{(\lambda)}$ of
$\mathfrak H_1^{(\lambda)}$ with $\mathfrak H_2^{(m(\lambda))}$, that is
$B(m(\lambda))=g^{(\lambda)}A(\lambda)(g^{(\lambda)})^{-1}$.}

One verifies this statement immediately, by substituting it into condition (3)
of Definition~\ref{def:3}.  Observe that the factor of the equivalence has no
effect on the $B(\lambda)$.

\vnpart{III}{Decomposition of rings}

\section{\texorpdfstring{$\sigma(\lambda)$}{sigma(lambda)}-measurable
dependence of \texorpdfstring{$\lambda$}{lambda}: for
\texorpdfstring{$f(\lambda)$, $A(\lambda)$, $M(\lambda)$}{f(lambda), A(lambda),
M(lambda)}. Discussion of \texorpdfstring{$f(\lambda)$, $A(\lambda)$}{f(lambda),
A(lambda)}}
\label{sec:15}

In what follows we will have to consider quantities of various sorts, which
depend on the parameter $\lambda$ of a generalized direct sum.  It will be of
importance to have a notion of measurability for such a dependence.  It is
given by the definition which follows:

\begin{definition}
\label{def:5}
Let $\mathfrak H$ be the generalized direct sum of the
$\mathfrak H^{(\lambda)}$ with the weight-function $\sigma(\lambda)$.  We will
consider an element $f(\lambda)$ of $\mathfrak H^{(\lambda)}$ or a linear
bounded operator $A(\lambda)$ in $\mathfrak H^{(\lambda)}$, or a ring
(containing $1$) $M(\lambda)$ in $\mathfrak H^{(\lambda)}$, depending on the
parameter $\lambda$.  The dependence on $\lambda$ is
$\sigma(\lambda)$-measurable if the following conditions are fulfilled:

\medskip\noindent
\emph{For $f(\lambda)$:}\quad For every $\sigma(\lambda)$-summable
$g(\lambda)$ the (numerical) function $(f(\lambda),g(\lambda))$ is
$\sigma(\lambda)$-measurable.

\medskip\noindent
\emph{For $A(\lambda)$:}\quad For every $\sigma(\lambda)$-measurable
$f(\lambda)$ the (element of $\mathfrak H^{(\lambda)}$-) function
$A(\lambda)f(\lambda)$ is $\sigma(\lambda)$-measurable.

\medskip\noindent
\emph{For $M(\lambda)$:}\quad A (finite or infinite) sequence of
$\sigma(\lambda)$-measurable (operator in $\mathfrak H^{(\lambda)}$-) functions
$A_n(\lambda)$, $n=1,2,\ldots$, exists, so that $M(\lambda)$ is the ring
generated by the $A_n(\lambda)$, $n=1,2,\ldots$ (for that $\lambda$, in
$\mathfrak H^{(\lambda)}$).
\end{definition}

\medskip\noindent\textbf{Remark 1.}\phantomsection\label{rem:15.1}
\emph{Ad $f(\lambda)$:} Our definition is clearly the condition (i)$_1$ in the
characterization (i) of the $\sigma(\lambda)$-summability in
Definition~\ref{def:1}.  If $\varphi_m(\lambda)$ is a measurable family, it
suffices to postulate the $\sigma(\lambda)$-measurability of all
$(f(\lambda),\varphi_m(\lambda))$, $m=1,2,\ldots$.  This is clearly necessary,
and it is sufficient too, because if $g(\lambda)$ is $\sigma(\lambda)$-summable,
then
\[
 (f(\lambda),g(\lambda))
 =\sum_m(f(\lambda),\varphi_m(\lambda))
   \overline{(g(\lambda),\varphi_m(\lambda))}
\]
will also be $\sigma(\lambda)$-measurable.

The above formula shows, also, that if $f(\lambda),g(\lambda)$ are
$\sigma(\lambda)$-measurable, then the (numerical) function
$(f(\lambda),g(\lambda))$ has that property too, and with it
$\lVert f(\lambda)\rVert=\sqrt{(f(\lambda),f(\lambda))}$.

It is obvious, that if $f(\lambda),g(\lambda)$ are
$\sigma(\lambda)$-measurable, then $f(\lambda)+g(\lambda)$ has that property
too; and if the (numerical) function $\alpha(\lambda)$ is
$\sigma(\lambda)$-measurable, then $\alpha(\lambda)f(\lambda)$ has that
property too.

\medskip\noindent\textbf{Remark 2.}\phantomsection\label{rem:15.2}
\emph{Ad $A(\lambda)$:} Considering \hyperref[rem:15.1]{Remark 1}, we may as well
require the $\sigma(\lambda)$-measurability of the (numerical) function
$(A(\lambda)f(\lambda),g(\lambda))$ for all $\sigma(\lambda)$-measurable
$f(\lambda),g(\lambda)$.  If $\varphi_m(\lambda)$ is a measurable family, it
suffices to postulate the $\sigma(\lambda)$-measurability of all
$(A(\lambda)\varphi_m(\lambda),\varphi_n(\lambda))$, $m,n=1,2,\ldots$.  This
is clearly necessary, and it is sufficient too, because if
$f(\lambda),g(\lambda)$ are $\sigma(\lambda)$-measurable, then
\[
 (A(\lambda)f(\lambda),g(\lambda))
 =\sum_n\left[\sum_m
 (A(\lambda)\varphi_m(\lambda),\varphi_n(\lambda))
 (f(\lambda),\varphi_m(\lambda))
 \overline{(g(\lambda),\varphi_n(\lambda))}\right]
\]

% p. 446
will also be $\sigma(\lambda)$-measurable.  This form of our definition is
clearly the condition (i) in the criterium for the existence of an
$A\sim\sum A(\lambda)$ in Theorem~\ref{thm:V}.

It is obvious that if $A(\lambda),B(\lambda)$ are
$\sigma(\lambda)$-measurable, then $A(\lambda)+B(\lambda)$ has that property
too; and if the (numerical) function $\alpha(\lambda)$ is
$\sigma(\lambda)$-measurable, then $\alpha(\lambda)A(\lambda)$ has that
property too.  As
\[
 (A^*(\lambda)f(\lambda),g(\lambda))
 =\overline{(A(\lambda)g(\lambda),f(\lambda))},
\]
the
$\sigma(\lambda)$-measurability of $A(\lambda)$ implies that one of
$A^*(\lambda)$.  If $A(\lambda),B(\lambda)$ are
$\sigma(\lambda)$-measurable, then $A(\lambda)B(\lambda)$ has that property
too: Let $f(\lambda),g(\lambda)$ be $\sigma(\lambda)$-measurable, then
$B(\lambda)f(\lambda)$ is $\sigma(\lambda)$-measurable, and hence
$A(\lambda)B(\lambda)f(\lambda)$ is $\sigma(\lambda)$-measurable too.

If $A_1(\lambda),A_2(\lambda),\ldots$ converge strongly (or even only weakly)
to $A(\lambda)$, then their $\sigma(\lambda)$-measurability implies that one of
$A(\lambda)$, because
$\lim_{n\to\infty}(A_n(\lambda)f(\lambda),g(\lambda))
=(A(\lambda)f(\lambda),g(\lambda))$.  All these facts permit us to conclude in
the usual way that if $A(\lambda)$ is $\sigma(\lambda)$-measurable, then any
operator-function $\varphi(A(\lambda))$ of it ($A(\lambda)$ Hermitian and
$\varphi(x)$ a Baire-function) is $\sigma(\lambda)$-measurable too.

If $A(\lambda)$ is $\sigma(\lambda)$-measurable, so is the (numerical)
function $\lvert\!\lvert\!\lvert A(\lambda)\rvert\!\rvert\!\rvert$: We know
from the considerations preceding Definition~\ref{def:4}, that
\[
 \lvert\!\lvert\!\lvert A(\lambda)\rvert\!\rvert\!\rvert
 =\sqrt{
 \mathop{\mathrm{l.u.b.}}_{\nu=1,2,\ldots}
 \sum_m\left|\sum_n a_{mn}(\lambda)x_n^{(\nu)}\right|^2},
\]
where the $x_1^{(\nu)},x_2^{(\nu)},\ldots$, $\nu=1,2,\ldots$, are the
sequences described there, and
$a_{mn}(\lambda)=(A(\lambda)\varphi_n(\lambda),\varphi_m(\lambda))$.  All these
expressions are $\sigma(\lambda)$-measurable.

Our next task is to discuss the consequences of $\sigma(\lambda)$-measurability
for rings (containing $1$) $M(\lambda)$, in $\mathfrak H^{(\lambda)}$.  Here
however a more elaborate argument is necessary.

\section{The operators of norm at most 1 in the weak topology form a complete
and separable metric space. A choice lemma for analytic sets}
\label{sec:16}

We need a lemma about the sets of M. Suslin and W. Lusin (cf. the monograph
\cite{ref14}, and the exposition in \cite[pp.~90--93, 179--196, 208--212,
276--277]{ref15}), which we will call ``analytic sets'' following the prevalent
usage.  But first we make some remarks of a general topological nature.

Observe that if $S$ is the set of all linear operators $A$ of a unitary space
$\mathfrak H$ with
$\lvert\!\lvert\!\lvert A\rvert\!\rvert\!\rvert\leq1$, then there exists a
metric in $S$ which is equivalent to its weak topology.  If
$\varphi_1,\varphi_2,\ldots$ is a complete normalized orthogonal set in
$\mathfrak H$ define
\[
 (\operatorname{Dist}(A,B))^2
 =\sum_{m,n=1}^{\infty}\frac{1}{2^{m+n}}
   |(A\varphi_m,\varphi_n)-(B\varphi_m,\varphi_n)|^2.
\]
This is clearly a metric, and as the weak topology in $S$ is separable
(cf. \cite[p.~386, footnote~38]{ref2}), we need only to prove, in order to
establish its equivalence to the weak topology, that
$\lim_{p\to\infty}\operatorname{Dist}(A_p,A)=0$ is equivalent to weak
$\lim_{p\to\infty}A_p=A$.  The latter means
$\lim_{p\to\infty}(A_p\varphi_m,\varphi_n)=(A\varphi_m,\varphi_n)$ for all
$m,n=1,2,\ldots$ (cf. \cite[pp.~385--386]{ref2}) and this is equivalent to the
former, as all series for $(\operatorname{Dist}(A,B))^2$ converge uniformly:
\[
 \frac{1}{2^{m+n}}|(A\varphi_m,\varphi_n)-(B\varphi_m,\varphi_n)|^2
 \leq\frac{4}{2^{m+n}},
 \qquad
 \sum_{m,n=1}^{\infty}\frac{4}{2^{m+n}}=4.
\]

% p. 447
This metric space is separable, as $S$ in the weak topology is separable
(cf. loc. cit. above), and it is complete too.  If
$\lim_{p,q\to\infty}\operatorname{Dist}(A_p,A_q)=0$, then
$\lim_{p,q\to\infty}|(A_p\varphi_m,\varphi_n)-(A_q\varphi_m,\varphi_n)|=0$,
and so $\lim_{p\to\infty}(A_p\varphi_m,\varphi_n)=a_{mn}$ exists.  If
$\sum_{m=1}^{\infty}|x_m|^2\leq1$ then
\[
 \begin{aligned}
 \sum_{n=1}^M\left|\sum_{m=1}^N a_{mn}x_m\right|^2
 &=\lim_{p\to\infty}\sum_{n=1}^M
      \left|\sum_{m=1}^N(A_p\varphi_m,\varphi_n)x_m\right|^2\\
 &\leq\limsup_{p\to\infty}\sum_{n=1}^{\infty}
      \left|\sum_{m=1}^N(A_p\varphi_m,\varphi_n)x_m\right|^2\\
 &=\limsup_{p\to\infty}\sum_{n=1}^{\infty}
      \left|\left(A_p\left\{\sum_{m=1}^N x_m\varphi_m\right\},
      \varphi_n\right)\right|^2\\
 &=\limsup_{p\to\infty}
      \left\lVert A_p\left\{\sum_{m=1}^N x_m\varphi_m\right\}\right\rVert^2\\
 &\leq\limsup_{p\to\infty}
      \left\lVert\sum_{m=1}^N x_m\varphi_m\right\rVert^2\\
 &=\limsup_{p\to\infty}\sum_{m=1}^N|x_m|^2\leq1,
 \end{aligned}
\]
and $N\to\infty$ and then $M\to\infty$ give
$\sum_{n=1}^{\infty}|\sum_{m=1}^{\infty}a_{mn}x_m|^2\leq1$.

Thus
\[
 A\left(\sum_{m=1}^{\infty}x_m\varphi_m\right)
 =\sum_{m=1}^{\infty}y_m\varphi_m,
 \qquad
 y_n=\sum_{m=1}^{\infty}a_{mn}x_m
\]
defines a linear operator $A$ with
$\lvert\!\lvert\!\lvert A\rvert\!\rvert\!\rvert\leq1$, so $A\in S$.  Now as
$\lim_{p\to\infty}(A_p\varphi_m,\varphi_n)=(A\varphi_m,\varphi_n)$, therefore
weak $\lim_{p\to\infty}A_p=A$, $\lim_{p\to\infty}\operatorname{Dist}(A_p,A)=0$.
So we see that $S$ in the weak topology may be looked at as a complete and
separable metric space.

Observe finally that if $S_1,S_2$ are two complete and separable metric spaces,
then their direct product $S_1\times S_2$ is another one.
\footnote{The direct product $S_1\times S_2$ is the set of all (ordered) pairs
$(a^{(1)},a^{(2)})$, $a^{(1)}\in S_1$, $a^{(2)}\in S_2$, the neighborhoods of
$(a_0^{(1)},a_0^{(2)})$ being the sets of all $(a^{(1)},a^{(2)})$ with
$a^{(1)}\in O_1$, $a^{(2)}\in O_2$ where $O_1$ is a neighborhood of
$a_0^{(1)}$, and $O_2$ one of $a_0^{(2)}$.  If a metric
$\operatorname{Dist}_1(a^{(1)},b^{(1)})$ exists in $S_1$ and another one
$\operatorname{Dist}_2(a^{(2)},b^{(2)})$ in $S_2$, then the topology in the
direct product $S_1\times S_2$ can be described by the metric
$\operatorname{Dist}((a^{(1)},a^{(2)}),(b^{(1)},b^{(2)}))
=\max(\operatorname{Dist}_1(a^{(1)},b^{(1)}),
\operatorname{Dist}_2(a^{(2)},b^{(2)}))$ in $S_1\times S_2$.\label{fn:11}}
This generalizes by repetition to any finite number of factors.

Our lemma on analytic sets is this:

\begin{lemma}
\label{lem:5}
Let $S$ be a complete and separable metric space, $A$ an analytic set in $S$
(cf. \cite[pp.~91, 211]{ref15}), and $F(a)$ a function which is defined and
continuous for all $a\in A$ and has (real) numerical values.  Let $K$ be the
 set of all values, which $F(a)$, $a\in S$, assumes.  Let $\sigma(\lambda)$ be
an arbitrary $N$-function.

Under these assumptions we have:
\begin{enumerate}
\renewcommand{\labelenumi}{(\roman{enumi})}
\item $K$ is a $\sigma(\lambda)$-measurable set.

% p. 448
\item There exists a function $f(\lambda)$, which is defined for all
$\lambda\in K$ and assumes values in $S$, such that:
\begin{enumerate}
\renewcommand{\labelenumii}{(\roman{enumi})$_{\arabic{enumii}}$.}
\item If $O$ is an open subset of $S$, then the set of all $\lambda\in K$ with
$f(\lambda)\in O$ is $\sigma(\lambda)$-measurable.
\item $F(f(\lambda))=\lambda$ for every $\lambda\in K$.
\end{enumerate}
\end{enumerate}
\end{lemma}

\begin{proof}
Denote the points of discontinuity of $\sigma(\lambda)$ by
$\lambda_1,\lambda_2,\ldots$ (their number is finite or countably infinite),
and the ``jumps'' of $\sigma(\lambda)$ in these points by resp.
$a_n=\sigma(\lambda_n)-\sigma(\lambda_n-0)$.  Form
$\tau_1(\lambda)=\sum_{\lambda_n\leq\lambda}a_n$.  We know from the discussion
at the beginning of \secref{sec:4}, that every set is
$\tau_1(\lambda)$-measurable.  As $\sigma(\lambda)-\tau_1(\lambda)$ is
monotonous (if $\lambda\leq\mu$, then
\[
 [\sigma(\mu)-\tau_1(\mu)]-[\sigma(\lambda)-\tau_1(\lambda)]
 =\sigma(\mu)-\sigma(\lambda)-\sum_{\lambda<\lambda_n\leq\mu}a_n\geq0),
\]
therefore we see that $\sigma(\lambda)-\tau_1(\lambda)$-measurability implies
$(\sigma(\lambda)-\tau_1(\lambda))+\tau_1(\lambda)=\sigma(\lambda)$-
measurability.  Besides $\sigma(\lambda)-\tau_1(\lambda)$ is obviously
continuous.  Let now $\tau_2(\lambda)$ be a continuous and nowhere constant
$N$-function (for instance: $\tau_2(\lambda)=\lambda/(1+|\lambda|)$).  Then
$\sigma(\lambda)-\tau_1(\lambda)+\tau_2(\lambda)=\tau(\lambda)$ is a continuous
and nowhere constant $N$-function too.  It is clear from ($\gamma$) in
\secref{sec:2} that $\sigma(\lambda)-\tau_1(\lambda)$ is
$\sigma(\lambda)-\tau_1(\lambda)+\tau_2(\lambda)=\tau(\lambda)$-totally
continuous.  Therefore $\tau(\lambda)$-measurability implies
$\sigma(\lambda)-\tau_1(\lambda)$-measurability, and thus
$\sigma(\lambda)$-measurability too.  Thus we may prove our lemma for
$\tau(\lambda)$ instead of $\sigma(\lambda)$.

Since $\tau(\lambda)$ is a continuous and nowhere constant $N$-function,
therefore $\lambda\rightleftarrows\mu=\tau(\lambda)$ is a topological mapping
of the infinite interval $-\infty<\lambda<+\infty$ on the finite interval
$a<\mu<b$, where $a=\tau(-\infty)$, $b=\tau(+\infty)$.  Denote the inverse
mapping by $\mu\rightleftarrows\lambda=\tau^{-1}(\mu)$.

Let us now replace $F(a),f(\lambda)$ by
$\tau(F(a)),f(\tau^{-1}(\mu))$, and then write again $\lambda$ for $\mu$.  This
leaves the structure of our assertions unaffected, but there are these changes
regarding $\lambda$.  The domain of $\lambda$ is now $a<\lambda<b$ and not
$-\infty<\lambda<+\infty$.  In this new domain we have to use the notion of
common (Lebesgue)-measurability, and not that one of
$\tau(\lambda)$-measurability.  (Our mapping
$\lambda\rightleftarrows\tau(\lambda)$ transformed the latter into the
former.)

We now introduce the ``0-space of Baire'', $S_0$, which is the set of all
sequences
$\alpha=[n_1,n_2,\ldots]$, $n_1,n_2,\ldots=1,2,\ldots$.  It is metrized by the
definition
\[
 \begin{aligned}
 &\operatorname{Dist}([m_1,m_2,\ldots],[n_1,n_2,\ldots])\\
 &\qquad=
 \begin{cases}
  0,
    &\begin{aligned}
      &m_\nu=n_\nu\text{ for all}\\[-1mm]
      &\nu=1,2,\ldots,
     \end{aligned}\\[2mm]
  \dfrac1{\nu_0},
    &\begin{aligned}
      &\nu_0\text{ is the smallest }\nu=1,2,\ldots\\[-1mm]
      &\text{for which }m_\nu\ne n_\nu.
     \end{aligned}
 \end{cases}
 \end{aligned}
\]
Thus $S_0$ is a complete and separable metric space.  (The countable set of all
those $\alpha=[n_1,n_2,\ldots]$, for which $n_\nu\ne1$ occurs for a finite
number of $\nu$'s only, is obviously everywhere dense in $S_0$.)

% p. 449
On the other hand the definition
\[
 [m_1,m_2,\ldots]<[n_1,n_2,\ldots]
\quad\begin{cases}
 \text{if a $\nu_0$ exists, such that }m_{\nu_0}<n_{\nu_0},\\
 \text{and }m_\nu=n_\nu\text{ for }\nu<\nu_0
\end{cases}
\]
establishes an ``ordering'' of $S_0$.  If $\alpha$ is fixed, then the set of all
$\beta<\alpha$ is open.  If $\beta<\alpha$, and $\nu_0$ is defined as above,
then we have $\gamma<\alpha$ in the entire sphere
$\operatorname{Dist}(\gamma,\beta)<1/\nu_0$.  Every sphere in $S_0$ is a set of
the form $\alpha_1\leq\beta<\alpha_2$ (with fixed $\alpha_1,\alpha_2$).  If its
center is $[m_1,m_2,\ldots]$ and its radius $\epsilon>0$, then
$[n_1,n_2,\ldots]=\beta$ belongs to it, if and only if,
$\operatorname{Dist}([m_1,m_2,\ldots],[n_1,n_2,\ldots])<\epsilon$, that is
$m_\nu=n_\nu$ for $\nu=1,\ldots,\nu_1$, where $\nu_1$ is the greatest integer
$\leq1/\epsilon$.  This may be written in the form
$\alpha_1\leq\beta<\alpha_2$ if we put
\[
 \alpha_1=[m_1,\ldots,m_{\nu_1-1},m_{\nu_1},1,1,\ldots],
 \quad
 \alpha_2=[m_1,\ldots,m_{\nu_1-1},m_{\nu_1}+1,1,1,\ldots].
\]

Apply now \cite[p.~211]{ref15}.  As $A$ is an analytic set in a complete and
separable metric space, therefore there exists a function $\varphi(\alpha)$,
which is defined and continuous for every $\alpha\in S_0$, and which maps
$S_0$ on $A$.  (The mapping may be many-to-one.)  Then $K$ is clearly the set
of all $F(\varphi(\alpha))$, $\alpha\in S_0$.  As $F(\varphi(\alpha))$ also is
defined and continuous for every $\alpha\in S_0$, and the real numbers also
form a complete and separable metric space, we can therefore now apply
\cite[p.~211]{ref15}, in the opposite sense: $K$ is an analytic set, and
therefore measurable by \cite[pp.~151--152]{ref14}.

Consider now a $\lambda\in K$.  Denote the set of all $\alpha\in S_0$ with
$F(\varphi(\alpha))=\lambda$ by $T(\lambda)$.  As $\lambda\in K$,
$T(\lambda)$ is not empty, and as $F(\varphi(\alpha))$ is continuous,
$T(\lambda)$ is closed.  Denote the smallest $n_1$ with
$[n_1,n_2,n_3,\ldots]\in T(\lambda)$ (for any $n_2,n_3,\ldots$) by $n_1^0$,
the smallest $n_2$ with
$[n_1^0,n_2,n_3,\ldots]\in T(\lambda)$ (for any $n_3,n_4,\ldots$) by $n_2^0$,
the smallest $n_3$ with
$[n_1^0,n_2^0,n_3,\ldots]\in T(\lambda)$ (for any $n_4,n_5,\ldots$) by
$n_3^0$, etc.  Put $\alpha(\lambda)=[n_1^0,n_2^0,\ldots]$.  It is clearly a
condensation point of $T(\lambda)$, and therefore $\alpha(\lambda)\in
T(\lambda)$; besides its definition guarantees
$\alpha(\lambda)\leq\beta$ for every $\beta\in T(\lambda)$.  In other words:
$T(\lambda)$ has a first element, and this is $\alpha(\lambda)$.

Choose now a fixed $\alpha_0$ and consider the set $M(\alpha_0)$ of all
$\lambda$'s with $\alpha(\lambda)<\alpha_0$.  $\lambda\in M(\alpha_0)$, that is
$\alpha(\lambda)<\alpha_0$, is clearly equivalent to the existence of a
$\beta<\alpha_0$ with $\beta\in T(\lambda)$, that is of a $\beta<\alpha_0$ with
$F(\varphi(\beta))=\lambda$.  In other words: $M(\alpha_0)$ is the
$\beta\rightleftarrows F(\varphi(\beta))$-image of the open set
$\beta<\alpha_0$ in $S_0$.  As $F(\varphi(\beta))$ is continuous, we can infer
by \cite[p.~209, Theorem III]{ref15}, that $M(\alpha_0)$ is an analytic set,
and therefore measurable by \cite{ref14}.

If $T$ is any subset of $S_0$, denote the set of all $\lambda$'s with
$\alpha(\lambda)\in T$ by $N(T)$.  If $T$ is a sphere, then it has the form
$\alpha_1\leq\beta<\alpha_2$, therefore $N(T)=M(\alpha_2)-M(\alpha_1)$.  Our
above result proves therefore, that $N(T)$ is measurable.  If $T$ is merely an
open set, then, as $S_0$ is separable, we can write $T$ as the sum of a sequence
of spheres: $T=T_1+T_2+\cdots$.  Now $N(T)=N(T_1)+N(T_2)+\cdots$, and as
$N(T_1),N(T_2),\ldots$ are all measurable, so is $N(T)$.

Let finally $O$ be an open subset of $S$, and denote by $T(O)$ the set of all
$\alpha\in S_0$

% p. 450
with $\varphi(\alpha)\in O$.  As $O$ is open and $\varphi(\alpha)$ is
continuous, therefore $T(O)$ is open too.  Thus $N(T(O))$ is measurable.
$\lambda\in N(T(O))$ means $\alpha(\lambda)\in T(O)$, that is
$\varphi(\alpha(\lambda))\in O$.

Observe finally, that as $\alpha(\lambda)\in T(\lambda)$ (assuming
$\lambda\in K$), therefore $F(\varphi(\alpha(\lambda)))=\lambda$.

The measurability of $K$ being already established, we have proved statement
(i) of our lemma.  And in order to prove statement (ii), it suffices to define
$f(\lambda)=\varphi(\alpha(\lambda))$.  Thus the proof is completed.
\end{proof}

\section{\texorpdfstring{$\sigma(\lambda)$}{sigma(lambda)}-measurable
dependence and the operations on rings; measurable ring relations and
measurable separating projections}
\label{sec:17}

We now return to the discussion of the $\sigma(\lambda)$-measurable rings
$M(\lambda)$:

\begin{lemma}
\label{lem:6}
Let the $M(\lambda)$ be rings containing $1$, depending
$\sigma(\lambda)$-measurably on $\lambda$.  Let $k(\lambda)$,
$\Delta_p$, $p=1,2,\ldots,\infty$, be as in Definition~\ref{def:2}, let
$\varphi_m(\lambda)$ be a measurable family.  Choose a $p=1,2,\ldots,\infty$
and a $p$-dimensional unitary space $\mathfrak H'_p$, and in it a complete
normalized orthogonal system $\psi'_1,\psi'_2,\ldots$.  For each
$\lambda\in\Delta_p$ denote the isomorphism of
$\mathfrak H^{(\lambda)}$ and $\mathfrak H'_p$, in which
$\varphi_m(\lambda)$ corresponds to $\psi'_m$, by $g_p^{(\lambda)}$.

Then there exists a set $N_p\subset\Delta_p$ of $\sigma(\lambda)$-measure $0$,
such that $\Delta_p-N_p$ is a Borel set and that the set of all pairs
$(\lambda,A)$ with $\lambda\in\Delta_p-N_p$,
$(g_p^{(\lambda)})^{-1}Ag_p^{(\lambda)}\in M(\lambda)'$ is a Borel-set in the
direct product space $r\times S_p$.  ($r$ is the set of all real numbers
$\lambda$ in their usual topology, $S_p$ the set of linear operators $A$ in
$\mathfrak H'_p$ with
$\lvert\!\lvert\!\lvert A\rvert\!\rvert\!\rvert\leq1$ in their weak topology,
cf. the beginning of \secref{sec:16}.)
\end{lemma}

\begin{proof}
Choose the $A_i(\lambda)$, $i=1,2,\ldots$, in the sense of
Definition~\ref{def:5}: $A_i(\lambda)$ is $\sigma(\lambda)$-measurable in
$\lambda$, and the $A_i(\lambda)$, $i=1,2,\ldots$, determine $M(\lambda)$.
Each function
$(A_i(\lambda)\varphi_m(\lambda),\varphi_n(\lambda))$, $i,m,n=1,2,\ldots$ (of
course $m,n\leq p$) is $\sigma(\lambda)$-measurable.  Therefore there exists a
set $N_1^{(i,m,n,p)}$ of $\sigma(\lambda)$-measure $0$, such that
$(A_i(\lambda)\varphi_m(\lambda),\varphi_n(\lambda))$ agrees with a
Baire-function for all $\lambda\notin N_1^{(i,m,n,p)}$
(cf. \cite[p.~171]{ref4}).
$N_2^{(p)}=\sum_{i,m,n}N_1^{(i,m,n,p)}$ has clearly the same properties, and is
even independent of $i,m,n$.  As $\Delta_p$ and $N_2^{(p)}$ are both
$\sigma(\lambda)$-measurable, $\Delta_p-\Delta_pN_2^{(p)}$ is
$\sigma(\lambda)$-measurable too, and so a Borel-set
$\Delta'_p\subset\Delta_p-\Delta_pN_2^{(p)}$ exists, for which
$(\Delta_p-\Delta_pN_2^{(p)})-\Delta'_p$ has the
$\sigma(\lambda)$-measure $0$.  Thus
$N_p=\Delta_p-\Delta'_p\subset
((\Delta_p-\Delta_pN_2^{(p)})-\Delta'_p)+N_2^{(p)}$ also has the
$\sigma(\lambda)$-measure $0$, and we have $\Delta'_p=\Delta_p-N_p$.  As
$\Delta'_p$ is a Borel-set, the function
\[
 1_{\Delta'_p}(\lambda)=
 \begin{cases}
  1,&\lambda\in\Delta'_p,\\
  0,&\lambda\notin\Delta'_p
 \end{cases}
\]
is obviously a Baire-function of $\lambda$.

Assume $\lambda\in r$, $A\in S_p$, we wish to find an analytic expression for
$\lambda\in\Delta'_p$,
$(g_p^{(\lambda)})^{-1}Ag_p^{(\lambda)}\in M(\lambda)'$.  The first requirement
means
\begin{equation}
 1_{\Delta'_p}(\lambda)-1=0.
 \label{eq:17.1}
\end{equation}
The second means that $(g_p^{(\lambda)})^{-1}Ag_p^{(\lambda)}$ commute with all
elements of $M(\lambda)$.  In order that this be the case, it suffices for
$(g_p^{(\lambda)})^{-1}Ag_p^{(\lambda)}$ to commute with the
$A_i(\lambda)$, $i=1,2,\ldots$.  In other words:
\[
 (g_p^{(\lambda)})^{-1}Ag_p^{(\lambda)}A_i(\lambda)
 =A_i(\lambda)(g_p^{(\lambda)})^{-1}Ag_p^{(\lambda)}
\]
for $i=1,2,\ldots$, or:
\[
 \bigl((g_p^{(\lambda)})^{-1}Ag_p^{(\lambda)}A_i(\lambda)
       \varphi_m(\lambda),\varphi_n(\lambda)\bigr)
 =\bigl(A_i(\lambda)(g_p^{(\lambda)})^{-1}Ag_p^{(\lambda)}
       \varphi_m(\lambda),\varphi_n(\lambda)\bigr)
\]
for $i,m,n=1,2,\ldots$.

% p. 451
We have however:
\[
 \begin{aligned}
 &\bigl((g_p^{(\lambda)})^{-1}Ag_p^{(\lambda)}A_i(\lambda)
       \varphi_m(\lambda),\varphi_n(\lambda)\bigr)\\
 &\quad=(A_i(\lambda)\varphi_m(\lambda),
       (g_p^{(\lambda)})^{-1}A^*g_p^{(\lambda)}\varphi_n(\lambda))\\
 &\quad=\sum_{l=1}^p(A_i(\lambda)\varphi_m(\lambda),\varphi_l(\lambda))
       (\varphi_l(\lambda),(g_p^{(\lambda)})^{-1}
       A^*g_p^{(\lambda)}\varphi_n(\lambda))\\
 &\quad=\sum_{l=1}^p(A_i(\lambda)\varphi_m(\lambda),\varphi_l(\lambda))
       (Ag_p^{(\lambda)}\varphi_l(\lambda),
        g_p^{(\lambda)}\varphi_n(\lambda))\\
 &\quad=\sum_{l=1}^p(A_i(\lambda)\varphi_m(\lambda),\varphi_l(\lambda))
       (A\psi'_l,\psi'_n),
 \end{aligned}
\]
and similarly:
\[
 \begin{aligned}
 &\bigl(A_i(\lambda)(g_p^{(\lambda)})^{-1}Ag_p^{(\lambda)}
       \varphi_m(\lambda),\varphi_n(\lambda)\bigr)\\
 &\quad=\sum_{l=1}^p(A\psi'_m,\psi'_l)
       (A_i(\lambda)\varphi_l(\lambda),\varphi_n(\lambda)).
 \end{aligned}
\]
So we obtain the following conditions:
\begin{equation}
 \begin{aligned}
 &\sum_{l=1}^p
   (A_i(\lambda)\varphi_m(\lambda),\varphi_l(\lambda))
   (A\psi'_l,\psi'_n)\\
 &\qquad-
   \sum_{l=1}^p(A\psi'_m,\psi'_l)
   (A_i(\lambda)\varphi_l(\lambda),\varphi_n(\lambda))=0,\\
 &\hspace{8em}\text{for }i,m,n=1,2,\ldots.
 \end{aligned}
 \label{eq:17.2}
\end{equation}
The left side of \eqref{eq:17.1} is clearly a Baire-function in $\lambda$, and
independent of $A$.  In the left side of \eqref{eq:17.2} the factors
$(A_i(\lambda)\varphi_m(\lambda),\varphi_l(\lambda))$ are Baire-functions in
$\lambda$, and independent of $A$; while the factors $(A\psi'_m,\psi'_l)$ are
(weakly) continuous in $A$, and independent of $\lambda$.  Thus all our
equations have Baire-functions in $\lambda,A$ as left sides.  Therefore each of
them has a Borel-set of pairs $(\lambda,A)$ in $r\times S_p$ as its domain of
validity (cf. \cite[p.~260]{ref15}).  As they are countably infinite in number,
the intersection of their domains of validity is a Borel-set too.  But this is
the set of all pairs $(\lambda,A)$ with $\lambda\in\Delta'_p$,
$(g_p^{(\lambda)})^{-1}Ag_p^{(\lambda)}\in M(\lambda)'$.

Thus the statement of our lemma holds for our present choice of $N_p$.
\end{proof}

\begin{lemma}
\label{lem:7}
Let the $M(\lambda)$ be rings containing $1$, depending
$\sigma(\lambda)$-measurably on $\lambda$.  Then the $M'(\lambda)$ too depend
$\sigma(\lambda)$-measurably on $\lambda$.
\end{lemma}

\begin{proof}
Introduce $k(\lambda)$, $p=1,2,\ldots,\infty$, $\Delta_p$,
$\mathfrak H'_p$, $g_p^{(\lambda)}$, $S_p$, as in the \hyperref[lem:6]{previous lemma}, and form
the corresponding sets $N_p$.  Let $O_{p1},O_{p2},\ldots$ be an equivalent set
of neighborhoods in the separable space $S_p$.  Denote the set of all pairs
$(\lambda,A)$ with $\lambda\in r$, $A\in O_{pj}$ by $O^*_{pj}$.  As $O_{pj}$
is open, therefore $O^*_{pj}$ is open too.  Let $M_p^*$ be the set of all pairs
$(\lambda,A)$ with $\lambda\in\Delta_p-N_p$,
$(g_p^{(\lambda)})^{-1}Ag_p^{(\lambda)}\in M(\lambda)'$; we know from
Lemma~\ref{lem:6}, that $M_p^*$ is a Borel-set.  Thus $O^*_{pj}M_p^*$ is a
Borel-set too.

As $O^*_{pj}M_p^*$ is a Borel-set in $r\times S_p$, it is \emph{a fortiori} an
analytic set.  Therefore Lemma~\ref{lem:5} applies to it: Define
$F(\lambda,A)=\lambda$, which is obviously a continuous

% p. 452
function.  We denote the set $K$ and the function $f(\lambda)$, as defined
there, by $K_{pj}$, $f_{pj}(\lambda)$ (because their dependence on $p,j$ is now
essential).  Clearly $K_{pj}\subset\Delta_p-N_p$.  $f_{pj}(\lambda)$ is a pair
$(\mu,A)$, as $F(f_{pj}(\lambda))=\lambda$, there is $\mu=\lambda$.  Thus we can
put: $f_{pj}(\lambda)=(\lambda,B^0_{pj}(\lambda))$.

Let $U$ be an open set in the complex-number-plane.  The set of all $A\in S_p$
with $(A\psi'_m,\psi'_n)\in U$ is obviously open, and thus the set of all
$(\lambda,A)\in r\times S_p$ with $(A\psi'_m,\psi'_n)\in U$ is open too.
Therefore the set of all those $\lambda\in K_{pj}$, for which
$f_{pj}(\lambda)=(\lambda,B^0_{pj}(\lambda))$ belongs to this set, is
$\sigma(\lambda)$-measurable by Lemma~\ref{lem:5}.  This set is characterized
by $(B^0_{pj}(\lambda)\psi'_m,\psi'_n)\in U$.  As this is true for all open
$U$'s, it means that the function
$(B^0_{pj}(\lambda)\psi'_m,\psi'_n)$ is $\sigma(\lambda)$-measurable.  As
\[
 \begin{aligned}
 (B^0_{pj}(\lambda)\psi'_m,\psi'_n)
 &=\bigl(B^0_{pj}(\lambda)g_p^{(\lambda)}\varphi_m(\lambda),
          g_p^{(\lambda)}\varphi_n(\lambda)\bigr)\\
 &=\bigl((g_p^{(\lambda)})^{-1}B^0_{pj}(\lambda)g_p^{(\lambda)}
          \varphi_m(\lambda),\varphi_n(\lambda)\bigr),
 \end{aligned}
\]
we can apply \hyperref[rem:15.2]{Remark 2} after Definition~\ref{def:5}: The operator
$(g_p^{(\lambda)})^{-1}B^0_{pj}(\lambda)g_p^{(\lambda)}$ (in
$\mathfrak H^{(\lambda)}$) is $\sigma(\lambda)$-measurable.

Now define
\[
 C_{pj}(\lambda)=
 \begin{cases}
  (g_p^{(\lambda)})^{-1}B^0_{pj}(\lambda)g_p^{(\lambda)},
       &\lambda\in K_{pj},\\
  0,&\lambda\notin K_{pj}.
 \end{cases}
\]
as $K_{pj}$ is $\sigma(\lambda)$-measurable, the operator $C_{pj}(\lambda)$
(in $\mathfrak H^{(\lambda)}$) is $\sigma(\lambda)$-measurable, too.

$\lambda\in K_{pj}$ means that an $a\in O^*_{pj}M_p^*$ with $F(a)=\lambda$
exists, that is:
\[
 a=(\lambda,A),\quad \lambda\in\Delta_p-N_p,\quad A\in O_{pj},\quad
 (g_p^{(\lambda)})^{-1}Ag_p^{(\lambda)}\in M(\lambda)',
\]
in other words: $\lambda\in\Delta_p-N_p$, $A\in O_{pj}$, and $M(\lambda)'$
possesses elements in $(g_p^{(\lambda)})^{-1}O_{pj}g_p^{(\lambda)}$.  If this
is the case, then $B^0_{pj}(\lambda)$ is defined, and
$C_{pj}(\lambda)=(g_p^{(\lambda)})^{-1}B^0_{pj}(\lambda)g_p^{(\lambda)}$.
Thus $f(\lambda)=(\lambda,B^0_{pj}(\lambda))
=(\lambda,g_p^{(\lambda)}C_{pj}(\lambda)(g_p^{(\lambda)})^{-1})
\in O^*_{pj}M_p^*$, that is:
$g_p^{(\lambda)}C_{pj}(\lambda)(g_p^{(\lambda)})^{-1}\in O_{pj}$,
$C_{pj}(\lambda)\in M(\lambda)'$, in other words: $C_{pj}(\lambda)$ is an
element of $M(\lambda)'$ in
$(g_p^{(\lambda)})^{-1}O_{pj}g_p^{(\lambda)}$.  If this is not the case, then
$C_{pj}(\lambda)=0$.  Assuming $\lambda\in\Delta_p-N_p$ we have furthermore
$C_{qj}(\lambda)=0$ for all $q\ne p$, as there cannot be a
$\lambda\in\Delta_q-N_q$.

Assume now
$\lambda\in\sum_{p=1,2,\ldots,\infty}(\Delta_p-N_p)
=r-\sum_{p=1,2,\ldots,\infty}N_p$.  Then we have $C_{qj}(\lambda)=0$, except if
$q=p$ and $M(\lambda)'$ possesses elements in
$(g_p^{(\lambda)})^{-1}O_{pj}g_p^{(\lambda)}$, in which case
$C_{qj}(\lambda)$ itself is such an element.  Remembering that the $O_{pj}$,
$j=1,2,\ldots$, form an equivalent set of neighborhoods for the $A$ with
$\lvert\!\lvert\!\lvert A\rvert\!\rvert\!\rvert\leq1$ in $\mathfrak H'_p$,
and thus the $(g_p^{(\lambda)})^{-1}O_{pj}g_p^{(\lambda)}$, $j=1,2,\ldots$, do
the same in $\mathfrak H^{(\lambda)}$, this proves: The $C_{qj}(\lambda)$,
$q,j=1,2,\ldots$, are everywhere dense in that part of $M(\lambda)'$ which
lies in $\lvert\!\lvert\!\lvert A\rvert\!\rvert\!\rvert\leq1$.

The ring generated by the $C_{qj}(\lambda)$ is contained in $M(\lambda)'$,
because the $C_{qj}(\lambda)$ are.  But as it contains all (weak) condensation
points of the $C_{qj}(\lambda)$, it must contain all $A\in M(\lambda)'$ with
$\lvert\!\lvert\!\lvert A\rvert\!\rvert\!\rvert\leq1$.  If
$A\in M(\lambda)'$, $\lvert\!\lvert\!\lvert A\rvert\!\rvert\!\rvert>1$, then
this ring will contain
$1/\lvert\!\lvert\!\lvert A\rvert\!\rvert\!\rvert\cdot A$ and thus $A$ too.
Thus it coincides with $M(\lambda)'$.  In other words: The
$C_{qj}(\lambda)$, $q,j=1,2,\ldots$, generate the ring $M(\lambda)'$.

Write $N=\sum_{p=1,2,\ldots,\infty}N_p$, and put the $C_{qj}(\lambda)$,
$q,j=1,2,\ldots$, in some way in a simple sequence:
$D_i(\lambda)$, $i=1,2,\ldots$.  Then we have proved: $N$ has the
$\sigma(\lambda)$-measure $0$, each $D_i(\lambda)$ (in
$\mathfrak H^{(\lambda)}$) is $\sigma(\lambda)$-measurable, if
$\lambda\notin N$ then the $D_i(\lambda)$,

% p. 453
$i=1,2,\ldots$, generate the ring $M(\lambda)'$.  Now change the
$D_i(\lambda)$, $i=1,2,\ldots$, for each $\lambda\in N$ in an arbitrary way
into a sequence which generates the ring $M(\lambda)'$.  (For instance again
into a [weakly] everywhere dense set in the part of $M(\lambda)'$ in
$\lvert\!\lvert\!\lvert A\rvert\!\rvert\!\rvert\leq1$ [which is separable].)
As the $\sigma(\lambda)$-measure of $N$ is $0$, this does not affect
$\sigma(\lambda)$-measurability of the $D_i(\lambda)$; but now the
$D_i(\lambda)$, $i=1,2,\ldots$, generate the ring $M(\lambda)'$ for each
$\lambda$.  Thus $M(\lambda)'$ is $\sigma(\lambda)$-measurable in the sense of
Definition~\ref{def:5}, and the proof is completed.
\end{proof}

\begin{lemma}
\label{lem:8}
Let the $M_s(\lambda)$, $s=1,2,\ldots$, (the sequence of the $s$ may be finite
or infinite) be rings containing $1$, depending for each $s$
$\sigma(\lambda)$-measurably on $\lambda$.  Let
$N(\lambda)=M_1(\lambda)M_2(\lambda)\cdots$ be their intersection,
$P(\lambda)=R(M_1(\lambda),M_2(\lambda),\ldots)$ the ring generated by them,
both rings containing $1$.  Then $N(\lambda),P(\lambda)$ depend
$\sigma(\lambda)$-measurably on $\lambda$ too.
\end{lemma}

\begin{proof}
Apply Definition~\ref{def:5} to $M_s(\lambda)$: Let
$A_{si}(\lambda)$, $i=1,2,\ldots$, be a sequence of
$\sigma(\lambda)$-measurable operators, such that $A_{si}(\lambda)$,
$i=1,2,\ldots$, generate the ring $M_s(\lambda)$.  Then all
$A_{si}(\lambda)$, $s,i=1,2,\ldots$, generate the ring
$P(\lambda)=R(M_1(\lambda),M_2(\lambda),\ldots)$.  Thus if we write the
$A_{si}(\lambda)$, $s,i=1,2,\ldots$, as a simple sequence:
$B_j(\lambda)$, $j=1,2,\ldots$, then the $B_j(\lambda)$ are
$\sigma(\lambda)$-measurable, and the $B_j(\lambda)$, $j=1,2,\ldots$,
generate the ring $P(\lambda)$.  So $P(\lambda)$ is
$\sigma(\lambda)$-measurable by Definition~\ref{def:5}.

This, together with Lemma~\ref{lem:7}, gives now our result for $N(\lambda)$
too, because
\[
 N(\lambda)=M_1(\lambda)M_2(\lambda)\cdots
 =\bigl(R(M_1(\lambda)',M_2(\lambda)',\ldots)\bigr)'.
\]
\footnote{The general relation
$M_1M_2\cdots=(R(M'_1,M'_2,\ldots))'$ is proved as follows: The operators which
commute with a given $A$ form a ring, so if they include
$M_1,M_2,\ldots$ they must also include $R(M_1,M_2,\ldots)$.  The converse is
obvious, as $M_1,M_2,\ldots$ are subsets of $R(M_1,M_2,\ldots)$.  Thus $A$
commutes with all of $M_1,M_2,\ldots$ if, and only if, it commutes with all of
$R(M_1,M_2,\ldots)$.  In other words:
$M'_1M'_2\cdots=(R(M_1,M_2,\ldots))'$.  If we now replace
$M_1,M_2,\ldots$ by $M'_1,M'_2,\ldots$ and apply $'$ to the resulting equation
then $(M''_1M''_2\cdots)'=(R(M'_1,M'_2,\ldots))''$ obtains.  Next $N''=N$ for
every ring $N$ containing $1$ by \cite[p.~372]{ref2}, and so the above equation
becomes $(M_1M_2\cdots)'=R(M'_1,M'_2,\ldots)$.  Thus $M'$ is an involutory
operation, which dualizes the operations $M_1M_2\cdots$ and
$R(M_1,M_2,\ldots)$.\label{fn:12}}
\end{proof}

\begin{lemma}
\label{lem:9}
Let $M(\lambda),N(\lambda)$ be rings containing $1$, both depending
$\sigma(\lambda)$-measurably on $\lambda$.  Assume $M(\lambda)\subset N(\lambda)$
always.  Let $\mathcal M$ be the set of all $\lambda$'s for which
$M(\lambda)\ne N(\lambda)$.

Under these assumptions we have:
\begin{enumerate}
\renewcommand{\labelenumi}{(\roman{enumi})}
\item $\mathcal M$ is a $\sigma(\lambda)$-measurable set.
\item There exists a $\sigma(\lambda)$-measurable $E(\lambda)$, which is a
projection for each $\lambda$, and for which $\lambda\in\mathcal M$ implies
$E(\lambda)\notin M(\lambda)$, $E(\lambda)\in N(\lambda)$.
\end{enumerate}
\end{lemma}

\begin{proof}
Apply Lemma~\ref{lem:6} to $M(\lambda)'$ and to $N(\lambda)'$ (in place of the
$M(\lambda)$ of that lemma).  Introduce $k(\lambda)$, $\Delta_p$,
$p=1,2,\ldots,\infty$, as there, and the sets of $\sigma(\lambda)$-measure $0$,
$\widetilde N_p$ and $\overline{\widetilde N}_p$ respectively.  Now
$\Delta_p-\widetilde N_p-\overline{\widetilde N}_p
=(\Delta_p-\widetilde N_p)(\Delta_p-\overline{\widetilde N}_p)$ is a Borel-set
in $r$.  Forming $r\times S_p$ again, the set $M_p^*$ of all $(\lambda,A)$ with
$\lambda\in\Delta_p-\widetilde N_p$,
$(g_p^{(\lambda)})^{-1}Ag_p^{(\lambda)}\in M(\lambda)''=M(\lambda)$ is a
Borel-set, and so is the set $M_p^{**}$ of all $(\lambda,A)$ with
$\lambda\in\Delta_p-\overline{\widetilde N}_p$,
$(g_p^{(\lambda)})^{-1}Ag_p^{(\lambda)}\in N_p(\lambda)''=N(\lambda)$.  The set
$N_p$ of all $(\lambda,A)$ with
$\lambda\in\Delta_p-\widetilde N_p-\overline{\widetilde N}_p$ is clearly a
Borel-set too.

% p. 454
The projections $E$ in $\mathfrak H'_p$ form a Borel-set in $S_p$.  They are
characterized by $E^*=E$, $E^2=E$, that is:
$(E^*\psi'_m,\psi'_n)=(E\psi'_m,\psi'_n)$,
$(E^2\psi'_m,\psi'_n)=(E\psi'_m,\psi'_n)$, that is (cf. the analogous
computations in the proof of Lemma~\ref{lem:6}):
\[
 \overline{(E\psi'_n,\psi'_m)}-(E\psi'_m,\psi'_n)=0,
\]
\[
 \sum_{l=1}^p(E\psi'_m,\psi'_l)(E\psi'_l,\psi'_n)
 -(E\psi'_m,\psi'_n)=0.
\]
As the left sides are Baire-functions of $E$ (in the weak topology of $S_p$,
cf. the considerations loc. cit. above), we conclude as loc. cit., that they
determine a Borel-set in $S_p$.  Now we can infer, that the set $\mathfrak E_p$
of all $(\lambda,E)$ where $E$ is a projection, is a Borel-set too (in
$r\times S_p$).

Thus the set
$D_p=N_p\cap(M_p^{**}\setminus M_p^*)\cap\mathfrak E_p$ is a Borel-set too (in
$r\times S_p$), its elements are all pairs $(\lambda,E)$ with
$\lambda\in\Delta_p-\widetilde N_p-\overline{\widetilde N}_p$,
$E\notin M(\lambda)$, $E\in N(\lambda)$, $E$ a projection.

As $D_p$ is a Borel-set in $r\times S_p$ it is \emph{a fortiori} an analytic
set.  Therefore Lemma~\ref{lem:5} applies to it: Define $F(\lambda,A)=\lambda$,
which is obviously a continuous function.  We denote the set $K$ and the
function $f(\lambda)$, as defined there, by $K_p,f_p(\lambda)$ (because their
dependence on $p$ is now essential).  Clearly
$K_p\subset\Delta_p-\widetilde N_p-\overline{\widetilde N}_p$.
$f_p(\lambda)$ is a pair $(\mu,A)$, as $F(f_p(\lambda))=\lambda$, there is
$\mu=\lambda$.  Thus we can put: $f_p(\lambda)=(\lambda,E_p^0(\lambda))$.

Literally as in the corresponding part of the proof of Lemma~\ref{lem:7}, we
prove that $E_p^0(\lambda)$ (which is defined in the
$\sigma(\lambda)$-measurable set $K_p$) is $\sigma(\lambda)$-measurable.

Now define
\[
 E(\lambda)=
 \begin{cases}
  E_p^0(\lambda),&\lambda\in K_p,\quad p=1,2,\ldots,\infty,\\
  0,&\text{otherwise}.
 \end{cases}
\]
(The $K_p\subset\Delta_p$ are mutually exclusive.)

$\lambda\in K_p$ means that an $a\in D_p$ with $F(a)=\lambda$ exists, that is:
\[
 a=(\lambda,E),\quad
 \lambda\in\Delta_p-\widetilde N_p-\overline{\widetilde N}_p,
 \quad E\notin M(\lambda),\quad E\in N(\lambda),
\]
$E$ a projection.  In other words, $\lambda\notin K_p$ means this: Either
$\lambda\in\widetilde N_p+\overline{\widetilde N}_p$, or
$\lambda\in\Delta_p-\widetilde N_p-\overline{\widetilde N}_p$, and every
projection $E\in N(\lambda)$ belongs to $M(\lambda)$ too.  Using the
terminology of \cite[p.~392]{ref2}, this latter statement means:
$N(\lambda)^P\subset M(\lambda)^P$.  By \cite[p.~393]{ref2}, this is equivalent
to $N(\lambda)\subset M(\lambda)$, and as $M(\lambda)\subset N(\lambda)$, it is
further equivalent to $M(\lambda)=N(\lambda)$.  Thus $K_p$ is the set of all
$\lambda\in\Delta_p-\widetilde N_p-\overline{\widetilde N}_p$ with
$M(\lambda)\ne N(\lambda)$.

$\lambda\in K_p$ implies $E(\lambda)=E_p^0(\lambda)$, and this is
$\notin M(\lambda)$, $\in N(\lambda)$, and a projection.  If
$\lambda\notin K_p$ for every $p=1,2,\ldots,\infty$, then $E(\lambda)=0$ is a
projection again.  Thus $E(\lambda)$ is always a projection.

As the $E_p^0(\lambda)$ are all $\sigma(\lambda)$-measurable, and the sets
$K_p$ are $\sigma(\lambda)$-measurable too, therefore $E(\lambda)$ is also
$\sigma(\lambda)$-measurable.

Put
$N=\widetilde N_1+\overline{\widetilde N}_1+\widetilde N_2
+\overline{\widetilde N}_2+\cdots+\widetilde N_\infty
+\overline{\widetilde N}_\infty$.  $N$ has the $\sigma(\lambda)$-measure $0$,
and
\[
 (\Delta_1-\widetilde N_1-\overline{\widetilde N}_1)
 +(\Delta_2-\widetilde N_2-\overline{\widetilde N}_2)+\cdots
 +(\Delta_\infty-\widetilde N_\infty-\overline{\widetilde N}_\infty)=r-N.
\]

% p. 455
Thus if $\lambda\notin N$, then $\lambda$ belongs to some
$\Delta_p-\widetilde N_p-\overline{\widetilde N}_p$,
$p=1,2,\ldots,\infty$, and thus $M(\lambda)\ne N(\lambda)$ implies
$\lambda\in K_p$, for some $p=1,2,\ldots,\infty$, and with it
$E(\lambda)\notin M(\lambda)$, $E(\lambda)\in N(\lambda)$.

Now change the $E(\lambda)$ for each $\lambda\in N$ with
$M(\lambda)\ne N(\lambda)$ in an arbitrary way in a projection
$E\notin M(\lambda)$, $E\in N(\lambda)$ (as to the existence of such an $E$,
cf. above).  As the $\sigma(\lambda)$-measure of $N$ is $0$, this does not
affect the $\sigma(\lambda)$-measurability of $E(\lambda)$; but now
$E(\lambda)\notin M(\lambda)$, $E(\lambda)\in N(\lambda)$ holds whenever
$M(\lambda)\ne N(\lambda)$.

The set $\mathcal M$ of the $\lambda$'s, for which
$M(\lambda)\ne N(\lambda)$, differs from the $\sigma(\lambda)$-measurable set
$K_1+K_2+\cdots$ only by a subset of $N$ (which therefore has the
$\sigma(\lambda)$-measure $0$).  Thus $\mathcal M$ is
$\sigma(\lambda)$-measurable.  This completes the proof.
\end{proof}

\medskip\noindent\textbf{Remark.}\phantomsection\label{rem:17}
As the $E(\lambda)$'s are projections,
$\lvert\!\lvert\!\lvert E(\lambda)\rvert\!\rvert\!\rvert\leq1$.  Thus they
fulfill condition (ii) in Theorem~\ref{thm:V}, while condition (i) in that
theorem follows from their $\sigma(\lambda)$-measurability
(cf. \hyperref[rem:15.2]{Remark 2} after Definition~\ref{def:5}).  Thus an operator
$E\sim\sum E(\lambda)$ in $\mathfrak H$ exists.  As all $E(\lambda)$'s are
projections, it follows from the results of \secref{sec:14}, that $E$ is a
projection too.

\begin{lemma}
\label{lem:10}
Let $M(\lambda),N(\lambda)$ be rings containing $1$, both depending
$\sigma(\lambda)$-measurably on $\lambda$.  Denote the set of all $\lambda$'s
for which $M(\lambda),N(\lambda)$ are respectively in the relation
$=,\ne,\subset,\supset,\varsubsetneqq,\varsupsetneqq$, respectively by
$K_1,\ldots,K_6$.  All these sets are $\sigma(\lambda)$-measurable.
\end{lemma}

\begin{proof}
$M(\lambda)N(\lambda)$ too depends $\sigma(\lambda)$-measurably on $\lambda$
by Lemma~\ref{lem:8}.  Thus we can apply Lemma~\ref{lem:9} to the rings
$M(\lambda)N(\lambda)$ and $N(\lambda)$ (as we have always
$M(\lambda)N(\lambda)\subset N(\lambda)$).  So we see that the set of all
$\lambda$'s with $M(\lambda)N(\lambda)=N(\lambda)$ is a
$\sigma(\lambda)$-measurable.  But this equation means
$M(\lambda)\subset N(\lambda)$, and so the set in question is $K_3$.

Interchanging $M(\lambda),N(\lambda)$ proves the
$\sigma(\lambda)$-measurability of $K_4$.  Now $K_1=K_3K_4$ too must be
$\sigma(\lambda)$-measurable, and similarly
$K_2=r-K_1$, $K_5=K_3-K_1$, $K_6=K_4-K_1$.
\end{proof}

There exists (in $\mathfrak H$) a smallest and a greatest ring containing $1$:
The set of all operators of the form $\alpha1$, resp. the set of all bounded
linear operators.  We will denote them respectively by $0$ and $1$.  In
$\mathfrak H^{(\lambda)}$ (instead of $\mathfrak H$) we will denote them by
$0(\lambda)$ and $1(\lambda)$ respectively.

Put $A_i(\lambda)=1$ for every $i=1,2,\ldots$, and every $\lambda$ with
$-\infty<\lambda<+\infty$.  These are clearly $\sigma(\lambda)$-measurable,
and so the rings $0(\lambda)$, which are generated by the $A_i(\lambda)$,
$i=1,2,\ldots$, are $\sigma(\lambda)$-measurable too.  By Lemma~\ref{lem:7}
the rings $1(\lambda)=0(\lambda)'$ too are $\sigma(\lambda)$-measurable.  So we
see:

\emph{The rings $0(\lambda)$ and $1(\lambda)$ depend both
$\sigma(\lambda)$-measurably on $\lambda$.}

\section{Definition and discussion of the relation
\texorpdfstring{$M\approx\sum M(\lambda)$}{M as a sum of M(lambda)}}
\label{sec:18}

The information obtained in the last section about rings $M(\lambda)$ which
depend $\sigma(\lambda)$-measurably on $\lambda$, permit us to establish and to
discuss the following definition:

\begin{definition}
\label{def:6}
Let $M$ be a ring in $\mathfrak H$ and $M(\lambda)$ a ring containing $1$ in
$\mathfrak H^{(\lambda)}$ for each $\lambda>-\infty$, $<+\infty$, depending
$\sigma(\lambda)$-measurably on $\lambda$.  We write
$M\approx\sum M(\lambda)$, if $A\in M$ is equivalent to
$A\sim\sum A(\lambda)$, $A(\lambda)\in M(\lambda)$.
\end{definition}

\origpage{456}
If the $M(\lambda)$ are given, it is clear that there cannot be
$M\approx\sum M(\lambda)$ for two different $M$'s; so there exists exactly
one such $M$, or none at all.

We prove:

\par\medskip\noindent\textbf{Lemma 11.}\phantomsection\label{lem:11}\quad
\begin{itshape}
Let $M(\lambda)$ depend $\sigma(\lambda)$-measurably on $\lambda$.
Then the operators $A_i(\lambda)$, $i=1,2,\ldots$, from \hyperref[def:5]{Definition~5} can be so
chosen that all $\lVert\!\lvert A_i(\lambda)\rvert\!\rVert\leq 1$. In this
case operators $A_i$ (in $\mathfrak H$) exist, so that
$A_i\sim\sum A_i(\lambda)$.
\end{itshape}

\par\smallskip\noindent\textit{Proof.}
By \hyperref[rem:15.2]{Remark~2} after \hyperref[def:5]{Definition~5} the numerical functions
$\lVert\!\lvert A_i(\lambda)\rvert\!\rVert$ are
$\sigma(\lambda)$-measurable along with $A_i(\lambda)$. Thus the operators
\[
A_i^0(\lambda)=
\begin{cases}
\dfrac{1}{\lVert\!\lvert A_i(\lambda)\rvert\!\rVert}\,A_i(\lambda),
  &\text{for }\lVert\!\lvert A_i(\lambda)\rvert\!\rVert\ne0,\\[6pt]
0,&\text{for }\lVert\!\lvert A_i(\lambda)\rvert\!\rVert=0
   \quad(\text{that is }A_i(\lambda)=0)
\end{cases}
\]
are $\sigma(\lambda)$-measurable too. Clearly the $A_i^0(\lambda)$,
$i=1,2,\ldots$, generate the same ring as the $A_i(\lambda)$,
$i=1,2,\ldots$, that is $M(\lambda)$; and
$\lVert\!\lvert A_i^0(\lambda)\rvert\!\rVert\leq1$. This proves the first
statement.

As to the second statement, observe that the $\sigma(\lambda)$-measurability
of $A_i(\lambda)$ means that it fulfills condition (i) in \hyperref[thm:V]{Theorem~V}, while
condition (ii) is satisfied owing to
$\lVert\!\lvert A_i(\lambda)\rvert\!\rVert\leq1$. Thus an
$A_i\sim\sum A_i(\lambda)$ exists.

\par\medskip\noindent\textbf{Lemma 12.}\phantomsection\label{lem:12}\quad
\begin{itshape}
Let $M(\lambda)$ depend $\sigma(\lambda)$-measurably on $\lambda$,
and form the $A_i\sim\sum A_i(\lambda)$, $i=1,2,\ldots$, of
\hyperref[lem:11]{Lemma 11}. Then the ring generated by
$P$\footnote{$P$ is the ring discussed in \secref{sec:11} and
\secref{sec:12}.} and by the $A_i$, $i=1,2,\ldots$, is
$M\approx\sum M(\lambda)$.
\end{itshape}

\par\smallskip\noindent\textit{Proof.}
Let $\widetilde M$ be the ring generated by $P$ and by the $A_i$,
$i=1,2,\ldots$. $P\subset P'$, and by \hyperref[thm:V]{Theorem~V} all $A_i\in P'$, so
$\widetilde M\subset P'$. Clearly $\widetilde M\supset P$, so
$P\subset\widetilde M\subset P'$. This implies
$P=P''\subset\widetilde M'\subset P'$\footnote{$N\subset L$ clearly implies
$L'\subset N'$; $N=N''$ by \cite[p.~388]{ref2}.}, so
$\widetilde M,\widetilde M'$ are both $\supset P$ and $\subset P'$. As the
space of all operators $A$ (in $\mathfrak H$) with
$\lVert\!\lvert A\rvert\!\rVert\leq1$ is separable
(cf. \cite[p.~386]{ref2}), therefore the part of $\widetilde M'$ in it is
separable too. Thus there exists a sequence of operators $B_j$,
$j=1,2,\ldots$, which is dense in it. As all $B_j\in\widetilde M'$, therefore
$R(B_1,B_2,\ldots)\subset\widetilde M'$; and by our construction
$R(B_1,B_2,\ldots)$ contains all $A\in\widetilde M'$,
$\lVert\!\lvert A\rvert\!\rVert\leq1$, thus all $A\in\widetilde M'$, that is
$R(B_1,B_2,\ldots)\supset\widetilde M'$. So
$R(B_1,B_2,\ldots)=\widetilde M'$, that is the $B_j$, $j=1,2,\ldots$,
determine the ring $\widetilde M'$. As $\widetilde M'\subset P'$, all
$B_j\in P'$, thus by \hyperref[thm:V]{Theorem~V}, $B_j\sim\sum B_j(\lambda)$.

As $1\in\widetilde M$, we have $\widetilde M''=\widetilde M$
(cf.\footnotemark[14]), so that $A\in\widetilde M$ means that $A$ commutes with
every element of $\widetilde M'=R(B_1,B_2,\ldots)$, or, which is equivalent,
with all $B_j,B_j^*$, $j=1,2,\ldots$. As $\widetilde M\subset P'$ we have
$A\sim\sum A(\lambda)$. Thus our condition is that
$A\sim\sum A(\lambda)$ must commute with all
$B_j\sim\sum B_j(\lambda)$, $B_j^*\sim\sum B_j(\lambda)^*$. By
\secref{sec:14}, this means that $A(\lambda)$ commutes with every
$B_j(\lambda)$ and $B_j(\lambda)^*$, except for a set of
$\sigma(\lambda)$-measure $0$. This set depends on $j$ and the choice of
$B_j(\lambda)$ or $B_j(\lambda)^*$, but by summing all these (countably many)
sets, we obtain a fixed set of exceptions which has
$\sigma(\lambda)$-measure $0$. On this $\lambda$-set we may replace
$A(\lambda)$ by $0$, so we see: The condition is
$A\sim\sum A(\lambda)$, where $A(\lambda)$ commutes with all
$B_j(\lambda),B_j(\lambda)^*$.

\origpage{457}
Denote the ring generated by $1,B_j(\lambda),B_j(\lambda)^*$ (in
$\mathfrak H^{(\lambda)}$) by $N(\lambda)$, then our condition is:
$A(\lambda)\in N(\lambda)'$. The $N(\lambda)$ are by definition
$\sigma(\lambda)$-measurable, and so by \hyperref[lem:7]{Lemma 7}
$N(\lambda)'$ are $\sigma(\lambda)$-measurable too. Our above result states,
remembering \hyperref[def:6]{Definition~6}, that
$\widetilde M\approx\sum N(\lambda)'$.

As $A_i\sim\sum A_i(\lambda)$, $A_i\in\widetilde M$, therefore
$A_i(\lambda)\in N(\lambda)'$, except for a set $N_1^i$ of
$\sigma(\lambda)$-measure $0$. $N_2=\sum_{i=1}^{\infty}N_1^i$ does the same,
and is independent of $i$. So for $\lambda\notin N_2$ all
$A_i(\lambda)\in N(\lambda)'$, and therefore
$M(\lambda)\subset N(\lambda)'$. As the behavior of $N(\lambda)$ on the set
of $\sigma(\lambda)$-measure $0$, $N_2$, does not matter, we may replace it
there by the set of all $\alpha1$, then $N(\lambda)'$ contains all operators. Thus
we can assume $M(\lambda)\subset N(\lambda)'$ for all $\lambda$'s. As
$M(\lambda),N(\lambda)'$ are both $\sigma(\lambda)$-measurable,
\hyperref[lem:9]{Lemma 9} applies. The set $\mathcal M$ of all $\lambda$ with
$M(\lambda)\ne N(\lambda)'$ is $\sigma(\lambda)$-measurable, and there exists
a $\sigma(\lambda)$-measurable $E(\lambda)$, which is a projection for every
$\lambda$, so that $\lambda\in\mathcal M$ implies
$E(\lambda)\notin M(\lambda)$, $E(\lambda)\in N(\lambda)'$. Replacing
$E(\lambda)$ by $0$ for $\lambda\notin\mathcal M$, which does not affect its
$\sigma(\lambda)$-measurability as $\mathcal M$ is
$\sigma(\lambda)$-measurable, we obtain furthermore
$E(\lambda)\in N(\lambda)'$ for all $\lambda$'s.

$R(E(\lambda))$ is $\sigma(\lambda)$-measurable and so by
\hyperref[lem:7]{Lemmas 7}, \hyperref[lem:8]{8},
$R(M(\lambda),E(\lambda))$, $R(M(\lambda),E(\lambda))'$ are
$\sigma(\lambda)$-measurable too. We have always
$R(M(\lambda),E(\lambda))\supset M(\lambda)$, and $\ne$ for
$\lambda\in\mathcal M$ (as $E(\lambda)\notin M(\lambda)$). Consequently we
have always $R(M(\lambda),E(\lambda))'\subset M(\lambda)'$, and $\ne$ for
$\lambda\in\mathcal M$. Apply \hyperref[lem:9]{Lemma 9} again. There exists
a $\sigma(\lambda)$-measurable $F(\lambda)$, which is a projection for every
$\lambda$, and $F(\lambda)\in M(\lambda)'$,
$F(\lambda)\notin R(M(\lambda),E(\lambda))'$ for
$\lambda\in\mathcal M$. We can again obtain
$F(\lambda)\in M(\lambda)'$ for all $\lambda$'s.

By the \hyperref[rem:17]{Remark} after \hyperref[lem:9]{Lemma 9} there exist two projections
$E,F$ with $E\sim\sum E(\lambda)$, $F\sim\sum F(\lambda)$. As
$E(\lambda)\in N(\lambda)'$ for all $\lambda$, we have
$E\in\widetilde M$. As $F(\lambda)\in M(\lambda)'$, $F(\lambda)$ commutes
with all $A_i(\lambda),A_i(\lambda)^*$, therefore $F$ commutes with all
$A_i,A_i^*$. By \hyperref[thm:V]{Theorem~V}, $F\in P'$, so $F$ commutes with all $P$ too. Thus
$F$ commutes with $\widetilde M$, hence with $E$, and therefore $F(\lambda)$
with $E(\lambda)$, except for a set of $\sigma(\lambda)$-measure $0$. For the
non-exceptional $\lambda$, $F(\lambda)$ commutes with $E(\lambda)$, and as
$F(\lambda)\in M(\lambda)'$, with all $M(\lambda)$ too, thus $F(\lambda)$
commutes with the ring generated by $M(\lambda)$ and $E(\lambda)$ too. This
means $F(\lambda)\in R(M(\lambda),E(\lambda))'$, and thus implies
$\lambda\notin\mathcal M$. Therefore $\mathcal M$ must have the
$\sigma(\lambda)$-measure $0$.

Hence $M(\lambda)=N(\lambda)'$, except for a set of
$\sigma(\lambda)$-measure $0$, and therefore
$\widetilde M\approx\sum N(\lambda)'$ implies
$\widetilde M\approx\sum M(\lambda)$. Thus the proof is completed.

We are now in the position to prove:

\par\medskip\noindent\textbf{Theorem VI.}\phantomsection\label{thm:VI}\quad
\begin{itshape}
If the $M(\lambda)$ are rings containing $1$, depending
$\sigma(\lambda)$-measurably on $\lambda$, then there exists a unique ring
$M\approx\sum M(\lambda)$. This $M$ is a ring containing $1$,
$P\subset M\subset P'$.

Conversely, if $M$ is a ring containing $1$, $P\subset M\subset P'$, then
there exists a system of rings containing $1$, $M(\lambda)$, depending
$\sigma(\lambda)$-measurably on $\lambda$, with
$M\approx\sum M(\lambda)$. This system is unique, except for an arbitrary
$\lambda$-set of $\sigma(\lambda)$-measure $0$.
\end{itshape}

\par\smallskip\noindent\textit{Proof.}
The existence of $M$, if the $M(\lambda)$ are given, follows from
\hyperref[lem:11]{Lemmas 11}, \hyperref[lem:12]{12}, the uniqueness of $M$
has been observed immediately after \hyperref[def:6]{Definition~6}.

As every $A\sim\sum A(\lambda)$ belongs to $P'$, necessarily
$M\subset P'$. If $A\in P$, then we know from \secref{sec:14}, that
$A\sim\sum \alpha(\lambda)1$, and as $1\in M(\lambda)$,
$\alpha(\lambda)1\in M(\lambda)$, this implies $A\in M$. Thus
$P\subset M$. As $1\in P\subset M$, we have $1\in M$ and
$P\subset M\subset P'$.

\origpage{458}
Assume now that $M$, $P\subset M\subset P'$, is given. Choose a sequence of
operators $A_i$, $i=1,2,\ldots$, which generate the ring $M$ (cf. the
corresponding construction in the proof of \hyperref[lem:12]{Lemma 12}, for
$\widetilde M'$). As $A_i\in M\subset P'$, we can write
$A_i\sim\sum A_i(\lambda)$. Denote the ring generated by the
$A_i(\lambda)$, $i=1,2,\ldots$, by $M(\lambda)$. (We may assume
$1\in M(\lambda)$, as we can include $1\sim\sum 1$ in the system of the
$A_i$.) The $M(\lambda)$ depend thus measurably on $\lambda$, and by
\hyperref[lem:12]{Lemma 12},
$R(P,A_1,A_2,\ldots)\approx\sum M(\lambda)$. Now
$R(P,A_1,A_2,\ldots)\supset R(A_1,A_2,\ldots)=M$, and as $P\subset M$ and
all $A_i\in M$, $R(P,A_1,A_2,\ldots)\subset M$. Thus
$R(P,A_1,A_2,\ldots)=M$, and so $M\approx\sum M(\lambda)$, proving the
existence of the desired $M(\lambda)$.

That the change of the $M(\lambda)$ on a set of $\sigma(\lambda)$-measure
$0$ does not affect $M\approx\sum M(\lambda)$, is obvious. Thus it remains
to prove only this: If $M\approx\sum M(\lambda)$ and
$M\approx\sum N(\lambda)$, then $M(\lambda)=N(\lambda)$, except for a
$\sigma(\lambda)$-set of measure $0$. Consider $M(\lambda)N(\lambda)$ and
$N(\lambda)$, both are $\sigma(\lambda)$-measurable
(cf. \hyperref[lem:8]{Lemma 8}). As there is always
$M(\lambda)N(\lambda)\subset N(\lambda)$, we can apply
\hyperref[lem:9]{Lemma 9}. There exists a $\sigma(\lambda)$-measurable
$E(\lambda)$, which is a projection for every $\lambda$, and
$E(\lambda)\notin M(\lambda)N(\lambda)$, $E(\lambda)\in N(\lambda)$ whenever
$M(\lambda)N(\lambda)\ne N(\lambda)$. As in the proof of
\hyperref[lem:12]{Lemma 12}, we can obtain $E(\lambda)\in N(\lambda)$ for all
$\lambda$'s. By the \hyperref[rem:17]{Remark} after \hyperref[lem:9]{Lemma 9}, a projection
$E\sim\sum E(\lambda)$ exists. As $E(\lambda)\in N(\lambda)$ for all
$\lambda$'s, therefore $M\approx\sum N(\lambda)$ necessitates $E\in M$.
Now $M\approx\sum M(\lambda)$ shows, that $E(\lambda)\in M(\lambda)$,
except for a set $N_1$ of $\sigma(\lambda)$-measure $0$. Thus
$\lambda\notin N_1$ implies $E(\lambda)\in M(\lambda)$, and as
$E\in N(\lambda)$, therefore $E\in M(\lambda)N(\lambda)$ too. This
necessitates, as we know, $M(\lambda)N(\lambda)=N(\lambda)$, and thus
$M(\lambda)\supset N(\lambda)$.

Similarly there exists a set $N_2$ of $\sigma(\lambda)$-measure $0$, so that
$\lambda\notin N_2$ implies $M(\lambda)\subset N(\lambda)$. Now
$N=N_1+N_2$ has $\sigma(\lambda)$-measure $0$ too, and
$\lambda\notin N$ implies obviously $M(\lambda)=N(\lambda)$.

Thus all parts of our theorem are proved.

\section{Algebra of the decomposition}\label{sec:19}
Our next task is to discuss the algebra of the decomposition
$M\approx\sum M(\lambda)$. This is done in the three lemmas which follow.

\par\medskip\noindent\textbf{Lemma 13.}\phantomsection\label{lem:13}\quad
\begin{itshape}
$M\approx\sum M(\lambda)$ implies
$M'\approx\sum M(\lambda)'$.
\end{itshape}

\par\smallskip\noindent\textit{Proof.}
The questions of $\sigma(\lambda)$-measurability are settled by
\hyperref[lem:7]{Lemma 7}. Form the $A_i\sim\sum A_i(\lambda)$,
$i=1,2,\ldots$, in the sense of \hyperref[lem:11]{Lemma 11}, for the
$M(\lambda)$. We know that $A_i\in M$.

If $A\in M'$ then as $M\supset P$, $M'\subset P'$, $A\in P'$ we can write by
\hyperref[thm:V]{Theorem~V}: $A\sim\sum A(\lambda)$. As $A\in M'$, $A_i\in M$, therefore $A$
commutes with $A_i,A_i^*$, so by \secref{sec:14}, $A(\lambda)$ commutes with
$A_i(\lambda),A_i(\lambda)^*$, for all $i=1,2,\ldots$, except for a
$\lambda$-set of $\sigma(\lambda)$-measure $0$. By adding the (countably
many) exceptional sets for all $i=1,2,\ldots$, we obtained a fixed
exceptional set of $\sigma(\lambda)$-measure $0$. In it we may replace
$A(\lambda)$ by $0$, thus we can assume that $A(\lambda)$ commutes with
$A_i(\lambda),A_i(\lambda)^*$ for all $i=1,2,\ldots$, and all $\lambda$'s.
Thus it commutes with $R(A_1(\lambda),A_2(\lambda),\ldots)=M(\lambda)$,
that is: $A(\lambda)\in M(\lambda)'$.

Conversely: Assume $A\sim\sum A(\lambda)$,
$A(\lambda)\in M(\lambda)'$. For any $B\in M$ we have
$B\sim\sum B(\lambda)$, $B(\lambda)\in M(\lambda)$, therefore we see:
$A(\lambda),B(\lambda)$ commute for every $\lambda$, so $A,B$ commute. As
this is true for every $B\in M$, it implies $A\in M'$. Thus we have proved
$M'\approx\sum M(\lambda)'$.

\par\medskip\noindent\textbf{Lemma 14.}\phantomsection\label{lem:14}\quad
\begin{itshape}
$M_s\approx\sum M_s(\lambda)$, $s=1,2,\ldots$ (the sequence of the
$s$ may be finite

\origpage{459}
or infinite) implies
\[
M_1\cdot M_2\cdot\ldots\approx
\sum M_1(\lambda)\cdot M_2(\lambda)\cdot\ldots
\]
and
\[
R(M_1,M_2,\ldots)\approx
\sum R(M_1(\lambda),M_2(\lambda),\ldots).
\]
\end{itshape}

\par\smallskip\noindent\textit{Proof.}
The questions of $\sigma(\lambda)$-measurability are settled by
\hyperref[lem:8]{Lemma 8}. The statement
$M_1\cdot M_2\cdot\ldots\approx
\sum M_1(\lambda)\cdot M_2(\lambda)\cdot\ldots$ is obvious, if we only
remember \hyperref[def:6]{Definition~6}. Now this gives, together with
\hyperref[lem:13]{Lemma 13}, the statement
\[
R(M_1,M_2,\ldots)'
 \approx\sum R(M_1(\lambda),M_2(\lambda),\ldots)',
\]
because of
$R(M_1,M_2,\ldots)'=M_1'\cdot M_2'\cdot\ldots$,
$R(M_1(\lambda),M_2(\lambda),\ldots)'=
M_1(\lambda)'\cdot M_2(\lambda)'\cdot\ldots$
(cf. footnote~12 in the proof of \hyperref[lem:8]{Lemma 8}).

\par\medskip\noindent\textbf{Lemma 15.}\phantomsection\label{lem:15}\quad
\begin{itshape}
If $M\approx\sum M(\lambda)$, $N\approx\sum N(\lambda)$, then
$M\subset N$ is equivalent to $M(\lambda)\subset N(\lambda)$ for all
$\lambda$'s, except perhaps for a set of $\sigma(\lambda)$-measure $0$.
\end{itshape}

\par\smallskip\noindent\textit{Proof.}
By \hyperref[lem:14]{Lemma 14}, we have
$M\cdot N\approx\sum M(\lambda)\cdot N(\lambda)$. As
$M\approx\sum M(\lambda)$, we see by \hyperref[thm:VI]{Theorem VI}, that
$M\cdot N=M$ means $M(\lambda)\cdot N(\lambda)=M(\lambda)$ for all
$\lambda$'s, except perhaps for a $\lambda$-set of
$\sigma(\lambda)$-measure $0$. But $M\cdot N=M$ is equivalent to
$M\subset N$, and $M(\lambda)\cdot N(\lambda)=M(\lambda)$ to
$M(\lambda)\subset N(\lambda)$. This gives the desired statement.

The behavior of the decomposition $M\approx\sum M(\lambda)$ under
equivalence is this:

\textit{If $M\approx\sum M(\lambda)$, and the representation of
$\mathfrak H$ as a generalized direct sum is replaced by an equivalent one,
with the mapping $\lambda\rightleftarrows m(\lambda)$, then the decomposition
of $M$ in this new representation will be $M\approx\sum N(\lambda)$, where
$N(\lambda)$ is defined as follows: $N(m(\lambda))$ is the ring which
corresponds to $M(\lambda)$ under the isomorphism of
$\mathfrak H_1^{(\lambda)}$ with $\mathfrak H_2^{(m(\lambda))}$.}

One verifies this statement immediately with the help of \hyperref[def:6]{Definition~6}, and
of the corresponding statement for $A\sim\sum A(\lambda)$ at the end of
\secref{sec:14}.

\vnpart{IV}{Decomposition of rings with respect to their center}

\section{Central decomposition into factors}\label{sec:20}
We can now apply the fundamental operation of the theory of finite-dimensional
matrices, which is indicated in the title of this part. The results of our
foregone discussions, particularly those of \hyperref[thm:VI]{Theorem VI}
and \hyperref[lem:13]{Lemmas 13--15} in \secref{sec:18}--\secref{sec:19},
have reduced the general problem to the same level. The meaning and the effect
of this operation are best expressed by the following theorem:

\par\medskip\noindent\textbf{Theorem VII.}\phantomsection\label{thm:VII}\quad
\begin{itshape}
Let $M$ be a ring in $\mathfrak H$ which contains $1$. Those
(Abelian) rings $P$ for which a decomposition
$M\approx\sum M(\lambda)$ in the sense of \hyperref[def:6]{Definition~6} is possible, are the
rings (containing $1$) $P\subset M\cdot M'$ and only these. (These rings are
all automatically Abelian.)

The choice $P=M\cdot M'$ is characterized by that property, that for this and
only for this $P$ all rings $M(\lambda)$ of the decomposition
$M\approx\sum M(\lambda)$ (except for a $\lambda$-set of
$\sigma(\lambda)$-measure $0$) are factors in the sense of \cite{ref10}. In
other words: $M(\lambda)\cdot M(\lambda)'=0(\lambda)$ (cf. at the end of
\secref{sec:17}).
\end{itshape}

\par\smallskip\noindent\textit{Proof.}
$M$ can be decomposed with respect to $P$ if and only if
$P\subset M\subset P'$ and $P$ is Abelian. (Of course always $1\in P$.) Now
$M\subset P'$ means that $M,P$ commute, which is equivalent to
$P\subset M'$ too. Thus $P\subset M\subset P'$ means
$P\subset M$ and $P\subset M'$, that is $P\subset M\cdot M'$. This implies
the Abelian character too, because $M\cdot M'$ is Abelian:
$M\cdot M'\subset M$ so $(M\cdot M')'\supset M'\supset M\cdot M'$.

By \hyperref[lem:13]{Lemmas 13}, \hyperref[lem:14]{14},
$M\cdot M'\approx\sum M(\lambda)\cdot M(\lambda)'$ and by the second
statement of

\origpage{460}
\secref{sec:14}, $P\approx\sum 0(\lambda)$. Thus
\hyperref[thm:VI]{Theorem VI} shows that
$M(\lambda)\cdot M(\lambda)'=0(\lambda)$ (except for a $\lambda$-set of
measure $0$) is equivalent to $M\cdot M'=P$. Thus all parts of our theorem
are proved.---

\hyperref[thm:VII]{Theorem VII} makes it possible to apply the results of
\cite{ref10} for \textit{factors} to all rings. The most important tools which
will be used in this connection are these: The classification of all factors
(\cite[Theorem~VIII]{ref10}), the notions of relative dimensionality
(\cite[Theorem~VII]{ref10}) and of the relative trace
(\cite[Theorem~XIII]{ref10}). The first steps towards a classification and
discussion of all rings on this basis will be made in the sections which
follow.

\section{Weight functions and relative dimensionality}\label{sec:21}
We defined loc. cit. above the notion of a relative dimensionality function
$d_M(E)$ for factors $M$, and proved, that for every $M$ there exists one and
(up to a constant factor $>0$, $<+\infty$) only one such $d_M(E)$. We now wish
to introduce a similar, although somewhat more general notion which covers
all rings $M$ (containing $1$):

\par\medskip\noindent\textbf{Definition 7.}\phantomsection\label{def:7}\quad
\begin{itshape}
Let $M$ be a ring (containing $1$). $\Delta_M(E)$ is a weight-function (with
respect to $M$) if it has the following properties:
\end{itshape}

\begin{enumerate}[label=(\roman*)]
\item $\Delta_M(E)$ is defined for all projections $E\in M$ and only these.
Its values are real numbers $\geq0$, $\leq+\infty$.
$\Delta_M(E)=0$ if and only if $E=0$.

\item $E_1,E_2,\ldots\in M$, all $E_j$'s projections,
$E_jE_k=0$ for $j\ne k$, imply
\[
\Delta_M(E_1+E_2+\cdots)
 =\Delta_M(E_1)+\Delta_M(E_2)+\cdots .
\]

\item $E,U\in M$, $E$ a projection, $U$ unitary, imply
$\Delta_M(UEU^{-1})=\Delta_M(E)$.
\end{enumerate}

In order to establish a first connection between the relative dimensionality
functions and the weight functions, we now prove:

\par\medskip\noindent\textbf{Lemma 16.}\phantomsection\label{lem:16}\quad
\begin{itshape}
If $M$ is a factor, then the notions of relative dimensionality functions and
of weight functions are equivalent, with this one exception:
\[
\Delta_M^{\infty}(E)=
\begin{cases}
0&\text{for }E=0,\\
\infty&\text{for }E\ne0
\end{cases}
\]
is a relative dimensionality function in the case III only
(cf. \cite[Theorem~VIII et seq.]{ref10}), but it is always a weight function.
\end{itshape}

\par\smallskip\noindent\textbf{Remark:}\phantomsection\label{rem:21}
It will appear from the proof, that this lemma remains true if we replace the
condition (ii) in \hyperref[def:7]{Definition 7}, by the weaker condition
(ii)$'$, which arises from (ii) by permitting two addends $E_j$ ($j=1,2$)
only. Thus it would not affect \hyperref[def:7]{Definition 7} essentially if
we replaced in it (ii) by (ii)$'$, provided that $M$ is a factor. But for
general rings $M$ this replacement is essential, as we will see in $(\gamma)$
in \secref{sec:24}.

\par\smallskip\noindent\textit{Proof.}
The statements concerning $\Delta_M^{\infty}(E)$ are obvious. Every relative
dimensionality function $d_M(E)$ is at the same time a weight function: It
fulfills (i) by definition, that is, by \cite[Definition~8.2.1]{ref10}, it
fulfills (ii) by \cite[Lemma~8.3.2]{ref10}; and (iii) follows from
\cite[Definition~8.2.1, condition (ii)]{ref10}, which states
$d_M(V^*V)=d_M(VV^*)$ for every partially isometric $V$
(cf. \cite[Definition~6.1.1 and Lemma~4.3.1]{ref10}), and we need only put
$V=UE$ (then $V^*=E^*U^*=EU^{-1}$,

\origpage{461}
and so $V^*V=E$, $VV^*=UEU^{-1}$). Thus we must only prove the converse
statement. If a weight function $d_M(E)$ is not $\Delta_M^{\infty}(E)$, that
is, if it does not assume only the values $0,+\infty$, then it is at the same
time a relative dimensionality function. By \cite[Lemma~8.3.5]{ref10}, this
is proved, if we can show that
\[
E\sim F\vnmod{M}\quad\hbox{implies}\quad\Delta_M(E)=\Delta_M(F).
\]

Assume therefore $E\sim F\vnmod{M}$. If
$1-E\sim1-F\vnmod{M}$, then let $V,W$ be partially isometric operators, with
the initial projections $E$ resp. $1-E$, and the final ones $F$ resp. $1-F$.
Then $U=V+W$ is unitary and maps $E$ on $F$, that is $UEU^{-1}=F$.
(Cf. \cite[Lemmas~4.3.1 and 4.3.3]{ref10}.) Now (iii) gives
$d_M(E)=d_M(F)$. So we need only to prove
$1-E\sim1-F\vnmod{M}$. This follows from $E\sim F\vnmod{M}$, if $E$, and with it
$F$, is finite, as one sees by using a relative dimensionality function
$d_M(E)$ and arguing as in \cite[Lemma~8.4.1, proof of (iii)]{ref10}.
Therefore we may assume that $E$, and with it $F$, is infinite.

In this case $E=E_1+E_2$, $E\sim E_1\sim E_2\vnmod{M}$, and
$F=F_1+F_2$, $F\sim F_1\sim F_2\vnmod{M}$.
(Cf. \cite[Definition~6.1.1]{ref10}.) Thus $E_1,E_2$ are infinite, and so
$1-E_1\geq E_2$ is infinite too. Similarly $F_1,F_2$, and
$1-F_1\geq F_2$ are infinite. Thus $E_1\sim F_1\vnmod{M}$,
$1-E_1\sim1-F_1\vnmod{M}$, which implies, as have already been seen,
$\Delta_M(E_1)=\Delta_M(F_1)$. Similarly
$\Delta_M(E_2)=\Delta_M(F_2)$, and so
\[
\Delta_M(E)=\Delta_M(E_1)+\Delta_M(E_2)
=\Delta(F_1)+\Delta_M(F_2)=\Delta_M(F).
\]
This completes the proof.---

Let $M$ be a ring containing $1$, put $P=M\cdot M'$, and decompose $M$ in
the sense of \hyperref[thm:VII]{Theorem VII} with respect to this $P$:
$M\approx M(\lambda)$. As each $M(\lambda)$ is a factor, we can select in
each one a weight function $\Delta_{M(\lambda)}(E)$, or even a relative
dimensionality function $d_{M(\lambda)}(E)$. It is clear, that if a definite
$d^*_{M(\lambda)}(E)$ has been chosen (for each $\lambda> -\infty$,
$<+\infty$), then all $\Delta_{M(\lambda)}(E)$ are obtained by forming
\begin{equation}
\Delta_{M(\lambda)}(E)=a(\lambda)d^*_{M(\lambda)}(E)
\tag{21.1}\label{eq:21.1}
\end{equation}
for all possible functions $a(\lambda)$ the values of which are real numbers
$>0$, $\leq+\infty$. (We use the convention $0\cdot\infty=0$. If a
$\Delta_{M(\lambda)}(E)$ is $\Delta_{M(\lambda)}^{\infty}(E)$, then we may
put $a(\lambda)=+\infty$. If we wish to obtain the
$d_{M(\lambda)}(E)$'s only, we must require $a(\lambda)>0$, $<+\infty$.)

Observe that \eqref{eq:21.1}, that is $\Delta_{M(\lambda)}(E)$, determines
$a(\lambda)$ uniquely, if an $E\in M(\lambda)$ with
$d^*_{M(\lambda)}(E)>0$, $<+\infty$ can be found, that is if we do not have
case III in the sense of \cite[Theorem~VIII]{ref10}. In case III,
\eqref{eq:21.1} holds identically for any $a(\lambda)>0$, $<+\infty$
(remember \hyperref[lem:16]{Lemma 16}). So we see: $a(\lambda)$ is uniquely
determined or perfectly arbitrary accordingly to whether we do not or do have
case III.

It is clear, that this $a(\lambda)$ may be a non
$\sigma(\lambda)$-measurable function. Therefore we cannot expect that
$\Delta_{M(\lambda)}(E(\lambda))$ will be a $\sigma(\lambda)$-measurable
(numerical) function for all $\sigma(\lambda)$-measurable $E(\lambda)$'s. We
now proceed to eliminate this inconvenience.

\par\medskip\noindent\textbf{Lemma 17.}\phantomsection\label{lem:17}\quad
\begin{itshape}
Let $d_{M(\lambda)}(E)$ be a relative dimensionality function, let the
projection $E(\lambda)\in M(\lambda)$ depend
$\sigma(\lambda)$-measurably on $\lambda$. Denote the sets of all
$\lambda$'s where

\origpage{462}
$d_{M(\lambda)}(E)=0$, or $=\infty$, or $>0$, $<\infty$, by
$M_0$ or $M_\infty$ or $M_r$, respectively. These sets are
$\sigma(\lambda)$-measurable.
\end{itshape}

\par\smallskip\noindent\textit{Proof.}
Let $\varphi_m(\lambda)$ be a measurable family.

$\lambda\in M_0$ means $E(\lambda)=0$ that is
$E(\lambda)\varphi_m(\lambda)=0$ from $m=1,2,\ldots$, that is
$\sum_{m=1}^{\infty}\lVert E(\lambda)\varphi_m(\lambda)\rVert^2=0$. As the
left side is a $\sigma(\lambda)$-measurable function of $\lambda$, this proves
the $\sigma(\lambda)$-measurability of $M_0$.

Consider now $M_\infty$: $\lambda\in M_\infty$ means that $E(\lambda)$ is
infinite (with respect to $M(\lambda)$) in the sense of
\cite[Definition~7.1.1]{ref10}, in other words: the existence of an
$F\sim E(\lambda)\vnmod{M(\lambda)}$, $F<E(\lambda)$, or: the existence of a
$V\in M(\lambda)$ with $V^*V=E(\lambda)$,
$VV^*E(\lambda)=VV^*\ne E(\lambda)$.\footnote{$V$ must be partially
isometric, but this follows from $V^*V=E(\lambda)$, because thus $V^*V$ is a
projection. (Cf. \cite[Lemma~4.3.2]{ref10}.)} As $M(\lambda)$ is
$\sigma(\lambda)$-measurable, $M(\lambda)'$ is also
$\sigma(\lambda)$-measurable, and so $\sigma(\lambda)$-measurable operators
$A_q(\lambda)$, $q=1,2,\ldots$, exist, such that
$1,A_q(\lambda)$, $q=1,2,\ldots$, generate the ring $M(\lambda)'$. Now
$V\in M(\lambda)$ means that $V$ commutes with all $M(\lambda)'$, that is
with every $A_q(\lambda),A_q(\lambda)^*$, $q=1,2,\ldots$.

Proceed now as in the proof of \hyperref[lem:6]{Lemma 6}. Introduce
$k(\lambda)$, and then the set $\Delta_p$ of all $\lambda$ with
$k(\lambda)=p$, $p=1,2,\ldots,\infty$. Choose there a set $N_p$ of
$\sigma(\lambda)$-measure $0$, such that $\Delta_p'=\Delta_p-N_p$ is a
Borel-set, and that all
$(E(\lambda)\varphi_m(\lambda),\varphi_n(\lambda))$ and
$(A_q(\lambda)\varphi_m(\lambda),\varphi_n(\lambda))$,
$q,m,n=1,2,\ldots$, are Baire functions in $\Delta_p'$. Introduce a
$p$-dimensional unitary space $\mathfrak H_p'$, in it a complete normalized
orthogonal set $\psi_1',\psi_2',\ldots$, and let $J^{(\lambda)}$ be the
isomorphic mapping of $\mathfrak H_p'$ on $\mathfrak H^{(\lambda)}$
($\lambda\in\Delta_p$), by which $\psi_m'$ corresponds to
$\varphi_m(\lambda)$.

Consider now the above condition for $\lambda\in M_\infty$. Replace the
operator $V$ in $\mathfrak H^{(\lambda)}$ by the $J^{(\lambda)}$-isomorphic
$\widetilde V=J^{(\lambda)-1}VJ^{(\lambda)}$ in $\mathfrak H_p'$. Rewriting
our conditions for this case, we see: $\lambda\in M_\infty$ means, if
$\lambda\in\Delta_p$, that a $\widetilde V$ in $\mathfrak H_p'$ exists, for
which
\begin{equation}
\left\{
\begin{aligned}
\widetilde VJ^{(\lambda)-1}A_q(\lambda)J^{(\lambda)}
 &=J^{(\lambda)-1}A_q(\lambda)J^{(\lambda)}\widetilde V,\\
\widetilde VJ^{(\lambda)-1}A_q(\lambda)^*J^{(\lambda)}
 &=J^{(\lambda)-1}A_q(\lambda)^*J^{(\lambda)}\widetilde V,
 &&q=1,2,\ldots,\\
\widetilde V^*\widetilde V
 &=J^{(\lambda)-1}E(\lambda)J^{(\lambda)},\\
\widetilde V\widetilde V^*J^{(\lambda)-1}E(\lambda)J^{(\lambda)}
 &=\widetilde V\widetilde V^*
 \ne J^{(\lambda)-1}E(\lambda)J^{(\lambda)}.
\end{aligned}
\right.
\tag{21.2}\label{eq:21.2}
\end{equation}
Observe that a $\widetilde V$ fulfilling \eqref{eq:21.2} is partially
isometric, and therefore automatically
$\lVert\!\lvert\widetilde V\rvert\!\rVert\leq1$.

Denote again the set $-\infty<\lambda<+\infty$ by $r$, and the set of all
operators $\widetilde A$ in $\mathfrak H_p'$ with
$\lVert\!\lvert\widetilde A\rvert\!\rVert\leq1$ (in its weak topology) by
$S_p$. Let the set of all pairs $(\lambda,\widetilde V)$ in $r\times S_p$
with $\lambda\in\Delta_p'$ and fulfilling \eqref{eq:21.2} be $D_p$. Proceed
as in the proof of \hyperref[lem:6]{Lemma 6}: Replace the operator equations
$X=Y$ of \eqref{eq:21.2} by the equivalent numerical equations
$(X\psi_m',\psi_n')=(Y\psi_m',\psi_n')$, $m,n=1,2,\ldots$. Now these
expressions are convergent, infinite sums, the terms of which are polynomials
of the $(\widetilde V\psi_m',\psi_n')$ and of the
\[
(J^{(\lambda)-1}A_q(\lambda)J^{(\lambda)}\psi_m',\psi_n')
 =(A_q(\lambda)\varphi_m(\lambda),\varphi_n(\lambda)),
\]
\[
(J^{(\lambda)-1}E(\lambda)J^{(\lambda)}\psi_m',\psi_n')
 =(E(\lambda)\varphi_m(\lambda),\varphi_n(\lambda)).
\]
The first expression is a (weakly) continuous function of $\widetilde V$, the
two others are (owing to $\lambda\in\Delta_p'$) Baire functions of $\lambda$.
Thus we can argue as in the proof of \hyperref[lem:6]{Lemma 6} (using
\cite[p.~260]{ref15}): $D_p$ is a Borel-set in $r\times S_p$.

\origpage{463}
Now we argue as in the proof of \hyperref[lem:7]{Lemma 7}: As $D_p$ is a
Borel-set in $S_p$, it is \textit{a fortiori} an analytic set. Therefore,
\hyperref[lem:5]{Lemma~5} applies to it: Define $F((\lambda,\widetilde V))=\lambda$, which is
obviously a continuous function. $\lambda\in\Delta_p'\cdot M_\infty$ means
that a $\widetilde V$ with $(\lambda,\widetilde V)\in D_p$ exists, i.e. that
an $a\in D_p$ with $F(a)=\lambda$ exists. Hence
$\Delta_p'\cdot M_\infty$ is $\sigma(\lambda)$-measurable. Next
$(\Delta_p-\Delta_p')\cdot M_\infty=N_p\cdot M_\infty\subset N_p$ has the
$\sigma(\lambda)$-measure $0$, thus $\Delta_p\cdot M_\infty$ is
$\sigma(\lambda)$-measurable too, and with it
$M_\infty=\sum_{p=1,2,\ldots,\infty}\Delta_p\cdot M_\infty$.

Thus $M_0,M_\infty$ are $\sigma(\lambda)$-measurable, and with them
$M_r=r-M_0-M_\infty$ is $\sigma(\lambda)$-measurable too. This completes the
proof.

\par\medskip\noindent\textbf{Lemma 18.}\phantomsection\label{lem:18}\quad
\begin{itshape}
Let $d_{M(\lambda)}(E)$ be a relative dimensionality function, and let the
projections $E(\lambda),F(\lambda)\in M(\lambda)$ depend
$\sigma(\lambda)$-measurably on $\lambda$. Then the (numerical) function
\[
\frac{d_{M(\lambda)}(E(\lambda))}{d_{M(\lambda)}(F(\lambda))}
\]
is $\sigma(\lambda)$-measurable.
\end{itshape}

\par\smallskip\noindent\textit{Proof.}
Form the sets $M_0,M_\infty,M_r$ in the sense of
\hyperref[lem:17]{Lemma 17} for $E(\lambda)$, and similarly
$P_0,P_\infty,P_r$ for $F(\lambda)$. We have
\[
r=M_0+M_r+M_\infty
=P_0+P_r+P_\infty
 =(M_0+M_r+M_\infty)(P_0+P_r+P_\infty)
\]
\[
=M_0P_0+M_rP_0+M_\infty P_0
 +M_0P_r+M_rP_r+M_\infty P_r
 +M_0P_\infty+M_rP_\infty+M_\infty P_\infty.
\]
All terms of this sum are $\sigma(\lambda)$-measurable by
\hyperref[lem:17]{Lemma 17}. In terms No.~1 and 9
$d_{M(\lambda)}(E(\lambda))/d_{M(\lambda)}(F(\lambda))$ has the form $0/0$
resp. $\infty/\infty$, and is thus nowhere defined; in all others it is
everywhere defined. In terms No.~2, 3, 4, 6, 7, 8, it is constant, being
respectively $\infty,\infty,0,\infty,0,0$; in term No.~5 it is variable and
$>0$, $<+\infty$. Thus we need only to prove the
$\sigma(\lambda)$-measurability of
$d_{M(\lambda)}(E(\lambda))/d_{M(\lambda)}(F(\lambda))$ in term No.~5, that
is in $M_rP_r$.

Choose two numbers $u,v=1,2,\ldots$, and define $M_{u,v}$ as the set of those
$\lambda$'s, for which projections
$G_1,\ldots,G_u,H_1,\ldots,H_v\in M(\lambda)$ exist, which have the following
properties: $G_1\sim\cdots\sim G_u\sim H_1\sim\cdots\sim H_v\vnmod{M(\lambda)}$,
$G_jG_k=H_jH_k=0$ for $j\ne k$,
$G_1+\cdots+G_u\leq E(\lambda)$,
$H_1+\cdots+H_v\geq F(\lambda)$. Our first objective is to prove that
$M_{u,v}$ is $\sigma(\lambda)$-measurable.

$\lambda\in M_{u,v}$ means that the above described
$G_1,\ldots,G_u,H_1,\ldots,H_v$ exist. Introducing another projection
$K\in M(\lambda)$, which is $\sim$ to all of them mod $M(\lambda)$, we may
describe the entire situation by speaking of the partially isometric
operators in $M(\lambda)$, which have all the initial projection $K$, and the
respective final projections $G_1,\ldots,G_u,H_1,\ldots,H_v$. We denote them
by $V_1,\ldots,V_{u+v}$. Now our conditions are: Each
$V_s\in M(\lambda)$, that is $V_s$ commutes with all $M(\lambda)'$, that is
with all $A_q(\lambda)$, $q=1,2,\ldots$ (the $A_q(\lambda)$,
$q=1,2,\ldots$, generate $M(\lambda)'$ and are
$\sigma(\lambda)$-measurable, as in the proof of
\hyperref[lem:17]{Lemma 17}); $V_sV_s^*V_s=V_s$ (this expresses the partial
isometry, cf. \cite[Lemma~4.3.2]{ref10});
$V_1^*V_1=\cdots=V_{u+v}^*V_{u+v}$;
$V_jV_j^*V_kV_k^*=0$ if $j\ne k$ and both $j,k=1,\ldots,u$ or both
$=u+1,\ldots,u+v$;
$(V_1V_1^*+\cdots+V_uV_u^*)E(\lambda)=V_1V_1^*+\cdots+V_uV_u^*$,
$(V_{u+1}V_{u+1}^*+\cdots+V_{u+v}V_{u+v}^*)F(\lambda)=F(\lambda)$.

Proceed now as in the proofs of \hyperref[lem:6]{Lemmas 6},
\hyperref[lem:17]{17}: Introduce $k(\lambda)$, the set $\Delta_p$ of all
$\lambda$'s with $k(\lambda)=p$, $p=1,2,\ldots,\infty$, choose
$N_p$, $\Delta_p'=\Delta_p-N_p$ as there, with respect to the operators
$E(\lambda),F(\lambda),A_q(\lambda),A_q(\lambda)^*$, $q=1,2,\ldots$.
Introduce $\mathfrak H_p'$, the $\psi_1',\psi_2',\ldots$, and the isomorphism
$J^{(\lambda)}$ as there. Replace now in our

\origpage{464}
above conditions the $V_s$ in $\mathfrak H^{(\lambda)}$ by their
$J^{(\lambda)}$-isomorphic
$\widetilde V_s=J^{(\lambda)-1}V_sJ^{(\lambda)}$ in $\mathfrak H_p'$. Then
our conditions become:
\begingroup\small
\begin{equation}
\left\{
\begin{aligned}
\widetilde V_sJ^{(\lambda)-1}A_q(\lambda)J^{(\lambda)}
 &=J^{(\lambda)-1}A_q(\lambda)J^{(\lambda)}\widetilde V_s,\\
\widetilde V_sJ^{(\lambda)-1}A_q(\lambda)^*J^{(\lambda)}
 &=J^{(\lambda)-1}A_q(\lambda)^*J^{(\lambda)}\widetilde V_s,
 &&\substack{s=1,\ldots,u+v,\\q=1,2,\ldots,}\\
\widetilde V_s\widetilde V_s^*\widetilde V_s&=\widetilde V_s,
 &&s=1,\ldots,u+v,\\
\widetilde V_1^*\widetilde V_1&=\cdots=
 \widetilde V_{u+v}^*\widetilde V_{u+v},\\
\widetilde V_j\widetilde V_j^*\widetilde V_k\widetilde V_k^*&=0,
 &&\substack{j\ne k,\quad j,k=1,\ldots,u,\\
 \text{or both }j,k=u+1,\ldots,u+v,}\\
(\widetilde V_1\widetilde V_1^*+\cdots+
 \widetilde V_u\widetilde V_u^*)J^{(\lambda)-1}E(\lambda)J^{(\lambda)}
 &=\widetilde V_1\widetilde V_1^*+\cdots+
 \widetilde V_u\widetilde V_u^*,\\
(\widetilde V_{u+1}\widetilde V_{u+1}^*+\cdots+
 \widetilde V_{u+v}\widetilde V_{u+v}^*)J^{(\lambda)-1}F(\lambda)J^{(\lambda)}
 &=J^{(\lambda)-1}F(\lambda)J^{(\lambda)}.
\end{aligned}
\right.
\tag{21.3}\label{eq:21.3}
\end{equation}
\endgroup
Observe that \eqref{eq:21.3} implies for each $\widetilde V_s$ partial
isometricity, and thus
$\lVert\!\lvert\widetilde V_s\rvert\!\rVert\leq1$.

Consider now again $r$ and $S_p$, but now form the direct product
$r\times S_p\times\cdots\times S_p$ with $u+v$ factors $S_p$. Let the set of
all $u+v+1$-uplets
$(\lambda,\widetilde V_1,\ldots,\widetilde V_{u+v})$ in
$r\times S_p\times\cdots\times S_p$, with $\lambda\in\Delta_p'$ and
fulfilling \eqref{eq:21.3} be $C_p$. We prove as in
\hyperref[lem:17]{Lemma 17} that $C_p$ is a Borel-set in
$r\times S_p\times\cdots\times S_p$. $\lambda\in\Delta_p'M_{u,v}$ means
that there exists a subsystem $\widetilde V_1,\ldots,\widetilde V_{u+v}$ with
$(\lambda,\widetilde V_1,\ldots,\widetilde V_{u+v})\in C_p$. From this we
infer, as in \hyperref[lem:17]{Lemma 17}, that
$\Delta_p'M_{u,v}$ is $\sigma(\lambda)$-measurable.
$(\Delta_p-\Delta_p')M_{u,v}=N_pM_{u,v}\subset N_p$ has the
$\sigma(\lambda)$-measure $0$; hence $\Delta_pM_{u,v}$ and
$M_{u,v}=\sum_{p=1,2,\ldots,\infty}\Delta_pM_{u,v}$ are
$\sigma(\lambda)$-measurable, too.

Let us now return to our function
$d_{M(\lambda)}(E(\lambda))/d_{M(\lambda)}(F(\lambda))$ in $M_rP_r$. We wish
to prove that it is $\sigma(\lambda)$-measurable, that is, that the set
$M_\alpha$ of all $\lambda\in M_rP_r$ for which
\[
\frac{d_{M(\lambda)}(E(\lambda))}
     {d_{M(\lambda)}(F(\lambda))}>\alpha
\]
is $\sigma(\lambda)$-measurable for every (real) $\alpha$. For $\alpha\leq0$
this set is $M_rP_r$ and our statement holds, thus we may assume $\alpha>0$.

Now if $\lambda\in M_rP_r$, then the last inequality is clearly equivalent
to the existence of two $u,v=1,2,\ldots$, such that $u/v>\alpha$, and
$\lambda\in M_{u,v}$.\footnote{The sufficiency of this condition is obvious.
In order to prove its necessity, consider the classification of
\cite[Theorem~VIII]{ref10} (for this $\lambda$). As
$\lambda\in M_rP_r$, therefore $d_{M(\lambda)}(E(\lambda))>0$, $<+\infty$,
hence case III cannot hold. If case I holds, choose $K$ so as to minimize
$d_{M(\lambda)}(K)$ ($>0$), denoting this minimum by $\varepsilon_\lambda$.
Then we can put
\[
u=\frac{1}{\varepsilon_\lambda}d_{M(\lambda)}(E(\lambda)),\qquad
v=\frac{1}{\varepsilon_\lambda}d_{M(\lambda)}(F(\lambda))
\]
as these quantities must be positive and integers, that is $=1,2,\ldots$.
If case II holds, then $u/v$ may be any rational number
$\leq d_{M(\lambda)}(E(\lambda))/d_{M(\lambda)}(F(\lambda))$ and
$>\alpha$. Then we can put
\[
d_{M(\lambda)}(K)=\frac{1}{v}d_{M(\lambda)}(F(\lambda))
\quad\text{and hence}\quad
\leq\frac{1}{u}d_{M(\lambda)}(E(\lambda)).
\]} Therefore
\[
Q_\alpha=M_rP_r\sum_{\substack{u,v=1,2,\ldots\\u/v>\alpha}}M_{u,v},
\]
and so it is $\sigma(\lambda)$-measurable along with $M_r,P_r$, and the
$M_{u,v}$. Thus the proof is completed.

\origpage{465}
\par\medskip\noindent\textbf{Lemma 19.}\phantomsection\label{lem:19}\quad
\begin{itshape}
There exists one and (up to a $\lambda$-set of
$\sigma(\lambda)$-measure $0$) only one $\sigma(\lambda)$-measurable set
$\overline M$ with the following properties.
\begin{enumerate}[label=(\roman*)]
\item There exists a projection $\overline E(\lambda)\in M(\lambda)$,
depending $\sigma(\lambda)$-measurably on $\lambda$, for which
$M_0=r-\overline M$, $M_\infty=0$, $M_r=\overline M$.
\item For every projection $E(\lambda)\in M(\lambda)$ which depends
$\sigma(\lambda)$-measurably on $\lambda$ we have
$\mu_{\sigma(\lambda)}(M_r-M_r\cdot\overline M)=0$.
\end{enumerate}
In other words: Disregarding sets of $\sigma(\lambda)$-measure $0$,
$\overline M$ is the greatest $M_r$.
\end{itshape}

\par\smallskip\noindent\textit{Proof.}
We proceed stepwise.

\begin{enumerate}[label=(\Roman*)]
\item Consider all $E(\lambda)$ depending $\sigma(\lambda)$-measurably on
$\lambda$, such that each $E(\lambda)$ is a projection in $M(\lambda)$.
Form their sets $M_r$ and denote the set of all $M_r$ by $T$.

\item If $M\in T$, then an $E(\lambda)$ with
$M_0=r-M$, $M_\infty=0$, $M_r=M$ exists.

\item If $M\in T$, then every $\sigma(\lambda)$-measurable $M'\subset M$ is
$M'\in T$ too.

We prove (II), (III) together, by constructing an $E(\lambda)$ with
$M_0=r-M'$, $M_\infty=0$, $M_r=M'$. Choose an $\widetilde E(\lambda)$ with
$\widetilde M_r=M$. Define now
\[
E(\lambda)=
\begin{cases}
\widetilde E(\lambda)&\text{for }\lambda\in M',\\
0&\text{for }\lambda\notin M'.
\end{cases}
\]
Then $E(\lambda)$ is $\sigma(\lambda)$-measurable along with
$\widetilde E(\lambda)$ and $M'$, and it obviously meets our requirements.

\item If $M_1,M_2,\ldots\in T$, then $M_1+M_2+\cdots\in T$ too.

By (III) we may replace $M_1,M_2,M_3,\ldots$ by
$M_1,M_2-M_1M_2,M_3-(M_1+M_2)M_3,\ldots$. Thus we may assume
$M_jM_k=0$ for $j\ne k$. Now choose $E_j(\lambda)$ for $M_j$ by (II), then
we have for every fixed $\lambda$ at most one $E_j(\lambda)\ne0$. Thus
$E(\lambda)=\sum_{j=1}^{\infty}E_j(\lambda)$ is a projection, and it is clear
that it belongs to $M(\lambda)$, that it depends $\sigma(\lambda)$-measurably
on $\lambda$, and that its $M_r=M_1+M_2+\cdots$.

\item Put
$\alpha=\operatorname{gr.l.b.}_{M\in T}\mu_{\sigma(\lambda)}(M)$. For each
$j=1,2,\ldots$ choose an $M_j\in T$ with
$\mu_{\sigma(\lambda)}(M_j)\geq\alpha-1/j$. By (IV)
$\overline M=\sum_{j=1}^{\infty}M_j\in T$, so
$\mu_{\sigma(\lambda)}(\overline M)\leq\alpha$. But
$\mu_{\sigma(\lambda)}(\overline M)\geq\mu_{\sigma(\lambda)}(M_j)
\geq\alpha-1/j$ for every $j=1,2,\ldots$, thus necessarily
$\mu_{\sigma(\lambda)}(\overline M)=\alpha$. Now if $M\in T$, then by (IV)
$\overline M+M\in T$ too; thus
$\mu_{\sigma(\lambda)}(\overline M+M)=\alpha$. Thus
$\mu_{\sigma(\lambda)}(M-\overline M\cdot M)
=\mu_{\sigma(\lambda)}(\overline M+M)
-\mu_{\sigma(\lambda)}(\overline M)=0$. Thus $\overline M$ has the property
(ii) of our lemma.

\item Applying (II) to $\overline M\in T$ proves that $\overline M$ has the
property (i) of our lemma. Its unicity (up to a $\lambda$-set of
$\sigma(\lambda)$-measure $0$) is an obvious consequence of (i) and (ii).
\end{enumerate}

Thus the proof is completed.

Observe that the sets $M_0,M_\infty,M_r$ (for a given $E(\lambda)$) and
$\overline M$ depend only apparently on the choice of
$d_{M(\lambda)}(E)$: The equations $d_{M(\lambda)}(E)=0$ and $\infty$ which
enter in their definitions mean the same thing for every choice of
$d_{M(\lambda)}(E)$: $E=0$ and $E$ infinite (with respect to $M(\lambda)$),
respectively.

We now proceed as follows:

\par\medskip\noindent\textbf{Definition 8.}\phantomsection\label{def:8}\quad
\begin{itshape}
A weight function $\Delta_{M(\lambda)}(E)$ is
$\sigma(\lambda)$-measurable (in its dependence on $\lambda$) if the
(numerical) function $\Delta_{M(\lambda)}(E(\lambda))$ is
$\sigma(\lambda)$-measurable whenever

\origpage{466}
$E(\lambda)$ is $\sigma(\lambda)$-measurable. (This nomenclature applies to
the relative dimensionality functions $d_{M(\lambda)}(E)$ too, as they
constitute a special case of the $\Delta_{M(\lambda)}(E)$.)
\end{itshape}

\par\medskip\noindent\textbf{Theorem VIII.}\phantomsection\label{thm:VIII}\quad
\begin{itshape}
The following statements are true:
\begin{enumerate}[label=(\roman*)]
\item A $\Delta_{M(\lambda)}(E)$, and even a $d_{M(\lambda)}(E)$, which is
$\sigma(\lambda)$-measurable, exists.
\item If $d^*_{M(\lambda)}(E)$ is $\sigma(\lambda)$-measurable, then all
$\sigma(\lambda)$-measurable $\Delta_{M(\lambda)}(E)$'s obtain by the formula
\eqref{eq:21.1}:
\begin{equation}
\Delta_{M(\lambda)}(E)=a(\lambda)d^*_{M(\lambda)}(E),
\qquad 0<a(\lambda)\leq+\infty,
\tag{21.1}\label{eq:21.1-repeat}
\end{equation}
$a(\lambda)$ being restricted to those functions which are
$\sigma(\lambda)$-measurable in $\overline M$
(cf. \hyperref[lem:19]{Lemma 19}).
\end{enumerate}
\end{itshape}

\par\smallskip\noindent\textit{Proof:}
Observe first that for every $\Delta_{M(\lambda)}(E)$ the (numerical) function
$\Delta_{M(\lambda)}(E(\lambda))$ is $\sigma(\lambda)$-measurable in
$r-\overline M$, if $E(\lambda)$ is $\sigma(\lambda)$-measurable. Choose any
$d_{M(\lambda)}(E(\lambda))$. Then
$\Delta_{M(\lambda)}(E)=a(\lambda)d_{M(\lambda)}(E)$ for some
$a(\lambda)>0$, $\leq+\infty$, and thus
$\Delta_{M(\lambda)}(E(\lambda))>0$, $<+\infty$ implies
$d_{M(\lambda)}(E(\lambda))>0$, $<+\infty$, and therefore for
$\lambda\in r-\overline M$ it implies
$\lambda\in M_r-M_r\cdot\overline M$. Thus we deduce from
\hyperref[lem:19]{Lemma 19}: $r-\overline M$ is
$\sigma(\lambda)$-measurable, and in it
$\Delta_{M(\lambda)}(E(\lambda))>0$, $<+\infty$ occurs only on a
$\lambda$-set of $\sigma(\lambda)$-measure $0$. Thus we must only prove that
the set of all $\lambda\in r-\overline M$,
$\Delta_{M(\lambda)}(E(\lambda))=0$, is $\sigma(\lambda)$-measurable.

Now $\Delta_{M(\lambda)}(E(\lambda))=0$ means $E(\lambda)=0$, and this
determines a $\sigma(\lambda)$-measurable set (cf. the first part of the proof
of \hyperref[lem:17]{Lemma 17}), and its intersection with
$r-\overline M$ will therefore be $\sigma(\lambda)$-measurable too.

Thus we can restrict ourselves from now on to the $\lambda\in\overline M$.

Choose an arbitrary $d_{M(\lambda)}(E)$, an $\overline E(\lambda)$ by
\hyperref[lem:19]{Lemma 19}, and define
\[
d^*_{M(\lambda)}(E)=a^*(\lambda)d_{M(\lambda)}(E),\qquad
a^*(\lambda)=
\begin{cases}
\dfrac{1}{d_{M(\lambda)}(\overline E(\lambda))},
   &\text{for }\lambda\in\overline M,\\[6pt]
1,&\text{for }\lambda\notin\overline M.
\end{cases}
\]
If $\lambda\in\overline M$, then
\[
d^*_{M(\lambda)}(E(\lambda))
=\frac{d_{M(\lambda)}(E(\lambda))}
       {d_{M(\lambda)}(\overline E(\lambda))},
\]
which is a $\sigma(\lambda)$-measurable function along with $E(\lambda)$, by
\hyperref[lem:18]{Lemma 18}. Thus $d^*_{M(\lambda)}(E(\lambda))$ is
$\sigma(\lambda)$-measurable in $\overline M$, and therefore in all $r$, in
other words $d^*_{M(\lambda)}(E)$ is $\sigma(\lambda)$-measurable, proving
(i).

Any $\Delta_{M(\lambda)}(E)$ can be represented by \eqref{eq:21.1} with some
$a(\lambda)>0$, $\leq+\infty$. Its $\sigma(\lambda)$-measurability means,
that $\Delta_{M(\lambda)}(E(\lambda))$ is $\sigma(\lambda)$-measurable in all
$r$ for every $\sigma(\lambda)$-measurable $E(\lambda)$, but we know that it
suffices to establish this in $\overline M$. Thus the
$\sigma(\lambda)$-measurability of $a(\lambda)$ in $\overline M$ is
sufficient. But it is necessary too: Put $E(\lambda)=\overline E(\lambda)$ in
\eqref{eq:21.1}, for $\lambda\in\overline M$ we have
$d^*_{M(\lambda)}(E(\lambda))>0$, $<+\infty$, and thus we can divide by it,
obtaining
\[
a(\lambda)=
\frac{\Delta_{M(\lambda)}(\overline E(\lambda))}
     {d^*_{M(\lambda)}(\overline E(\lambda))},
\]
which proves the $\sigma(\lambda)$-measurability of $a(\lambda)$ in
$\overline M$. Thus (ii) too is established, and the proof is completed.---

It is worth emphasizing that \hyperref[thm:VIII]{Theorem VIII} shows that the
$a(\lambda)$ of a $\sigma(\lambda)$-measurable
$\Delta_{M(\lambda)}(E)$ need only to be $\sigma(\lambda)$-measurable in
$\overline M$, and not in all $r$! Now we remarked in the discussion which
preceded \hyperref[lem:17]{Lemma 17} that \eqref{eq:21.1} determines
$a(\lambda)$ in $r-C_{III}$ uniquely, but that it leaves $a(\lambda)$
completely arbitrary in $C_{III}$. ($C_{III}$ is the set of all $\lambda$'s,
for which $M(\lambda)$ belongs to the class III, cf.
\cite[Theorem~VIII]{ref10}. Cf. also the notations introduced in
\secref{sec:22}, and in particular

\origpage{467}
in \hyperref[def:9]{Definition 9}.) Besides
$\overline M\subset r-C_{III}$, because $\lambda\in\overline M$ implies
$d_{M(\lambda)}(\overline E(\lambda))>0$, $<+\infty$, and therefore
$\lambda\notin C_{III}$. Thus our result would be very natural, if
$\overline M$ coincided with $r-C_{III}$ up to a set of
$\sigma(\lambda)$-measure $0$. But this, that is
$\mu_{\sigma(\lambda)}(r-C_{III}-\overline M)=0$, is unproved! In fact we
cannot prove at present even the $\sigma(\lambda)$-measurability of $C_{III}$
(which is equivalent to that one of $r-C_{III}-\overline M$). (Cf.
\hyperref[thm:IX]{Theorem IX} and the observations following it.).

A natural attempt to prove
$\mu_{\sigma(\lambda)}(r-C_{III}-\overline M)=0$ would be this: If
$\lambda\in r-C_{III}$, then a projection $E(\lambda)\in M(\lambda)$ with
$d_{M(\lambda)}(E(\lambda))>0$, $<+\infty$ exists. If we could select such
an $E(\lambda)$ for every $\lambda\in r-C_{III}$, and an arbitrary
$E(\lambda)$ for every $\lambda\in C_{III}$, and if this could be done so as
to make the resulting $E(\lambda)$ depend $\sigma(\lambda)$-measurably on
$\lambda$, then we would have for this $E(\lambda)$,
$M_r=r-C_{III}$. Then $M_r-M_r\cdot\overline M=r-C_{III}-\overline M$, and
therefore \hyperref[lem:19]{Lemma 19} gives the desired
$\mu_{\sigma(\lambda)}(r-C_{III}-\overline M)=0$. But at present we are not
able to construct such an $E(\lambda)$: The ``axiom of choice'' does not
guarantee the $\sigma(\lambda)$-measurability, while the application of
\hyperref[lem:5]{Lemma~5} (which would guarantee it) is impossible because the resulting set
$A$ turns out to be the complement of an analytic set, and not an analytic
set. (Cf. \hyperref[thm:IX]{Theorem IX} and its proof. As to the importance
of this difference cf. \cite[pp.~194--195, 197--202, 267--273]{ref14}.)

\section{Classification of the factor fibres}\label{sec:22}
The classification of the factors will now be applied:

\par\medskip\noindent\textbf{Definition 9.}\phantomsection\label{def:9}\quad
\begin{itshape}
A classification of factors was established in
\cite[Theorem~VIII]{ref10}, distributing the factors $M^*$ in a unitary
space over the following classes:
\[
(I_m,I_n),\quad m,n=1,2,\ldots,\infty;
\qquad (II_\alpha,II_\beta),\quad \alpha,\beta=1,\infty;
\qquad (III,III).
\]
(This is really a classification for the ``coupled factorizations''
$M^*,M^{*'}$ of $\mathfrak H^*$, cf. \cite[Definition~3.1.3]{ref10}, but we
wish to use it for the factors $M^*$ themselves.) Denote the set of these
classes by $\mathfrak C$, and the general element of $\mathfrak C$ by $c$.

Let $M$ be a ring (containing $1$) in $\mathfrak H$. Decompose $M$ with
respect to $P=M\cdot M'$ in the sense of \hyperref[thm:VII]{Theorem VII}:
$M\approx\sum M(\lambda)$. Each $M(\lambda)$ is a factor in its
$\mathfrak H^{(\lambda)}$, and so it belongs to precisely one class
$c\in\mathfrak C$. For a given $c\in\mathfrak C$, denote by $C_c$ the set of
all $\lambda$'s for which $M(\lambda)$ is of class $c$.
\end{itshape}

As a decomposition $M\approx\sum M(\lambda)$ is only determined up to an
arbitrary $\lambda$-set of $\sigma(\lambda)$-measure $0$, therefore we have:

\textit{Each $\lambda$-set $C_c$ ($c\in\mathfrak C$) is only determined up
to an arbitrary $\lambda$-set of $\sigma(\lambda)$-measure $0$.}

The last statement of \secref{sec:19} implies:

\textit{If the representation of $\mathfrak H$ as a generalized direct sum
(with the ring $P=M\cdot M'$) is replaced by an equivalent one, with the
mapping $\lambda\rightleftarrows m(\lambda)$, then each set $C_c$ is replaced
by its image under this mapping. The same is true for the set $\overline M$
of \hyperref[lem:19]{Lemma 19}.}

The following abbreviations suggest themselves:
\[
C_{I_m}=\sum_{n=1,2,\ldots,\infty}C_{(I_m,I_n)},\qquad
C_I=\sum_{m=1,2,\ldots,\infty}C_{I_m}
=\sum_{m,n=1,2,\ldots,\infty}C_{(I_m,I_n)},
\]
\[
\begin{aligned}
C_{II_\alpha}&=\sum_{\beta=1,\infty}C_{(II_\alpha,II_\beta)},\qquad
C_{II}=\sum_{\alpha=1,\infty}C_{II_\alpha}
=\sum_{\alpha,\beta=1,\infty}C_{(II_\alpha,II_\beta)},\\
C_{III}&=C_{(III,III)}.
\end{aligned}
\]

\origpage{468}
We define furthermore two sets $C_f,C_i$:

\begin{itshape}
$C_f$ is the set of all $\lambda$'s for which $1$ (in
$\mathfrak H^{(\lambda)}$) is finite with respect to $M(\lambda)$; $C_i$ is
the set of all $\lambda$ for which it is infinite.
\end{itshape}

Let $d_{M(\lambda)}(E)$ be $\sigma(\lambda)$-measurable, as the constant
operator $1$ (in $\mathfrak H^{(\lambda)}$) is
$\sigma(\lambda)$-measurable, the (numerical) function
$d_{M(\lambda)}(1)$ is $\sigma(\lambda)$-measurable. Thus the set of all
$\lambda$'s with $d_{M(\lambda)}(1)=\infty$ is
$\sigma(\lambda)$-measurable, but this is $C_i$, and with it
$C_f=r-C_i$ is $\sigma(\lambda)$-measurable too. So we have proved:

\begin{itshape}
$C_f$ and $C_i$ are both $\sigma(\lambda)$-measurable.
\end{itshape}

The relations
\[
C_f=\sum_{m=1,2,\ldots}C_{I_m}+C_{II_1},\qquad
C_i=C_{I_\infty}+C_{II_\infty}+C_{III}
\]
are obvious.

Our next objective is to discuss the $\sigma(\lambda)$-measurability of the
sets $C_c$ for all $c\in\mathfrak C$. It is a remarkable feature of this
theory that these $\sigma(\lambda)$-measurabilities are not at all
obvious---in spite of the fact that each set $C_c$ had been obtained by an
effective construction. The reason for this strange phenomenon is that a
direct discussion of these sets would characterize $C_c$, as well as
$C_I+C_{II}=r-C_{III}$ as the projection of a complement of an analytic set
(that is: a ``projective set of the second order''), and the measurability
properties of such sets are not known as yet.\footnote{\phantomsection\label{fn:17}This was written in
1938. In the meantime K. G\"odel has shown that the (Lebesgue-) measurability
of all ``projective sets of the second order'' is indemonstrable in the
conventional system of mathematics. (Oral communication to the author.) The
result is implicit in G\"odel's work on ``The consistency of the Axiom of
choice and of the generalized continuum-Hypothesis with the Axioms of Set
Theory'', Princeton Lecture notes, 1938--1939. The non-measurability in
question follows from the relation ``$V=L$'', cf. loc. cit., which G\"odel
proves to be irrefutable in his system.} (Cf.
\cite[pp.~267--273, 290--292]{ref14}.) We will be able, however, to prove by
an indirect method the $\sigma(\lambda)$-measurability of all $C_c$,
$c\ne(II_\infty,II_\infty),(III,III)$, and of
$C_{(II_\infty,II_\infty)}+C_{(III,III)}$, but we are not able at present to
separate $C_{(II_\infty,II_\infty)}$ and $C_{(III,III)}$!

The auxiliary considerations, which we must carry out for this purpose are
as follows.

\par\medskip\noindent\textbf{Lemma 20.}\phantomsection\label{lem:20}\quad
\begin{itshape}
Let $M\approx\sum M(\lambda)$ be the decomposition of $M$ with respect to
any Abelian $P$. (Cf. \hyperref[thm:VII]{Theorem VII}.) Let $M^*$ be a ring
(containing $1$) in a unitary space $\mathfrak H^*$. Denote by $D_{M^*}$ the
set of all those $\lambda$'s for which an isomorphic mapping $J^*$ of
$\mathfrak H^*$ on $\mathfrak H^{(\lambda)}$ exists, which carries $M^*$
into $M(\lambda)$. Then the set $D_{M^*}$ is
$\sigma(\lambda)$-measurable.
\end{itshape}

\par\smallskip\noindent\textit{Proof.}
Let $p$ be the dimensionality of $\mathfrak H^*$,
$p=1,2,\ldots,\infty$. Then $\lambda\in D_{M^*}$ necessitates
$k(\lambda)=p$, $\lambda\in\Delta_p$. Proceed now as in the proofs of
\hyperref[lem:6]{Lemmas 6}, \hyperref[lem:17]{17}, and
\hyperref[lem:18]{18}. Let $\psi_1^*,\psi_2^*,\ldots$ be a complete
normalized orthogonal set in $\mathfrak H^*$, $\varphi_m^{(\lambda)}$ a
measurable family, $J^{(\lambda)}$ an isomorphism of
$\mathfrak H^{(\lambda)}$ on $\mathfrak H^*$ in which
$\varphi_m^{(\lambda)}$ corresponds to
$\psi_m^*$. Let the $A_q(\lambda)$, $q=1,2,\ldots$, be
$\sigma(\lambda)$-measurable, such that $A_q(\lambda)$, $q=1,2,\ldots$,
generate the ring $M(\lambda)$ and the $B_q(\lambda)$,
$q=1,2,\ldots$, the same for $M(\lambda)'$. Let the
$\overline A_1,\overline A_2,\ldots$ be operators in $\mathfrak H^*$
generating the ring $M^*$, and

\origpage{469}
$\overline B_1,\overline B_2,\ldots$ the ring $M^{*'}$. Choose a set
$N_p\subset\Delta_p$ of $\sigma(\lambda)$-measure $0$, as in the proofs
quoted above, such that $\Delta_p'=\Delta_p-N_p$ is a Borel-set, and the
$(A_q(\lambda)\varphi_m(\lambda),\varphi_n(\lambda))$,
$(B_q(\lambda)\varphi_m(\lambda),\varphi_n(\lambda))$,
$q,m,n=1,2,\ldots$, are Baire-functions.

Any isomorphic mapping $J^*$ of $\mathfrak H^*$ on
$\mathfrak H^{(\lambda)}$ may be put in the form $J^{(\lambda)-1}U$ where
$U$ is an isomorphic mapping of $\mathfrak H^*$ on itself, that is a
unitary operator in $\mathfrak H^*$.

That $J^*$ carries $M^*$ into $M(\lambda)$ means that it carries $M^*$ into
a subset of $M(\lambda)$ and at the same time $M^{*'}$ into a subset of
$M(\lambda)'$ (because the latter means that $M^*$ is mapped in a set
containing $M(\lambda)$), or: that the $J^*$-image of $M^*$ commutes with
$M(\lambda)'$ and that the $J^*$-image of $M^{*'}$ commutes with
$M(\lambda)$. This is equivalent to this: that the $J^*$-image of every
$\overline A_r$ commutes with every $B_q(\lambda)$, and similarly for every
$\overline B_r$ and $A_q(\lambda)$, $r,q=1,2,\ldots$. If we now replace
$J^*$ by $J^{(\lambda)-1}U$ then this means:
$U\overline A_rU^{-1}$ and
$J^{(\lambda)}B_q(\lambda)J^{(\lambda)-1}$ commute,
$U\overline B_rU^{-1}$ and
$J^{(\lambda)}A_q(\lambda)J^{(\lambda)-1}$ commute,
$r,q=1,2,\ldots$. Thus $\lambda\in D_{M^*}$ is equivalent to the following
equations (the last line expresses the unitarity of $U$, and therefore it was
permissible to replace $U^{-1}$ by $U^*$ in the previous ones):
\begin{equation}
\left\{
\begin{aligned}
U\overline A_rU^*J^{(\lambda)}B_q(\lambda)J^{(\lambda)-1}
 &=J^{(\lambda)}B_q(\lambda)J^{(\lambda)-1}U\overline A_rU^*,\\
U\overline B_rU^*J^{(\lambda)}A_q(\lambda)J^{(\lambda)-1}
 &=J^{(\lambda)}A_q(\lambda)J^{(\lambda)-1}U\overline B_rU^*,
 &&r,q=1,2,\ldots,\\
U^*U&=UU^*=1.
\end{aligned}
\right.
\tag{22.1}\label{eq:22.1}
\end{equation}
Observe that such a $U$ is unitary, and so
$\lVert\!\lvert U\rvert\!\rVert\leq1$.

Consider now the set $C_p$ of all pairs $(\lambda,U)$ in $r\times S_p$, for
which $\lambda\in\Delta_p'$ and $\lambda,U$ fulfill \eqref{eq:22.1}. One
proves by the same computations as in the proofs of
\hyperref[lem:17]{Lemmas 17} and \hyperref[lem:18]{18}, that $C_p$ is a
Borel-set in $r\times S_p$. $\lambda\in\Delta_p'D_{M^*}$ means that there
exists a $U$ with $(\lambda,U)\in C_p$. From this we infer as in
\hyperref[lem:17]{Lemmas 17} and \hyperref[lem:18]{18} that
$\Delta_p'D_{M^*}$ is $\sigma(\lambda)$-measurable.

As $D_{M^*}\subset\Delta_p$, therefore
$D_{M^*}-\Delta_p'D_{M^*}=N_pD_{M^*}\subset N_p$ has the
$\sigma(\lambda)$-measure $0$, therefore $D_{M^*}$ is
$\sigma(\lambda)$-measurable along with $\Delta_p'D_{M^*}$. Thus the proof
is completed.

We are now in the position to prove:

\par\medskip\noindent\textbf{Theorem IX.}\phantomsection\label{thm:IX}\quad
\begin{itshape}
The following sets are $\sigma(\lambda)$-measurable:
$C_{(I_m,I_n)}$, $m,n=1,2,\ldots,\infty$;
$C_{(II_\alpha,II_\beta)}$, $\alpha,\beta=1,\infty$, excepting however,
$\alpha=\beta=\infty$;
$C_{(II_\infty,II_\infty)}+C_{(III,III)}$. The question remains, however,
open for $C_{(II_\infty,II_\infty)}$ and $C_{(III,III)}$ separately.
\end{itshape}

\par\smallskip\noindent\textit{Proof.}
Choose two $m,n=1,2,\ldots,\infty$, an $mn$-dimensional unitary space
$\mathfrak H^*$, and in it a coupled factorization $M^*,M^{*'}$ of class
$(I_m,I_n)$. (Cf. \cite[Definitions~3.1.1--3.1.3 and
Theorem~VIII]{ref10}.) We know from \cite[Lemma~8.6.1]{ref10} that
$M(\lambda)$ in $\mathfrak H^{(\lambda)}$ is of class $(I_m,I_n)$ if and
only if an isomorphic mapping $J^*$ of $\mathfrak H^*$ on
$\mathfrak H^{(\lambda)}$ exists, which carries $M^*$ into $M(\lambda)$.
Thus we have by \hyperref[lem:20]{Lemma 20}
$C_{(I_m,I_n)}=D_{M^*}$, and this proves that it is
$\sigma(\lambda)$-measurable.

As we proved at the beginning of this section, $C_f$ is
$\sigma(\lambda)$-measurable too. As
\[
C_f=\sum_{m,n=1,2,\ldots,\infty}C_{(I_m,I_n)}
 +C_{(II_1,II_1)}+C_{(II_1,II_\infty)},
\]
these results imply the $\sigma(\lambda)$-measurability of
$C_{(II_1,II_1)}+C_{(II_1,II_\infty)}$. If we replace $M$ by $M'$, then

\origpage{470}
$M(\lambda)$ and $M(\lambda)'$ are interchanged, and so
$C_{(II_1,II_1)},C_{(II_1,II_\infty)}$ become
$C_{(II_1,II_1)},C_{(II_\infty,II_1)}$. Thus
$C_{(II_1,II_1)}+C_{(II_\infty,II_1)}$ is
$\sigma(\lambda)$-measurable, too. Now we see that
\[
\bigl(C_{(II_1,II_1)}+C_{(II_1,II_\infty)}\bigr)
\bigl(C_{(II_1,II_1)}+C_{(II_\infty,II_1)}\bigr)
=C_{(II_1,II_1)}
\]
is $\sigma(\lambda)$-measurable, and with it
\[
\bigl(C_{(II_1,II_1)}+C_{(II_1,II_\infty)}\bigr)-C_{(II_1,II_1)}
=C_{(II_1,II_\infty)},
\]
\[
\bigl(C_{(II_1,II_1)}+C_{(II_\infty,II_1)}\bigr)-C_{(II_1,II_1)}
=C_{(II_\infty,II_1)}.
\]

We have established the $\sigma(\lambda)$-measurability of all $C_c$ with
$c\ne(II_\infty,II_\infty)$ and $(III,III)$. This implies the
$\sigma(\lambda)$-measurability of
\[
r-\sum_{\substack{c\in\mathfrak C,\;
c\ne(II_\infty,II_\infty),(III,III)}}C_c
=C_{(II_\infty,II_\infty)}+C_{(III,III)}
\]
too. Thus we have proved all statements of our theorem.

A direct treatment of $C_{(II_\infty,II_\infty)}$ or $C_{(III,III)}$ fails
because it leads to projections of complements of analytic sets (that is:
projections of complements of projections of Borel-sets), for which the
problem of measurability is yet unsolved. (Cf. the remarks preceding
\hyperref[lem:20]{Lemma 20}.)

The difficulties in proving the $\sigma(\lambda)$-measurability of
$C_{(II_\infty,II_\infty)}$ and $C_{(III,III)}$ are remarkable under several
aspects. We do not know how real they are: $(III,III)$ is the only class,
which might be empty (cf. \cite[Theorem~XII]{ref10}), and if this were the
case, we would have
$C_{(II_\infty,II_\infty)}=
C_{(II_\infty,II_\infty)}+C_{(III,III)}$ and $C_{(III,III)}=0$, so that both
sets would be $\sigma(\lambda)$-measurable.\footnote{\phantomsection\label{fn:18}This was written in
1938. In the meantime the existence of factors belonging to the class
$(III,III)$ has been established. Cf. J. v. Neumann: ``On rings of
operators. III'', \textit{Annals of Math.}, Vol.~41, pp.~94--161, 1940. Cf.
in particular Theorem~IX and \S4.4, loc. cit.} On the other hand, if we had
simple normal forms for factors of class $(II_\infty,II_\infty)$ or
$(III,III)$, as \cite[Lemma~8.6.1]{ref10} provided some for the classes
$(I_m,I_n)$, then we could use \hyperref[lem:20]{Lemma 20}, to prove the
$\sigma(\lambda)$-measurability of $C_{(II_\infty,II_\infty)}$ resp.
$C_{(III,III)}$, and thus of both. This connects with
\cite[Problem~6 (following Theorem~VIII)]{ref10}.

The connection of these sets with the set $\overline M$ of
\hyperref[lem:19]{Lemma 19} is of interest, too. We have already proved
$\overline M\subset r-C_{III}=r-C_{(III,III)}$. Consider now
$E(\lambda)=1$, for this $M_r=C_f$, and so by \hyperref[lem:19]{Lemma 19}
$\mu_{\sigma(\lambda)}(C_f-C_f\overline M)=0$. In other words: If we
disregard a set of $\sigma(\lambda)$-measure $0$, then
$\overline M\supset C_f$. As $\overline M$ is only defined up to such a set,
we may assume $\overline M\supset C_f$. One can prove furthermore that
$\overline M\supset C_{(I_\infty,I_\infty)}$, and that $\overline M$ is
unchanged, if we replace $M$ by $M'$. We will not give the details here. Now
we have $\overline M\supset C_f\supset C_{(I_m,I_\infty)}$,
$m=1,2,\ldots$, and $\supset C_{(II_1,II_\infty)}$; and as we can replace
$M$ by $M'$ without affecting $\overline M$, but thereby changing
$C_{(I_m,I_\infty)}$ into $C_{(I_\infty,I_m)}$, and
$C_{(II_1,II_\infty)}$ into $C_{(II_\infty,II_1)}$, this implies
$\overline M\supset C_{(I_\infty,I_m)}$, $m=1,2,\ldots$, and
$\supset C_{(II_\infty,II_1)}$. Thus we have
\[
\begin{aligned}
\overline M\supset{}&C_f+C_{(I_\infty,I_1)}+C_{(I_\infty,I_2)}+\cdots
+C_{(I_\infty,I_\infty)}+C_{(II_\infty,II_1)}\\
={}&r-\bigl(C_{(II_\infty,II_\infty)}+C_{(III,III)}\bigr).
\end{aligned}
\]
So we have
\[
r-\bigl(C_{(II_\infty,II_\infty)}+C_{(III,III)}\bigr)
\subset\overline M\subset r-C_{(III,III)}.
\]

\origpage{471}
The two first sets are $\sigma(\lambda)$-measurable, while we do not know
this about the third one; but we saw that there are reasons to surmise that
the two last sets are identical, and they are at any rate similar under
several aspects.

\section{Integration of weight functions}\label{sec:23}
Our information about the weight functions $\Delta_{M(\lambda)}(E)$ of the
rings $M(\lambda)$ is now sufficiently complete to make it possible to
discuss the weight function $\Delta_M(E)$ of $M$. We do, however, need one
more lemma concerning the $M(\lambda)$, before we attack the main problem
directly.

\par\medskip\noindent\textbf{Lemma 21.}\phantomsection\label{lem:21}\quad
\begin{itshape}
Let $M$ be a ring (containing $1$) in $\mathfrak H$. Decompose $M$ with
respect to $P=M\cdot M'$ in the sense of \hyperref[thm:VII]{Theorem VII}:
$M\approx\sum M(\lambda)$ so that each $M(\lambda)$ is a factor in its
$\mathfrak H^{(\lambda)}$.

Consider two projections $E,F\in M$ and their decompositions:
$E\sim\sum E(\lambda)$, $F\sim\sum F(\lambda)$. Then there exist three
partially isometric operators $V_j\in M$, $j=1,2,3$, with the decompositions
$V_j\sim\sum V_j(\lambda)$, and with the following properties:
\begin{enumerate}[label=(\roman*)]
\item If $\lambda>-\infty$, $<+\infty$, and if there exists for this
$\lambda$ a partially isometric $V_1\in M(\lambda)$ with the initial
projection $E(\lambda)$ and with the final projection $F(\lambda)$, then
$V_1(\lambda)$ is such a $V_1$. Otherwise $V_1(\lambda)=0$.

\item If $\lambda>-\infty$, $<+\infty$, and if there exists for this
$\lambda$ a partially isometric $V_2\in M(\lambda)$ with the initial
projection $E(\lambda)$ and with a final projection $\leq F(\lambda)$, then
$V_2(\lambda)$ is such a $V_2$. Otherwise $V_2(\lambda)=0$.

\item If $\lambda>-\infty$, $<+\infty$, and if there exist for this
$\lambda$ two partially isometric $V_3,V_3'\in M(\lambda)$, both with the
initial projection $E(\lambda)$, and with final projections which have the
sum $E(\lambda)$, then $V_3(\lambda)$ is such a $V_3$. Otherwise
$V_3(\lambda)=0$.
\end{enumerate}
\end{itshape}

\par\smallskip\noindent\textit{Proof.}
We first construct $\sigma(\lambda)$-measurable $V_j(\lambda)$'s,
$j=1,2,3$, which satisfy respectively (i), (ii), (iii). It suffices, however,
to do this within each set $\Delta_p$, defined by $k(\lambda)=p$, for each
$p=1,2,\ldots$---as we can then put together all these solutions, owing to
$\sum_{p=1}^{\infty}\Delta_p=r$. Furthermore we may replace $\Delta_p$ by
any $\Delta_p'=\Delta_p-N_p$ if $\mu_{\sigma(\lambda)}(N_p)=0$: Because we
can define the $V_j(\lambda)$ for $\lambda\in N_p$ by the axiom of choice, as
$\mu_{\sigma(\lambda)}(N_p)=0$, this does not affect their
$\sigma(\lambda)$-measurability.

Now choose $\sigma(\lambda)$-measurable operators $A_q(\lambda)$,
$q=1,2,\ldots$, generating $M(\lambda)'$. For each
$p=1,2,\ldots,\infty$ proceed as in the proofs of
\hyperref[lem:17]{Lemmas 17}, \hyperref[lem:18]{18}, and
\hyperref[lem:20]{20}: Choose $N_p$ as a set with
$\mu_{\sigma(\lambda)}(N_p)=0$, such that $\Delta_p'=\Delta_p-N_p$ is a
Borel-set, and the (numerical) functions
$(E(\lambda)\varphi_m(\lambda),\varphi_n(\lambda))$,
$(F(\lambda)\varphi_m(\lambda),\varphi_n(\lambda))$,
$(A_q(\lambda)\varphi_m(\lambda),\varphi_n(\lambda))$,
$q,m,n=1,2,\ldots$, are Baire-functions in $\Delta_p'$. As there, introduce
a $p$-dimensional unitary space $\mathfrak H_p$, in it a complete normalized
orthogonal set $\psi_1',\psi_2',\ldots$, and (for
$\lambda\in\Delta_p$) an isomorphism $J^{(\lambda)}$ of
$\mathfrak H^{(\lambda)}$ on $\mathfrak H_p$ which carries
$\varphi_m(\lambda)$ into $\psi_m'$.

Under these conditions our $V_j(\lambda)$'s are characterized by
\begin{equation}
\begin{gathered}
(j=1)\qquad V^*V=E(\lambda),\quad VV^*=F(\lambda),\\
VA_q(\lambda)=A_q(\lambda)V,\quad
VA_q(\lambda)^*=A_q(\lambda)^*V\quad(q=1,2,\ldots),
\end{gathered}
\tag{23.1$_1$}\label{eq:23.1.1}
\end{equation}
and
\begin{equation}
\begin{gathered}
(j=2)\qquad V^*V=E(\lambda),\quad VV^*F(\lambda)=VV^*,\\
VA_q(\lambda)=A_q(\lambda)V,\quad
VA_q(\lambda)^*=A_q(\lambda)^*V\quad(q=1,2,\ldots),
\end{gathered}
\tag{23.1$_2$}\label{eq:23.1.2}
\end{equation}
and

\origpage{472}
\begin{equation}
\left\{
\begin{aligned}
(j=3)\qquad&V^*V=V'^*V'=E(\lambda),\quad
VV^*+V'V'^*=F(\lambda),\\
&VA_q(\lambda)=A_q(\lambda)V,\quad
VA_q(\lambda)^*=A_q(\lambda)^*V,\\
&V'A_q(\lambda)=A_q(\lambda)V',\quad
V'A_q(\lambda)^*=A_q(\lambda)^*V',
&&q=1,2,\ldots,
\end{aligned}
\right.
\tag{23.1$_3$}\label{eq:23.1.3}
\end{equation}
respectively. (We write $V$ for $V_j(\lambda)$, $j=1,2,3$, and $V'$ for
$V_3'(\lambda)$.)

Put now $\widetilde V=J^{(\lambda)}VJ^{(\lambda)-1}$,
$\widetilde V'=J^{(\lambda)}V'J^{(\lambda)-1}$, then
$\widetilde V,\widetilde V'$ are operators in $\mathfrak H_p$, and if they
satisfy the equivalent of any of \eqref{eq:23.1.1},
\eqref{eq:23.1.2}, and \eqref{eq:23.1.3}, then they are partially
isometric, implying
$\lVert\!\lvert\widetilde V\rvert\!\rVert\leq1$,
$\lVert\!\lvert\widetilde V'\rvert\!\rVert\leq1$. Denote the sets of all
pairs $(\lambda,\widetilde V)$ and triplets
$(\lambda,\widetilde V,\widetilde V')$ fulfilling
$\lambda\in\Delta_p'$ and \eqref{eq:23.1.1}, \eqref{eq:23.1.2}, and
\eqref{eq:23.1.3}, respectively by $G_p^1,G_p^2$ and $G_p^3$. We prove as in
the proofs of \hyperref[lem:17]{Lemmas 17}, \hyperref[lem:18]{18} and
\hyperref[lem:20]{20} that these are respectively Borel-sets in
$r\times S_p$ and $r\times S_p\times S_p$. Now apply \hyperref[lem:5]{Lemma~5} to
$S=r\times S_p$ and $r\times S_p\times S_p$, $A=G_p^1,G_p^2$ and $G_p^3$
and $F((\lambda,V))=\lambda$ and
$F((\lambda,V,V'))=\lambda$ respectively. $K_j$ and $f_j(\lambda)$,
$j=1,2,3$, will result. For $j=1,2$,
$f_j(\lambda)=(\lambda,\widetilde V_j(\lambda))\in r\times S_p$; for $j=3$,
$f_3(\lambda)=(\lambda,\widetilde V_3(\lambda),
\widetilde V_3'(\lambda))\in r\times S_p\times S_p$, considering \hyperref[lem:5]{Lemma~5},
(ii)$_2$. The set of all $\widetilde V$ with
$\RePart(\widetilde V\psi_m',\psi_n')>\alpha$ and similarly that one with
$\ImPart(\widetilde V\psi_m',\psi_n')>\alpha$ is clearly an open set in
$S=r\times S_p$ resp. $r\times S_p\times S_p$. Thus the sets of all
$\lambda$'s with
$\RePart(\widetilde V_j(\lambda)\psi_m',\psi_n')>\alpha$ or
$\ImPart(\widetilde V_j(\lambda)\psi_m',\psi_n')>\alpha$ are
$\sigma(\lambda)$-measurable, owing to \hyperref[lem:5]{Lemma~5}, (ii)$_1$. Put
$V_j^0(\lambda)=J^{(\lambda)-1}\widetilde V_j(\lambda)J^{(\lambda)}$, then
the sets of all $\lambda$'s with
$\RePart(V_j^0(\lambda)\varphi_m(\lambda),\varphi_n(\lambda))>\alpha$ or with
$\ImPart(V_j^0(\lambda)\varphi_m(\lambda),\varphi_n(\lambda))>\alpha$ are
$\sigma(\lambda)$-measurable, and thus $V_j^0(\lambda)$ is
$\sigma(\lambda)$-measurable. Define now
\[
V_j(\lambda)=
\begin{cases}
V_j^0(\lambda)&\text{for }\lambda\in K_j,\\
0&\text{for }\lambda\notin K_j,
\end{cases}
\qquad j=1,2,3.
\]
As $K_j$ is $\sigma(\lambda)$-measurable by \hyperref[lem:5]{Lemma~5}, (i), the
$V_j(\lambda)$ are $\sigma(\lambda)$-measurable too.

One verifies immediately that the $V_j(\lambda)$, $j=1,2,3$, fulfill the
corresponding conditions (i), (ii), (iii) of our Lemma. Thus we have the
desired $V_j(\lambda)$'s in each $\Delta_p'$, and therefore in $r$ too. As
$V_j(\lambda)\in M(\lambda)$ and as it depends
$\sigma(\lambda)$-measurably on $\lambda$, we can form
$V_j\sim\sum V_j(\lambda)$ by \hyperref[def:6]{Definition~6}. As all $V_j(\lambda)$'s are
partially isometric, the $V_j$ are also partially isometric. Thus the proof
of our lemma is completed.---

Consider now a weight function $\Delta_M(E)$ with respect to $M$. For every
$\sigma(\lambda)$-measurable set $K$ form the function
\[
1_K(\lambda)=
\begin{cases}
1&\text{for }\lambda\in K,\\
0&\text{for }\lambda\notin K,
\end{cases}
\]
and the operator $1_K$ (cf. \secref{sec:7}). As
$1_K(\lambda)^2=1_K(\lambda)=\overline{1_K(\lambda)}$, therefore $1_K$ is a
projection, and by definition (in \secref{sec:11})
$1_K\in P=M\cdot M'$. Thus for every $E\in M$, $E$ and $1_K$ commute, and
therefore $1_K E$ is a projection along with $E$. So we can form
$\Delta_M(1_KE)$. We write provisionally
\[
m(K)=\Delta_M(1_KE).
\]
We enumerate these properties of $m(K)$:
\begin{enumerate}[label=(\alph*)]
\item $m(K)$ is defined for every $\sigma(\lambda)$-measurable $K$, its
values being $\geq0$, $\leq+\infty$.

\origpage{473}
\item If a sequence $K_1,K_2,\ldots$ with $K_jK_k=0$ for $j\ne k$ is given,
then $1_{K_1+K_2+\cdots}(\lambda)=1_{K_1}(\lambda)+1_{K_2}(\lambda)+\cdots$,
and so by \hyperref[def:7]{Definition 7}, (ii),
\[
m(K_1+K_2+\cdots)=m(K_1)+m(K_2)+\cdots.
\]

\item If $\mu_{\sigma(\lambda)}(K)=0$ then $1_K(\lambda)$ differs from $0$
only on a $\lambda$-set of $\sigma(\lambda)$-measure $0$, so $1_K=0$, and
thus $m(K)=0$.
\end{enumerate}

Consider those sets $K$ for which $m(K)$ is finite, and let $\alpha$ be the
l.u.b. of their $\mu_{\sigma(\lambda)}(K)$'s. For each $j=1,2,\ldots$ choose
such a $K_j$ with $\mu_{\sigma(\lambda)}(K_j)\geq\alpha-1/j$. If $m(K)$ is
finite and $K'\subset K$, then $(\beta)$ gives
$m(K')+m(K-K')=m(K)$, and so $m(K')$ is finite too; if $m(K')$, $m(K'')$ are
finite, then $m(K'-K'K'')$ is finite, and so by $(\beta)$
\[
\begin{aligned}
m(K'+K'')&=m(K'+(K''-K'K''))\\
&=m(K')+m(K''-K'K'').
\end{aligned}
\]
is finite too. Thus if we define
$L_j=K_j-(K_1+\cdots+K_{j-1})K_j$, then $m(L_j)$ is finite too; besides
$L_jL_k=0$ for $j\ne k$, and $L_1+L_2+\cdots=K_1+K_2+\cdots$. If $m(K)$ is
finite, then $m(K_j+K)$ is finite too, so
$\mu_{\sigma(\lambda)}(K_j+K)\leq\alpha$. Therefore
\[
\begin{aligned}
\mu_{\sigma(\lambda)}(K-K_jK)
 &=\mu_{\sigma(\lambda)}((K_j+K)-K_j)\\
 &=\mu_{\sigma(\lambda)}(K_j+K)-\mu_{\sigma(\lambda)}(K_j)\\
 &\leq\alpha-(\alpha-1/j)=1/j.
\end{aligned}
\]
and \textit{a fortiori}
\[
\begin{aligned}
\mu_{\sigma(\lambda)}(K-(L_1+L_2+\cdots)K)
 &=\mu_{\sigma(\lambda)}(K-(K_1+K_2+\cdots)K)\\
 &\leq\mu_{\sigma(\lambda)}(K-K_jK)\leq1/j.
\end{aligned}
\]
As this holds for every $j=1,2,\ldots$, it gives
$\mu_{\sigma(\lambda)}(K-(L_1+L_2+\cdots)K)=0$. In particular: if
$K\subset r-(L_1+L_2+\cdots)=L_\infty$ and $m(K)$ finite, then
$\mu_{\sigma(\lambda)}(K)=0$, and thus $m(K)=0$. In other words:

\noindent$(\delta)$ For $K\subset L_\infty$ either
$\mu_{\sigma(\lambda)}(K)=0,m(K)=0$, or
$\mu_{\sigma(\lambda)}(K)>0,m(K)=+\infty$.

Consider now $m(L_jK)$ for $j=1,2,\ldots$. It again fulfills $(\alpha)$,
$(\beta)$, $(\gamma)$, but as
$m(L_jK)\leq m(L_jK)+m(L_j-L_jK)=m(L_j)<+\infty$, more than $(\alpha)$ holds:

\noindent$(\alpha')$ $m(L_jK)$ is defined for every
$\sigma(\lambda)$-measurable $K$, its values being $\geq0$,
$\leq m(L_j)<+\infty$.

As $m(L_jK)$ fulfills $(\alpha')$, $(\beta)$, $(\gamma)$, it is a
non-negative and finite totally additive set-function, and absolutely
continuous with respect to $\mu_{\sigma(\lambda)}(K)$. Therefore the theorem
which holds for such functions applies to it: There exists a
$\sigma(\lambda)$-measurable function $w_j(\lambda)\geq0$, $<+\infty$, such
that
\begin{equation}
m(L_jK)=\int_K w_j(\lambda)\,d\sigma(\lambda)
\tag{23.2}\label{eq:23.2}
\end{equation}
for every $\sigma(\lambda)$-measurable $K$. (Cf.
\cite[Chapter~IX]{ref13}. The theorem is formulated there for common
Lebesgue-measure only, but this restriction is unessential.) Equation
\eqref{eq:23.2} is proved for $j=1,2,\ldots$, but it holds for $j=\infty$
too, if we put
\[
w_\infty(\lambda)=
\begin{cases}
+\infty&\text{for }\lambda\in L_\infty,\\
0&\text{for }\lambda\notin L_\infty,
\end{cases}
\]
considering that we know by $(\delta)$ that
\[
m(L_\infty K)=
\begin{cases}
+\infty&\text{for }\mu_{\sigma(\lambda)}(L_\infty K)>0,\\
0&\text{for }\mu_{\sigma(\lambda)}(L_\infty K)=0.
\end{cases}
\]
Now put
$\delta_M^{(\lambda)}(E)=\sum_{j=1}^{\infty}w_j(\lambda)$ (we

\origpage{474}
indicate in this way the dependence on $\lambda,M,E$), then summation of
\eqref{eq:23.2} over $j=1,2,\ldots,\infty$ gives, considering $(\beta)$:
\begin{equation}
m(K)=\Delta_M(1_K(\lambda)E)
=\int_K\delta_M^{(\lambda)}(E)\,d\sigma(\lambda).
\tag{23.3}\label{eq:23.3}
\end{equation}
For $E=0$ we may redefine $\delta_M^{(\lambda)}(0)=0$, without violating
\eqref{eq:23.3}.

We now apply \hyperref[def:7]{Definition 7}, (ii), (iii) to
\eqref{eq:23.3}. We use (ii) for two addends $E_1=E,E_2=F$ only (while
$E_3=E_4=\cdots=0$), and remember in (iii) that $1_K(\lambda)\in M\cdot M'$,
$U\in M$ imply $U1_K(\lambda)U^{-1}=1_K(\lambda)$. Thus we obtain:
\[
\int_K\delta_M^{(\lambda)}(E+F)\,d\sigma(\lambda)
=\int_K\delta_M^{(\lambda)}(E)\,d\sigma(\lambda)
 +\int_K\delta_M^{(\lambda)}(F)\,d\sigma(\lambda)
\quad\text{for }EF=0,
\]
\[
\int_K\delta_M^{(\lambda)}(UEU^{-1})\,d\sigma(\lambda)
=\int_K\delta_M^{(\lambda)}(E)\,d\sigma(\lambda),
\]
for every $\sigma(\lambda)$-measurable set $K$. Therefore:
\begin{equation}
\left\{
\begin{aligned}
\delta_M^{(\lambda)}(E+F)
 &=\delta_M^{(\lambda)}(E)+\delta_M^{(\lambda)}(F),&&EF=0,\\
\delta_M^{(\lambda)}(UEU^{-1})&=\delta_M^{(\lambda)}(E),
\end{aligned}
\right.
\tag{23.4}\label{eq:23.4}
\end{equation}
except for certain $\lambda$-sets of $\sigma(\lambda)$-measure $0$, which we
will denote by $N_1(E,F)$ resp. $N_2(E,U)$.

We need this additional Lemma\footnote{It is possible to avoid the use of
this Lemma in the present context, but the Lemma has a certain interest of
its own.}:

\par\medskip\noindent\textbf{Lemma 22:}\phantomsection\label{lem:22}\quad
\begin{itshape}
Given two projections $E,F$, form
\[
G_n=F\text{ or }E\cdots EFE
\]
($n$ factors, the first factor is $F$ or $E$ if $n$ is even or odd,
respectively), $n=1,2,\ldots$. Then
\begin{equation}
G=\operatorname*{strong\,lim}_{n\to\infty}G_n
\tag{23.5}\label{eq:23.5}
\end{equation}
exists. This $G$ has the following properties:
\begin{enumerate}[label=(\roman*)]
\item \phantomsection\label{lem:22.i}$G$ is a projection.
\item \phantomsection\label{lem:22.ii}For any projection $H$, $H\leq G$ is equivalent to the conjunction of
$H\leq E$ and $H\leq F$.
\item \phantomsection\label{lem:22.iii}The closed linear set of $G$ is the intersection of the closed linear
sets of $E$ and of $F$.
\end{enumerate}
Denote this $G$ by $E\mathbin{\#}F$.
\end{itshape}

\par\smallskip\noindent\textit{Proof:}
We note:
\[
\tag{I}\label{eq:23.I} G_n\text{ is Hermitean if $n$ is odd.}
\]
\[
\tag{II}\label{eq:23.II} G_m^*G_n=G_p,\quad\text{where $p=m+n$ or $m+n-1$, whichever is odd.}
\]
\[
\tag{III}\label{eq:23.III}G_{n+1}=
\begin{cases}
EG_n&\text{for $n$ even},\\
FG_n&\text{for $n$ odd}.
\end{cases}
\]

\origpage{475}
From \hyperref[eq:23.II]{(II)} with $m=n$
\[
\tag{IV}\label{eq:23.IV}\lVert G_nf\rVert^2=(G_{2n-1}f,f).
\]
\hyperref[eq:23.III]{(III)} implies $\lVert G_{n+1}f\rVert\leq\lVert G_nf\rVert$, hence \hyperref[eq:23.IV]{(IV)} gives
$(G_{2n+1}f,f)\leq(G_{2n-1}f,f)\geq0$, i.e.
\[
\tag{V}\label{eq:23.V}(G_1f,f)\geq(G_3f,f)\geq(G_5f,f)\geq\cdots\geq0.
\]
Consequently,
\[
\tag{VI}\label{eq:23.VI}\lim_{n\to\infty}(G_{2n-1}f,f)\text{ exists and is finite.}
\]
Now by \hyperref[eq:23.II]{(II)}
\[
\begin{aligned}
\lVert G_mf-G_nf\rVert^2
 &=\lVert(G_m-G_n)f\rVert^2\\
 &=((G_m-G_n)^*(G_m-G_n)f,f)\\
 &=((G_{2m-1}-2G_p+G_{2n-1})f,f)
\end{aligned}
\]
($p$ according to \hyperref[eq:23.II]{(II)}). Hence $m,n\to\infty$ and \hyperref[eq:23.VI]{(VI)} give
\[
\tag{VII}\label{eq:23.VII}\lim_{m,n\to\infty}\lVert G_mf-G_nf\rVert=0.
\]
This proves \eqref{eq:23.5}.

We now prove \hyperref[lem:22.i]{(i)}--\hyperref[lem:22.iii]{(iii)}:

$n\to\infty$ through even values and through odd values in $n$ gives
\[
\tag{VIII}\label{eq:23.VIII}G=EG=FG.
\]
If $A=EA=FA$, then $A=G_nA$, hence $A=GA$:
\[
\tag{IX}\label{eq:23.IX}A=EA=FA\quad\hbox{implies}\quad A=GA.
\]
Now $n\to\infty$ in \hyperref[eq:23.I]{(I)} shows that $G$ is Hermitean and \hyperref[eq:23.VIII]{(VIII)} with \hyperref[eq:23.IX]{(IX)}
for $A=G$ gives $G^2=G$, hence $G$ is a projection. This proves \hyperref[lem:22.i]{(i)}.

Now \hyperref[eq:23.VIII]{(VIII)} with \hyperref[eq:23.IX]{(IX)} for projections $A$ states precisely \hyperref[lem:22.ii]{(ii)}.

Stating \hyperref[lem:22.ii]{(ii)} for the closed linear sets of the operators involved coincides
with \hyperref[lem:22.iii]{(iii)}.---

We infer immediately from \hyperref[lem:22]{Lemma 22}:

\noindent(a) $E,F\in M$ imply $E\mathbin{\#}F\in M$,

and replacing $\mathfrak H,M$ by $\mathfrak H^{(\lambda)},M(\lambda)$:

\noindent(b) $E(\lambda),F(\lambda)\in M(\lambda)$ imply
$E(\lambda)\mathbin{\#}F(\lambda)\in M(\lambda)$,

\noindent(c) if $E(\lambda),F(\lambda)$ are $\sigma(\lambda)$-measurable,
then $E(\lambda)\mathbin{\#}F(\lambda)$ is also
$\sigma(\lambda)$-measurable.

Now put $E\sim\sum E(\lambda)$, $F\sim\sum F(\lambda)$, and form
$E_1\sim\sum E(\lambda)\mathbin{\#}F(\lambda)$. Applying
\hyperref[lem:22]{Lemma 22} to $E(\lambda),F(\lambda)$ and to $E,F$ then
gives this: Consider a projection $H\in M$. Put
$H\sim\sum H(\lambda)$. Then $H\leq E\mathbin{\#}F$ means $H\leq E$ and
$H\leq F$, hence $H(\lambda)\leq E(\lambda)$ and
$H(\lambda)\leq F(\lambda)$ (except for a $\lambda$-set of
$\sigma(\lambda)$-measure $0$), hence
$H(\lambda)\leq E(\lambda)\mathbin{\#}F(\lambda)$ (except as above), hence
$H\leq E_1$. Now $E_1\in M$ by

\origpage{476}
its definition and $E\mathbin{\#}F\in M$ by (a). Therefore, the equivalence
of $H\leq E\mathbin{\#}F$ and of $H\leq E_1$ for all projections $H\in M$
implies $E\mathbin{\#}F=E_1$. We restate this

\noindent(d) $E\mathbin{\#}F\sim\sum
E(\lambda)\mathbin{\#}F(\lambda)$.

The following further construction is necessary:

For every $A\in M$ form an operator $A_u\in M$ as follows: Form
$A\sim\sum A(\lambda)$, then $A(\lambda)\in M(\lambda)$ and
$\sigma(\lambda)$-measurable. Thus $1-A(\lambda)^*A(\lambda)$ and
$1-A(\lambda)A(\lambda)^*$ are $\sigma(\lambda)$-measurable too, and
therefore the set $H_A$ where both are $=0$ is
$\sigma(\lambda)$-measurable, too. $\lambda\in H_A$ means that $A(\lambda)$
is unitary. Now define
\[
A_u(\lambda)=
\begin{cases}
A(\lambda)&\text{for }\lambda\in H_A,\\
1&\text{for }\lambda\notin H_A,
\end{cases}
\]
then $A_u(\lambda)$ is $\sigma(\lambda)$-measurable, unitary (hence
$\lVert\!\lvert A_u(\lambda)\rvert\!\rVert=1$) and in $M(\lambda)$. Thus we
can define $A_u\sim\sum A_u(\lambda)$, and we have: $A_u\in M$, $A_u$
unitary, and $A_u(\lambda)=A(\lambda)$ for those $\lambda$'s for which
$A(\lambda)$ is unitary.

After these preparations we undertake the following construction: Choose an
$E_0\in M$. Take furthermore the projection $\overline E\in M$ from
\hyperref[lem:19]{Lemma 19}. Now form a sequence of operators
$A_1,A_2,\ldots$ with the following properties:
\begin{enumerate}[label=(\alph*)]
\item $1,\overline E,E_0$ are $A_m$'s.
\item $A_m^*$, $(A_m)_u$, $A_m\pm A_n$, $A_mA_n$ are $A_p$'s again.
\item If $A_m,A_n$ are projections, then the $V_1,V_2,V_3$ of
\hyperref[lem:21]{Lemma 21} (with $E=A_m,F=A_n$), and the
$A_m\mathbin{\#}A_n$ of \hyperref[lem:22]{Lemma 22} are $A_p$'s again.
\end{enumerate}
Choose for each $A_m$ a fixed decomposition
$A_m\sim\sum A_m(\lambda)$. Denote the projections and the unitary operators
in the sequence $A_1,A_2,\ldots$ respectively by $E_1,E_2,\ldots$ and
$U_1,U_2,\ldots$, and if $A_m$ is $E_n$ or $U_n$, denote $A_m(\lambda)$ by
$E_n(\lambda)$ or $U_n(\lambda)$, respectively.

Let $A_p$ be successively equal to each one of
$1,A_m^*,(A_m)_u,A_m\pm A_n,A_mA_n,A_m\mathbin{\#}A_n$, then we have
correspondingly for all $\lambda$
$A_p(\lambda)=1,A_m(\lambda)^*,(A_m)_u(\lambda),
A_m(\lambda)\pm A_n(\lambda),A_m(\lambda)A_n(\lambda),
A_m(\lambda)\mathbin{\#}A_n(\lambda)$, except for certain $\lambda$-sets of
$\sigma(\lambda)$-measure $0$, say
$N_3,N_4(m),N_5(m),N_6(m,n),N_7(m,n),N_8(m,n),N_9(m,n)$.
($N_1,N_2$, or to be more specific, $N_1(E,F),N_2(E,U)$, were defined
earlier, in this section, in connection with \eqref{eq:23.4} further above.)

Compare now two projections $E_m,E_n$. Denote the set of all $\lambda$'s
with $E_m(\lambda)=E_n(\lambda)$,
$\delta_M^{(\lambda)}(E_m)>\delta_M^{(\lambda)}(E_n)$, by $K_{m,n}$; it is
clearly $\sigma(\lambda)$-measurable. Now
$1_{K_{m,n}}(\lambda)E_m(\lambda)=1_{K_{m,n}}(\lambda)E_n(\lambda)$ for all
$\lambda$, so $1_{K_{m,n}}(\lambda)E_m=1_{K_{m,n}}(\lambda)E_n$,
\[
\begin{aligned}
\Delta_M(1_{K_{m,n}}(\lambda)E_m)
 &=\Delta_M(1_{K_{m,n}}(\lambda)E_n),\\
\int_{K_{m,n}}\delta_M^{(\lambda)}(E_m)\,d\sigma(\lambda)
 &=\int_{K_{m,n}}\delta_M^{(\lambda)}(E_n)\,d\sigma(\lambda).
\end{aligned}
\]
As $\delta_M^{(\lambda)}(E_m)>\delta_M^{(\lambda)}(E_n)$ for all
$\lambda\in K_{m,n}$, this necessitates
$\mu_{\sigma(\lambda)}(K_{m,n})=0$. Put
$N_{10}(m,n)=K_{m,n}+K_{n,m}$, then
$\mu_{\sigma(\lambda)}(N_{10}(m,n))=0$, and if
$\lambda\notin N_{10}(m,n)$, then $E_m(\lambda)=E_n(\lambda)$ implies
$\delta_M^{(\lambda)}(E_m)=\delta_M^{(\lambda)}(E_n)$.

Now put
\[
\begin{aligned}
N_{11}={}&\sum_{m,n=1}^{\infty}N_1(E_m,E_n)
+\sum_{m,n=1}^{\infty}N_2(E_m,U_n)+N_3\\
&+\sum_{m=1}^{\infty}\bigl(N_4(m)+N_5(m)\bigr)
+\sum_{m,n=1}^{\infty}\bigl(N_6(m,n)+\cdots+N_{10}(m,n)\bigr).
\end{aligned}
\]

\origpage{477}
Then $\mu_{\sigma(\lambda)}(N_{11})=0$ and
$\lambda\notin N_{11}$ has the following consequences:
\begin{enumerate}[label=(\alph*$'$)]
\item $\delta_M^{(\lambda)}(E_m+E_n)
=\delta_M^{(\lambda)}(E_m)+\delta_M^{(\lambda)}(E_n)$ for $E_mE_n=0$,
\item $\delta_M^{(\lambda)}(U_nE_mU_n^{-1})
=\delta_M^{(\lambda)}(E_m)$,
\item If $A_p$ is successively equal to each one of
$1,A_m^*,(A_m)_u,A_m\pm A_n,A_mA_n,A_m\mathbin{\#}A_n$, then
correspondingly for all $\lambda$
$A_p(\lambda)=1,A_m(\lambda)^*,(A_m)_u(\lambda),
A_m(\lambda)\pm A_n(\lambda),A_m(\lambda)A_n(\lambda),
A_m(\lambda)\mathbin{\#}A_n(\lambda)$.
\item If $E_m(\lambda)=E_n(\lambda)$ (for one particular $\lambda$!), then
$\delta_M^{(\lambda)}(E_m)=\delta_M^{(\lambda)}(E_n)$.
\end{enumerate}

Now introduce a ($\sigma(\lambda)$-measurable) relative dimensionality
function $d_{M(\lambda)}(E)$. Then we prove:

\begin{itshape}
If $E_m,E_n$ are given, then there exists a $U_p$ with the following
property. For every $\lambda\notin N_{11}$ for which
$d_{M(\lambda)}(E_m(\lambda))$ is finite and
$\leq d_{M(\lambda)}(E_n(\lambda))$, the projection
$U_p(\lambda)E_m(\lambda)U_p(\lambda)^{-1}$ is $\leq E_n(\lambda)$.
\end{itshape}

Indeed if for a certain $\lambda^*$,
$d_{M(\lambda)}(E_m(\lambda^*))\leq d_{M(\lambda)}(E_n(\lambda^*))$, then a
partially isometric $V\in M(\lambda)$ exists, which has the initial
projection $E_m(\lambda)$ and a final projection $\leq E_n(\lambda)$. Thus
if $A_r=V_2$ (in the sense of \hyperref[lem:21]{Lemma 21} for $E=E_m$ and
$F=E_n$), then $A_r(\lambda)$ is such a $V$, and its final projection is
$A_r(\lambda^*)A_r(\lambda^*)^*$. As $A_r$ is partially isometric, therefore
$A_rA_r^*$ is a projection. It is an $A_l$, and hence an $E_q$. Thus
$A_r(\lambda^*)A_r(\lambda^*)^*=E_q(\lambda^*)$. Now
$d_{M(\lambda)}(E_q(\lambda))=d_{M(\lambda)}(E_m(\lambda))$ is finite, so
$d_{M(\lambda)}(1-E_q(\lambda))=d_{M(\lambda)}(1-E_m(\lambda))$. Therefore a
partially isometric $W\in M(\lambda)$ exists, which has the initial and final
projections $1-E_m(\lambda)$ resp. $1-E_q(\lambda)$. Therefore if $A_s=V_1$
(in the sense of \hyperref[lem:21]{Lemma 21} for $E=1-E_m$, $F=1-E_q$), then
$A_s(\lambda)$ is such a $W$. Thus if $A_r+A_s=A_t$, then
$A_t(\lambda)=A_r(\lambda)+A_s(\lambda)$ is a unitary operator, which maps
$E_m(\lambda)$ on $E_q(\lambda)\leq E_n(\lambda)$. Now consider $(A_t)_u$,
which is unitary, say $U_p$. As $A_t(\lambda)$ is unitary we have
$U_p(\lambda)=(A_t)_u(\lambda)=A_t(\lambda)$. Thus $U_p,E_q$ meet our
requirements.

Next we prove:

\begin{itshape}
$\lambda\notin N_{11}$ and $E_m(\lambda)\leq E_n(\lambda)$ imply
$\delta_M^{(\lambda)}(E_m)\leq\delta_M^{(\lambda)}(E_n)$.
\end{itshape}

Indeed: Put $E_m\mathbin{\#}E_n=E_p$, then
$E_p(\lambda)=E_m(\lambda)\mathbin{\#}E_n(\lambda)=E_m(\lambda)$,
$\delta_M^{(\lambda)}(E_p)=\delta_M^{(\lambda)}(E_m)$. But
$E_p\leq E_n$, so $E_n-E_p$ is a projection, say $E_q$. Thus
$\delta_M^{(\lambda)}(E_n)=
\delta_M^{(\lambda)}(E_p+E_q)=
\delta_M^{(\lambda)}(E_p)+\delta_M^{(\lambda)}(E_q)
\geq\delta_M^{(\lambda)}(E_p)$, proving
$\delta_M^{(\lambda)}(E_m)\leq\delta_M^{(\lambda)}(E_n)$.

We now proceed to discuss the following question:

\begin{itshape}
For which (real) numbers $\theta>0$ does $\lambda\notin N_{11}$ and
$d_{M(\lambda)}(E_m(\lambda))\leq$ or
$\geq\theta d_{M(\lambda)}(E_n(\lambda))$ ($d_{M(\lambda)}(E_m)$ or
$d_{M(\lambda)}(E_n)$ being finite for $\leq$ or $\geq$, respectively)
imply $\delta_M^{(\lambda)}(E_m)\leq$ resp.
$\geq\theta\delta_M^{(\lambda)}(E_n)$?
\end{itshape}

We proceed stepwise:
\begin{enumerate}[label=(\Roman*)]
\item $\theta=1$ is such. As $\leq$ and $\geq$ differ only by an interchange
of $m$ and $n$, we consider only $\leq$. Then this follows by combining our
two above results.

\item $1/\theta$ is such, along with $\theta$. Obvious, as interchanging $m$
and $n$, as well as $\leq$ and $\geq$, replaces $\theta$ by $1/\theta$.

\item $\theta+1$ is such along with $\theta$. We have
$d_{M(\lambda)}(E_m(\lambda))\leq$ or
$\geq(\theta+1)d_{M(\lambda)}(E_n(\lambda))$, the finiteness restriction as
above and wish to prove
$\delta_M^{(\lambda)}(E_m)\leq$ resp.
$\geq(\theta+1)\delta_M^{(\lambda)}(E_n)$.

Assume first $d_{M(\lambda)}(E_m(\lambda))<d_{M(\lambda)}(E_n(\lambda))$.
This excludes the $\geq$ in the assumption, so it can only occur in the case
with $\leq$, and then $\delta_M^{(\lambda)}(E_m)$ is finite.

\origpage{478}
Now we have
$d_{M(\lambda)}(E_m(\lambda))\leq d_{M(\lambda)}(E_n(\lambda))$, so by (I)
$\delta_M^{(\lambda)}(E_m)\leq\delta_M^{(\lambda)}(E_n)$, and
\textit{a fortiori} $\leq(\theta+1)\delta_M^{(\lambda)}(E_n)$, disposing of
this case.

Assume next
$d_{M(\lambda)}(E_m(\lambda))\geq d_{M(\lambda)}(E_n(\lambda))$, which makes
$d_{M(\lambda)}(E_n(\lambda))$ finite in any event. Applying our previous
results, this implies the existence of a $U_p$ with
$U_p(\lambda)E_n(\lambda)U_p(\lambda)^{-1}\leq E_m(\lambda)$. (For this only:
We have interchanged the roles of $m$ and $n$.) Now form successively
$E_r=U_pE_nU_p^{-1}$, $E_s=E_m\mathbin{\#}E_r\leq E_m$,
$E_t=E_m-E_s$. Then
$d_{M(\lambda)}(E_m(\lambda))
=d_{M(\lambda)}(E_s(\lambda))+d_{M(\lambda)}(E_t(\lambda))$, and therefore
$\delta_M^{(\lambda)}(E_m)=
\delta_M^{(\lambda)}(E_s)+\delta_M^{(\lambda)}(E_t)$. But as
$E_s(\lambda)=E_r(\lambda)\mathbin{\#}E_m(\lambda)=E_r(\lambda)$ (owing to
$E_r(\lambda)=U_p(\lambda)E_n(\lambda)U_p(\lambda)^{-1}\leq E_m(\lambda)$),
we may replace $s$ by $r$ in these equations; and as
$E_r(\lambda)=U_p(\lambda)E_n(\lambda)U_p(\lambda)^{-1}$,
$E_r=U_pE_nU_p^{-1}$, even by $n$. Thus we have
\[
d_{M(\lambda)}(E_m(\lambda))
=d_{M(\lambda)}(E_t(\lambda))+d_{M(\lambda)}(E_n(\lambda)),
\quad
\delta_M^{(\lambda)}(E_m)
=\delta_M^{(\lambda)}(E_t)+\delta_M^{(\lambda)}(E_n).
\]
As $d_{M(\lambda)}(E_n(\lambda))$ is finite, the first equation leads
respectively to
$d_{M(\lambda)}(E_t(\lambda))\leq$ and
$\geq\theta d_{M(\lambda)}(E_n(\lambda))$. As our question is assumed to
have an affirmative answer for $\theta$, this implies respectively
$\delta_M^{(\lambda)}(E_t)\leq$ and
$\geq\theta\delta_M^{(\lambda)}(E_n)$ (apply to $t,n$ instead of $m,n$, the
finiteness of $d_{M(\lambda)}(E_m(\lambda))$ implies clearly that one of
$d_{M(\lambda)}(E_t(\lambda))$ for $\leq$), and now the second equation gives
respectively $\delta_M^{(\lambda)}(E_m)\leq$ and
$\geq(\theta+1)\delta_M^{(\lambda)}(E_n)$, as desired.

\item All rational $\theta$'s are such. By (I), (III) every
$\theta=k=1,2,\ldots$ is such; by (II), (III) $k+1/\theta$,
$k=0,1,2,\ldots$, is such, along with $\theta$. Thus all $\theta$ equal to a
finite continued fraction
\[
k_0+1/k_1+1/k_2+\cdots+1/k_\nu
\quad(k_0=0,1,2,\ldots,\quad k_1,\ldots,k_\nu=1,2,\ldots)
\]
have the property in question. That is: All rational $\theta$'s have it.

\item By continuity we can extend the property to all real $\theta>0$,
considering that (IV) stated it for all rational $\theta>0$. Thus if
$d_{M(\lambda)}(E_m(\lambda)),d_{M(\lambda)}(E_n(\lambda))>0$, $<+\infty$,
we can put
$\theta=d_{M(\lambda)}(E_m(\lambda))/d_{M(\lambda)}(E_n(\lambda))$ and apply
$\leq$ and $\geq$ simultaneously: This gives
\[
\delta_M^{(\lambda)}(E_m)
=\frac{d_{M(\lambda)}(E_m(\lambda))}
       {d_{M(\lambda)}(E_n(\lambda))}\delta_M^{(\lambda)}(E_n),
\]
\[
\frac{\delta_M^{(\lambda)}(E_m)}{d_{M(\lambda)}(E_m(\lambda))}
=\frac{\delta_M^{(\lambda)}(E_n)}{d_{M(\lambda)}(E_n(\lambda))}.
\]
\end{enumerate}

In other words:

\begin{itshape}
For all $\lambda\notin N_{11}$ we have
\begin{equation}
\frac{\delta_M^{(\lambda)}(E_m)}{d_{M(\lambda)}(E_m(\lambda))}
=\gamma(\lambda)\qquad\text{(independent of $m$!)},
\tag{23.6}\label{eq:23.6}
\end{equation}
whenever $d_{M(\lambda)}(E_m(\lambda))>0$, $<+\infty$.
\end{itshape}

$\gamma(\lambda)$ is $\geq0$, $\leq+\infty$ and
$\sigma(\lambda)$-measurable by its construction. Let $\overline K_0$ be the
set of those $\lambda$, where $\gamma(\lambda)=0$. If the projection
$\overline E$ of \hyperref[lem:19]{Lemma 19} is substituted for $E_m$, then
we have for $\lambda\in\overline M$
$d_{M(\lambda)}(\overline E(\lambda))>0$, $<+\infty$, so for
$\lambda\in\overline K_0\cdot\overline M$, we have
$\delta_M^{(\lambda)}(\overline E)=0$.

Thus
\[
\Delta_M(1_{\overline K_0\cdot\overline M}(\lambda)\overline E)
=\int_{\overline K_0\cdot\overline M}
 \delta_M^{(\lambda)}(\overline E)\,d\sigma(\lambda)=0,
\]
and therefore $1_{\overline K_0\cdot\overline M}(\lambda)\overline E=0$.
But
$1_{\overline K_0\cdot\overline M}(\lambda)\overline E
\sim\sum1_{\overline K_0\cdot\overline M}(\lambda)\overline E(\lambda)$, so
$1_{\overline K_0\cdot\overline M}(\lambda)\overline E(\lambda)=0$, except
for a $\lambda$-set of $\sigma(\lambda)$-measure $0$. Now for
$\lambda\in\overline K_0\cdot\overline M$,

\origpage{479}
$1_{\overline K_0\cdot\overline M}(\lambda)=1$, and as
$\lambda\in\overline M$, $\overline E(\lambda)\ne0$, therefore
$1_{\overline K_0\cdot\overline M}(\lambda)\overline E(\lambda)\ne0$. All
of this implies
$\mu_{\sigma(\lambda)}(\overline K_0\cdot\overline M)=0$. Putting
$N_{12}=N_{11}+\overline K_0\cdot\overline M$, we still have
$\mu_{\sigma(\lambda)}(N_{12})=0$, and we can state: For
$\lambda\notin N_{12}$, $\lambda\in\overline M$ implies
$\gamma(\lambda)\notin\overline K_0$, that is $\gamma(\lambda)>0$. So
$0<\gamma(\lambda)\leq+\infty$, and as
$0<d_{M(\lambda)}(E_m(\lambda))\leq+\infty$ too (if
$E_m(\lambda)\ne0$), we can write for \eqref{eq:23.6}:
\begin{equation}
\delta_M^{(\lambda)}(E_m)
=\gamma(\lambda)d_{M(\lambda)}(E_m(\lambda))
\quad\text{if }\lambda\notin N_{12},\ \lambda\in\overline M,
\qquad0<\gamma(\lambda)\leq+\infty.
\tag{23.7}\label{eq:23.7}
\end{equation}
Here $d_{M(\lambda)}(E_m(\lambda))=0$ or $=+\infty$ is to be excluded. But
$d_{M(\lambda)}(E_m(\lambda))=0$ means $E_m(\lambda)=0$, and so
$\delta_M^{(\lambda)}(E_m)=\delta_M^{(\lambda)}(0)=0$, so that
\eqref{eq:23.7} holds in this case too, if we use again the convention
$0\cdot+\infty=0$. If $d_{M(\lambda)}(E_m(\lambda))=+\infty$, then we have
$d_{M(\lambda)}(E_m(\lambda))\geq
\theta d_{M(\lambda)}(\overline E(\lambda))$ for all $\theta>0$, so
\[
\delta_M^{(\lambda)}(E_m)
\geq\theta\delta_M^{(\lambda)}(\overline E)
=\theta\gamma(\lambda)d_{M(\lambda)}(\overline E(\lambda))
\]
for all $\theta>0$. As the second and third factors are $>0$, this gives
$\delta_M^{(\lambda)}(E_m)=+\infty$, and thus \eqref{eq:23.7} holds again.
Thus \eqref{eq:23.7} holds for all $E_m$'s.

\begin{enumerate}[label=(\Roman*),start=6]
\item We can use \eqref{eq:23.7} to express $\gamma(\lambda)$ in
$\overline M$, if we substitute the projection $\overline E$ of
\hyperref[lem:19]{Lemma 19} for $E_m$. Then we see: $\gamma(\lambda)$ is a
$\sigma(\lambda)$-measurable function in $\overline M$, and, disregarding a
$\lambda$-set of $\sigma(\lambda)$-measure $0$, we have:
\begin{equation}
\gamma(\lambda)=
\frac{\delta_M^{(\lambda)}(\overline E)}
     {d_{M(\lambda)}(\overline E(\lambda))}
\qquad\text{for }\lambda\in\overline M.
\tag{23.7$_1$}\label{eq:23.7.1}
\end{equation}

\item We next discuss the behavior of $\delta_M^{(\lambda)}(E_m)$ for
$\lambda\notin\overline M$. We again assume $\lambda\notin N_{12}$, too.
We know from \hyperref[lem:19]{Lemma 19} that the set of all
$\lambda\notin\overline M$ with
$d_{M(\lambda)}(E_m(\lambda))>0$, $<+\infty$, say $N_{13}(m)$, has the
$\sigma(\lambda)$-measure $0$. So if we put
$N_{14}=N_{12}+\sum_{m=1}^{\infty}N_{13}(m)$, then
$\mu_{\sigma(\lambda)}(N_{14})=0$, and we see: If
$\lambda\notin N_{14}$ and $\lambda\notin\overline M$, then either
$E_m(\lambda)=0$ (that is $d_{M(\lambda)}(E_m(\lambda))=0$) or
$E_m(\lambda)$ is infinite (with respect to $M(\lambda)$, that is
$d_{M(\lambda)}(E_m(\lambda))=+\infty$). In both cases there are two
partially isometric $V,V'\in M(\lambda)$, which have both the initial
projection $E_m(\lambda)$, and final projections with the sum $E_m(\lambda)$:
For $E_m(\lambda)=0$ these are $0,0$; for $E_m(\lambda)$ infinite apply
\cite[Lemma~7.2.3]{ref10}. Now form $A_p=V_3$ (in the sense of
\hyperref[lem:21]{Lemma 21} for $E=E_m$). As $A_p$ is partially isometric,
$A_pA_p^*$ is a projection, say $E_q$. Thus
$E_q(\lambda)=A_p(\lambda)A_p(\lambda)^*$ is
$\sim E_m(\lambda)\vnmod{M(\lambda)}$ and so is $E_m(\lambda)-E_q(\lambda)$.
Furthermore $E_q\mathbin{\#}E_m$ is a projection say $E_r$, and
$\leq E_m$, thus $E_m-E_r$ too is a projection, say $E_s$. Now
$E_q(\lambda)\leq E_m(\lambda)$, therefore
$E_q(\lambda)\mathbin{\#}E_m(\lambda)=E_q(\lambda)$, therefore
$E_s(\lambda)=E_m(\lambda)-E_q(\lambda)$. Thus $E_s(\lambda)$ too is
$\sim E_m(\lambda)\vnmod{M(\lambda)}$.

So we see:
\[
E_m=E_q+E_s,\qquad
E_m(\lambda)\sim E_q(\lambda)\sim E_s(\lambda)\vnmod{M(\lambda)}.
\]
Similarly we can construct for another $E_n$, two $E_t,E_w$ with
\[
E_n=E_t+E_w,\qquad
E_n(\lambda)\sim E_t(\lambda)\sim E_w(\lambda)\vnmod{M(\lambda)}.
\]
Thus if $E_m(\lambda),E_n(\lambda)\ne0$, then both are infinite, and with
them $E_q(\lambda),E_s(\lambda),E_t(\lambda),E_w(\lambda)$ as well as
$1-E_q(\lambda)\geq E_s(\lambda)$ and
$1-E_t(\lambda)\geq E_w(\lambda)$ are infinite, too. Thus partially
isometric $W,W'$'s with the initial projections $E_q(\lambda)$,

\origpage{480}
$1-E_q(\lambda)$ and the final projections $E_t(\lambda)$,
$1-E_t(\lambda)$ exist. If $A_x,A_y$ are the resp. $V_1$ (in the sense of
\hyperref[lem:21]{Lemma 21}, with $E=E_q$ resp. $1-E_q$, and $F=E_t$ resp.
$1-E_t$), then $A_x(\lambda),A_y(\lambda)$ are such $W,W'$'s, and thus if
$A_z=A_x+A_y$, then $A_z(\lambda)=A_x(\lambda)+A_y(\lambda)$ is unitary, and
maps $E_q(\lambda)$ on $E_t(\lambda)$. $(A_z)_u$ is unitary, say $U_a$, and
as $A_z(\lambda)$ is unitary, $U_a(\lambda)=(A_z)_u(\lambda)=A_z(\lambda)$.
So we have $U_a(\lambda)E_q(\lambda)U_a(\lambda)^{-1}=E_t(\lambda)$. Put
$U_aE_qU_a^{-1}=E_b$, then $E_b(\lambda)=E_t(\lambda)$. So we have:
$\delta_M^{(\lambda)}(E_q)=\delta_M^{(\lambda)}(E_b)
=\delta_M^{(\lambda)}(E_t)$. Similarly
$\delta_M^{(\lambda)}(E_s)=\delta_M^{(\lambda)}(E_w)$. Thus
\[
\begin{aligned}
\delta_M^{(\lambda)}(E_m)
 &=\delta_M^{(\lambda)}(E_q+E_s)
  =\delta_M^{(\lambda)}(E_q)+\delta_M^{(\lambda)}(E_s)\\
 &=\delta_M^{(\lambda)}(E_t)+\delta_M^{(\lambda)}(E_w)
  =\delta_M^{(\lambda)}(E_t+E_w)
  =\delta_M^{(\lambda)}(E_n).
\end{aligned}
\]
In other words: $\delta_M^{(\lambda)}(E_m)$ has the same value for all
$E_m(\lambda)\ne0$, say $\delta$. Now this applies to $E_m,E_q,E_s$ also.
Then $\delta_M^{(\lambda)}(E_m)
=\delta_M^{(\lambda)}(E_q)+\delta_M^{(\lambda)}(E_s)$,
$\delta=2\delta$, so that is: $\delta=0$ or $+\infty$. As $1\ne0$, these
alternatives are equivalent to $\delta_M^{(\lambda)}(1)=0$ or $+\infty$.

Denote the set of all $\lambda$'s where $\delta_M^{(\lambda)}(1)=0$ by
$\overline K_1$, it is $\sigma(\lambda)$-measurable along with
$\delta_M^{(\lambda)}(1)$. Then \eqref{eq:23.3} gives
\[
\Delta_M(1_{\overline K_1}(\lambda))
=\int_{\overline K_1}\delta_M^{(\lambda)}(1)\,d\sigma(\lambda)=0,
\qquad1_{\overline K_1}(\lambda)=0,
\]
and therefore $\mu_{\sigma(\lambda)}(\overline K_1)=0$. Put
$N_{15}=N_{14}+\overline K_1$, then
$\mu_{\sigma(\lambda)}(N_{15})=0$, and we have:
\[
\delta_M^{(\lambda)}(E_m)=
\begin{cases}
+\infty&\text{for }E_m(\lambda)\ne0,\\
0&\text{for }E_m(\lambda)=0,
\end{cases}
\qquad\text{if }\lambda\notin N_{15},\ \lambda\notin\overline M.
\]
In other words: \eqref{eq:23.7} holds again (with $N_{15}$ instead of
$N_{12}$), but now we have
\begin{equation}
\gamma(\lambda)=+\infty\qquad\text{for }\lambda\notin\overline M.
\tag{23.7$_2$}\label{eq:23.7.2}
\end{equation}
\end{enumerate}

Putting \eqref{eq:23.6} and \eqref{eq:23.3} together we obtain
\begin{equation}
\tag{A}\label{eq:23.A}
\Delta_M(1_K(\lambda)E_m)
=\int_K\gamma(\lambda)d_{M(\lambda)}(E_m(\lambda))\,d\sigma(\lambda).
\end{equation}
And substituting $E_0$ for $E_m$, and specializing $K$ to $r$ (as
$1_r(\lambda)=1$):
\begin{equation}
\tag{B}\label{eq:23.B}
\Delta_M(E_0)
=\int\gamma(\lambda)d_{M(\lambda)}(E_0(\lambda))\,d\sigma(\lambda).
\end{equation}

The entire construction which lead to \hyperref[eq:23.B]{(B)} depends on the projection
$E_0\in M$. But the result does not: \hyperref[eq:23.B]{(B)} expresses the given weight-function
$\Delta_M(E)$ (for $E=E_0$) by means of the arbitrarily (but
$\sigma(\lambda)$-measurably) prescribed $\sigma(\lambda)$-measurable
relative dimensionality function $d_{M(\lambda)}(E)$, with the sole help of
the $\sigma(\lambda)$-measurable function $\gamma(\lambda)>0$,
$\leq+\infty$. And this $\gamma(\lambda)$ is defined by
\eqref{eq:23.7.1} and \eqref{eq:23.7.2}, in which $E_0$ does not enter. Thus
$\gamma(\lambda)$ is a fixed function, and \hyperref[eq:23.B]{(B)} holds for all projections
$E_0\in M$ (with the same $\gamma(\lambda)$).

We can now prove easily:

\par\medskip\noindent\textbf{Theorem X.}\phantomsection\label{thm:X}\quad
\begin{itshape}
If $\Delta_{M(\lambda)}(E)$ is a weight function with respect to
$M(\lambda)$, depending $\sigma(\lambda)$-measurably on $\lambda$, then
\begin{equation}
\Delta_M(E)=\int\Delta_{M(\lambda)}(E(\lambda))\,d\sigma(\lambda),
\qquad\text{where }E\sim\sum E(\lambda),
\tag{23.8}\label{eq:23.8}
\end{equation}
is a weight function with respect to $M$. Conversely, all weight functions
$\Delta_M(E)$ (with respect to $M$) can be obtained by this process.
\end{itshape}

\par\smallskip\noindent\textit{Proof:}
One verifies easily the conditions (i)--(iii) of
\hyperref[def:7]{Definition 7} for

\origpage{481}
$\Delta_M(E)$ (and $M$) if they do hold for all
$\Delta_{M(\lambda)}(E)$ (and $M(\lambda)$). This proves the first statement.

The second statement follows from the considerations which preceded this
theorem, particularly from the final formula \hyperref[eq:23.B]{(B)} and the comments following
it. We need only to put
$\Delta_{M(\lambda)}(E)=\gamma(\lambda)d_{M(\lambda)}(E)$, and remember
\hyperref[thm:VIII]{Theorem VIII}. Thus the proof is complete.

\section{Consequences for weights and equivalence}\label{sec:24}
\hyperref[thm:X]{Theorem X} permits some simple inferences which have an
interest of their own, and which we are going to discuss now.

\noindent$(\alpha)$ Combine the results of \hyperref[thm:VIII]{Theorems VIII}
and \hyperref[thm:X]{X}. Then we find: The formula
\begin{equation}
\Delta_M(E)=\int a(\lambda)d_{M(\lambda)}(E(\lambda))\,d\sigma(\lambda)
\qquad\text{for }E\sim\sum E(\lambda)
\tag{24.1}\label{eq:24.1}
\end{equation}
expresses all weight functions $\Delta_M(E)$ (with respect to $M$), if
$a(\lambda)$ runs over all those functions which fulfill
$0<a(\lambda)\leq+\infty$ everywhere, and which are
$\sigma(\lambda)$-measurable in $\overline M$. It is natural to ask, how far
$\Delta_M(E)$ determines $a(\lambda)$.

If $K\subset\overline M$ and $\sigma(\lambda)$-measurable, then
\[
\Delta_M(1_K(\lambda)E)
=\int_Ka(\lambda)d_{M(\lambda)}(E(\lambda))\,d\sigma(\lambda),
\]
thus
$\int_Ka(\lambda)d_{M(\lambda)}(\overline E(\lambda))\,d\sigma(\lambda)$ is
determined for all these sets $K$. Thus
$a(\lambda)d_{M(\lambda)}(\overline E(\lambda))$ is uniquely determined in
$\overline M$, up to a set of $\sigma(\lambda)$-measure $0$; and as
$d_{M(\lambda)}(\overline E(\lambda))>0$, $<+\infty$ for all
$\lambda\in\overline M$, therefore $a(\lambda)$ is determined to the same
extent. In $r-\overline M$ we have by \hyperref[lem:19]{Lemma 19} always
$d_{M(\lambda)}(\overline E(\lambda))=0$ or $+\infty$, except for a set of
$\sigma(\lambda)$-measure $0$, and therefore the numerical value of
$a(\lambda)$ (which must be $>0$, $\leq+\infty$) does not matter at all. So
we see:

\begin{itshape}
$a(\lambda)$ may be any function which fulfills
$0<a(\lambda)\leq+\infty$ and which is $\sigma(\lambda)$-measurable in
$\overline M$. $\Delta_M(E)$ determines $a(\lambda)$ uniquely in
$\overline M$, up to a set of $\sigma(\lambda)$-measure $0$, while it leaves
it entirely arbitrary in $r-\overline M$.
\end{itshape}

The same observation, as the one we made after
\hyperref[thm:VIII]{Theorem VIII} in \secref{sec:21}, applies here: While it
is obvious that the numerical value of $a(\lambda)$ is irrelevant if
$\lambda\in C_{III}=C_{(III,III)}$ (because then
$d_{M(\lambda)}(E)=0$ or $+\infty$ for all $E\in M(\lambda)$), it is
surprising, that this is the case even for $\lambda\in r-\overline M$, where
$r-\overline M\subset C_{(III,III)}$. We discussed the relationship of
$r-\overline M$ and $C_{(III,III)}$ loc. cit.

\noindent$(\beta)$ We have for every partially isometric
$V\in M(\lambda)$
$d_{M(\lambda)}(V^*V)=d_{M(\lambda)}(VV^*)$, by
\cite[Definition~8.2.1 and Theorem~VIII]{ref10}. Formula \eqref{eq:24.1}
extends this at once to all partially isometric $V\in M$: Then
$V\sim\sum V(\lambda)$, $V(\lambda)\in M(\lambda)$, $V(\lambda)$ partially
isometric, thus
$d_{M(\lambda)}(V(\lambda)^*V(\lambda))
=d_{M(\lambda)}(V(\lambda)V(\lambda)^*)$, and so by \eqref{eq:24.1}
$\Delta_M(V^*V)=\Delta_M(VV^*)$. Thus the condition

\begin{itshape}
(iii)$'$ for every partially isometric $V\in M$,
$\Delta_M(V^*V)=\Delta_M(VV^*)$,
\end{itshape}

follows from \hyperref[def:7]{Definition 7}, (i)--(iii). It clearly implies
(iii). It suffices to put, as at the beginning of the proof of
\hyperref[lem:16]{Lemma 16}, $V=UE$. Then
$V^*=E^*U^*=EU^{-1}$, and so $V^*V=E$, $VV^*=UEU^{-1}$, giving (iii). So we
see:

\begin{itshape}
Condition (iii) in \hyperref[def:7]{Definition 7} is equivalent to (iii)$'$
(conserving (i), (ii)).
\end{itshape}

Another obvious way to formulate (iii)$'$ is this:

\begin{itshape}
(iii)$''$ If $E,F$ are the initial resp. final projections of a partially
isometric

\origpage{482}
$V\in M$ (that is: if $V$ maps their resp. linear, closed sets on each other;
or, using the notations of \cite[Definition~6.1.1]{ref10}: if
$E\sim F\vnmod{M}$), then $\Delta_M(E)=\Delta_M(F)$.
\end{itshape}

\noindent$(\gamma)$ We observed in the \hyperref[rem:21]{Remark} following
\hyperref[lem:16]{Lemma 16}, that for factors the condition (ii) of
\hyperref[def:7]{Definition 7} can be replaced by the weaker condition
(ii)$'$: (ii)$'$ arises from (ii) by permitting two addends $E_j$ ($j=1,2$)
only. For a general $M$ this is not the case, that is: (ii)$'$ is then
essentially weaker than (ii). This appears if we consider the opposite
extreme to a factor $M$: an Abelian ring $M$.

The Abelian character of $M$ implies: $M\subset M'$,
$M=M\cdot M'=P\approx\sum0(\lambda)$ (cf. the remark in the proof of
\hyperref[thm:VII]{Theorem VII}); conversely if
$M\approx\sum0(\lambda)$, then $M$ is Abelian as all $0(\lambda)$'s are
Abelian. $0(\lambda)$ is obviously the only factor of class $C_{I_1}$, so we
see:

\begin{itshape}
$M$ is Abelian, if and only if $C_{I_1}=r$ (up to a set of
$\sigma(\lambda)$-measure $0$, as usual).
\end{itshape}

Now the general element of $M$ is (as $M=P$) $\alpha(\lambda)$, and this is
a projection if $\alpha(\lambda)^2=\alpha(\lambda)=\overline{\alpha(\lambda)}$,
that is if $\alpha(\lambda)=0$ or $1$. In other words:
\[
\alpha(\lambda)=1_K(\lambda)=
\begin{cases}
1&\text{for }\lambda\in K,\\
0&\text{for }\lambda\notin K,
\end{cases}
\]
where $K$ is any $\sigma(\lambda)$-measurable set. Thus a $\Delta_M(E)$,
$E\in M$, is defined for the $E=1_K(\lambda)$ only, and so we may write:
$\Delta_M(1_K(\lambda))=n(K)\geq0$, $\leq+\infty$. Now the conditions of
\hyperref[def:7]{Definition 7} become:

\noindent(i): $n(K)$ depends on $1_K(\lambda)$ only, that is, it is
unchanged, if $K$ is changed by a set of $\sigma(\lambda)$-measure $0$, and
$n(K)=0$ is equivalent to $1_K(\lambda)=0$, that is, to
$\mu_{\sigma(\lambda)}(K)=0$.

\noindent(ii) $n(K_1+K_2+\cdots)=n(K_1)+n(K_2)+\cdots$ if $K_jK_k=0$ for
$j\ne k$.

\noindent(ii)$'$ $n(K_1+K_2)=n(K_1)+n(K_2)$ for $K_1K_2=0$.

\noindent(iii) Void.

Thus the problem of finding a weight-function coincides in the case of an
Abelian $M$ with the problem of finding a Lebesgue-measure, which is
absolutely continuous with respect to the $\sigma(\lambda)$-measure, and
conversely. (Cf. the notion of equivalence in \secref{sec:2}, $(\delta)$.)
Now it is well known that the ``countable additivity'' postulated by (ii) is
an essential feature of Lebesgue-measure, and cannot be replaced by the
``finite additivity'' which (ii)$'$ expresses.\footnote{This is even not
possible in the totally discontinuous case, discussed in \secref{sec:4},
where the points $\lambda_1,\lambda_2,\ldots$ (a finite or infinite sequence)
have the $\sigma(\lambda)$-measures $a_1,a_2,\ldots$ (all $>0$,
$\sum a_\nu$ finite), and $r-(\lambda_1,\lambda_2,\ldots)$ has the measure
$0$: We assume that the sequence $\lambda_1,\lambda_2,\ldots$ is infinite.
Now a construction of Banach permits to define a function
$n'(L)\geq0$, $\leq1$ for all sets $L\subset(1,2,\ldots)$ with the following
properties:
\[
n'(L)=0\text{ if $L$ is finite};\qquad n'(1,2,\ldots)=1;\qquad
n'(L_1+L_2)=n'(L_1)+n'(L_2)\text{ if }L_1L_2=0.
\]
Then denote for each $K\subset r$ the set of all $\nu=1,2,\ldots$ with
$\lambda_\nu\in K$ by $L_K$, and define
\[
n(K)=\sum_{\gamma\in L_K}\frac1{2^\gamma}+n'(L_K).
\]
This $n(K)$ fulfills (i), (ii)$'$ ((iii) is void), but it violates (ii),
because
\[
\sum_{\nu=1}^{\infty}n(\lambda_\nu)
=\sum_{\nu=1}^{\infty}\frac1{2^\nu}=1,
\qquad n(\lambda_1,\lambda_2,\ldots)=2.
\]}

\origpage{483}
\noindent$(\delta)$ We can answer the following question:

\begin{itshape}
If $E,F$ are projections in $M$, when does a partially isometric $V\in M$
with the initial projection $E$ and a final projection $\leq F$ exist?
\end{itshape}

In the notation of \cite[Definition~6.3.1]{ref10}, this means: When is
$E\lesssim F\vnmod{M}$? Loc. cit. the problem was solved for factors. Then the
relative dimensionality function $d_M(G)$ gives the answer,
$E\lesssim F\vnmod{M}$ being equivalent to $d_M(E)\leq d_M(F)$. We will now
find the answer for a general $M$.

Our question is: When does a $V\in M$ with $V^*V=E$, $VV^*\leq F$ exist?
Putting $E\sim\sum E(\lambda)$, $F\sim\sum F(\lambda)$,
$V\sim\sum V(\lambda)$, we see that we need a partially isometric
$V(\lambda)\in M(\lambda)$, depending $\sigma(\lambda)$-measurably on
$\lambda$, with $V(\lambda)^*V(\lambda)=E(\lambda)$,
$V(\lambda)V(\lambda)^*\leq F(\lambda)$ for each $\lambda$ (except, of
course, for a $\lambda$-set of $\sigma(\lambda)$-measure $0$). Thus
$d_{M(\lambda)}(E(\lambda))\leq d_{M(\lambda)}(F(\lambda))$ (with the same
exceptions) is necessary. But it is sufficient, too, because it involves
for each $\lambda$ (with the same exceptions as above) the existence of a
partially isometric $V'\in M(\lambda)$ with $V'^*V'=E(\lambda)$,
$V'V'^*\leq F(\lambda)$, and then the $V_2$ of
\hyperref[lem:21]{Lemma 21} has the desired properties:
$V=V_2$, $V\sim\sum V(\lambda)$.

So we have proved:

\begin{itshape}
The necessary and sufficient condition is this: For all $\lambda$, except
for a set of $\sigma(\lambda)$-measure $0$,
$d_{M(\lambda)}(E(\lambda))\leq d_{M(\lambda)}(F(\lambda))$, where
$E\sim\sum E(\lambda)$, $F\sim\sum F(\lambda)$.
\end{itshape}

We saw in $(\gamma)$ above that the projections $\widetilde E\in P$ coincide
with the $\widetilde E=1_K(\lambda)$, $K$ any
$\sigma(\lambda)$-measurable set,
\[
1_K(\lambda)=
\begin{cases}
1&\text{for }\lambda\in K,\\
0&\text{for }\lambda\notin K.
\end{cases}
\]
(The considerations there assumed that $M$ is Abelian, but that part which
dealt with $P$ only made no use of this.) Hence
$E\sim\sum E(\lambda)$ implies
$\widetilde EE\sim\sum1_K(\lambda)E(\lambda)$. Given any weight function
$\Delta_M(G)$ and using the formula \eqref{eq:24.1}, we have therefore
\begin{equation}
\Delta_M(\widetilde EE)
=\int_Ka(\lambda)d_{M(\lambda)}(E(\lambda))\,d\sigma(\lambda).
\tag{24.2}\label{eq:24.2}
\end{equation}
Applying \eqref{eq:24.2} to $E$ and to $F$, our above condition, i.e.
$d_{M(\lambda)}(E(\lambda))\leq d_{M(\lambda)}(F(\lambda))$ for all
$\lambda$ excepting a set of $\sigma(\lambda)$-measure $0$, is seen to imply
\begin{equation}
\left\{
\begin{gathered}
\Delta_M(\widetilde EE)\leq\Delta_M(\widetilde EF)\\
\text{for all weight functions $\Delta_M(G)$ and all projections
$\widetilde E\in P$.}
\end{gathered}
\right.
\tag{24.3}\label{eq:24.3}
\end{equation}

Consider the converse: \eqref{eq:24.3} for all projections
$\widetilde E\in P$, and we may even restrict it to one $\Delta_M(G)$,
provided that the $a(\lambda)$ of \eqref{eq:24.1} is always finite. Then
\eqref{eq:24.2} for $E$ and $F$ gives
\[
\int_Ka(\lambda)d_{M(\lambda)}(E(\lambda))\,d\sigma(\lambda)
\leq\int_Ka(\lambda)d_{M(\lambda)}(F(\lambda))\,d\sigma(\lambda)
\]
for every $\sigma(\lambda)$-measurable set $K$. Now let $K$ be the set of all
$\lambda$ with
\[
d_{M(\lambda)}(E(\lambda))>d_{M(\lambda)}(F(\lambda)).
\]

\origpage{484}
Since $a(\lambda)$ is always finite, this implies
\[
\int_Ka(\lambda)d_{M(\lambda)}(E(\lambda))\,d\sigma(\lambda)
>\int_Ka(\lambda)d_{M(\lambda)}(F(\lambda))\,d\sigma(\lambda)
\]
provided that $\mu_{\sigma(\lambda)}(K)>0$. This, however, is a
contradiction, hence $\mu_{\sigma(\lambda)}(K)=0$. That is,
$d_{M(\lambda)}(E(\lambda))\leq d_{M(\lambda)}(F(\lambda))$, except for a
set of $\sigma(\lambda)$-measure $0$. This is our above condition.

We assumed that $a(\lambda)$ is always finite. We know, however, that
changing $a(\lambda)$ on a set of $\sigma(\lambda)$-measure $0$ or outside
$\overline M$ does not matter. (Concerning the latter, cf. $(\alpha)$ above.)

So we have obtained:

\begin{itshape}
Another necessary and sufficient condition is this \eqref{eq:24.3} above.
It is even sufficient to postulate \eqref{eq:24.3} for one given weight
function $\Delta_M(G)$ only, provided that the $a(\lambda)$ of
\eqref{eq:24.2} is finite for all $\lambda\in\overline M$, except for a
$\lambda$-set of $\sigma(\lambda)$-measure $0$.
\end{itshape}

\noindent$(\epsilon)$ A similar question is this:

\begin{itshape}
When does a partially isometric operator $V\in M$ with the initial and final
projections $E$ resp. $F$ exist?
\end{itshape}

In the notation of \cite[Definition~6.1.1]{ref10}: When is
$E\sim F\vnmod{M}$? The same procedure as in $(\delta)$ (only using the $V_1$
of \hyperref[lem:21]{Lemma 21} instead of the $V_2$) gives two necessary and
sufficient conditions, which obtain from those of $(\delta)$ by replacing
the relation $\leq$ in each case by the relation $=$. Another way to derive
this result obtains by observing that $E\sim F\vnmod{M}$ is equivalent to the
conjunction of $E\lesssim F\vnmod{M}$ and of $F\lesssim E\vnmod{M}$ and then by
applying $(\delta)$ to $E,F$ and $F,E$.


\end{document}
